% Options for packages loaded elsewhere
\PassOptionsToPackage{unicode}{hyperref}
\PassOptionsToPackage{hyphens}{url}
\PassOptionsToPackage{dvipsnames,svgnames,x11names}{xcolor}
%
\documentclass[
  11pt,
  11pt,
  a4paper,
  openany]{book}
\usepackage{amsmath,amssymb}
\usepackage{setspace}
\usepackage{iftex}
\ifPDFTeX
  \usepackage[T1]{fontenc}
  \usepackage[utf8]{inputenc}
  \usepackage{textcomp} % provide euro and other symbols
\else % if luatex or xetex
  \usepackage{unicode-math} % this also loads fontspec
  \defaultfontfeatures{Scale=MatchLowercase}
  \defaultfontfeatures[\rmfamily]{Ligatures=TeX,Scale=1}
\fi
\usepackage{lmodern}
\ifPDFTeX\else
  % xetex/luatex font selection
  \setmainfont[]{TeX Gyre Pagella}
  \setsansfont[]{TeX Gyre Heros}
  \setmonofont[]{DejaVu Sans Mono}
  \setmathfont[]{Latin Modern Math}
\fi
% Use upquote if available, for straight quotes in verbatim environments
\IfFileExists{upquote.sty}{\usepackage{upquote}}{}
\IfFileExists{microtype.sty}{% use microtype if available
  \usepackage[]{microtype}
  \UseMicrotypeSet[protrusion]{basicmath} % disable protrusion for tt fonts
}{}
\makeatletter
\@ifundefined{KOMAClassName}{% if non-KOMA class
  \IfFileExists{parskip.sty}{%
    \usepackage{parskip}
  }{% else
    \setlength{\parindent}{0pt}
    \setlength{\parskip}{6pt plus 2pt minus 1pt}}
}{% if KOMA class
  \KOMAoptions{parskip=half}}
\makeatother
\usepackage{xcolor}
\usepackage[margin=2.4cm,top=2.6cm,bottom=2.6cm]{geometry}
\usepackage{color}
\usepackage{fancyvrb}
\newcommand{\VerbBar}{|}
\newcommand{\VERB}{\Verb[commandchars=\\\{\}]}
\DefineVerbatimEnvironment{Highlighting}{Verbatim}{commandchars=\\\{\}}
% Add ',fontsize=\small' for more characters per line
\newenvironment{Shaded}{}{}
\newcommand{\AlertTok}[1]{\textcolor[rgb]{1.00,0.00,0.00}{\textbf{#1}}}
\newcommand{\AnnotationTok}[1]{\textcolor[rgb]{0.38,0.63,0.69}{\textbf{\textit{#1}}}}
\newcommand{\AttributeTok}[1]{\textcolor[rgb]{0.49,0.56,0.16}{#1}}
\newcommand{\BaseNTok}[1]{\textcolor[rgb]{0.25,0.63,0.44}{#1}}
\newcommand{\BuiltInTok}[1]{\textcolor[rgb]{0.00,0.50,0.00}{#1}}
\newcommand{\CharTok}[1]{\textcolor[rgb]{0.25,0.44,0.63}{#1}}
\newcommand{\CommentTok}[1]{\textcolor[rgb]{0.38,0.63,0.69}{\textit{#1}}}
\newcommand{\CommentVarTok}[1]{\textcolor[rgb]{0.38,0.63,0.69}{\textbf{\textit{#1}}}}
\newcommand{\ConstantTok}[1]{\textcolor[rgb]{0.53,0.00,0.00}{#1}}
\newcommand{\ControlFlowTok}[1]{\textcolor[rgb]{0.00,0.44,0.13}{\textbf{#1}}}
\newcommand{\DataTypeTok}[1]{\textcolor[rgb]{0.56,0.13,0.00}{#1}}
\newcommand{\DecValTok}[1]{\textcolor[rgb]{0.25,0.63,0.44}{#1}}
\newcommand{\DocumentationTok}[1]{\textcolor[rgb]{0.73,0.13,0.13}{\textit{#1}}}
\newcommand{\ErrorTok}[1]{\textcolor[rgb]{1.00,0.00,0.00}{\textbf{#1}}}
\newcommand{\ExtensionTok}[1]{#1}
\newcommand{\FloatTok}[1]{\textcolor[rgb]{0.25,0.63,0.44}{#1}}
\newcommand{\FunctionTok}[1]{\textcolor[rgb]{0.02,0.16,0.49}{#1}}
\newcommand{\ImportTok}[1]{\textcolor[rgb]{0.00,0.50,0.00}{\textbf{#1}}}
\newcommand{\InformationTok}[1]{\textcolor[rgb]{0.38,0.63,0.69}{\textbf{\textit{#1}}}}
\newcommand{\KeywordTok}[1]{\textcolor[rgb]{0.00,0.44,0.13}{\textbf{#1}}}
\newcommand{\NormalTok}[1]{#1}
\newcommand{\OperatorTok}[1]{\textcolor[rgb]{0.40,0.40,0.40}{#1}}
\newcommand{\OtherTok}[1]{\textcolor[rgb]{0.00,0.44,0.13}{#1}}
\newcommand{\PreprocessorTok}[1]{\textcolor[rgb]{0.74,0.48,0.00}{#1}}
\newcommand{\RegionMarkerTok}[1]{#1}
\newcommand{\SpecialCharTok}[1]{\textcolor[rgb]{0.25,0.44,0.63}{#1}}
\newcommand{\SpecialStringTok}[1]{\textcolor[rgb]{0.73,0.40,0.53}{#1}}
\newcommand{\StringTok}[1]{\textcolor[rgb]{0.25,0.44,0.63}{#1}}
\newcommand{\VariableTok}[1]{\textcolor[rgb]{0.10,0.09,0.49}{#1}}
\newcommand{\VerbatimStringTok}[1]{\textcolor[rgb]{0.25,0.44,0.63}{#1}}
\newcommand{\WarningTok}[1]{\textcolor[rgb]{0.38,0.63,0.69}{\textbf{\textit{#1}}}}
\usepackage{longtable,booktabs,array}
\usepackage{calc} % for calculating minipage widths
% Correct order of tables after \paragraph or \subparagraph
\usepackage{etoolbox}
\makeatletter
\patchcmd\longtable{\par}{\if@noskipsec\mbox{}\fi\par}{}{}
\makeatother
% Allow footnotes in longtable head/foot
\IfFileExists{footnotehyper.sty}{\usepackage{footnotehyper}}{\usepackage{footnote}}
\makesavenoteenv{longtable}
\setlength{\emergencystretch}{3em} % prevent overfull lines
\providecommand{\tightlist}{%
  \setlength{\itemsep}{0pt}\setlength{\parskip}{0pt}}
\setcounter{secnumdepth}{5}
\ifLuaTeX
\usepackage[bidi=basic]{babel}
\else
\usepackage[bidi=default]{babel}
\fi
\babelprovide[main,import]{french}
\ifPDFTeX
\else
\babelfont{rm}[]{TeX Gyre Pagella}
\fi
% get rid of language-specific shorthands (see #6817):
\let\LanguageShortHands\languageshorthands
\def\languageshorthands#1{}
\input{entete.tex}
\ifLuaTeX
  \usepackage{selnolig}  % disable illegal ligatures
\fi
\IfFileExists{bookmark.sty}{\usepackage{bookmark}}{\usepackage{hyperref}}
\IfFileExists{xurl.sty}{\usepackage{xurl}}{} % add URL line breaks if available
\urlstyle{same}
\hypersetup{
  pdftitle={Comparer les programmes présidentiels de 2027},
  pdfauthor={Projet « Présidentielle » --- modèle macroéconométrique ouvert de l'économie française},
  pdflang={fr-FR},
  colorlinks=true,
  linkcolor={blue!50!black},
  filecolor={Maroon},
  citecolor={Blue},
  urlcolor={blue!50!black},
  pdfcreator={LaTeX via pandoc}}

\title{Comparer les programmes présidentiels de 2027}
\usepackage{etoolbox}
\makeatletter
\providecommand{\subtitle}[1]{% add subtitle to \maketitle
  \apptocmd{\@title}{\par {\large #1 \par}}{}{}
}
\makeatother
\subtitle{Un modèle macroéconométrique de la France, ses fondements, ses
résultats et ses limites --- étude scientifique exploratoire}
\author{Projet « Présidentielle » --- modèle macroéconométrique ouvert
de l'économie française}
\date{Version provisoire du 25 novembre 2026 (corpus des programmes
arrêté au 24 novembre 2026)}

\begin{document}
\frontmatter
\maketitle

{
\hypersetup{linkcolor=}
\setcounter{tocdepth}{2}
\tableofcontents
}
\setstretch{1.12}
\mainmatter
\hypertarget{avant-propos}{%
\chapter*{Avant-propos}\label{avant-propos}}
\addcontentsline{toc}{chapter}{Avant-propos}

Cet ouvrage présente un modèle macroéconométrique de l'économie
française construit à partir de données publiques, et l'utilise pour
comparer les programmes économiques de trois candidats à l'élection
présidentielle de 2027. Il a une double ambition : documenter de manière
complète un outil de simulation, et permettre au lecteur non spécialiste
mais motivé de comprendre les mécanismes économiques en jeu, les
désaccords entre écoles de pensée et ce qu'un modèle peut --- et ne peut
pas --- dire des choix politiques.

\textbf{Statut du document.} Il s'agit d'une \textbf{version
provisoire}, clairement datée. Les programmes des candidats ne sont pas
définitifs : La France insoumise a annoncé un chiffrage complet pour le
début de 2027, Édouard Philippe un « volume 2 » de son programme pour
janvier 2027, et le Rassemblement national n'a ventilé qu'une partie de
ses économies. Le corpus est arrêté au 24 novembre 2026. La procédure de
mise à jour (corpus, simulations, régénération du document) est décrite
en annexe \ref{sec:annexe-j}, et la liste des éléments à actualiser
figure au chapitre \ref{sec:conclusion}.

\textbf{Ce que ce document n'est pas.} Ce n'est ni une prévision, ni un
classement des candidats, ni une prise de position. Les simulations sont
des scénarios conditionnels : elles décrivent les conséquences des
programmes \emph{selon ce modèle et sous des hypothèses explicites}.
Aucun « vainqueur » n'est désigné ; les conclusions distinguent les
critères (activité, emploi, finances publiques, distribution,
contraintes institutionnelles) et ce qui résiste aux variantes.

\textbf{Comment le lire.} Le lecteur pressé peut lire l'introduction
(chapitre \ref{sec:intro}), la présentation des programmes (chapitre
\ref{sec:programmes}), les résultats (chapitre \ref{sec:resultats}), la
robustesse (chapitre \ref{sec:robustesse}) et la conclusion (chapitre
\ref{sec:conclusion}). Le lecteur qui souhaite apprendre l'économie qui
sous-tend le modèle lira le chapitre \ref{sec:theorie} (fondements
théoriques), le chapitre \ref{sec:transmission} (mécanismes de
transmission) et le chapitre \ref{sec:comportements} (limites des
hypothèses comportementales). Le lecteur technique trouvera
l'architecture, les équations, les données et la validation aux
chapitres \ref{sec:architecture} à \ref{sec:donnees} et dans les
annexes.

\textbf{Conventions.} Statuts des éléments : \textbf{{[}D{]}} donnée
observée ou règle publiée ; \textbf{{[}E{]}} paramètre estimé par le
projet (ou repris d'une estimation publiée, mention « littérature ») ;
\textbf{{[}C{]}} paramètre calibré ; \textbf{{[}H{]}} hypothèse de
modélisation ou de scénario ; \textbf{{[}X{]}} exploratoire. Les nombres
sont écrits à la française (virgule décimale) ; les montants en
milliards d'euros (Md€) ; les effets en écart à une trajectoire de
référence sans le programme.

\textbf{Reproductibilité.} Tous les résultats chiffrés des chapitres de
résultats sont produits automatiquement à partir des simulations
archivées ; le code, les données et les commandes de reproduction sont
dans le dépôt du projet (annexe \ref{sec:annexe-j}). L'annexe
\ref{sec:annexe-k} relie chaque conclusion aux simulations et aux
sources qui la fondent.

\hypertarget{sec:intro}{%
\chapter{Introduction générale et problématique}\label{sec:intro}}

\hypertarget{sec:contexte}{%
\section{Le contexte économique et budgétaire français à l'automne
2026}\label{sec:contexte}}

L'élection présidentielle d'avril-mai 2027 s'ouvre dans une
configuration macroéconomique que la France n'avait pas connue depuis la
création de l'euro : un déficit public durablement supérieur à 5 \% du
produit intérieur brut (PIB), une dette publique qui dépasse 115 \% du
PIB et continue de croître, une charge d'intérêts qui a doublé en
quelques années sous l'effet de la remontée des taux, une croissance
faible (de l'ordre de 1 \% par an ou moins), un taux de chômage remonté
au-dessus de 7 \% et une procédure européenne pour déficit excessif
ouverte depuis 2024. Les chiffres de référence utilisés dans cette étude
sont ceux des comptes nationaux publiés par l'Insee et Eurostat jusqu'en
2025, prolongés par la trajectoire « à politique inchangée » décrite au
chapitre \ref{sec:architecture} : déficit de 5,1 \% du PIB en 2025,
dette de 115,7 \%, chômage de 7,4 \%, taux de l'obligation assimilable
du Trésor (OAT) à dix ans de 4,89 \% en octobre 2026 (dernier point
mensuel disponible, Eurostat) contre 3,82 \% en moyenne sur douze mois.

Trois éléments de ce contexte conditionnent toute évaluation de
programme.

\begin{enumerate}
\def\labelenumi{\arabic{enumi}.}
\tightlist
\item
  \textbf{L'effet « boule de neige »} de la dette. Lorsque le taux
  d'intérêt apparent de la dette publique devient supérieur au taux de
  croissance nominal du PIB, le ratio d'endettement augmente
  mécaniquement, même avec un solde primaire (hors intérêts) à
  l'équilibre. Cette condition, presque oubliée pendant la décennie de
  taux très bas (2012-2021), est redevenue d'actualité. Elle signifie
  que l'horizon long des simulations est dominé par la dynamique de la
  dette et par l'hypothèse de taux, beaucoup plus que par le contenu des
  programmes.
\item
  \textbf{Le vieillissement}. Les générations nombreuses de
  l'après-guerre sont à la retraite ; la masse des pensions de droit
  direct représente de l'ordre de 13 à 14 \% du PIB, et le rapport entre
  cotisants et retraités continue de diminuer. Les trois candidats
  étudiés proposent des réformes des retraites de sens opposés (retour à
  62 ans, voire 60 ans, d'un côté ; 65 ans et 45 annuités de l'autre),
  ce qui en fait l'un des terrains de comparaison les plus nets.
\item
  \textbf{La contrainte européenne}. La France est soumise au cadre
  budgétaire réformé en 2024 (trajectoire de dépenses nettes, procédure
  de déficit excessif), à la politique monétaire unique de la Banque
  centrale européenne (BCE) et au droit de l'Union (TVA, aides d'État,
  politique commerciale, libre circulation des capitaux). Plusieurs
  mesures des programmes se heurtent à ce cadre ; leur faisabilité et
  leur calendrier en dépendent (chapitre \ref{sec:europe}).
\end{enumerate}

\hypertarget{sec:problematique}{%
\section{Problématique : comparer rigoureusement des programmes
différents}\label{sec:problematique}}

La question centrale de l'étude est la suivante :

\begin{quote}
\textbf{Comment comparer rigoureusement les conséquences économiques et
sociales de programmes présidentiels qui diffèrent par leur philosophie,
leurs instruments, leur calendrier et la précision de leur chiffrage ?}
\end{quote}

Elle se décompose en trois difficultés.

\textbf{Hétérogénéité des instruments.} Un programme combine des
dizaines de mesures : impôts, prestations, dépenses publiques, règles du
marché du travail, retraites, normes, politique commerciale. Leurs
effets passent par des canaux différents (demande, offre, prix,
incitations, finances publiques) et à des horizons différents. Les
comparer exige un cadre commun où chaque mesure est traduite dans les
mêmes variables.

\textbf{Hétérogénéité de la documentation.} Les programmes ne sont pas
chiffrés avec la même précision. À la date de l'étude (novembre 2026),
La France insoumise (LFI) a publié un « budget rectificatif » pour 2027
détaillant 106,4 milliards d'euros (Md€) de recettes et 81,7 Md€ de
dépenses nouvelles ; le Rassemblement national (RN) a présenté une
enveloppe de 140 Md€ d'économies nettes d'ici 2032, partiellement
ventilée ; Horizons (Édouard Philippe) annonce près de 150 Md€
d'économies et un déficit sous 2,5 \% du PIB en fin de quinquennat, avec
peu de mesures chiffrées une à une. Une comparaison honnête ne doit pas
donner artificiellement le même niveau de précision à des programmes
dont le chiffrage diffère. C'est pourquoi l'étude propose, pour chaque
candidat, \textbf{deux lectures} : une lecture « par mesures documentées
» (lecture A) et une lecture « par objectifs budgétaires » (lecture B),
qui reconstruit l'effort nécessaire pour atteindre la trajectoire de
solde annoncée (chapitre \ref{sec:resultats}).

\textbf{Incertitude des mécanismes.} Les effets macroéconomiques d'une
mesure dépendent de paramètres mal connus (multiplicateurs, élasticités,
réactions des marchés et de la BCE, anticipations) et de choix de
modélisation sur lesquels les économistes ne s'accordent pas (chapitre
\ref{sec:theorie}). Une comparaison crédible doit montrer ce qui résiste
à ces incertitudes et ce qui en dépend.

\hypertarget{sec:statique}{%
\section{Pourquoi les chiffrages statiques ne suffisent
pas}\label{sec:statique}}

Le débat public s'appuie surtout sur des \textbf{chiffrages statiques} :
le coût d'une mesure est calculé en multipliant une assiette par un
taux, à comportements et à activité inchangés. Ces chiffrages sont
indispensables (ils sont le point de départ de tout le reste) mais ils
ne répondent pas à la question posée, pour quatre raisons.

\begin{enumerate}
\def\labelenumi{\arabic{enumi}.}
\tightlist
\item
  \textbf{Les comportements réagissent.} Une hausse de l'imposition des
  revenus du capital réduit l'assiette (les dividendes distribués
  diminuent, certains contribuables optimisent ou partent) ; une baisse
  de TVA n'est que partiellement répercutée sur les prix ; un
  allongement de la durée d'assurance ne se traduit que partiellement
  par des départs plus tardifs. Le rendement « net » d'une mesure peut
  être très inférieur à son rendement statique (section
  \ref{sec:incidence}).
\item
  \textbf{L'activité réagit.} Une dépense publique supplémentaire
  augmente la demande, donc la production, les revenus et les recettes
  fiscales ; une économie budgétaire produit l'effet inverse. Le coût
  budgétaire final d'une mesure diffère de son coût direct d'un montant
  qui dépend du multiplicateur et de l'élasticité des recettes à
  l'activité (encadré ci-dessous).
\item
  \textbf{Les prix et les taux réagissent.} Les hausses de salaires se
  diffusent aux prix ; les déficits peuvent faire monter la prime de
  risque souveraine ; la BCE réagit à l'inflation et à l'activité de la
  zone euro.
\item
  \textbf{Le temps compte.} Les effets de demande sont rapides, les
  effets d'offre (capital, productivité, participation au marché du
  travail) sont lents ; une mesure peut être favorable à cinq ans et
  défavorable à vingt, ou l'inverse.
\end{enumerate}

\begin{quote}
\textbf{Encadré 1.1 --- Coût budgétaire direct et effet macroéconomique
final : un exemple chiffré.} Considérons une baisse de 10 Md€ de la
consommation publique, dans une économie où le PIB vaut 3 000 Md€, où le
multiplicateur de court terme est \(m\) et où les recettes publiques
représentent \(\tau = 52\,\%\) du PIB. L'économie directe est de 10 Md€
(0,33 point de PIB). Le PIB diminue de \(m \times 10\) Md€ ; les
recettes diminuent de \(\tau \times m \times 10\) Md€. Le gain
budgétaire net est donc \(10 \times (1 - \tau m)\) Md€ : 4,8 Md€ si
\(m = 1\), 2,8 Md€ si \(m = 1{,}4\). Quant au ratio de dette
\(d = D/Y\), il varie de
\(\Delta d \approx (-\Delta D)/Y - d \cdot \Delta Y / Y\) : la baisse du
dénominateur peut l'emporter sur la baisse du numérateur. Avec
\(d = 1{,}2\), la consolidation \textbf{augmente} le ratio de dette la
première année dès que \(m > 1/(d + \tau) \approx 0{,}58\). Le modèle
utilisé ici a un multiplicateur de consommation publique d'environ 1,2
la première année (section \ref{sec:multiplicateurs}) : une
consolidation par la dépense dégrade le ratio de dette pendant plusieurs
années avant de l'améliorer. Ce n'est pas un paradoxe du modèle, c'est
une propriété arithmétique bien documentée (Blanchard et Leigh, 2013 ;
FMI, 2010).
\end{quote}

\hypertarget{sec:approche}{%
\section{L'intérêt d'une approche dynamique, macroéconomique et
microéconomique}\label{sec:approche}}

Pour dépasser le chiffrage statique, l'étude s'appuie sur un modèle
macroéconométrique de la France construit pour ce projet, à partir de
données ouvertes uniquement, et couplé à une microsimulation des
ménages. Ses principes :

\begin{itemize}
\tightlist
\item
  \textbf{Macroéconomie bouclée.} Les comptes de tous les secteurs
  (ménages, sociétés non financières, administrations publiques, reste
  du monde) sont tenus de manière cohérente ; un revenu versé par l'un
  est reçu par l'autre ; les identités comptables sont vérifiées à
  chaque simulation (chapitre \ref{sec:donnees}).
\item
  \textbf{Multisectoriel.} La production est décrite par 25 branches
  marchandes reliées par le tableau des entrées-sorties (TES) de
  l'Insee, ce qui permet de suivre la diffusion des prix et des contenus
  en importations.
\item
  \textbf{Microéconomie des ménages.} Les mesures fiscales et sociales
  sont évaluées ménage par ménage sur un échantillon représentatif
  construit à partir du recensement et des données fiscales, avec le
  moteur de calcul public OpenFisca-France ; leurs effets distributifs
  sont mesurés par décile de niveau de vie (chapitre
  \ref{sec:distribution}).
\item
  \textbf{Régimes de politique économique explicites.} La règle
  budgétaire suivie après le programme, la réaction de la BCE, la
  formation des taux souverains et les réactions institutionnelles
  européennes sont des choix de scénario, documentés et variables, et
  non des hypothèses cachées.
\end{itemize}

L'architecture distingue trois objets (chapitre \ref{sec:architecture})
: l'\textbf{environnement} (l'économie à politique inchangée, calibrée
une fois), le \textbf{régime} (les règles de politique budgétaire et
monétaire), et le \textbf{programme} (les mesures du candidat). L'effet
d'un programme est toujours mesuré comme la différence entre deux
simulations dans le même environnement et sous le même régime, l'une
avec et l'autre sans le programme.

\hypertarget{sec:candidats}{%
\section{Les candidats retenus}\label{sec:candidats}}

L'étude porte sur trois candidats représentatifs de trois philosophies
économiques distinctes :

\begin{itemize}
\tightlist
\item
  \textbf{Jean-Luc Mélenchon} (La France insoumise), dont le programme
  \emph{L'Avenir en commun} (5\textsuperscript{e} édition, novembre
  2026, 1 069 mesures) privilégie la relance par les revenus des ménages
  modestes et la dépense publique, financée par une hausse de
  l'imposition du capital et des hauts patrimoines ;
\item
  \textbf{Marine Le Pen} (Rassemblement national), dont le programme
  combine des baisses ciblées d'impôts (TVA sur l'énergie, jeunes,
  familles), un retour partiel sur la réforme des retraites et des
  économies concentrées sur l'immigration, la contribution européenne et
  l'organisation de l'État ;
\item
  \textbf{Édouard Philippe} (Horizons), retenu comme candidat du centre,
  dont le programme privilégie la réduction des dépenses publiques (près
  de 150 Md€), l'allongement de la vie active (65 ans, 45 annuités), la
  réduction de 250 000 postes de fonctionnaires et un « pacte » de
  baisse des impôts de production financée par la baisse des aides aux
  entreprises.
\end{itemize}

\textbf{Justification du choix du candidat du centre.} À la date du
corpus (24 novembre 2026), le programme d'Édouard Philippe est le seul
programme du centre et du centre droit publié sous une forme structurée
(quinze thèmes sur son site, un « volume 1 » de 142 mesures présenté le
18 novembre 2026, un volet de finances publiques avec un objectif
chiffré). Pour Gabriel Attal, les sources consultées ne donnaient qu'une
enveloppe d'économies (120 à 150 Md€ selon les sources). Ce choix sera
réévalué lorsque d'autres programmes seront suffisamment détaillés ; le
corpus et le code permettent d'ajouter un candidat sans modification du
modèle (annexe \ref{sec:annexe-j}).

Le choix de trois candidats ne préjuge ni de leur poids électoral ni
d'un jugement sur les autres programmes ; il vise à couvrir la diversité
des philosophies économiques présentes dans le débat.

\hypertarget{sec:questions}{%
\section{Questions de recherche, hypothèses et
périmètre}\label{sec:questions}}

L'étude cherche à répondre aux questions suivantes.

\begin{itemize}
\tightlist
\item
  \textbf{Q1.} Quels sont les effets macroéconomiques (activité, emploi,
  prix, revenus, finances publiques) de chaque programme, mesure par
  mesure et dans son ensemble, à 5 et 10 ans, et sous quelles hypothèses
  à 20 ans ?
\item
  \textbf{Q2.} Quelle part de ces effets est robuste aux incertitudes de
  modélisation (paramètres, régimes, environnement, réactions
  européennes) ?
\item
  \textbf{Q3.} Quels mécanismes de transmission expliquent les écarts
  entre programmes ?
\item
  \textbf{Q4.} Qui gagne et qui perd, au moins au premier ordre, parmi
  les ménages (déciles, actifs et retraités, générations) ?
\item
  \textbf{Q5.} Quelles mesures se heurtent aux contraintes européennes,
  et avec quelles conséquences possibles ?
\end{itemize}

Hypothèses de travail : les programmes sont pris tels qu'ils sont
publiés, sans les compléter par des intentions supposées ; lorsqu'une
traduction dans le modèle exige une hypothèse (par exemple la
composition d'une enveloppe d'économies), cette hypothèse est
explicitée, marquée {[}H{]} et accompagnée de variantes. Les hypothèses
de traduction retenues par l'auteur du modèle ne sont jamais présentées
comme des engagements des candidats.

Périmètre : économie française, horizon 2027-2046, pas annuel. Hors
périmètre : modélisation complète de l'économie européenne et mondiale,
effets environnementaux (émissions), effets sur la qualité des services
publics, effets politiques et sociaux non économiques.

\begin{quote}
\textbf{Avertissement méthodologique.} Les simulations présentées dans
cet ouvrage sont des \textbf{scénarios conditionnels} : elles répondent
à la question « que se passerait-il, selon ce modèle et sous ces
hypothèses, si ce programme était appliqué tel qu'il est décrit ? ». Ce
ne sont pas des prévisions. Elles ne disent pas ce qui va se passer, ni
ce qu'un gouvernement ferait effectivement, ni comment évolueraient les
chocs extérieurs. Leur valeur tient à la cohérence comptable, à la
transparence des hypothèses et à l'analyse de sensibilité, non à une
précision qu'aucun modèle ne possède.
\end{quote}

\hypertarget{sec:plan}{%
\section{Plan de l'étude}\label{sec:plan}}

Le chapitre \ref{sec:theorie} présente les fondements théoriques et les
désaccords entre écoles qui justifient --- ou contestent --- les
équations du modèle. Les chapitres \ref{sec:architecture} à
\ref{sec:donnees} décrivent le modèle : architecture, équations,
données, estimations et validation. Le chapitre \ref{sec:programmes}
présente en détail les trois programmes et leur traduction. Le chapitre
\ref{sec:transmission} explique la transmission macroéconomique des
grandes familles de mesures, et le chapitre \ref{sec:europe} les
contraintes européennes. Le chapitre \ref{sec:resultats} présente les
résultats des simulations comparatives, le chapitre
\ref{sec:distribution} leurs effets distributifs, le chapitre
\ref{sec:comportements} les limites des hypothèses comportementales, et
le chapitre \ref{sec:robustesse} la robustesse et les incertitudes. Le
chapitre \ref{sec:discussion} compare les philosophies économiques des
programmes et le chapitre \ref{sec:conclusion} conclut. Les annexes (A à
K) rassemblent les registres, les équations complètes, les résultats
détaillés, les tests, le glossaire et le guide de reproduction.

\textbf{Conventions de lecture.} Tout au long de l'ouvrage, les statuts
suivants sont utilisés : \textbf{{[}D{]}} donnée observée ou règle
publiée ; \textbf{{[}E{]}} paramètre estimé par le projet ;
\textbf{{[}C{]}} paramètre calibré sur la littérature ou sur une cible ;
\textbf{{[}H{]}} hypothèse de modélisation ou de scénario ;
\textbf{{[}X{]}} exploratoire. Les effets sont exprimés en écart à la
trajectoire de référence (même environnement, même régime, sans le
programme) : en pourcentage pour les volumes et les prix, en points pour
les taux (chômage, solde public, dette rapportés au PIB), en milliers
pour l'emploi. Les horizons « 5 ans », « 10 ans » et « 20 ans »
désignent les années 2031, 2036 et 2046 (le programme commence en 2027).

\hypertarget{sec:theorie}{%
\chapter{Fondements théoriques de l'économie
modélisée}\label{sec:theorie}}

Un modèle macroéconométrique n'est pas une théorie unique : c'est un
assemblage d'équations dont chacune résume un mécanisme, et dont la
forme reflète des choix théoriques. Ce chapitre présente les grandes
écoles de pensée qui justifient --- ou contestent --- ces choix. Pour
chacune, on expose les intuitions, les mécanismes, les équations
fondamentales, les principaux résultats empiriques, les controverses, et
ce que le modèle de l'étude en reprend, en simplifie ou en écarte.
L'objectif est double : permettre au lecteur de comprendre pourquoi le
modèle produit les résultats qu'il produit, et lui donner les moyens de
juger ce qui changerait sous d'autres hypothèses.

Un fil conducteur traverse tout le chapitre : \textbf{la question de
savoir si, et pendant combien de temps, la demande détermine la
production}. Les écoles qui répondent « longtemps » (keynésiennes,
post-keynésiennes) attribuent de grands effets aux politiques
budgétaires ; celles qui répondent « peu ou pas » (classiques, nouvelle
macroéconomie classique) attribuent l'essentiel aux conditions de
l'offre. Les modèles institutionnels contemporains, dont celui de
l'étude, adoptent une position intermédiaire : la demande compte à court
terme, l'offre à long terme, et la vitesse du passage de l'un à l'autre
est l'un des paramètres les plus incertains et les plus influents.

\hypertarget{sec:classique}{%
\section{L'économie classique et néoclassique}\label{sec:classique}}

\textbf{Intuitions.} Pour les économistes classiques (Smith, Ricardo,
Say, puis Walras et Marshall pour la tradition néoclassique), les prix
flexibles ajustent l'offre et la demande sur chaque marché. Sur le
marché du travail, le salaire réel s'ajuste jusqu'à ce que toute
personne souhaitant travailler au salaire courant trouve un emploi ; le
chômage durable ne peut être que volontaire ou dû à des rigidités
(salaire minimum, syndicats, réglementations). Sur le marché des biens,
la production est déterminée par les facteurs disponibles (travail,
capital) et la technologie : c'est la « loi de Say », selon laquelle
l'offre crée sa propre demande, car tout revenu versé est dépensé,
consommé ou épargné, et l'épargne finance l'investissement par
l'intermédiaire du taux d'intérêt.

\textbf{Mécanismes et équations.} La production est donnée par une
fonction de production

\[Y = F(K, AL),\]

où \(Y\) est la production, \(K\) le capital, \(L\) le travail et \(A\)
l'efficience du travail (la technologie). Les entreprises embauchent
jusqu'à ce que la productivité marginale du travail égale le salaire
réel, \(F_L(K, AL) = W/P\), et investissent jusqu'à ce que la
productivité marginale du capital égale le coût réel du capital. Sur le
marché des fonds prêtables, le taux d'intérêt réel \(r\) égalise
l'épargne \(S(r)\) et l'investissement \(I(r)\). Une hausse de la
dépense publique financée par emprunt réduit l'épargne nationale, fait
monter \(r\) et évince l'investissement privé d'un montant égal : le
multiplicateur budgétaire est nul (« éviction totale »). La monnaie est
neutre : la théorie quantitative \(MV = PY\) implique qu'une hausse de
la masse monétaire \(M\) se traduit par une hausse proportionnelle des
prix \(P\).

\textbf{Résultats empiriques.} À long terme, plusieurs prédictions
classiques sont bien vérifiées : la neutralité de la monnaie sur
l'activité réelle, le rôle dominant de la productivité dans la
croissance du niveau de vie, l'ajustement des salaires réels à la
productivité. À court terme, elles le sont beaucoup moins : les
récessions s'accompagnent de hausses massives du chômage involontaire,
que les prix et salaires ne résorbent que lentement.

\textbf{Controverses.} La critique keynésienne porte sur la flexibilité
des prix et sur le lien entre épargne et investissement : si les
salaires nominaux baissent, la demande des salariés baisse aussi, et la
déflation peut aggraver le chômage en alourdissant les dettes (Fisher,
1933). La critique de la loi de Say porte sur la possibilité de
thésauriser : un revenu épargné sous forme de monnaie n'est pas dépensé.

\textbf{Ce que le modèle reprend.} Le cadre classique gouverne le
\textbf{long terme} du modèle : la production tend vers un PIB potentiel
déterminé par la fonction de production (CES, section
\ref{sec:eq-production}), la population active et le chômage structurel
; la monnaie n'y a pas d'effet réel durable. Ce que le modèle écarte :
l'ajustement instantané des prix et l'éviction totale à court terme.

\hypertarget{sec:keynes}{%
\section{La théorie keynésienne et le multiplicateur
budgétaire}\label{sec:keynes}}

\textbf{Intuitions.} Keynes (1936) part d'un constat : pendant la Grande
Dépression, le chômage massif persiste alors que les salaires baissent.
Son explication tient en une idée : \textbf{la production est déterminée
par la demande effective}. Les entreprises produisent ce qu'elles
anticipent de vendre ; si les ménages et les entreprises réduisent leurs
dépenses (par précaution, pessimisme, contrainte de crédit), la
production et l'emploi baissent, ce qui réduit les revenus, donc la
dépense, dans un enchaînement cumulatif. Rien ne garantit que l'économie
revienne spontanément au plein emploi, car les salaires et les prix sont
rigides à la baisse et parce qu'une baisse des salaires réduit aussi la
demande.

\textbf{Le multiplicateur.} Dans le modèle le plus simple, la
consommation dépend du revenu disponible, \(C = c_0 + c\,(Y - T)\), où
\(c\) est la propension marginale à consommer. L'équilibre
\(Y = C + I + G\) donne

\[Y = \frac{c_0 - cT + I + G}{1 - c}, \qquad \frac{\partial Y}{\partial G} = \frac{1}{1-c}.\]

Avec \(c = 0{,}6\), un euro de dépense publique supplémentaire produit
2,5 euros de PIB. En économie ouverte, une part \(\mu\) de la dépense
est importée et une part \(\tau\) du revenu est prélevée : le
multiplicateur devient \(1/[1 - c(1-\tau) + \mu]\), soit environ 1,1
avec \(c = 0{,}6\), \(\tau = 0{,}45\) et \(\mu = 0{,}3\). Le
multiplicateur d'une baisse d'impôt est plus faible que celui d'une
dépense (\(c\) fois plus faible au premier tour), car une partie du
revenu supplémentaire est épargnée.

\textbf{Le modèle IS-LM.} Hicks (1937) formalise Keynes en deux courbes
: la courbe IS (équilibre du marché des biens, production décroissante
avec le taux d'intérêt, qui réduit l'investissement) et la courbe LM
(équilibre du marché monétaire, taux croissant avec la production pour
une offre de monnaie donnée). Une relance budgétaire déplace IS ;
l'éviction passe par la hausse du taux d'intérêt, partielle si LM n'est
pas verticale. Lorsque le taux est bloqué à sa borne inférieure («
trappe à liquidité ») ou que la banque centrale le maintient fixe,
l'éviction disparaît et le multiplicateur est maximal.

\textbf{Résultats empiriques.} L'estimation des multiplicateurs est l'un
des terrains les plus travaillés de la macroéconomie empirique. Trois
grandes familles d'approches existent :

\begin{itemize}
\tightlist
\item
  les \textbf{vecteurs autorégressifs structurels} (VAR) identifiant les
  chocs budgétaires par le délai de décision (Blanchard et Perotti,
  2002) trouvent des multiplicateurs de dépense proches de 1 aux
  États-Unis ;
\item
  les approches \textbf{narratives} (chocs de dépenses militaires,
  Ramey, 2011 ; chocs fiscaux, Romer et Romer, 2010) trouvent des
  multiplicateurs de dépense entre 0,6 et 1,2 et des effets d'impôt plus
  élevés ;
\item
  les études \textbf{conditionnelles à l'état de l'économie} (Auerbach
  et Gorodnichenko, 2012) trouvent des multiplicateurs plus élevés en
  récession (supérieurs à 1,5) qu'en expansion (proches de 0,5), un
  résultat discuté (Ramey et Zubairy, 2018).
\end{itemize}

La revue de Ramey (2019) situe la plupart des estimations du
multiplicateur de dépense entre 0,6 et 1. L'erreur de prévision des
organisations internationales pendant la crise de la zone euro
(Blanchard et Leigh, 2013) a montré que les multiplicateurs implicites
de la consolidation de 2010-2012 avaient été sous-estimés (supérieurs à
1 au lieu de 0,5), ce qui a joué un rôle important dans la réévaluation
de l'austérité.

\textbf{Controverses.} La critique monétariste et classique porte sur
l'éviction et les anticipations (sections \ref{sec:monetarisme} et
\ref{sec:anticipations}) ; la critique « ricardienne » (Barro, 1974)
soutient que des ménages prévoyants épargnent la baisse d'impôt financée
par la dette en prévision des impôts futurs. À l'opposé, la thèse de «
l'austérité expansionniste » (Giavazzi et Pagano, 1990 ; Alesina et
Ardagna, 2010) soutenait qu'une consolidation crédible pouvait stimuler
l'activité par la confiance ; les travaux ultérieurs du FMI (Guajardo,
Leigh et Pescatori, 2014) ont montré que ce résultat tenait largement à
l'identification des épisodes et que les consolidations sont en moyenne
récessives à court terme.

\textbf{Ce que le modèle reprend.} Le cœur de court terme du modèle est
keynésien : la production résulte de la demande (consommation,
investissement, dépenses publiques, exportations nettes des
importations), elle-même fonction des revenus. Les multiplicateurs du
modèle (section \ref{sec:multiplicateurs}) se situent dans le haut de la
fourchette empirique pour la consommation publique (environ 1,2 la
première année, 0,75 à cinq ans, avec un taux d'intérêt exogène), parce
que le modèle n'a pas d'effet d'éviction par le taux dans son régime
principal et que la fermeture vers le potentiel est lente. Ce point est
central pour l'interprétation : sous le régime avec réaction de la BCE
et prime de risque, ou avec une fermeture plus rapide, les
multiplicateurs sont plus faibles (chapitre \ref{sec:robustesse}).

\hypertarget{sec:nouvelle-keynesienne}{%
\section{Synthèse néoclassique et nouvelle économie
keynésienne}\label{sec:nouvelle-keynesienne}}

\textbf{Intuitions.} La « synthèse néoclassique » (Samuelson,
Modigliani, Tobin) réconcilie Keynes et les classiques : keynésienne à
court terme parce que les prix sont rigides, classique à long terme
quand ils s'ajustent. La \textbf{nouvelle économie keynésienne} (années
1980-2000) donne des fondements microéconomiques à cette rigidité :
coûts de changement des prix (« coûts de menu »), contrats échelonnés
(Taylor, 1980 ; Calvo, 1983), concurrence monopolistique, salaires
d'efficience, rigidités réelles.

\textbf{Équations fondamentales.} Le modèle canonique (Clarida, Galí et
Gertler, 1999 ; Woodford, 2003 ; Galí, 2015) tient en trois équations
log-linéarisées :

\[\tilde y_t = \mathbb{E}_t \tilde y_{t+1} - \frac{1}{\sigma}\left(i_t - \mathbb{E}_t \pi_{t+1} - r^n_t\right) \quad\text{(IS dynamique)}\]
\[\pi_t = \beta\,\mathbb{E}_t \pi_{t+1} + \kappa\,\tilde y_t \quad\text{(courbe de Phillips nouvelle keynésienne)}\]
\[i_t = \rho + \phi_\pi \pi_t + \phi_y \tilde y_t \quad\text{(règle de Taylor)}\]

où \(\tilde y\) est l'écart de production, \(\pi\) l'inflation, \(i\) le
taux nominal, \(r^n\) le taux d'intérêt naturel. La courbe de Phillips
relie l'inflation aux tensions sur les capacités de production ; la
règle de Taylor (1993) décrit une banque centrale qui relève son taux
plus que l'inflation (\(\phi_\pi > 1\)), stabilisant l'économie.

\textbf{Résultats empiriques.} Les modèles DSGE (dynamiques,
stochastiques, d'équilibre général) de cette famille (Smets et Wouters,
2007) sont devenus l'outil standard des banques centrales. Ils
reproduisent correctement les corrélations conjoncturelles, mais ont mal
anticipé la crise financière de 2008 (absence de secteur financier) et
la faiblesse de l'inflation de la décennie 2010. La pente de la courbe
de Phillips estimée sur données récentes est faible (courbe « plate »),
ce qui signifie que de fortes variations de chômage modifient peu
l'inflation --- jusqu'à l'épisode inflationniste de 2021-2023, dont
l'interprétation (chocs d'offre et non-linéarités) reste débattue.

\textbf{Controverses.} Les critiques portent sur les anticipations
rationnelles (section \ref{sec:anticipations}), sur l'agent
représentatif (section \ref{sec:heterogenes}), et sur les effets de «
guidage prospectif » irréalistes (l'« énigme du guidage prospectif » :
une promesse de taux bas dans dix ans a, dans ces modèles, un effet
aujourd'hui plus fort qu'une baisse immédiate).

\textbf{Ce que le modèle reprend.} Le modèle de l'étude n'est pas un
DSGE : il appartient à la famille des modèles macroéconométriques
institutionnels (comme Mésange, de l'Insee et de la DG Trésor, ou FR-BDF
de la Banque de France pour son bloc d'équations de comportement), dont
les équations sont estimées en correction d'erreur. Il reprend de la
nouvelle synthèse : la rigidité des prix et des salaires (ajustements
partiels vers des cibles de long terme), une courbe de Phillips sur les
salaires (faible, calibrée) et une pente des prix à l'écart de
production, et une règle de Taylor pour la BCE dans le régime monétaire
endogène (section \ref{sec:eq-monetaire}). Il écarte les anticipations
rationnelles généralisées (sauf en option pour le taux allemand).

\hypertarget{sec:monetarisme}{%
\section{Le monétarisme et le rôle de la politique
monétaire}\label{sec:monetarisme}}

\textbf{Intuitions.} Friedman (1968) et Phelps (1967) affirment qu'il
n'existe pas d'arbitrage durable entre inflation et chômage : si la
banque centrale cherche à maintenir le chômage sous son niveau « naturel
», l'inflation accélère, car les salariés intègrent l'inflation passée
dans leurs revendications. La politique monétaire a des effets réels à
court terme, mais avec des délais « longs et variables » ; elle devrait
suivre une règle simple plutôt que chercher à piloter finement la
conjoncture.

\textbf{Équations.} La courbe de Phillips augmentée des anticipations
s'écrit

\[\pi_t = \pi^e_t - \alpha\,(u_t - u^*),\]

où \(u^*\) est le chômage d'équilibre (NAIRU, \emph{non-accelerating
inflation rate of unemployment}). Avec des anticipations adaptatives,
\(\pi^e_t = \pi_{t-1}\), maintenir \(u_t < u^*\) entraîne une inflation
qui accélère.

\textbf{Résultats empiriques.} La stagflation des années 1970 a validé
l'essentiel du diagnostic de Friedman et Phelps et conduit à
l'indépendance des banques centrales et au ciblage d'inflation. Le lien
quantitatif entre monnaie et inflation s'est en revanche distendu à
court terme (instabilité de la demande de monnaie), d'où l'abandon du
ciblage d'agrégats monétaires au profit des règles de taux.

\textbf{Controverses.} Le niveau de \(u^*\) est incertain et varie dans
le temps ; les estimations françaises vont de 7 à 9 \% selon les
méthodes et les périodes. Les hystérèses (Blanchard et Summers, 1986)
suggèrent qu'un chômage élevé durablement peut augmenter \(u^*\)
lui-même.

\textbf{Ce que le modèle reprend.} Le chômage structurel \(u^*\) est une
hypothèse de la référence (7,5 \% puis 7,0 \%, cohérente avec les
hypothèses du FMI et du Conseil d'orientation des retraites) {[}C{]} ;
il peut être modifié par certaines mesures (assurance chômage,
formation) ; la courbe de Phillips sur les salaires est faible. Dans la
zone euro, la politique monétaire est commune : la BCE réagit à
l'inflation et à l'activité \textbf{de la zone}, dans lesquelles la
France pèse environ un cinquième ; une relance française n'appelle donc
qu'une faible réaction de taux (section \ref{sec:eq-monetaire}).

\hypertarget{sec:offre}{%
\section{Théorie de l'offre, fiscalité et incitations}\label{sec:offre}}

\textbf{Intuitions.} L'économie de l'offre met l'accent sur les effets
incitatifs des prélèvements : un impôt modifie les prix relatifs (entre
travail et loisir, entre consommation présente et future, entre
investissement en France et à l'étranger) et donc les comportements. Le
coût d'un impôt pour l'économie dépasse la recette qu'il rapporte :
c'est la « charge excédentaire » (ou perte sèche), proportionnelle, en
première approximation, au carré du taux et à l'élasticité de
l'assiette.

\textbf{Équations.} L'élasticité du revenu imposable au taux net
d'impôt, \(e = \frac{\partial \ln z}{\partial \ln(1-\tau)}\), résume les
réactions (heures, effort, optimisation, évasion). Saez, Slemrod et
Giertz (2012) montrent que le taux marginal qui maximise la recette sur
les hauts revenus vaut

\[\tau^* = \frac{1}{1 + a\,e},\]

où \(a\) est le paramètre de Pareto de la distribution des hauts revenus
(environ 1,5 à 2). Avec \(e = 0{,}25\) et \(a = 2\),
\(\tau^* \approx 67\,\%\) ; avec \(e = 0{,}5\), \(\tau^* = 50\,\%\). La
« courbe de Laffer » est l'illustration générale de cette relation entre
taux et recette.

\textbf{Résultats empiriques.} Les élasticités sont faibles pour les
heures travaillées des salariés à temps plein, plus élevées pour la
participation des femmes mariées et des seniors (marge extensive) et
pour les très hauts revenus (par l'optimisation et la requalification
des revenus plutôt que par l'effort). Pour la France, les estimations
sur réformes (Cabannes, Houdré et Landais, 2014) trouvent des
élasticités du revenu imposable de l'ordre de 0,1 à 0,2 pour l'ensemble
et plus élevées en haut de la distribution. Pour l'impôt sur les
sociétés, les effets sur l'investissement passent par le coût du capital
(section \ref{sec:investissement}) ; l'incidence sur les salaires a été
estimée à environ la moitié pour l'impôt local allemand sur les
entreprises (Fuest, Peichl et Siegloch, 2018). Pour la TVA, l'incidence
est asymétrique : les hausses sont répercutées plus que les baisses
(Benzarti et al., 2020).

\textbf{Controverses.} Les élasticités de court terme (optimisation,
décalage temporel des revenus) surestiment souvent les effets de long
terme ; la mobilité internationale des hauts patrimoines est discutée
(les départs liés à l'ancien impôt de solidarité sur la fortune sont
estimés faibles en nombre mais concentrés en patrimoine). Les promesses
d'autofinancement des baisses d'impôt sont rarement vérifiées aux taux
observés dans les pays de l'OCDE.

\textbf{Ce que le modèle reprend.} Le module M13 applique aux mesures
sur les hauts revenus et les patrimoines des élasticités issues de la
littérature, avec des fourchettes basse, centrale et haute (revenu
imposable 0,1-0,2-0,4 ; capital 0,2-0,5-1,0 ; pertes d'assiette de
l'impôt sur la fortune de 25 \% puis 35 \% à cinq ans en valeur
centrale) {[}C{]}. Le module M15 traduit les réformes des prestations en
effets sur la participation au marché du travail {[}B{]}. L'incidence
des prélèvements sur les entreprises (impôt sur les sociétés, impôts de
production, cotisations) est répartie entre prix, salaires et marges
selon des parts calibrées (chapitre \ref{sec:transmission}).

\hypertarget{sec:croissance}{%
\section{Théories de la croissance exogène et
endogène}\label{sec:croissance}}

\textbf{Intuitions.} Le modèle de Solow (1956) explique la croissance de
long terme par l'accumulation du capital et le progrès technique. Le
capital ayant des rendements décroissants, l'accumulation seule ne
suffit pas : à long terme, le niveau de vie croît au rythme du progrès
technique, supposé exogène. Les théories de la croissance endogène
(Romer, 1990 ; Lucas, 1988 ; Aghion et Howitt, 1992) rendent ce progrès
dépendant de choix économiques : recherche et développement, capital
humain, concurrence, diffusion des connaissances. Les politiques
publiques (éducation, recherche, infrastructures) peuvent alors modifier
le \textbf{taux} de croissance et pas seulement le niveau du PIB.

\textbf{Équations.} Dans le modèle de Solow,
\(\dot k = s\,f(k) - (n + g + \delta)\,k\), où \(k\) est le capital par
unité de travail efficace, \(s\) le taux d'épargne, \(n\) la croissance
démographique, \(g\) le progrès technique et \(\delta\) la dépréciation.
À l'état stationnaire, le PIB par tête croît au taux \(g\). La
comptabilité de la croissance décompose
\(\Delta \ln Y = \alpha\,\Delta \ln K + (1-\alpha)\,\Delta \ln L + \Delta \ln \text{PGF}\),
la productivité globale des facteurs (PGF) étant le résidu.

\textbf{Résultats empiriques.} La convergence conditionnelle entre pays
est confirmée (Mankiw, Romer et Weil, 1992). Les rendements de la R\&D
sont élevés et largement sociaux (Hall, Mairesse et Mohnen, 2010) ;
l'élasticité de la production au capital public est estimée autour de
0,08 en moyenne (Bom et Ligthart, 2014). Le ralentissement de la
productivité dans les pays avancés depuis le milieu des années 2000
reste mal expliqué.

\textbf{Controverses.} L'ampleur des effets de long terme des politiques
publiques sur la croissance est très incertaine ; les effets de la
fiscalité sur la croissance de long terme sont discutés (effets de
niveau plutôt que de taux dans la plupart des estimations).

\textbf{Ce que le modèle reprend.} La croissance de long terme du modèle
est \textbf{exogène} : la productivité tendancielle de chaque branche
est un paramètre (dérivé des huit dernières années observées). Deux
canaux endogènes exploratoires existent : l'investissement public
accroît le capital public, qui augmente la productivité avec une
élasticité de 0,08 {[}C{]} ; l'aide à la R\&D augmente le stock de R\&D
(élasticité 0,08, statut exploratoire) {[}X{]}. Les effets de
l'éducation, de la qualité des services publics et des réformes
structurelles sur la productivité ne sont pas représentés : c'est une
limite qui pénalise les programmes dont les effets attendus passent par
ces canaux (chapitre \ref{sec:comportements}).

\hypertarget{sec:travail}{%
\section{Marché du travail, chômage structurel et loi
d'Okun}\label{sec:travail}}

\textbf{Intuitions.} Le chômage résulte de plusieurs causes :
insuffisance de la demande (chômage keynésien), coût du travail trop
élevé par rapport à la productivité (chômage classique), frictions
d'appariement entre offres et demandes d'emploi (chômage frictionnel),
inadéquation des qualifications. Les théories modernes (Mortensen et
Pissarides, 1994) décrivent le marché du travail comme un processus de
recherche : le chômage d'équilibre dépend des flux de créations et de
destructions d'emplois, de l'efficacité de l'appariement et des
conditions de négociation salariale (dont l'assurance chômage, qui
augmente le salaire de réserve).

\textbf{Équations.} La loi d'Okun (1962) relie les variations du chômage
à l'écart de croissance :

\[\Delta u_t \approx -\beta_O\,(g_t - g^*) \quad\text{ou}\quad u_t - u^* \approx -\frac{1}{\beta_O'}\,\tilde y_t,\]

avec pour la France un coefficient de l'ordre de 0,3 à 0,5 point de
chômage par point d'écart de production. La courbe de Beveridge relie le
taux d'emplois vacants au chômage. Dans les modèles de négociation, le
salaire d'équilibre augmente avec la productivité, l'indemnisation du
chômage et le pouvoir de négociation, et diminue avec le taux de
chômage.

\textbf{Résultats empiriques.} La relation d'Okun est robuste en
direction mais instable en coefficient. Les effets de l'assurance
chômage sur la durée du chômage sont bien établis (une indemnisation
plus généreuse ou plus longue allonge les épisodes), mais leurs effets
agrégés sur le chômage d'équilibre sont plus incertains ; pour la
France, Le Barbanchon, Rathelot et Roulet (2019) trouvent un effet
faible de la générosité sur le salaire de réserve. Les allègements de
cotisations sur les bas salaires ont des effets sur l'emploi peu
qualifié documentés, plus faibles et plus incertains au-delà de 1,6
SMIC.

\textbf{Controverses.} Le débat oppose ceux qui voient dans le chômage
européen surtout un problème d'institutions (coût du travail, protection
de l'emploi, générosité des allocations) et ceux qui y voient surtout un
problème de demande persistante et d'hystérèse.

\textbf{Ce que le modèle reprend.} L'emploi de chaque branche marchande
est déterminé par une demande de travail à production donnée (CES,
élasticité de substitution capital-travail \(\sigma\)) ; le chômage est
la différence entre la population active (démographie, retraites, offre
de travail) et l'emploi. Le chômage structurel est exogène sauf pour
l'assurance chômage (0,063 point par 10 \% d'allocations, fourchette
0-0,13) {[}B{]} et la formation {[}X{]}. Les salaires suivent les prix
avec une indexation partielle, la productivité et une courbe de Phillips
faible. Un point de vigilance, documenté au chapitre
\ref{sec:robustesse} : dans certains chocs, l'écart de chômage du modèle
réagit plus fortement que ne l'implique la loi d'Okun, parce que
l'emploi s'ajuste à la production branche par branche alors que la
population active est en partie exogène.

\hypertarget{sec:consommation}{%
\section{Théories de la consommation et de
l'épargne}\label{sec:consommation}}

\textbf{Intuitions.} Keynes faisait dépendre la consommation du revenu
courant. Friedman (1957, revenu permanent) et Modigliani et Brumberg
(1954, cycle de vie) la font dépendre des ressources de toute la vie :
un ménage lisse sa consommation, épargne pendant la vie active et
désépargne à la retraite ; un revenu transitoire est surtout épargné, un
revenu permanent surtout consommé. Ces théories impliquent des
propensions marginales à consommer faibles pour les transferts
temporaires --- sauf si les ménages sont contraints dans leur accès au
crédit.

\textbf{Équations.} Avec un taux d'intérêt égal au taux de préférence
pour le présent, l'hypothèse du revenu permanent donne
\(C_t = \frac{r}{1+r}\left[A_t + \sum_{k \ge 0} \frac{\mathbb{E}_t Y_{t+k}}{(1+r)^k}\right]\)
: la consommation réagit aux révisions du revenu permanent, pas aux
variations anticipées. L'équation d'Euler,
\(u'(C_t) = \beta(1+r)\,\mathbb{E}_t u'(C_{t+1})\), relie la croissance
de la consommation au taux réel avec une élasticité de substitution
intertemporelle (EIS) dont la méta-analyse de Havranek (2015) situe la
valeur centrale autour de 0,3 à 0,5 une fois corrigé le biais de
publication.

\textbf{Résultats empiriques.} Les propensions marginales à consommer
observées sont beaucoup plus élevées que ne le prédit l'hypothèse du
revenu permanent pure, surtout pour les ménages modestes ou peu liquides
(Jappelli et Pistaferri, 2010) ; l'épargne de précaution (Carroll, 1997)
et les contraintes de liquidité (Deaton, 1991) expliquent cet écart. Les
ménages réagissent aux chocs permanents de revenu bien plus qu'aux chocs
transitoires (Blundell, Pistaferri et Preston, 2008). Pendant la crise
sanitaire, l'épargne forcée a montré la sensibilité de la consommation à
l'incertitude.

\textbf{Controverses.} L'importance relative des ménages « non
ricardiens » (qui consomment leur revenu courant) varie selon les études
et les périodes ; l'effet de la dette publique sur l'épargne des ménages
(équivalence ricardienne) est jugé faible empiriquement, mais non nul.

\textbf{Ce que le modèle reprend.} La consommation agrégée suit une
équation à correction d'erreur estimée sur le revenu disponible réel
(élasticité de long terme unitaire), complétée par une épargne de
précaution (0,3 \% de consommation par point de chômage, fourchette
0,1-0,6) et un effet du taux réel du crédit à l'habitat (0,2, fourchette
0-0,5) {[}C{]}. Pour les mesures socio-fiscales, des propensions
marginales à consommer par décile (module M14) font que les transferts
aux ménages modestes stimulent davantage la consommation que les baisses
d'impôt pour les ménages aisés. L'équivalence ricardienne est écartée.

\hypertarget{sec:investissement}{%
\section{Théories de l'investissement et du coût du
capital}\label{sec:investissement}}

\textbf{Intuitions.} Trois familles d'explications coexistent. Le
\textbf{principe d'accélérateur} fait dépendre l'investissement de la
croissance anticipée de la demande. La théorie \textbf{néoclassique}
(Jorgenson, 1963 ; Hall et Jorgenson, 1967) le fait dépendre du coût
d'usage du capital, qui combine taux d'intérêt, dépréciation, fiscalité
et prix relatif des biens d'équipement. La théorie du \textbf{q de
Tobin} (1969) le fait dépendre du rapport entre la valeur de marché des
entreprises et le coût de remplacement de leur capital. S'y ajoutent les
contraintes financières (l'autofinancement compte quand l'accès au
crédit est imparfait) et l'incertitude (option d'attendre, Bloom, 2009).

\textbf{Équations.} Le coût d'usage s'écrit, en notation simplifiée,

\[c_K = \frac{p_I}{p}\,\frac{(r + \delta)(1 - z\,\tau)}{1 - \tau},\]

où \(p_I/p\) est le prix relatif de l'investissement, \(r\) le taux
d'intérêt réel, \(\delta\) la dépréciation, \(\tau\) le taux de l'impôt
sur les sociétés et \(z\) la valeur actualisée des amortissements
fiscaux. Le capital désiré \(K^* = \alpha Y / c_K\) (production
Cobb-Douglas) ou \(K^* \propto Y\,c_K^{-\sigma}\) (CES).

\textbf{Résultats empiriques.} L'accélérateur est l'explication la plus
robuste à court terme ; l'élasticité de l'investissement au coût du
capital est significative mais modeste à moyen terme, de l'ordre de −0,5
à −1 selon les études microéconomiques récentes. L'effet du CICE sur
l'investissement a été jugé faible par les évaluations de France
Stratégie ; la trésorerie supplémentaire a été surtout utilisée pour les
marges et le désendettement.

\textbf{Ce que le modèle reprend.} L'investissement de chaque branche
suit un accélérateur en correction d'erreur {[}E{]}, avec un coût
d'usage qui intègre l'impôt sur les sociétés et, depuis la phase 5, le
taux des crédits aux sociétés non financières transmis par le bloc
monétaire {[}E{]}. Un canal de trésorerie (0,10 de l'épargne
supplémentaire des entreprises investie, fourchette 0-0,3) {[}C{]} relie
profits et investissement. Les effets de l'incertitude politique sur
l'investissement ne sont pas représentés (chapitre
\ref{sec:comportements}).

\hypertarget{sec:anticipations}{%
\section{Anticipations rationnelles, adaptatives et contraintes de
liquidité}\label{sec:anticipations}}

\textbf{Intuitions.} La façon dont les agents forment leurs
anticipations change radicalement l'effet des politiques. Avec des
\textbf{anticipations adaptatives}, les agents extrapolent le passé ;
une politique nouvelle les surprend et produit des effets tant qu'ils ne
l'ont pas apprise. Avec des \textbf{anticipations rationnelles} (Muth,
1961 ; Lucas, 1972 ; Sargent et Wallace, 1975), les agents connaissent
le modèle et utilisent toute l'information disponible : une politique
annoncée et crédible est intégrée immédiatement, et seules les surprises
ont des effets réels (proposition d'inefficacité de Sargent et Wallace).

\textbf{La critique de Lucas.} Lucas (1976) montre que les paramètres
des équations estimées (propension à consommer, réaction des salaires)
dépendent eux-mêmes de la politique en vigueur ; un modèle estimé sous
un régime ne permet donc pas, en toute rigueur, de prévoir les effets
d'un changement de régime. C'est une critique directe des modèles
macroéconométriques du type de celui de l'étude.

\textbf{Contraintes de liquidité.} Même des agents parfaitement
rationnels ne peuvent lisser leur consommation s'ils ne peuvent
emprunter ; une part significative des ménages se comporte donc « à la
Keynes ». Les modèles modernes combinent anticipations rationnelles et
hétérogénéité des contraintes (section \ref{sec:heterogenes}).

\textbf{Résultats empiriques.} Les enquêtes d'anticipations (inflation,
revenus) montrent des anticipations ni parfaitement rationnelles ni
purement adaptatives : elles sont hétérogènes, réagissent avec retard,
et dépendent de l'expérience vécue. Les marchés financiers, en revanche,
intègrent rapidement l'information sur les politiques monétaires
annoncées.

\textbf{Ce que le modèle reprend.} Les anticipations sont adaptatives
(implicites dans les dynamiques de correction d'erreur), avec une option
d'anticipations rationnelles pour la transmission du taux directeur au
taux allemand à dix ans. La critique de Lucas s'applique : les résultats
doivent être lus comme les effets d'une politique dans un environnement
dont les comportements restent ceux observés de 1997 à 2019, ce qui est
d'autant moins vrai que la politique est éloignée de l'expérience
historique (chapitre \ref{sec:comportements}). Les \textbf{effets de
crédibilité} (une règle d'or constitutionnelle qui réduirait la prime de
risque, une politique jugée insoutenable qui l'augmenterait) ne sont
représentés que par la prime souveraine et les variantes de tension.

\hypertarget{sec:ouverte}{%
\section{Économie ouverte, commerce international et
compétitivité}\label{sec:ouverte}}

\textbf{Intuitions.} Dans une économie ouverte, une partie de la demande
intérieure est satisfaite par des importations (fuite du
multiplicateur), et la demande étrangère soutient la production
nationale par les exportations. La compétitivité dépend des prix
relatifs (coûts salariaux, productivité, change) et de facteurs hors
prix (qualité, spécialisation, innovation). Le modèle de Mundell-Fleming
(1962-1963) montre qu'en changes fixes et mobilité des capitaux, la
politique budgétaire est efficace et la politique monétaire nationale
inopérante --- la situation d'un pays de la zone euro vis-à-vis de ses
partenaires.

\textbf{Équations.} Les exportations en volume dépendent de la demande
mondiale adressée et de la compétitivité-prix,
\(\ln X = \ln D^* - \eta_X \ln (p_X / p^*)\) ; les importations
dépendent de la demande pondérée par les contenus en importations et des
prix relatifs, \(\ln M = \ln D - \eta_M \ln(p_M / p)\). La condition de
Marshall-Lerner donne les conditions sous lesquelles une dépréciation
améliore le solde commercial.

\textbf{Résultats empiriques.} Les élasticités-prix du commerce sont
modestes à court terme et plus élevées à long terme ; la compétitivité
hors prix explique une part importante de la dégradation des parts de
marché françaises depuis 2000. Les droits de douane sont en grande
partie supportés par les acheteurs du pays importateur (épisode
américain de 2018-2019) et appellent des représailles.

\textbf{Controverses.} La pertinence des politiques de relocalisation et
de protection, l'ampleur des gains de l'intégration commerciale et leur
répartition (« choc chinois », Autor, Dorn et Hanson, 2013) restent
discutées.

\textbf{Ce que le modèle reprend.} Exportations et importations suivent
des équations estimées {[}E{]} ; le TES permet de suivre les contenus en
importations de chaque composante de la demande. Les droits de douane
(curseur hors programme) sont modélisés de manière stylisée (part hors
UE des importations, répercussion, riposte), en rappelant qu'ils
relèvent de la compétence européenne. Il n'y a pas de taux de change
propre à la France.

\hypertarget{sec:institutions}{%
\section{Économie institutionnelle et contraintes
budgétaires}\label{sec:institutions}}

\textbf{Intuitions.} Les institutions --- règles budgétaires,
indépendance de la banque centrale, droit européen, organisation de la
protection sociale --- encadrent les politiques possibles et influencent
les anticipations. North (1990) en fait des « règles du jeu » qui
déterminent la performance économique de long terme. En matière
budgétaire, deux questions dominent : la soutenabilité de la dette et
l'interaction entre politiques budgétaire et monétaire.

\textbf{Équations.} La dynamique du ratio de dette \(d\) s'écrit

\[d_t = \frac{1 + i_t}{1 + g_t}\,d_{t-1} - s_t,\]

où \(i\) est le taux d'intérêt apparent, \(g\) la croissance nominale et
\(s\) le solde primaire (en \% du PIB). Le solde primaire qui stabilise
le ratio est \(s^* = \frac{i - g}{1 + g}\,d\). Bohn (1998) propose un
critère de soutenabilité : la dette est soutenable si le solde primaire
réagit positivement au ratio de dette. Leeper (1991) distingue les
régimes où la politique budgétaire est « passive » (elle stabilise la
dette) de ceux où elle est « active » (la politique monétaire doit alors
s'adapter, ou l'inflation ajuste la valeur réelle de la dette).

\textbf{Résultats empiriques.} Le seuil de dette au-delà duquel la
croissance serait durablement plus faible, avancé par Reinhart et Rogoff
(2010), n'a pas résisté à la vérification (erreurs de calcul relevées
par Herndon, Ash et Pollin, 2014) ; il n'existe pas de seuil universel.
Blanchard (2019) a souligné qu'avec des taux inférieurs à la croissance,
le coût budgétaire de la dette était faible ; la remontée des taux
depuis 2022 a inversé cette situation pour la France. Les primes de
risque souveraines réagissent à la dette et aux déficits, de manière non
linéaire et fortement dépendante du contexte (crise de la zone euro de
2010-2012).

\textbf{Ce que le modèle reprend.} Le module M23 représente les règles
budgétaires (stabilisation du ratio de dette, ajustement vers une cible,
clause de sauvegarde) ; M24 représente la réaction de la BCE et la prime
souveraine estimée sur un panel de la zone euro ; M26 représente le
cadre européen (procédure de déficit excessif, droit de la TVA,
compétences de l'Union). La dynamique de la dette est calculée par
cohortes d'émission (M17).

\hypertarget{sec:heterogenes}{%
\section{Approches stock-flux cohérentes et modèles à agents
hétérogènes}\label{sec:heterogenes}}

\textbf{Approches stock-flux cohérentes.} Héritières de Tobin et
développées par Godley et Lavoie (2007), elles imposent que chaque flux
financier ait une contrepartie et que les stocks (dettes, actifs)
résultent de l'accumulation des flux. Elles attirent l'attention sur les
déséquilibres de bilans (endettement privé, position extérieure) que les
modèles centrés sur les flux ignorent, et sur l'interdépendance des
soldes sectoriels : le solde public, le solde privé et le solde
extérieur somment à zéro.

\textbf{Modèles à agents hétérogènes.} Les modèles HANK
(\emph{Heterogeneous Agent New Keynesian}, Kaplan, Moll et Violante,
2018) remplacent l'agent représentatif par une distribution de ménages
différant par leurs revenus, leur patrimoine et leur accès au crédit.
Ils montrent que la transmission de la politique monétaire passe surtout
par les effets de revenu indirects (emploi, salaires) plutôt que par la
substitution intertemporelle, et que les effets des politiques
budgétaires dépendent de qui reçoit les transferts et de qui paie les
impôts.

\textbf{Ce que le modèle reprend.} Le modèle tient les comptes de tous
les secteurs institutionnels et vérifie les identités entre flux
(capacités et besoins de financement) à chaque simulation ; il suit la
dette publique par cohortes et la dette des sociétés non financières. Il
n'a pas de bilan complet des ménages ni de position extérieure
détaillée. L'hétérogénéité des ménages est représentée de manière «
semi-structurelle » : la microsimulation calcule l'effet de chaque
mesure ménage par ménage, et des propensions à consommer par décile
traduisent cet effet en demande. Ce n'est pas un modèle HANK : les
comportements individuels ne sont pas optimisés et la distribution ne
réagit pas dynamiquement à la conjoncture (sauf dans le module de
distribution dynamique M16, simplifié).

\hypertarget{sec:synthese-theorie}{%
\section{Synthèse : ce que les désaccords théoriques changent pour
l'évaluation des programmes}\label{sec:synthese-theorie}}

Le tableau suivant résume, pour les grandes familles de mesures des
programmes, le sens des effets selon les écoles, et la position du
modèle.

\begin{longtable}[]{@{}
  >{\raggedright\arraybackslash}p{(\columnwidth - 6\tabcolsep) * \real{0.2500}}
  >{\raggedright\arraybackslash}p{(\columnwidth - 6\tabcolsep) * \real{0.2500}}
  >{\raggedright\arraybackslash}p{(\columnwidth - 6\tabcolsep) * \real{0.2500}}
  >{\raggedright\arraybackslash}p{(\columnwidth - 6\tabcolsep) * \real{0.2500}}@{}}
\toprule\noalign{}
\begin{minipage}[b]{\linewidth}\raggedright
Famille de mesures
\end{minipage} & \begin{minipage}[b]{\linewidth}\raggedright
Lecture keynésienne
\end{minipage} & \begin{minipage}[b]{\linewidth}\raggedright
Lecture classique / offre
\end{minipage} & \begin{minipage}[b]{\linewidth}\raggedright
Position du modèle (régime principal)
\end{minipage} \\
\midrule\noalign{}
\endhead
\bottomrule\noalign{}
\endlastfoot
Hausse de dépense publique & effet positif fort à court terme
(multiplicateur \textgreater{} 1 en sous-emploi) & éviction, effet
faible ; hausse future des impôts & multiplicateur ≈ 1,2 à 1 an, 0,75 à
5 ans ; éviction par les taux seulement en régime BCE + prime \\
Hausse des transferts aux ménages modestes & effet positif (propension à
consommer élevée) & effet désincitatif possible sur le travail & effet
de demande selon les propensions par décile ; effet sur la participation
pour les prestations sous condition de ressources (M15) \\
Hausse des impôts sur le capital et les patrimoines & effet de demande
faible (propension faible des hauts revenus) & réduction de l'épargne,
de l'investissement, départs & érosion d'assiette (M13) ; effet de
demande faible ; pas d'effet sur l'investissement par le patrimoine \\
Baisse des impôts de production, des cotisations & effet de demande
modéré & baisse du coût, compétitivité, emploi & baisse des prix
(répercussion 50 \%), hausse des marges et des salaires selon
l'incidence \\
Réforme des retraites (âge) & relèvement : baisse de revenu des seniors,
chômage & hausse de l'offre de travail et du potentiel & M9 : population
active, pensions ; la fermeture absorbe progressivement l'offre
supplémentaire \\
Consolidation budgétaire & récessive, peut dégrader le ratio de dette à
court terme & confiance, baisse des taux & récessive ; amélioration du
ratio de dette à moyen terme seulement ; prime souveraine en option \\
\end{longtable}

Ces écarts ne se résolvent pas par une meilleure estimation d'un seul
paramètre : ils tiennent à la vitesse à laquelle l'économie revient vers
son potentiel, à la réaction des taux et aux anticipations. C'est
pourquoi l'étude présente systématiquement ses résultats sous plusieurs
régimes (règle budgétaire, politique monétaire, vitesse de fermeture,
ancrage des anticipations) et rend compte de ce qui en dépend.

\hypertarget{sec:architecture}{%
\chapter{Architecture scientifique du modèle}\label{sec:architecture}}

Ce chapitre décrit l'organisation du modèle : ses modules, la manière
dont ils s'articulent, le déroulement d'une simulation annuelle,
l'articulation avec la microsimulation, le traitement du calendrier des
mesures, les régimes de politique économique, la démographie et les
retraites, et la cohérence comptable entre secteurs. Les équations sont
détaillées au chapitre \ref{sec:equations}.

\hypertarget{sec:modules}{%
\section{Vue d'ensemble et modules}\label{sec:modules}}

Le modèle est écrit en Python ; il est organisé en modules numérotés (M1
à M27), conservés dans le dépôt du projet avec leurs tests. Le tableau
suivant en donne la liste fonctionnelle.

\begin{longtable}[]{@{}
  >{\raggedright\arraybackslash}p{(\columnwidth - 4\tabcolsep) * \real{0.3333}}
  >{\raggedright\arraybackslash}p{(\columnwidth - 4\tabcolsep) * \real{0.3333}}
  >{\raggedright\arraybackslash}p{(\columnwidth - 4\tabcolsep) * \real{0.3333}}@{}}
\toprule\noalign{}
\begin{minipage}[b]{\linewidth}\raggedright
Couche
\end{minipage} & \begin{minipage}[b]{\linewidth}\raggedright
Modules
\end{minipage} & \begin{minipage}[b]{\linewidth}\raggedright
Rôle
\end{minipage} \\
\midrule\noalign{}
\endhead
\bottomrule\noalign{}
\endlastfoot
Données et comptes & M1 démographie, M3 cadre comptable, M5 tableau
entrées-sorties et prix de Leontief & population par âge et sexe
(projections Insee 2026), comptes des secteurs institutionnels,
structure productive en 64 produits et 25 branches marchandes \\
Ménages (micro) & M4 population et OpenFisca, M10 couplage micro-macro,
M14 hétérogénéité, M16 distribution dynamique & impôt sur le revenu et
prestations ménage par ménage ; propensions à consommer par décile ;
transitions emploi-chômage en espérance \\
Comportements & M7 équations estimées, M13 réponses des hauts revenus et
patrimoines, M15 offre de travail, M9 activité des seniors et retraites
& emploi, prix, salaires, consommation, investissement, logement,
commerce ; élasticités fiscales ; participation ; départs à la retraite
et pensions \\
Cœur & M8 (avec M17 dette par cohortes, M18 prix, M19 salaires, M21
compte des sociétés) & résolution annuelle, fermeture d'offre, finances
publiques, comptes d'agents \\
Référence & M8\_reference & trajectoire à politique inchangée : calage
2023-2025, prévisions du FMI 2026-2031, hypothèses de long terme \\
Europe et monnaie & M11 (zone euro), M24 politique monétaire et
conditions financières & règle de la BCE, taux allemand, prime
souveraine, taux des crédits \\
Incertitude, validation & M12, M22 & variantes de paramètres, validation
prédictive (millésimes du FMI) \\
Régimes et programmes & M23 régimes, M25 traduction des programmes, M26
contraintes européennes, M27 corpus des programmes & moteur d'évaluation
environnement × régime × programme \\
\end{longtable}

La version de référence de cette étude est la configuration \textbf{P4}
(section \ref{sec:gel}), qui ajoute à la version 0.9 une équation des
prix de valeur ajoutée fondée sur une frontière des prix des facteurs et
un compte complet des sociétés non financières.

\hypertarget{sec:env-regime-prog}{%
\section{Environnement, régime et programme}\label{sec:env-regime-prog}}

Le principe d'organisation de l'évaluation distingue trois objets.

\begin{itemize}
\tightlist
\item
  \textbf{L'environnement} \(E\) décrit l'économie à politique inchangée
  : les équations, leurs paramètres, et une trajectoire de référence
  calibrée de 2022 à 2047 (exogènes et termes de calage). Il est calibré
  une seule fois par configuration et par convention de dépense, puis
  mis en cache avec une empreinte de tous ses paramètres. Un
  environnement est modifié seulement si l'on change une équation ou un
  paramètre structurel qui exige un recalage (par exemple l'élasticité
  de substitution \(\sigma\)).
\item
  \textbf{Le régime} \(R\) décrit les règles de politique économique
  appliquées pendant la projection : règle budgétaire (aucune,
  stabilisation du ratio de dette, ajustement vers une cible ; effort
  figé ou réactif), régime monétaire (taux exogènes ou règle de la BCE,
  prime souveraine), jeu de paramètres fragiles (ancrage des
  anticipations, vitesse de fermeture), paramètres structurels
  modifiables sans recalage.
\item
  \textbf{Le programme} \(P\) est l'ensemble des mesures d'un candidat,
  traduites en modifications des exogènes du modèle (taux d'imposition,
  dépenses, prestations, règles de retraite, etc.), chacune avec sa date
  de première application.
\end{itemize}

L'effet d'un programme est défini comme

\[\text{Effet}(P \mid E, R) = Y(E, R, P) - Y(E, R, \varnothing),\]

où \(Y\) désigne la trajectoire simulée. Dans le régime principal (« R2
figé »), l'effort budgétaire de la règle est celui de la trajectoire
sans programme : le programme ne déclenche pas d'effort supplémentaire,
de sorte que son effet n'est pas confondu avec la réaction de la règle.
Dans les régimes « actifs », la règle réagit au programme (chapitre
\ref{sec:robustesse}).

\begin{figure}[htbp]
\centering
\begin{tikzpicture}[node distance=1.2cm and 1.6cm, every node/.style={font=\small},
  boite/.style={draw, rounded corners, align=center, minimum width=3.2cm, minimum height=1cm, fill=gray!8},
  fleche/.style={-{Stealth[length=2.5mm]}, thick}]
\node[boite] (env) {\textbf{Environnement $E$}\\équations, paramètres,\\référence calibrée (cache)};
\node[boite, right=of env] (reg) {\textbf{Régime $R$}\\règle budgétaire, BCE,\\prime souveraine, jeux};
\node[boite, right=of reg] (prog) {\textbf{Programme $P$}\\mesures datées\\(corpus M27 → M25)};
\node[boite, below=1.6cm of reg, minimum width=5cm] (sim) {\textbf{Simulation annuelle M8}\\2023 → 2047};
\node[boite, below left=1.4cm and -0.6cm of sim] (yp) {$Y(E,R,P)$};
\node[boite, below right=1.4cm and -0.6cm of sim] (y0) {$Y(E,R,\varnothing)$};
\node[boite, below=3.4cm of sim, minimum width=5cm] (eff) {\textbf{Effet} $= Y(E,R,P) - Y(E,R,\varnothing)$};
\draw[fleche] (env) -- (sim); \draw[fleche] (reg) -- (sim); \draw[fleche] (prog) -- (sim);
\draw[fleche] (sim) -- (yp); \draw[fleche] (sim) -- (y0);
\draw[fleche] (yp) -- (eff); \draw[fleche] (y0) -- (eff);
\end{tikzpicture}
\caption{Organisation de l'évaluation : environnement, régime, programme.}
\end{figure}

\hypertarget{sec:annee}{%
\section{Fonctionnement d'une année de simulation}\label{sec:annee}}

Chaque année \(t\) de la projection est résolue par un \textbf{point
fixe amorti} : on cherche des valeurs de la production, des salaires et
du chômage telles que toutes les équations de l'année soient satisfaites
simultanément. L'algorithme itère au plus 60 fois, avec un amortissement
de 0,5 et un critère de convergence de \(10^{-7}\). L'ordre des calculs
dans une itération est le suivant :

\begin{enumerate}
\def\labelenumi{\arabic{enumi}.}
\tightlist
\item
  \textbf{Crochets de régime} (avant la résolution, avec l'information
  disponible jusqu'en \(t-1\)) : régime monétaire (taux directeur, taux
  allemand, écart souverain, taux des crédits), puis règle budgétaire
  (effort et répartition entre instruments), puis, le cas échéant,
  objectif budgétaire du candidat (lecture B) ;
\item
  \textbf{prix d'importation} (prix mondiaux, énergie) ;
\item
  \textbf{salaires, prix de la valeur ajoutée par branche, prix de
  Leontief} des produits et de la consommation ;
\item
  \textbf{emploi et investissement} par branche ;
\item
  \textbf{revenu disponible des ménages et demande finale}
  (consommation, logement, investissement, dépenses publiques,
  exportations) ;
\item
  \textbf{importations et production} par branche (système de Leontief).
\end{enumerate}

Après convergence, le modèle calcule le PIB potentiel et l'écart de
production, puis les comptes des administrations publiques (recettes,
dépenses, charge d'intérêts par cohortes, solde, dette) et le compte des
sociétés (intérêts, impôt sur les sociétés, dividendes, épargne,
capacité de financement, dette). Ces comptes ne rétroagissent sur
l'année qu'à travers des variables retardées (dividendes de l'exercice
précédent, dette de fin d'année précédente), ce qui correspond au
calendrier réel (les dividendes de l'année \(t\) rémunèrent l'exercice
\(t-1\)) et garantit qu'aucune information future n'est utilisée.

\begin{figure}[htbp]
\centering
\begin{tikzpicture}[node distance=0.55cm, every node/.style={font=\small},
  etape/.style={draw, rounded corners, align=center, minimum width=9cm, minimum height=0.75cm, fill=blue!4},
  fleche/.style={-{Stealth[length=2.5mm]}, thick}]
\node[etape] (a) {Crochets de régime : BCE, taux, règle budgétaire, objectif (info. $\le t-1$)};
\node[etape, below=of a] (b) {Prix d'importation};
\node[etape, below=of b] (c) {Salaires $\to$ prix de VA $\to$ prix de Leontief};
\node[etape, below=of c] (d) {Emploi et investissement par branche};
\node[etape, below=of d] (e) {Revenu disponible $\to$ demande finale};
\node[etape, below=of e] (f) {Importations $\to$ production (Leontief)};
\node[etape, below=0.9cm of f, fill=green!6] (g) {Potentiel, finances publiques, compte des sociétés, comptes d'agents};
\draw[fleche] (a)--(b); \draw[fleche] (b)--(c); \draw[fleche] (c)--(d); \draw[fleche] (d)--(e); \draw[fleche] (e)--(f);
\draw[fleche] (f.east) -- ++(1.0,0) |- node[pos=0.25, right, align=left]{point fixe\\(≤ 60 itérations)} (c.east);
\draw[fleche] (f)--(g);
\end{tikzpicture}
\caption{Déroulement d'une année de simulation.}
\end{figure}

\hypertarget{sec:fermeture}{%
\section{La fermeture d'offre : comment la demande rejoint le
potentiel}\label{sec:fermeture}}

Un modèle dont la production est déterminée par la demande ne revient
pas nécessairement à son potentiel : rien dans les équations de court
terme ne garantit que le chômage revienne à son niveau structurel. Les
modèles institutionnels résolvent ce problème par des mécanismes de
retour (effet des prix sur la demande réelle, réaction de la politique
monétaire, effets de richesse). Dans le modèle de l'étude, ces canaux
sont soit absents (pas de taux de change propre, pas d'effets de
richesse), soit faibles ; une \textbf{fermeture explicite} les
représente : un supplément de demande résorbe chaque année une fraction
\(\rho = 0{,}2\) de l'écart de production {[}B, inféré de la durée des
écarts de production 1982-2019{]}, et anticipe une part
\(\theta = 0{,}75\) des révisions du potentiel {[}C{]}. Cette fermeture
est l'hypothèse qui gouverne la transition du court au long terme : avec
\(\rho = 0{,}1\), les effets de demande durent deux fois plus longtemps
; avec \(\rho = 0{,}3\), ils s'estompent plus vite. Les deux variantes,
et une fermeture pilotée par le chômage, font partie des régimes de
robustesse.

\hypertarget{sec:micro}{%
\section{Microsimulation et articulation avec
OpenFisca}\label{sec:micro}}

Les mesures qui modifient l'impôt sur le revenu, la CSG ou les
prestations sous condition de ressources sont évaluées sur une
\textbf{population de ménages} construite par le module M4 à partir du
recensement de 2022 (fichier détail des individus) et des tableaux
fiscaux de la DGFiP (impôt sur le revenu par tranche, déclarations
nationales). L'impôt est calculé par OpenFisca-France (moteur public de
la législation) sur un échantillon de 2 \% des ménages ; les prestations
(RSA, prime d'activité, allocations familiales, aides au logement) sur
deux sous-échantillons de 0,5 \%, avec un taux de recours apparent
mesuré en législation actuelle (RSA 0,68, prime d'activité 0,66).

Le \textbf{couplage micro-macro} (M10) fonctionne dans un seul sens :

\begin{enumerate}
\def\labelenumi{\arabic{enumi}.}
\tightlist
\item
  la réforme est décrite par des paramètres OpenFisca modifiés ;
\item
  l'impôt et les prestations sont calculés ménage par ménage en
  législation actuelle et après réforme (effet statique, comportements
  inchangés, revenus de 2024) ;
\item
  les écarts sont ramenés aux montants versés (foyers DGFiP, taux de
  recours) et répartis par décile de niveau de vie ;
\item
  les montants agrégés (M€ de 2024) entrent dans le modèle bouclé,
  indexés sur le PIB nominal simulé ; des propensions marginales à
  consommer par décile (M14) traduisent la répartition en effet sur la
  consommation.
\end{enumerate}

Ce couplage a deux limites assumées : les revenus de la microsimulation
ne sont pas revalorisés avec les salaires, l'emploi et les prix simulés
(pas de retour macro → micro), et les comportements des ménages ne sont
pas microsimulés (l'offre de travail passe par des élasticités moyennes,
M15). Les mesures sur les revenus du capital et les patrimoines
(prélèvement forfaitaire unique, impôt sur la fortune, successions) sont
évaluées sur des assiettes agrégées, avec une répartition approchée sur
le haut de la distribution.

\textbf{Risque de double compte.} Une mesure ne doit pas être comptée à
la fois par la microsimulation et par un canal agrégé. Le pont des
programmes (M25) le garantit : les montants statiques de la
microsimulation remplacent les montants des canaux agrégés
correspondants ; les profits induits par les mesures passent par le
compte des sociétés (M21), et non plus par une hypothèse séparée de M14
(double compte supprimé en phase 4.3). Le chapitre \ref{sec:donnees}
présente les contrôles correspondants.

\hypertarget{sec:dates}{%
\section{Traitement des dates de lancement des
mesures}\label{sec:dates}}

Chaque mesure porte une \textbf{première année d'application}. Par
défaut, il s'agit de 2028, premier budget voté après l'élection de 2027
; le corpus retient 2027 lorsque le candidat annonce une application dès
l'année de l'élection (collectif budgétaire de l'été 2027, ordonnances
de juillet 2027) et des dates ultérieures lorsqu'il l'annonce (par
exemple la réforme des retraites du RN au 1er janvier 2028). Le pont des
programmes accepte des \textbf{listes de mesures datées}, de sorte que
deux mesures de même nature peuvent avoir des calendriers différents.
Une application au cours de l'année est approchée par l'année entière
suivante ou, si nécessaire, par un prorata (non utilisé dans les
résultats de base). Le pas annuel a été conservé après examen : les
mesures des programmes sont annuelles et les horizons d'analyse sont de
cinq ans et plus (chapitre \ref{sec:donnees}).

\hypertarget{sec:regimes-budgetaires}{%
\section{Régimes budgétaires}\label{sec:regimes-budgetaires}}

La règle budgétaire générique (M23) détermine chaque année un
\textbf{effort} \(e_t\) (en points de PIB), réparti entre quatre
instruments : consommation publique, prestations, impôts des ménages,
impôts sur les produits (TVA). Sa forme est

\[\bar s_t = \frac{\bar r - \bar g}{1 + \bar g}\,d_{t-1}, \qquad s^*_t = \bar s_t + \beta\,(d_{t-1} - d^*), \qquad
\Delta e_t = \text{clip}\left(\alpha\,(s^*_t - s^{ca}_{t-1}),\ 0,\ \bar e\right),\]

où \(\bar s\) est le solde primaire stabilisant la dette, \(\bar r\) et
\(\bar g\) les moyennes sur trois ans du taux apparent et de la
croissance nominale potentielle, \(s^{ca}\) le solde primaire corrigé du
cycle (semi-élasticité 0,5), \(d^*\) la cible de dette, \(\alpha\) la
vitesse d'ajustement et \(\bar e\) le plafond annuel. Une clause de
sauvegarde suspend l'effort quand l'écart de production est inférieur à
−1,5 point. Les régimes nommés sont : R0 (aucune règle), R1
(stabilisation, \(\beta = 0\)), R2 (ajustement, \(\alpha = 0{,}5\),
\(\beta = 0{,}02\), \(d^* = 90\,\%\), plafond 0,5 point, répartition
30/20/30/20) {[}H, choix du projet{]}, en mode figé ou actif, symétrique
ou non.

\textbf{Lecture par objectifs budgétaires (lecture B).} Pour représenter
la trajectoire de solde annoncée par un candidat, un second dispositif
s'ajoute à la règle du régime : chaque année, un effort supplémentaire
\(\Delta e^B_t = \text{clip}\left(c_t - s_{t-1}, -2, 2\right)\)
rapproche le solde effectif \(s\) de la cible \(c_t\) (interpolée
linéairement entre les points annoncés, tenue ensuite), et se répartit
selon la composition déduite des déclarations du candidat. Trois passes
de simulation corrigent la cible interne de l'écart réalisé, ce qui
revient à supposer que le gouvernement anticipe correctement la dérive
du solde {[}H{]}. Le chapitre \ref{sec:resultats} présente les écarts
résiduels à la cible.

\hypertarget{sec:bce}{%
\section{Politique monétaire et réaction de la BCE}\label{sec:bce}}

Dans le régime principal, les taux d'intérêt sont \textbf{exogènes} : le
taux court de référence est l'Euribor à trois mois observé (2,63 \% en
octobre 2026) et le taux long l'OAT à dix ans (dernière observation
mensuelle), maintenus constants. Ce choix isole l'effet propre des
programmes ; il revient à supposer que la BCE ne réagit pas à un choc
français (la France pèse environ 20 \% de la zone euro) et que la prime
souveraine est inchangée.

Dans le \textbf{régime monétaire endogène} (M24), la BCE suit une règle
de Taylor estimée en écart à la référence, sur l'inflation et l'activité
de la zone euro, auxquelles la France contribue selon son poids ; le
taux allemand à dix ans réagit au taux directeur ; l'écart entre l'OAT
et le taux allemand dépend de la dette et du solde public français en
écart à la référence (prime souveraine estimée sur un panel de la zone
euro hors 2010-2013, avec une variante incluant la crise) ; les taux des
crédits nouveaux aux ménages et aux sociétés s'ajustent partiellement
aux taux de marché (séries MIR de la BCE). Les taux des crédits entrent
dans l'arbitrage consommation-épargne des ménages et dans le coût
d'usage du capital des entreprises.

\begin{figure}[htbp]
\centering
\begin{tikzpicture}[node distance=0.8cm and 1.3cm, every node/.style={font=\small},
  b/.style={draw, rounded corners, align=center, minimum width=2.8cm, minimum height=0.9cm, fill=orange!6},
  f/.style={-{Stealth[length=2.5mm]}, thick}]
\node[b] (fr) {Inflation, activité\\France ($t-1$)};
\node[b, right=of fr] (ze) {Inflation, activité\\zone euro};
\node[b, right=of ze] (bce) {Taux directeur\\(règle de Taylor)};
\node[b, below=of bce] (bund) {Taux allemand\\10 ans};
\node[b, below=of fr] (finpub) {Dette, solde\\France ($t-1$)};
\node[b, right=of finpub] (prime) {Écart souverain\\(prime)};
\node[b, below=of prime] (oat) {OAT 10 ans};
\node[b, below=of bund] (cred) {Taux des crédits\\ménages, sociétés};
\node[b, below=of oat] (dem) {Consommation,\\investissement, intérêts};
\draw[f] (fr) -- (ze); \draw[f] (ze) -- (bce); \draw[f] (bce) -- (bund);
\draw[f] (finpub) -- (prime); \draw[f] (prime) -- (oat); \draw[f] (bund) -- (oat);
\draw[f] (bce) -- (cred); \draw[f] (oat) -- (cred); \draw[f] (cred) -- (dem); \draw[f] (oat) -- (dem);
\end{tikzpicture}
\caption{Transmission monétaire et financière dans le régime endogène (M24).}
\end{figure}

\hypertarget{sec:prime}{%
\section{Conditions financières et prime souveraine}\label{sec:prime}}

La prime souveraine (écart entre l'OAT et le taux allemand) est estimée
sur un panel de pays de la zone euro avec effets fixes : environ 1,3
point de base par point de dette publique (en \% du PIB) et 1,5 point de
base par point de déficit, hors période de crise ; en incluant
2010-2013, l'effet du déficit est beaucoup plus fort (11,8 points de
base par point) {[}E{]}. Ces coefficients sont des moyennes historiques
: ils ne disent rien des épisodes de crise de confiance, qui sont non
linéaires. Le module M26 permet donc d'ajouter des chocs d'écart
(variante « tension »), dont l'ampleur est une hypothèse de scénario.

\hypertarget{sec:contraintes-ue}{%
\section{Contraintes européennes}\label{sec:contraintes-ue}}

Le module M26 qualifie chaque mesure au regard du droit de l'Union
(niveau 0 : compétence nationale ; 1 : aménagement ; 2 : contentieux
probable ; 3 : effet sur la trajectoire budgétaire et la procédure de
déficit excessif ; 4 : incompatibilité avec les traités) et produit
trois variantes : \emph{conforme} (mesures mises en conformité, part non
conforme des mesures de niveau 4 retirée), \emph{contentieux} (en plus,
mesures de niveaux 1-2 reportées et partiellement retirées) et
\emph{tension} (en plus, amendes de la procédure de déficit excessif et
choc de prime souveraine). Le chapitre \ref{sec:europe} les détaille.

\hypertarget{sec:demo-retraites}{%
\section{Démographie et retraites}\label{sec:demo-retraites}}

La démographie (M1) reprend les projections de population de l'Insee
(2026, scénario central) par âge et sexe. Le module des retraites (M9)
relie la population, les règles de retraite, l'activité des seniors, le
nombre de retraités et la masse des pensions :

\begin{itemize}
\tightlist
\item
  l'\textbf{activité des seniors} de 55 à 74 ans est modélisée par âge,
  sexe et génération, en fonction de l'écart entre l'âge et un
  \textbf{âge pivot} qui dépend de l'âge d'ouverture des droits (AOD) ;
  le profil est estimé sur les taux d'activité de l'enquête Emploi
  (Eurostat, 2003-2019) et testé hors échantillon sur 2020-2025
  (relèvement de l'AOD de 62 ans à 62 ans et 9 mois) ;
\item
  depuis la finalisation, l'âge pivot dépend aussi de la \textbf{durée
  d'assurance requise} : chaque année de durée supplémentaire reporte
  l'âge effectif de cessation d'activité de 9 mois pour les hommes et de
  5 mois pour les femmes, d'après l'évaluation de la réforme de 1993
  (Bozio, 2011) {[}E, littérature{]} ; et une fraction de chaque
  génération peut relever de \textbf{conditions particulières}
  (carrières longues à un âge donné, ou maintien des conditions
  actuelles) ;
\item
  les \textbf{pensions} sont calculées par génération (salaire de
  référence, revalorisation sur les prix), calées sur la masse observée
  ; la durée requise modifie le niveau des pensions de la fraction des
  assurés qui ne reportent pas leur départ (proratisation et décote sur
  la part de base) {[}H{]} ;
\item
  les \textbf{liens avec le modèle bouclé} passent par la population
  active, l'emploi des seniors supplémentaires (part estimée sur
  2010-2019) et l'écart de masse des pensions.
\end{itemize}

La section \ref{sec:eq-retraites} donne les équations et le chapitre
\ref{sec:distribution} les effets générationnels.

\hypertarget{sec:distrib-archi}{%
\section{Distribution des revenus et effets sur les
ménages}\label{sec:distrib-archi}}

Trois niveaux de lecture distributive coexistent : (i) l'effet
\textbf{statique} des mesures socio-fiscales par décile de niveau de vie
(microsimulation) ; (ii) l'effet \textbf{macroéconomique moyen} sur le
pouvoir d'achat, les salaires, l'emploi et les prix ; (iii) un module de
\textbf{distribution dynamique} simplifié (M16), qui répartit les
variations d'emploi, de salaires et de transferts entre ménages en
espérance. Seuls (i) et (ii) sont utilisés dans les résultats de cette
étude ; (iii) est mentionné comme piste.

\hypertarget{sec:secteurs}{%
\section{Relations entre secteurs institutionnels}\label{sec:secteurs}}

Le modèle tient les comptes de quatre secteurs : ménages (S14-S15),
sociétés non financières (S11, compte complet depuis la phase 4.3 ; les
sociétés financières sont traitées de manière agrégée), administrations
publiques (S13) et reste du monde (S2). Les flux principaux sont
représentés sur la figure suivante.

\begin{figure}[htbp]
\centering
\begin{tikzpicture}[every node/.style={font=\scriptsize},
  s/.style={draw, rounded corners, align=center, minimum width=3.0cm, minimum height=1.0cm, fill=gray!10, font=\small},
  f/.style={-{Stealth[length=2.2mm]}, thick}]
\node[s] (men) at (-5.2,0) {\textbf{Ménages}};
\node[s] (snf) at (5.2,0) {\textbf{Sociétés}};
\node[s] (apu) at (0,4.2) {\textbf{Administrations}\\\textbf{publiques}};
\node[s] (rdm) at (0,-4.2) {\textbf{Reste du monde}};
\draw[f] (snf.165) -- node[above]{salaires, dividendes, intérêts} (men.15);
\draw[f] (men.345) -- node[below]{consommation, logement} (snf.195);
\draw[f] (men.70) to[bend left=12] node[above left, align=right]{impôts, cotisations,\\TVA} (apu.180);
\draw[f] (apu.210) to[bend left=12] node[below right, align=left, pos=0.55]{prestations,\\salaires publics} (men.45);
\draw[f] (snf.110) to[bend right=12] node[above right, align=left]{IS, impôts de production,\\cotisations} (apu.0);
\draw[f] (apu.330) to[bend right=12] node[below left, align=right, pos=0.55]{commande publique,\\subventions} (snf.135);
\draw[f] (snf.250) to[bend left=12] node[below right, align=left]{importations,\\dividendes versés} (rdm.0);
\draw[f] (rdm.30) to[bend left=12] node[above left, align=right, pos=0.45]{exportations} (snf.225);
\draw[f] (apu.270) -- node[right, pos=0.35]{intérêts versés aux non-résidents} (rdm.90);
\end{tikzpicture}
\caption{Principaux flux entre secteurs institutionnels représentés dans le modèle.}
\end{figure}

\textbf{Conventions comptables.} Les identités suivantes sont vérifiées
à chaque simulation : la somme des capacités de financement des secteurs
est nulle (aux erreurs de calage près, explicitées) ; le revenu national
brut est égal au PIB plus les revenus primaires nets reçus du reste du
monde ; la variation de la dette publique est égale au besoin de
financement (hors ajustements stock-flux observés jusqu'en 2025) ; les
prix de Leontief vérifient l'égalité entre coûts et prix pour chaque
produit. Le chapitre \ref{sec:donnees} présente les résultats de ces
contrôles.

\textbf{Risques de double comptabilisation identifiés et traités.}

\begin{longtable}[]{@{}
  >{\raggedright\arraybackslash}p{(\columnwidth - 2\tabcolsep) * \real{0.5000}}
  >{\raggedright\arraybackslash}p{(\columnwidth - 2\tabcolsep) * \real{0.5000}}@{}}
\toprule\noalign{}
\begin{minipage}[b]{\linewidth}\raggedright
Risque
\end{minipage} & \begin{minipage}[b]{\linewidth}\raggedright
Traitement
\end{minipage} \\
\midrule\noalign{}
\endhead
\bottomrule\noalign{}
\endlastfoot
Une mesure fiscale comptée par la microsimulation et par un canal agrégé
& les montants microsimulés remplacent le canal agrégé (M25) \\
Profits induits par une mesure comptés deux fois (M14 et compte des
sociétés) & supprimé en phase 4.3 : seul le compte des sociétés répartit
les profits \\
Investissement hors budget (pôle public bancaire) compté en dépense
publique & mesure \texttt{investissement\_hors\_apu} : effet de demande
sans coût pour le solde public ; distincte de l'investissement public \\
Composantes d'une enveloppe comptées en plus de l'enveloppe (RN :
collectivités et agences dans la « réorganisation de l'État ») &
composantes non additionnées (corpus) \\
Effort de la lecture B superposé à une mesure du programme & l'effort B
est calculé sur le solde résultant des mesures : il ne couvre que
l'écart restant \\
Amendes européennes comptées comme dépense intérieure & comptées en
dépense publique sans effet de demande (transfert à l'Union) \\
Coût des retraites compté dans M9 et dans une ligne de transferts & les
réformes de retraite passent par M9 seulement ; les revalorisations de
pensions annoncées séparément passent par les transferts \\
\end{longtable}

\hypertarget{sec:equations}{%
\chapter{Équations et mécanismes économiques}\label{sec:equations}}

Ce chapitre présente les équations du modèle dans la configuration de
référence P4. Pour chacune, on indique sa signification, ses variables
et unités, son origine théorique, sa méthode d'estimation ou de
calibration, l'interprétation de ses paramètres, la sensibilité des
résultats, ses limites et les formulations alternatives. Les notations
suivent une convention simple : les minuscules désignent des logarithmes
(\(y = \ln Y\)), \(\Delta\) la variation d'une année à l'autre, l'indice
\(b\) la branche, \(t\) l'année. Les coefficients estimés le sont sur
1997-2019 (hors années de crise sanitaire) sauf mention contraire.
L'annexe \ref{sec:annexe-b} reprend l'ensemble sous forme compacte.

\textbf{Comment lire une équation à correction d'erreur.} La plupart des
équations de comportement ont la forme

\[\Delta x_t = \underbrace{\sum_k a_k \,\Delta z_{k,t}}_{\text{dynamique de court terme}}
\;-\; \lambda \,\underbrace{\left(x_{t-1} - x^*_{t-1}\right)}_{\text{écart à la cible}}, \qquad
x^*_t = \sum_k b_k z_{k,t} + \text{constante},\]

où \(x^*\) est la valeur de long terme (la « cible ») déterminée par la
théorie, et \(\lambda\) la vitesse de retour vers cette cible (la
fraction de l'écart résorbée chaque année). Cette forme, popularisée par
Engle et Granger (1987), concilie une relation de long terme théorique
et une dynamique de court terme ajustée aux données. Elle est celle des
modèles Mésange (Insee, DG Trésor) et FR-BDF (Banque de France) pour
leurs blocs d'équations de comportement.

\hypertarget{sec:eq-production}{%
\section{Production et PIB potentiel}\label{sec:eq-production}}

\textbf{Fonction de production.} Le potentiel des branches marchandes
est représenté par une fonction CES (élasticité de substitution
constante) normalisée en heures :

\[Q^* = \left[\alpha\,(K)^{\frac{\sigma-1}{\sigma}} + (1-\alpha)\,(E\,H)^{\frac{\sigma-1}{\sigma}}\right]^{\frac{\sigma}{\sigma-1}},\]

où \(Q^*\) est la valeur ajoutée potentielle (volume, euros de l'année
de base), \(K\) le capital (volume net), \(H\) les heures travaillées au
chômage structurel, \(E\) l'efficience du travail (tendance),
\(\alpha = 0{,}35\) la part du capital {[}A, comptes nationaux{]} et
\(\sigma = 0{,}45\) l'élasticité de substitution capital-travail
{[}C{]}. Le potentiel non marchand (administrations, éducation, santé,
action sociale, loyers) est égal à sa valeur ajoutée tendancielle.

\emph{Origine et alternatives.} La CES généralise la Cobb-Douglas
(\(\sigma = 1\)) ; les estimations françaises et européennes de
\(\sigma\) sont en majorité inférieures à 1 (entre 0,3 et 0,7). Une
Cobb-Douglas donnerait une substitution plus forte et des effets du coût
du travail sur l'emploi plus importants.

\emph{Sensibilité.} Deux environnements calibrés (\(\sigma = 0{,}3\) et
\(\sigma = 0{,}6\)) montrent que \(\sigma\) change peu les effets sur le
PIB (au plus 0,05 point à 20 ans pour les stratégies de test) mais
modifie le contenu en emplois et, par là, la dette (jusqu'à 7 points de
PIB à 20 ans) (section \ref{sec:robustesse-env}).

\emph{Limites.} La productivité tendancielle est exogène ; les effets de
l'éducation, de l'innovation ou des réformes structurelles sur
l'efficience ne sont pas représentés (sauf capital public et R\&D,
exploratoires).

\textbf{Écart de production.}
\(\tilde y_t = 100 \times \ln(Y_t / Y^*_t)\) (en \%), où \(Y\) est le
PIB et \(Y^*\) le PIB potentiel.

\textbf{Fermeture d'offre.} Un supplément de demande \(S_t\) ramène la
production vers le potentiel :

\[S_t = S_{t-1} - \rho\,\tilde y_{t-1} + \theta\,\Delta y^*_t \cdot (\text{révision}),\]

avec \(\rho = 0{,}2\) {[}B{]} et \(\theta = 0{,}75\) {[}C{]}. Il
représente les canaux de retour non modélisés (politique monétaire avant
M24, crédit, confiance). C'est l'équation la plus influente pour le
passage du court au long terme (section \ref{sec:fermeture}).

\hypertarget{sec:eq-demande}{%
\section{Demande agrégée et multiplicateurs}\label{sec:eq-demande}}

La demande finale en volume est la somme de la consommation des ménages
\(C\), de leur investissement en logement \(I^m\), de l'investissement
des branches \(I^b\), des autres investissements (administrations,
institutions sans but lucratif) \(I^a\), de la consommation publique
\(G\), des exportations \(X\) et du supplément de fermeture \(S\) ; la
production de chaque produit est obtenue par l'inverse de Leontief :

\[q = (I - A_d)^{-1}\,d_d,\]

où \(A_d\) est la matrice des coefficients techniques domestiques (64
produits, TES de l'Insee) et \(d_d\) la demande finale adressée à la
production intérieure (après déduction des importations, réparties selon
les contenus en importations de chaque composante). Le multiplicateur
n'est pas un paramètre : il résulte de l'ensemble du système
(propensions, contenus en importations, prélèvements, réactions des prix
et de la fermeture). Le chapitre \ref{sec:transmission} en donne les
valeurs simulées.

\hypertarget{sec:eq-consommation}{%
\section{Consommation et épargne}\label{sec:eq-consommation}}

\[\Delta c_t = 0{,}85\,\Delta r_t - 0{,}22\,\left(c_{t-1} - r_{t-1} - \mu_c\right) - \psi_u\,\Delta u_t - \psi_r\,\Delta i^{m,r}_t + \text{calage},\]

où \(c\) est la consommation des ménages en volume (log), \(r\) le
revenu disponible brut réel (déflaté par le prix de la consommation),
\(u\) le taux de chômage, \(i^{m,r}\) le taux réel des crédits à
l'habitat (M24), \(\mu_c\) la constante de long terme (propension
moyenne). Élasticité de court terme au revenu 0,85 (t = 4,2) et vitesse
de retour 0,22 {[}E{]} ; épargne de précaution \(\psi_u = 0{,}3\,\%\)
par point de chômage (fourchette 0,1-0,6) et effet du taux réel
\(\psi_r = 0{,}2\) (fourchette 0-0,5) {[}C{]}.

\emph{Interprétation.} À long terme, la consommation est proportionnelle
au revenu disponible réel (élasticité unitaire, taux d'épargne stable) ;
à court terme, 85 \% d'une variation de revenu passe dans la
consommation la première année. Ce comportement est intermédiaire entre
la fonction de consommation keynésienne et le revenu permanent.

\emph{Hétérogénéité.} Pour les mesures socio-fiscales, la variation de
revenu de chaque décile est pondérée par une propension marginale à
consommer propre (plus élevée en bas de la distribution) ; l'écart à la
propension moyenne entre comme ajustement de consommation (M14).

\emph{Limites et alternatives.} Pas d'effet de richesse (patrimoine
immobilier et financier), pas d'anticipation des impôts futurs (non
ricardien), pas d'épargne de précaution liée à l'incertitude politique.
Un modèle à agents hétérogènes (HANK) ferait dépendre les propensions de
la liquidité des ménages et de la nature du choc.

\textbf{Logement.} \(\Delta i^m_t\) suit une équation estimée sur le
revenu disponible réel {[}E{]} ; le taux du crédit à l'habitat n'y entre
pas (testé et non retenu, section \ref{sec:donnees}).

\textbf{Taux d'épargne.} \(s_t = 1 - P_C\,C_t / R_t\) en résulte (avec
la correction des revenus du patrimoine et de l'investissement des
ménages).

\hypertarget{sec:eq-investissement}{%
\section{Investissement}\label{sec:eq-investissement}}

\[\Delta i^b_{t} = 0{,}67\,\Delta q^b_{t} + 0{,}26\,\Delta q^b_{t-1} - \lambda_I\left(i^b_{t-1} - k^{*b}_{t-1}\right) + \theta_T\,\frac{\Delta \text{Épargne}^{S11}_{t}}{I_{t-1}} + \text{calage},\]

\[k^{*}_t = q_t - \sigma_K \ln c_{K,t}, \qquad c_{K,t} = \frac{p_I}{p_Q}\,\frac{\left(i^{e,r}_t + \delta\right)\left(1 - z\,\tau^{IS}_t\right)}{1 - \tau^{IS}_t},\]

où \(i^b\) est l'investissement de la branche (log), \(q^b\) sa
production, \(k^*\) le capital désiré, \(c_K\) le coût d'usage du
capital, \(i^{e,r}\) le taux réel des crédits aux sociétés (M24),
\(\delta\) la dépréciation, \(\tau^{IS}\) le taux de l'impôt sur les
sociétés (semi-élasticité des bénéfices au taux pour les groupes : 0,8
{[}B{]}). Accélérateur estimé (0,67 ; 0,26) {[}E{]} ; élasticité au coût
d'usage −0,18 pour l'agrégat des sociétés non financières {[}E,
fragile{]} ; canal de trésorerie \(\theta_T = 0{,}10\) {[}C{]}.

\emph{Limites.} L'investissement des sociétés est l'équation la moins
bien validée (erreur de niveau de 11 \% en 2024 dans le test hors
échantillon) ; l'incertitude politique et réglementaire n'y figure pas.

\hypertarget{sec:eq-travail}{%
\section{Emploi, chômage, salaires et offre de
travail}\label{sec:eq-travail}}

\textbf{Demande de travail (branches marchandes).} À production donnée,
la condition du premier ordre de la CES donne

\[n^b_t = q^b_t - \sigma \ln\left(\frac{c_{h,t}}{p_{va,t}}\right) - (1-\sigma)\left(\ln E_t + \ln d_t\right),\]

où \(n\) est l'emploi (heures, puis personnes via la durée \(d\)),
\(c_h\) le coût horaire du travail (rémunérations moins subventions sur
salaires) et \(p_{va}\) le prix de la valeur ajoutée. La dynamique de
court terme (élasticité à la production 0,26, retour de 21 \% par an
dans l'équation estimée) lisse l'ajustement {[}E{]}. \textbf{Emploi non
marchand :} volume de valeur ajoutée divisé par une productivité
tendancielle propre.

\textbf{Population active.} Somme des actifs de 15 à 74 ans : taux
d'activité par âge et sexe observés jusqu'à 54 ans, modèle des seniors
au-delà (M9, section \ref{sec:eq-retraites}), effets de participation
des réformes des prestations (M15 : élasticité de participation 0,2,
d'heures 0,1 {[}B{]}).

\textbf{Chômage.} \(u_t = 1 - N_t / PA_t\) ; le chômage structurel
\(u^*\) est exogène (7,5 \% puis 7,0 \%) sauf effets de l'assurance
chômage (\(+0{,}063\) point par 10 \% d'allocations {[}B{]}) et de la
formation {[}X{]}.

\textbf{Salaires (secteur marchand, par tête).}

\[\Delta w_t = 0{,}39\,\Delta p^c_t + \text{(productivité)} - 0{,}09\left(w_{t-1} - p_{va,t-1} - \text{prod}_{t-1} - \mu_w\right) - \phi\,(u_t - u^*_t) + \text{calage},\]

où \(w\) est le salaire moyen par tête, \(p^c\) le prix de la
consommation, la cible de long terme étant une part des salaires dans la
valeur ajoutée stable (retour de 9 \% par an) {[}E{]} ;
\(\phi = 0{,}004\) par point d'écart de chômage {[}C, non identifié en
annuel et en trimestriel{]}, multiplié par 0,5 quand le chômage est
au-dessus de son niveau structurel (rigidité à la baisse) {[}C{]}. Le
SMIC suit sa règle légale (inflation des 20 \% de ménages les plus
modestes et moitié du gain de pouvoir d'achat du salaire horaire de base
ouvrier et employé), avec diffusion de 11 \% de ses hausses à la masse
salariale {[}H{]}.

\emph{Alternatives.} Une courbe de Phillips plus pentue (ancrage faible
: \(\times 2\)) ou une indexation plus rapide modifient surtout les prix
et la compétitivité ; elles sont testées en régime de robustesse.

\hypertarget{sec:eq-prix}{%
\section{Prix et inflation}\label{sec:eq-prix}}

\textbf{Prix de la valeur ajoutée (frontière des prix des facteurs,
P4).}

\[p^{*}_{va,b,t} = \theta_b \ln \text{CSU}_{b,t} + (1-\theta_b)\,\ln\left(\frac{p_{I,t}\,K_{b,t-1}}{Q_{b,t}}\right),\]

\[\Delta p_{va,b,t} = a_1\,\Delta \ln \text{CSU}_{b,t} + a_2\,\Delta q_{b,t} - \lambda_p\left(p_{va,b,t-1} - p^*_{va,b,t-1}\right) + \kappa\,\tilde y_t\,\mathbb{1}_{\text{hors industrie}},\]

où CSU est le coût salarial unitaire, \(\theta_b\) la part du travail de
la branche (médiane 1997-2019) {[}D{]}, et \(\kappa = 0{,}10\) la pente
sur l'écart de production hors industrie {[}C{]}. Les coefficients de
court terme et la vitesse de retour sont estimés en panel {[}E{]}. Une
variante (équation v0.9, marge vers sa norme, transmission du CSU 0,67)
dispose de son propre environnement calibré.

\textbf{Prix de production et de consommation (Leontief).} Pour chaque
produit \(j\),

\[p_j = \sum_i a_{ij}^d\,p_i + \sum_i a_{ij}^m\,p^m_i + v_j\,p_{va,j} + \tau_j,\]

où \(a^d\) et \(a^m\) sont les coefficients techniques domestiques et
importés, \(p^m\) les prix d'importation (transmission partielle des
prix mondiaux, la moitié l'année même {[}E{]}), \(v\) la part de la
valeur ajoutée et \(\tau\) les impôts nets sur les produits. Le prix de
la consommation est la moyenne pondérée par la structure de la
consommation.

\emph{Limites.} Pas de distinction des taux de TVA par produit
énergétique ; pas de substitution entre énergies ; les grands chocs de
prix (2022-2023) sont mal reproduits.

\hypertarget{sec:eq-productivite}{%
\section{Productivité}\label{sec:eq-productivite}}

La productivité tendancielle de chaque branche (croissance moyenne de la
productivité par tête sur les huit dernières années observées hors
2020-2021) est un paramètre {[}D/H{]} ; la convention de dépense de la
référence retient 0,7 \% par an pour la croissance potentielle
conventionnelle. Deux canaux endogènes : capital public (élasticité
0,08, capital public net de 55 \% du PIB, dépréciation 4 \% par an)
{[}C{]} et R\&D (additionnalité 0,7 €/€, élasticité au stock 0,08)
{[}X{]}. La durée du travail modifie la productivité horaire (70 \%
d'une baisse de durée absorbée sans embauche, calé sur l'ordre de
grandeur des 35 heures {[}H{]} ; appliqué symétriquement aux hausses).

\hypertarget{sec:eq-fiscalite}{%
\section{Fiscalité et transferts}\label{sec:eq-fiscalite}}

Les recettes publiques sont calculées par \textbf{taux apparents}
appliqués à des assiettes simulées :

\begin{itemize}
\tightlist
\item
  impôts sur les produits (TVA et accises) : taux apparent ×
  consommation des ménages (et autres emplois taxés) ; variation de taux
  répercutée à 80 \% sur les prix (95 \% pour les hausses et 50 \% pour
  les baisses du taux normal, d'après Benzarti et al., 2020) {[}B{]} ;
\item
  impôts des ménages (impôt sur le revenu, CSG, autres) : taux apparent
  × revenus ; mesures en points de PIB (\texttt{impots\_menages\_pib})
  ou par microsimulation ;
\item
  cotisations sociales : taux apparents × masse salariale ; transmission
  des cotisations employeurs aux salaires 0,5 à long terme {[}C{]} ;
\item
  impôt sur les sociétés : compte des sociétés (section
  \ref{sec:eq-societes}) ;
\item
  impôts de production : répercussion de 50 \% sur les prix de la valeur
  ajoutée, incidence de 20 \% sur les salaires {[}D, hypothèse{]} ;
\item
  autres recettes proportionnelles au PIB nominal.
\end{itemize}

Les \textbf{prestations} (hors pensions) sont indexées sur les prix de
l'année précédente et, dans la convention « part constante », croissent
avec le PIB potentiel conventionnel ; les allocations chômage sont
proportionnelles au nombre de chômeurs ; les pensions de droit direct
viennent de M9.

\hypertarget{sec:eq-depenses}{%
\section{Dépenses publiques}\label{sec:eq-depenses}}

\[\text{Dép}_t = \underbrace{W^{pub}_t\,N^{pub}_t}_{\text{rémunérations}} + \underbrace{P_{G,t}\,G^{autres}_t}_{\text{consommation hors rémunérations}} + \underbrace{\text{Prest}_t}_{\text{prestations}} + \underbrace{\text{Int}_t}_{\text{intérêts}} + \underbrace{\omega_t\,Y^{nom}_t}_{\text{autres dépenses}} + \text{calage},\]

où \(W^{pub}\) est le salaire public (instrument : point d'indice),
\(N^{pub}\) l'emploi public (non marchand), \(G^{autres}\) la
consommation publique hors rémunérations en volume, \(\omega\) la part
du PIB des autres dépenses (subventions, investissement public, autres
transferts). Convention de référence « part constante » : consommation
publique et prestations hors pensions croissent comme le PIB potentiel
conventionnel (population active × (1 − u*) × productivité de 0,7 \%)
{[}H, décision du projet{]}.

\hypertarget{sec:eq-retraites}{%
\section{Retraites}\label{sec:eq-retraites}}

\textbf{Âge d'ouverture des droits (AOD)} par génération \(g\) :
législation actuelle (loi de 2023 aménagée par la LFSS pour 2026 : 62
ans et 9 mois pour les générations 1963-1964, puis +3 mois par
génération jusqu'à 64 ans pour 1969) {[}D{]}, ou réforme « cible à
partir de l'année de début, progressive ou immédiate ».

\textbf{Durée d'assurance requise} \(D_g\) (trimestres) : 170 pour
1963-1964, 171 pour 1965, 172 à partir de 1966 en législation actuelle
{[}D{]} ; calendrier de la loi de 2014 (un trimestre toutes les trois
générations jusqu'à 172 pour 1973) pour l'abrogation de la réforme de
2023 ; cible d'une réforme (immédiate ou un trimestre par génération).

\textbf{Âge pivot et activité des seniors.} Pour le sexe \(s\) et la
génération \(g\),

\[\pi_{g,s} = \beta_s\,\text{AOD}_g + (1 - \beta_s)\cdot 62 + \varepsilon_s\,\frac{D_g - D^{act}_g}{4}, \qquad
a(x,t) = A_s(t)\cdot h\big(x + 0{,}5 - \pi_{g,s}\big),\]

\[h(d) = \begin{cases} 1 - e(t)\,\min\left(1, \frac{d+5}{5}\right)^+ & d < 0 \\ s\,e^{-k d}\,(1 - e(t)) & d \ge 0 \end{cases}\]

où \(a(x,t)\) est le taux d'activité à l'âge \(x\) relatif à celui des
50-54 ans \(A_s(t)\), \(e(t)\) la part des sorties précoces (tendance
arrêtée en 2025), \(s\) et \(k\) la marche et la décroissance après
l'âge pivot, \(\beta_s\) la part du profil qui se décale avec l'AOD ---
tous estimés par moindres carrés non linéaires sur les taux d'activité
par tranche quinquennale de 2003 à 2019 {[}E{]}.
\(\varepsilon_H = 0{,}75\) et \(\varepsilon_F = 5/12\) (report de 9 et 5
mois par année de durée supplémentaire, Bozio, 2011) {[}E, littérature,
appliqué symétriquement aux baisses : H{]}.

\textbf{Groupes particuliers.} Pour une fraction \(p\) de chaque
génération (carrières longues à l'âge \(a_{CL}\), ou maintien des
conditions actuelles), le profil devient, en écart à la législation
actuelle,

\[\tilde h = h(d_{reg}) + p\left[h(d_{grp}) - h(d_{reg})\right] - p\left[h(d^{act}_{grp}) - h(d^{act})\right],\]

ce qui reproduit exactement le profil estimé en législation actuelle.
L'âge actuel du groupe « carrière longue » est approché par
\(\max(60, \text{AOD} - 2)\) {[}H{]} ; \(p = 0{,}17\) pour les carrières
longues (part des départs anticipés pour carrière longue au régime
général en 2024, CNAV) {[}D, approchant l'éligibilité{]}.

\textbf{Pensions.} Pension moyenne de la génération à la liquidation
\(= \nu \times\) moyenne des 25 derniers salaires moyens revalorisés sur
les prix, puis revalorisée sur les prix ; masse \(= \sum_{x,s}\)
retraités \(\times\) pension. Le facteur de durée

\[f_{g,s} = 1 - \phi_{base}\,(1 - \varepsilon_s)\,\Delta D_g \left(\frac{4}{D^{act}_g} + 0{,}05\right)\]

applique la proratisation et la décote (1,25 \% par trimestre {[}D{]}) à
la part de base de la pension (\(\phi_{base} = 0{,}5\) {[}H{]}) pour la
fraction des assurés qui ne reportent pas leur départ. L'écart de masse
dû à une réforme d'âge est multiplié par 0,55 (calage sur les chiffrages
de l'étude d'impact de 2023 et de l'Institut Montaigne : les départs
déplacés concernent surtout des carrières incomplètes) {[}C{]}.

\emph{Limites.} Pas de distinction des régimes, de la réversion, de la
surcote, des minima de pension ; pas de piliers de capitalisation ; les
durées validées individuelles ne sont pas simulées (effet moyen par
génération et sexe).

\hypertarget{sec:eq-dette}{%
\section{Dette et intérêts}\label{sec:eq-dette}}

\[D_t = D_{t-1} - \text{Solde}_t + \text{AF}_t, \qquad \text{Int}_t = \sum_{c} r_c\,B_{c,t}, \]

où la dette \(D\) (au sens de Maastricht) évolue avec le besoin de
financement et des ajustements stock-flux \(\text{AF}\) observés
jusqu'en 2025 (nuls ensuite) ; la charge d'intérêts est la somme, sur
les cohortes d'émission \(c\), du taux à l'émission \(r_c\) (courbe des
taux) appliqué à l'encours restant \(B_{c,t}\) (module M17, validé sur
les données de l'Agence France Trésor) {[}D/E{]}. Le \textbf{taux
apparent} rejoint ainsi progressivement le taux de marché, au rythme du
renouvellement de la dette (durée moyenne de l'ordre de 8 ans). La
boucle intra-annuelle intègre le besoin de financement de l'année.

\hypertarget{sec:eq-monetaire}{%
\section{Politique monétaire}\label{sec:eq-monetaire}}

En écart à la référence (années \(\ge 2027\), information de \(t-1\)) :

\[\Delta i_t = 0{,}728\,\Delta i_{t-1} + 0{,}393\,\Delta\pi^{ZE}_{t-1} + 0{,}302\,\Delta y^{ZE}_{t-1},\]

\[\Delta\pi^{ZE} = \omega_{FR}\,\Delta\pi^{FR} + (1-\omega_{FR})\,\text{choc}^\pi, \qquad \Delta y^{ZE} = \omega_{FR}\,\Delta y^{FR} + (1-\omega_{FR})\,\text{choc}^y,\]

où \(i\) est le taux directeur, \(\pi\) l'inflation, \(y\) l'écart
d'activité (écart du PIB à la référence, ou loi d'Okun sur le chômage en
variante) et \(\omega_{FR}\) le poids de la France {[}E : règle estimée
sous forme contemporaine sur des données annuelles de la zone euro,
appliquée avec un an de décalage car la forme retardée estimée avait un
signe erroné{]}. Variante : anticipations rationnelles du taux allemand
sous la règle.

\emph{Lecture.} Avec un poids de 20 \% de la France, une hausse de 1
point de l'inflation française n'induit qu'environ 0,08 point de hausse
du taux directeur la première année. La politique monétaire commune
amortit peu les chocs nationaux, ce qui rapproche la France d'une
économie en change fixe pour son budget.

\hypertarget{sec:eq-credit}{%
\section{Taux de crédit}\label{sec:eq-credit}}

\[\Delta i^{m}_t = (1 - v_m)\,\Delta i^m_{t-1} + b^c_m\,\Delta i_t + b^l_m\,\Delta i^{OAT}_t, \qquad
\Delta i^{e}_t = (1 - v_e)\,\Delta i^e_{t-1} + b^c_e\,\Delta i_t + b^l_e\,\Delta i^{OAT}_t,\]

(en écart à la référence), avec pour l'habitat une vitesse
\(v_m = 0{,}43\) et une transmission de long terme de 0,90, et pour les
sociétés non financières \(v_e = 0{,}93\) et 0,78 {[}E, séries MIR de la
BCE{]}. Le crédit aux sociétés, majoritairement à taux variable ou
court, s'ajuste presque immédiatement ; le crédit à l'habitat, à taux
fixe, plus lentement.

\hypertarget{sec:eq-prime}{%
\section{Prime de risque souveraine}\label{sec:eq-prime}}

\[\Delta(\text{OAT} - \text{Bund})_t = \gamma_d\,\Delta d_{t-1} + \gamma_s\,\Delta s_{t-1} + \text{chocs}_t,\]

où \(\Delta d\) et \(\Delta s\) sont les écarts du ratio de dette et du
solde public à la référence ; \(\gamma_d = 0{,}013\) et
\(\gamma_s = -0{,}0148\) (points de taux par point de PIB) hors crise,
\(\gamma_d = 0{,}012\) et \(\gamma_s = -0{,}118\) en incluant 2010-2013
{[}E, panel de la zone euro, effets fixes par pays{]}. Le coefficient de
solde est fragile. Variante en niveau au-delà d'un seuil de dette
{[}H{]}.

\emph{Littérature.} Les études empiriques trouvent en général 2 à 5
points de base par point de dette en moyenne, avec de fortes
non-linéarités et une dépendance au contexte (liquidité, appartenance à
l'union monétaire, programme de la BCE).

\hypertarget{sec:eq-exterieur}{%
\section{Économie extérieure}\label{sec:eq-exterieur}}

\[\Delta x_t = 0{,}71\,\Delta d^*_t - 0{,}60\,\Delta\ln\left(\frac{p_X}{p^*}\right)_t + \text{calage},\]
\[\Delta m_t = 1{,}47\,\Delta \bar d_t - \lambda_M\left(m_{t-1} - 1{,}37\,\bar d_{t-1} - 0{,}006\,t + 0{,}16\,\ln\frac{p_M}{p}_{t-1}\right),\]

où \(x\) et \(m\) sont les exportations et importations en volume,
\(d^*\) la demande mondiale adressée, \(\bar d\) la demande intérieure
pondérée par les contenus en importations, \(p_X/p^*\) la
compétitivité-prix {[}E{]}. Le revenu national brut est le PIB augmenté
des revenus primaires nets du reste du monde (indexés sur le PIB
nominal, diminués de l'écart des dividendes versés aux non-résidents, 50
\% des dividendes des sociétés).

\hypertarget{sec:eq-societes}{%
\section{Compte des sociétés non financières}\label{sec:eq-societes}}

\[\text{Int}^{S11}_t = \rho_t\,D^{S11}_{t-1}, \qquad \rho_t = \rho_{t-1} + \mu\left(\rho^*_t - \rho_{t-1}\right), \qquad \rho^*_t = c_\rho + 1{,}29\left(0{,}77\,i^{c}_t + 0{,}23\,i^{l}_t\right),\]
\[\text{IS}_t = \tau^{app}\,\text{RCAI}_t\cdot f^{fisc}_t, \qquad \text{Div}_t = 0{,}364\,\Pi_{t-1}, \qquad
D^{S11}_t = D^{S11,base}_t + 1{,}06\sum_{k \le t}\left(\text{BF}_k - \text{BF}^{base}_k\right),\]

où \(\rho\) est le taux apparent de la dette des sociétés, qui rejoint à
la vitesse \(\mu = 0{,}60\) {[}E, 1999-2024{]} un taux de marché
\(\rho^*\) combinant taux court (poids 0,77) et long, avec une
transmission de 1,29 et un écart moyen de 3,4 points {[}E{]}, RCAI le
résultat courant avant impôt, \(\Pi\) le résultat après intérêts et
impôt, BF le besoin de financement. Les intérêts reçus valent 0,77 fois
les intérêts versés {[}C{]} ; les dividendes sont répartis entre ménages
(42 \%), reste du monde (50 \%) et administrations (8 \%) {[}D{]}.
L'équation de Lintner (lissage des dividendes) n'est pas identifiée : la
règle proportionnelle est retenue, avec un paramètre de lissage en
variante {[}H{]}.

\hypertarget{sec:eq-distribution}{%
\section{Distribution des revenus}\label{sec:eq-distribution}}

Pour une mesure socio-fiscale, la variation du revenu disponible du
ménage \(h\) est \(\Delta R_h = -\Delta T_h + \Delta B_h\) (impôts
\(T\), prestations \(B\)), calculée par OpenFisca sur les revenus de
2024 ; elle est agrégée par décile de niveau de vie \(q\) (revenu
disponible par unité de consommation, échelle de l'OCDE modifiée) :

\[\Delta R_q = \sum_{h \in q} w_h\,\Delta R_h, \qquad \Delta C^{M14} = \sum_q (\text{pmc}_q - \overline{\text{pmc}})\,\Delta R_q,\]

où \(w_h\) est le poids de sondage et pmc la propension marginale à
consommer du décile. Les indicateurs d'inégalité (rapports interdéciles,
Gini) sont calculés avant et après réforme. Les effets macroéconomiques
induits (emploi, salaires, prix) ne sont pas redistribués ménage par
ménage dans les résultats de cette étude.

\hypertarget{sec:donnees}{%
\chapter{Données, estimations, calibrations et
validation}\label{sec:donnees}}

Un modèle macroéconométrique vaut ce que valent ses données et la façon
dont ses paramètres ont été obtenus. Ce chapitre présente les sources,
la manière dont chaque paramètre a été fixé (estimation, calibration,
littérature, hypothèse), les méthodes d'estimation, les tests réalisés
et, surtout, ce que ces tests prouvent et ne prouvent pas.

\hypertarget{sec:sources}{%
\section{Sources statistiques}\label{sec:sources}}

\textbf{Principe : données ouvertes uniquement.} Toutes les données
utilisées sont publiques, accessibles sans habilitation, téléchargées
par script et identifiées par une empreinte numérique (SHA-256)
consignée dans un registre. Ce choix garantit la reproductibilité mais
exclut les données administratives individuelles (déclarations fiscales
individuelles, échantillons interrégimes de retraités, fichiers
d'entreprises), ce qui limite la microsimulation et l'estimation de
certains comportements.

\begin{longtable}[]{@{}
  >{\raggedright\arraybackslash}p{(\columnwidth - 8\tabcolsep) * \real{0.2000}}
  >{\raggedright\arraybackslash}p{(\columnwidth - 8\tabcolsep) * \real{0.2000}}
  >{\raggedright\arraybackslash}p{(\columnwidth - 8\tabcolsep) * \real{0.2000}}
  >{\raggedright\arraybackslash}p{(\columnwidth - 8\tabcolsep) * \real{0.2000}}
  >{\raggedright\arraybackslash}p{(\columnwidth - 8\tabcolsep) * \real{0.2000}}@{}}
\toprule\noalign{}
\begin{minipage}[b]{\linewidth}\raggedright
Domaine
\end{minipage} & \begin{minipage}[b]{\linewidth}\raggedright
Source et jeu de données
\end{minipage} & \begin{minipage}[b]{\linewidth}\raggedright
Période, fréquence
\end{minipage} & \begin{minipage}[b]{\linewidth}\raggedright
Usage
\end{minipage} & \begin{minipage}[b]{\linewidth}\raggedright
Qualité, limites
\end{minipage} \\
\midrule\noalign{}
\endhead
\bottomrule\noalign{}
\endlastfoot
Comptes nationaux par branche & Eurostat \texttt{nama\_10\_a64},
\texttt{nama\_10\_a64\_e}, \texttt{nama\_10\_a64\_p5} (comptes de
l'Insee) & 1995-2025, annuel & équations d'emploi, de prix,
d'investissement & comptes semi-définitifs pour 2023-2025, révisables \\
Comptes des secteurs & Eurostat \texttt{nasa\_10\_nf\_tr},
\texttt{nasa\_10\_f\_bs}, \texttt{nasa\_10\_f\_tr} & 1995-2025, annuel &
ménages, sociétés, administrations ; bilans & sociétés financières peu
détaillées \\
Finances publiques & Eurostat \texttt{gov\_10a\_main},
\texttt{gov\_10dd\_edpt1} (dette de Maastricht) & 1995-2025 & recettes,
dépenses, dette & recoupé avec l'Insee \\
Tableau entrées-sorties & Insee, TES symétrique 64 produits & année de
base & prix de Leontief, contenus en importations & structure figée sur
l'année de base \\
Chômage & Insee, BDM, série 001688526 (BIT, trimestriel) & 1975-2026 &
Phillips, Okun & --- \\
Population & Insee, projections 2026 (29 scénarios), recensement 2022
(fichier détail) & 1962-2070 & démographie, population de
microsimulation & projections conditionnelles \\
Activité par âge & Eurostat \texttt{lfsa\_argan} (enquête Emploi) &
2003-2025 & activité des seniors (M9) & échantillon d'enquête \\
Retraites & Eurostat \texttt{spr\_pns\_ben}, \texttt{spr\_exp\_pens} ;
LFSS 2026 (durées) ; DREES, CNAV (carrières longues) & 2000-2024 & M9 &
pas de données individuelles \\
Impôt sur le revenu & DGFiP, IRCOM 2025 (revenus 2024), déclarations
2042 nationales, déciles & 2006-2024 & population M4, calage & agrégats
par tranche \\
Taux d'intérêt & Eurostat \texttt{irt\_lt\_mcby\_a/m},
\texttt{irt\_st\_m} ; BCE, séries MIR (crédits nouveaux) & 1980-2026 &
M24 & --- \\
Prévisions & FMI, World Economic Outlook (millésimes) & 2010-2031 &
référence, validation ex ante & --- \\
Législation & OpenFisca-France 177.0.0 ; textes (CGI, CSS, LFSS 2026) &
2026 & microsimulation, règles & moteur communautaire \\
\end{longtable}

Données indisponibles et substitutions principales : potentiel en temps
réel (reconstruit par filtrage), patrimoine des ménages par ménage
(approché par les revenus du patrimoine), héritages reçus (assiette
pondérée par l'âge), transitions individuelles emploi-chômage
(probabilités moyennes), durées validées individuelles (effets moyens
par génération dans M9).

\hypertarget{sec:typologie}{%
\section{Typologie des paramètres}\label{sec:typologie}}

Le registre complet figure en annexe \ref{sec:annexe-a}. Le tableau
suivant classe les paramètres les plus influents.

\begin{longtable}[]{@{}
  >{\raggedright\arraybackslash}p{(\columnwidth - 6\tabcolsep) * \real{0.2500}}
  >{\raggedright\arraybackslash}p{(\columnwidth - 6\tabcolsep) * \real{0.2500}}
  >{\raggedright\arraybackslash}p{(\columnwidth - 6\tabcolsep) * \real{0.2500}}
  >{\raggedright\arraybackslash}p{(\columnwidth - 6\tabcolsep) * \real{0.2500}}@{}}
\toprule\noalign{}
\begin{minipage}[b]{\linewidth}\raggedright
Catégorie
\end{minipage} & \begin{minipage}[b]{\linewidth}\raggedright
Exemples
\end{minipage} & \begin{minipage}[b]{\linewidth}\raggedright
Nombre approximatif
\end{minipage} & \begin{minipage}[b]{\linewidth}\raggedright
Conséquence
\end{minipage} \\
\midrule\noalign{}
\endhead
\bottomrule\noalign{}
\endlastfoot
\textbf{Estimés économétriquement} {[}E{]} & accélérateur (0,67),
élasticité de la consommation au revenu (0,85) et vitesse de retour
(0,22), indexation des salaires (0,39), élasticités du commerce, règle
de la BCE, transmission aux crédits, prime souveraine, activité des
seniors, vitesse du taux apparent des sociétés (0,60) & plusieurs
dizaines (dont 25 effets fixes de branche) & intervalles d'incertitude
disponibles \\
\textbf{Calibrés sur la littérature} {[}C{]} & \(\sigma = 0{,}45\),
pente de Phillips (−0,004), pente des prix (0,10), épargne de précaution
(0,3), effet du taux réel (0,2), incidences fiscales, élasticités des
hauts revenus, capital public (0,08) & une trentaine & variantes basse
et haute systématiques \\
\textbf{Issus de la littérature, appliqués tels quels} {[}E,
littérature{]} & report de l'âge de départ par année de durée requise (9
mois / 5 mois, Bozio 2011) ; répercussion asymétrique de la TVA
(Benzarti et al., 2020) & quelques-uns & dépendent de la
transposabilité \\
\textbf{Hypothèses de scénario} {[}H{]} & règle budgétaire (α, β, d*),
convention de dépense, OAT de référence, durée des contentieux
européens, compositions des enveloppes non ventilées, additionnalité de
l'investissement hors budget & une vingtaine & présentées en
variantes \\
\textbf{Fragiles ou insuffisamment identifiés} & \(\sigma\), Phillips,
fermeture (ρ), coefficient de solde de la prime souveraine, élasticité
de l'investissement au coût d'usage (−0,18), effet de l'assurance
chômage sur u* & --- & priorité des analyses de sensibilité \\
\end{longtable}

\hypertarget{sec:methodes}{%
\section{Méthodes d'estimation}\label{sec:methodes}}

\textbf{Équations de comportement (M7).} Modèles à correction d'erreur
estimés en une étape (relation de long terme et dynamique de court terme
ensemble), sur 1997-2019 hors années de crise sanitaire. Les équations
de branche sont estimées en panel de 25 branches avec effets fixes, par
moindres carrés pondérés par la taille de la branche, avec des écarts
types groupés par année ; les équations agrégées par moindres carrés
ordinaires avec écarts types robustes à l'autocorrélation (Newey-West).
Lorsque la théorie l'impose et que la version libre est biaisée par la
colinéarité, des contraintes sont imposées (homogénéité de long terme
des prix, rendements constants de l'emploi), et le résultat libre est
publié.

\textbf{Pourquoi 1997-2019 ?} L'échantillon commence avec les séries
homogènes des comptes nationaux et s'arrête avant la crise sanitaire,
dont les mouvements (activité partielle, épargne forcée, aides massives)
ne relèvent pas des comportements ordinaires. Les équations de taux
d'intérêt font exception (estimées jusqu'en 2024, seul épisode de hausse
des taux de l'échantillon). Conséquence : le modèle décrit des
comportements d'une période de taux en baisse et d'inflation faible ;
leur stabilité dans un environnement de taux élevés n'est pas garantie
(critique de Lucas).

\textbf{Fermeture (ρ, θ).} Inférence indirecte : on choisit les valeurs
qui reproduisent la durée moyenne des écarts de production observée dans
les données AMECO et FMI (1982-2019), et la réaction de la consommation
à un choc permanent (Blundell, Pistaferri et Preston, 2008).

\textbf{Activité des seniors (M9).} Moindres carrés non linéaires sur
les taux d'activité par tranche quinquennale et par sexe (2003-2019),
pondérés par la population ; test hors échantillon sur 2020-2025.

\textbf{Politique monétaire (M24).} Règle de la BCE et transmission au
taux allemand : moindres carrés sur données annuelles de la zone euro ;
crédits : ajustement partiel sur les séries MIR ; prime souveraine :
panel de pays de la zone euro à effets fixes, avec et sans 2010-2013.

\textbf{Paramètres non identifiés.} Quand une estimation libre est non
significative, de signe opposé à la théorie, ou instable selon
l'échantillon, le paramètre est calibré au milieu de la littérature et
reçoit des variantes. Règle de méthode appliquée dans tout le projet :
\textbf{aucun paramètre n'est choisi pour améliorer une simulation
particulière}, et une spécification n'est pas retenue sur le seul
critère de l'erreur quadratique.

\hypertarget{sec:calibration}{%
\section{Calibration de la référence}\label{sec:calibration}}

La trajectoire de référence à politique inchangée est construite en
trois temps : (i) les constantes de branche sont recalées sur les
résidus moyens des dernières années observées (2023-2025) au lieu
d'extrapoler les effets fixes estimés ; (ii) les exogènes reproduisent
les prévisions du FMI jusqu'en 2031 (croissance, inflation, demande
mondiale) ; (iii) au-delà, les hypothèses de long terme (productivité,
chômage structurel, démographie centrale de l'Insee) sont prolongées. La
convention de dépense publique « part constante » fait croître la
consommation publique et les prestations hors pensions comme le PIB
potentiel conventionnel.

\textbf{Une référence explosive.} Avec l'OAT à 4,89 \% (dernier point
mensuel observé) maintenue, la référence voit le taux apparent de la
dette passer de 2,0 \% en 2025 à 4,0 \% en 2032 et 4,6 \% en 2047, alors
que la croissance nominale ralentit (de 2,9 \% à 1,8 \%) : le
différentiel taux-croissance devient positif et la dette atteint 139 \%
du PIB en 2032 et 262 \% en 2047 sans règle budgétaire (R0), le déficit
passant de 5,1 \% à 8,2 \% puis 16,3 \% du PIB. La charge d'intérêts
explique l'essentiel de la dérive (de 2,3 à 5,3 points de PIB en 2032 et
11,4 points en 2047) ; le solde primaire se dégrade beaucoup moins (de
−2,8 à −2,9 points en 2032, −4,9 points en 2047). \textbf{Cette
référence n'est pas une prévision} : elle décrit ce que produirait
l'absence de toute correction, avec des taux élevés maintenus. Elle a
deux conséquences pour l'étude : les effets des programmes se lisent en
écart à cette référence et sous règle budgétaire ; et la lecture par
objectifs budgétaires (lecture B) mesure un effort \emph{par rapport à
cette dérive}, ce qui la rend sensible à l'hypothèse de taux. Une
variante avec l'OAT à sa moyenne des douze derniers mois (3,82 \%) est
présentée systématiquement.

\textbf{Revalidation de l'environnement.} Chaque environnement calibré
est mis en cache avec une empreinte de tous ses paramètres (coefficients
estimés, constantes des modules, paramètres globaux, cibles de
calibration, exogènes calibrées) et une trajectoire de contrôle
d'environ 70 variables. À chaque chargement, si l'empreinte diffère, le
modèle resimule la référence ; il n'accepte le cache que si toutes les
trajectoires sont identiques à \(10^{-9}\) près et si les paramètres
calibrés n'ont pas changé ; sinon il recalibre. Un test vérifie qu'un
coefficient modifié est refusé.

\hypertarget{sec:validation}{%
\section{Validation : rétrosimulations et
comparaisons}\label{sec:validation}}

La validation d'un modèle macroéconomique se fait à plusieurs niveaux,
du plus faible au plus exigeant.

\textbf{Niveau 1 --- équation par équation.} Chaque équation est simulée
dynamiquement hors échantillon (estimation 1997-2012, simulation
2013-2019 ; estimation 1997-2019, simulation 2022-2024). Les équations
battent une prévision naïve (tendance prolongée) pour la consommation,
les salaires, les importations, les exportations et l'emploi (2022-2024)
; elles font moins bien pour les prix de la valeur ajoutée (2013-2019)
et pour l'investissement des sociétés (erreur de niveau de +11 \% en
2024).

\textbf{Niveau 2 --- rétrosimulation conditionnelle du modèle complet.}
Le modèle est simulé sur une période passée avec les exogènes observées.
Écarts en dernière année (simulé − observé), configuration P4 :

\begin{longtable}[]{@{}
  >{\raggedright\arraybackslash}p{(\columnwidth - 10\tabcolsep) * \real{0.1667}}
  >{\raggedright\arraybackslash}p{(\columnwidth - 10\tabcolsep) * \real{0.1667}}
  >{\raggedright\arraybackslash}p{(\columnwidth - 10\tabcolsep) * \real{0.1667}}
  >{\raggedright\arraybackslash}p{(\columnwidth - 10\tabcolsep) * \real{0.1667}}
  >{\raggedright\arraybackslash}p{(\columnwidth - 10\tabcolsep) * \real{0.1667}}
  >{\raggedright\arraybackslash}p{(\columnwidth - 10\tabcolsep) * \real{0.1667}}@{}}
\toprule\noalign{}
\begin{minipage}[b]{\linewidth}\raggedright
Période
\end{minipage} & \begin{minipage}[b]{\linewidth}\raggedright
PIB (\%)
\end{minipage} & \begin{minipage}[b]{\linewidth}\raggedright
Prix à la consommation (\%)
\end{minipage} & \begin{minipage}[b]{\linewidth}\raggedright
Emploi (\%)
\end{minipage} & \begin{minipage}[b]{\linewidth}\raggedright
Chômage (pt)
\end{minipage} & \begin{minipage}[b]{\linewidth}\raggedright
Dette (pt de PIB)
\end{minipage} \\
\midrule\noalign{}
\endhead
\bottomrule\noalign{}
\endlastfoot
2010 → 2019 & +1,64 & +1,36 & +1,66 & −1,52 & −4,64 \\
2014 → 2019 & −0,42 & −1,53 & +0,27 & −0,25 & −1,77 \\
2019 → 2024 & +0,90 & −3,20 & −1,65 & +1,53 & −5,93 \\
\end{longtable}

Les erreurs de croissance annuelle sont faibles ; les erreurs de niveau
s'accumulent sur dix ans, et l'épisode inflationniste de 2022-2023, hors
du domaine d'estimation, est mal reproduit (prix sous-estimés de 3,2
\%).

\textbf{Niveau 3 --- prévision ex ante.} Le modèle est utilisé comme il
l'aurait été à l'époque, avec les seules informations disponibles
(millésimes du FMI des années 2010-2016). Erreur absolue moyenne sur la
croissance : 0,39 point à un an et 0,25 point à cinq ans, contre 0,59 et
0,34 pour le FMI du même millésime et 0,84 et 0,38 pour une prévision
naïve. Sur l'inflation, le modèle fait moins bien que le FMI (0,72
contre 0,43 à un an). Sur la dette, il présente un biais de +10,6 points
à cinq ans, dû à la convention de législation inchangée (le gouvernement
corrige toujours une partie de la dérive).

\textbf{Niveau 4 --- choc de politique identifié.} Le crédit d'impôt
pour la compétitivité et l'emploi (CICE, 2013-2018) est simulé
rétrospectivement (avec et sans, coût budgétaire compté) : +71 000
emplois marchands en 2016 et +86 000 en 2017, contre environ 100 000
dans les évaluations microéconométriques (2013-2016) et environ 160 000
dans une évaluation macro-sectorielle (2013-2017), selon la synthèse de
France Stratégie (2020). Le modèle est dans l'ordre de grandeur, plutôt
en dessous, ce qui est cohérent avec une transmission aux prix plus
forte que celle mesurée ex post.

\hypertarget{sec:tests}{%
\section{Ce que prouvent les tests automatisés --- et ce qu'ils ne
prouvent pas}\label{sec:tests}}

Le projet comprend des suites de tests automatisés (annexe
\ref{sec:annexe-f}) : 32 tests pour les phases 4.3 à 6, 6 tests pour la
finalisation (retraites, lecture par objectifs, corpus), 142 contrôles
d'intégration comptable, et les suites antérieures (versions 0.8 et
0.9).

\textbf{Ce qu'ils prouvent.}

\begin{itemize}
\tightlist
\item
  \textbf{Cohérence comptable} : les identités entre secteurs, entre
  flux et stocks, entre coûts et prix sont respectées à la précision
  numérique (écart relatif de l'ordre de \(10^{-16}\) sur l'identité
  dette-solde des simulations de programmes).
\item
  \textbf{Reproductibilité} : les versions antérieures (0.8, 0.9) et la
  référence sans règle sont reproduites exactement ; un environnement
  mis en cache n'est accepté que si ses paramètres et ses trajectoires
  sont inchangés.
\item
  \textbf{Absence de régression} : une modification du code n'altère pas
  les résultats des configurations qu'elle ne vise pas.
\item
  \textbf{Conformité des traductions} : les conversions entre montants
  annoncés et paramètres du modèle inversent exactement les rendements
  statiques ; les totaux publiés par LFI (106,4 Md€ et 81,7 Md€) sont
  recoupés ; les calendriers légaux de durée d'assurance sont respectés
  ; les variantes européennes agissent comme spécifié.
\item
  \textbf{Sens des effets} : les réformes des retraites ont l'effet
  attendu sur la population active (durée plus courte, départs anticipés
  → moins d'actifs ; âge et durée plus élevés → plus d'actifs).
\end{itemize}

\textbf{Ce qu'ils ne prouvent pas.}

\begin{itemize}
\tightlist
\item
  Ils ne prouvent pas que \textbf{les paramètres sont justes} : un
  modèle peut être parfaitement cohérent et se tromper de
  multiplicateur.
\item
  Ils ne prouvent pas que \textbf{les comportements estimés sur
  1997-2019} resteront valables dans un environnement de taux élevés,
  après des réformes de grande ampleur (critique de Lucas).
\item
  Ils ne prouvent pas que \textbf{la traduction d'un programme} reflète
  les intentions du candidat : elle reflète un texte publié, complété
  d'hypothèses explicitées.
\item
  Ils ne constituent pas une \textbf{validation scientifique} au sens
  d'une confrontation réussie à des expériences indépendantes : seuls
  les niveaux 2 à 4 ci-dessus s'en approchent, et leurs résultats sont
  mitigés (bonne performance sur la croissance, moins bonne sur les prix
  et l'investissement).
\end{itemize}

Cette distinction est essentielle : la cohérence et la reproductibilité
sont des conditions nécessaires, pas suffisantes, de la crédibilité des
résultats.

\hypertarget{sec:ecartes}{%
\section{Choix techniques écartés}\label{sec:ecartes}}

\begin{itemize}
\tightlist
\item
  \textbf{Pas trimestriel.} Examiné en phase 5 et non retenu : les
  mesures sont annuelles, les horizons de 5 ans et plus, et les
  équations trimestrielles disponibles ne couvrent que les salaires et
  les prix marchands ; leur généralisation aurait demandé une
  recalibration complète sans gain démontré pour la comparaison.
\item
  \textbf{Contenu importé selon le niveau de vie, tension des capacités
  sur les importations, taux d'intérêt dans l'équation de logement} :
  testés, non retenus (effets non significatifs ou oscillations non
  observées).
\item
  \textbf{Équation de Lintner pour les dividendes} : non identifiée ;
  règle proportionnelle retenue.
\item
  \textbf{Effet du taux légal de l'impôt sur les sociétés sur
  l'assiette} : élasticité historique non significative (−0,10 ± 0,29) ;
  le taux apparent observé est utilisé en rétrosimulation.
\end{itemize}

\hypertarget{sec:programmes}{%
\chapter{Présentation détaillée des programmes
politiques}\label{sec:programmes}}

Ce chapitre présente les trois programmes tels qu'ils sont publiés, puis
leur traduction dans le modèle. Deux règles le gouvernent. D'une part,
\textbf{les montants, calendriers et effets attendus sont ceux annoncés
par les candidats} ou, à défaut, des chiffrages de tiers explicitement
cités ; aucun montant n'est inventé. D'autre part, \textbf{les
hypothèses de traduction de l'auteur du modèle sont distinguées des
engagements des candidats} : elles sont marquées {[}H{]} et accompagnées
de variantes. Le corpus est daté du 24 novembre 2026 ; il est structuré
dans le module M27 et régénérable (annexe \ref{sec:annexe-c}).

\hypertarget{sec:methode-corpus}{%
\section{Méthode de constitution du corpus}\label{sec:methode-corpus}}

Pour chaque candidat, les sources primaires (programme, documents
budgétaires, conférences de presse publiées sur les sites des candidats)
ont été privilégiées ; les articles de presse (Public Sénat, LCP,
Euronews, Moneyvox, Sud Radio) servent à dater et à compléter les
annonces orales. Chaque mesure est décrite par : la source, le montant
annoncé (signe du point de vue du solde public : positif pour une
recette ou une économie), le statut du chiffrage (par le candidat, par
un tiers, aucun), le calendrier, les bénéficiaires, le canal du modèle,
la règle de conversion, le degré de quantification (\emph{oui},
\emph{partiel}, \emph{non}) et la qualification au regard du droit
européen (chapitre \ref{sec:europe}).

\textbf{Conversion des montants en paramètres.} Quand le paramètre du
modèle n'est pas exprimé en euros (taux du prélèvement forfaitaire
unique, taux moyen de l'impôt sur la fortune, variation des droits de
succession, taux de TVA, taux de l'impôt sur les sociétés), sa valeur
est obtenue en inversant le rendement \textbf{statique} du simulateur à
partir du montant annoncé. On simule ainsi la mesure \emph{telle que le
candidat la chiffre} ; le modèle applique ensuite ses propres réactions
(érosion d'assiette, bouclage macroéconomique). L'écart entre le
rendement annoncé et le rendement simulé est donc un résultat, et non
une hypothèse.

\begin{longtable}[]{@{}
  >{\raggedright\arraybackslash}p{(\columnwidth - 12\tabcolsep) * \real{0.1429}}
  >{\raggedright\arraybackslash}p{(\columnwidth - 12\tabcolsep) * \real{0.1429}}
  >{\raggedright\arraybackslash}p{(\columnwidth - 12\tabcolsep) * \real{0.1429}}
  >{\raggedright\arraybackslash}p{(\columnwidth - 12\tabcolsep) * \real{0.1429}}
  >{\raggedright\arraybackslash}p{(\columnwidth - 12\tabcolsep) * \real{0.1429}}
  >{\raggedright\arraybackslash}p{(\columnwidth - 12\tabcolsep) * \real{0.1429}}
  >{\raggedright\arraybackslash}p{(\columnwidth - 12\tabcolsep) * \real{0.1429}}@{}}
\toprule\noalign{}
\begin{minipage}[b]{\linewidth}\raggedright
Candidat
\end{minipage} & \begin{minipage}[b]{\linewidth}\raggedright
Mesures du corpus
\end{minipage} & \begin{minipage}[b]{\linewidth}\raggedright
quantifiables
\end{minipage} & \begin{minipage}[b]{\linewidth}\raggedright
partielles
\end{minipage} & \begin{minipage}[b]{\linewidth}\raggedright
non quantifiables
\end{minipage} & \begin{minipage}[b]{\linewidth}\raggedright
montants chiffrés (Md€)
\end{minipage} & \begin{minipage}[b]{\linewidth}\raggedright
dont avec canal dans le modèle
\end{minipage} \\
\midrule\noalign{}
\endhead
\bottomrule\noalign{}
\endlastfoot
Mélenchon (LFI) & 36 & 18 & 15 & 3 & 188,1 & 188,1 \\
Le Pen (RN) & 16 & 4 & 9 & 3 & 97,6 & 80,1 \\
Philippe (Horizons) & 14 & 2 & 8 & 4 & 2,0 & 2,0 \\
\end{longtable}

\hypertarget{sec:prog-lfi}{%
\section{Jean-Luc Mélenchon (La France insoumise)}\label{sec:prog-lfi}}

\textbf{Vision économique et sociale.} \emph{L'Avenir en commun} se
présente comme un programme de « rupture avec l'ordre économique »
néolibéral, organisé en six « ruptures » dont la première est la «
bifurcation écologique ». La doctrine économique est explicitement
keynésienne et redistributive : relancer l'activité par les revenus des
classes populaires et les services publics, financer cette relance par
l'imposition du capital, des hauts patrimoines et des grandes
entreprises, et planifier la transition écologique en mobilisant un pôle
public bancaire. Le document budgétaire du 14 octobre 2026 précise : «
Si vous nous jugez d'après les critères du néolibéralisme, nos projets
sont inacceptables » (citation rapportée par LCP).

\textbf{Principaux objectifs chiffrés.} Pour 2027 : 106,4 Md€ de
recettes supplémentaires, 81,7 Md€ de dépenses nouvelles, soit une
réduction du déficit de 24,7 Md€ ; stabilisation du ratio de dette sur
un quinquennat ; « retour de l'activité et redressement des comptes »
dès 2028 (formulation du document). Le document précise lui-même qu'il «
ne constitue pas un chiffrage complet » du programme, promis pour le
début de 2027.

\textbf{Mesures fiscales.} Suppression du prélèvement forfaitaire unique
(9 Md€), rétablissement d'un impôt sur la fortune renforcé (10 Md€),
taxe sur les patrimoines supérieurs à 100 M€ (15 Md€), réforme des
droits de succession (20 Md€ la première année), impôt sur le revenu à
14 tranches (non chiffré), rétablissement de la CVAE sur les plus
grandes entreprises (3 Md€), réforme de l'impôt sur les sociétés (3
Md€), taxe universelle sur les multinationales (5 Md€), renforcement de
la taxe sur les services numériques (3,9 Md€), taxe sur les superprofits
(5 Md€), rabot des niches fiscales et sociales « les plus injustes » (22
Md€), taxes sur les transactions financières (5 Md€), les rachats
d'actions (2,5 Md€) et les crypto-actifs (1 Md€), lutte contre la fraude
(2 Md€).

\textbf{Dépenses publiques et services publics.} Dégel et hausse de 10
\% du point d'indice (17 Md€), 10 000 postes dans l'Éducation nationale
(0,5 Md€), gratuité de la cantine (3,6 Md€), 250 000 places à
l'université (3,5 Md€), hôpital public (3 Md€), lutte contre les
violences sexistes, sexuelles et faites aux enfants (4,6 Md€),
indexation des budgets des services supprimés en 2026 (7 Md€).

\textbf{Protection sociale et retraites.} Abrogation de la réforme de
2023 (9 Md€) et des réformes de l'assurance chômage (3 Md€) ;
revalorisation de 10 \% des pensions, allocations familiales et aides au
logement (8 Md€) ; garantie d'autonomie pour les jeunes (9,5 Md€) ;
objectif de mandat : retraite à 60 ans avec 40 annuités, pension
minimale alignée sur le SMIC pour une carrière complète (non chiffrés).

\textbf{Emploi et salaires.} SMIC à 1 700 € net (non chiffré ; soit
+15,0 \% par rapport au SMIC net de 1 477,93 € au 1er juin 2026),
conférence sur les salaires et indexation généralisée (non chiffrées).

\textbf{Entreprises, investissement et industrie ; énergie et
transition.} Planification écologique : 10 Md€ pour l'adaptation au
dérèglement climatique, 700 000 rénovations par an (non chiffré), et 100
Md€ par an d'investissements « portés par le pôle public bancaire » hors
budget de l'État ; blocage du prix des carburants à 1,70 €/l (non
chiffré).

\textbf{Commerce extérieur et souveraineté ; finance.} Contrôle des
capitaux, fin de la cotation continue en Bourse, relèvement du plafond
du livret A à 30 000 € et de la centralisation de l'épargne réglementée,
création d'un « livret B » (non chiffrés).

\textbf{Politique budgétaire et financement.} Le financement repose sur
les recettes nouvelles ; 25 Md€ par an sont annoncés pour la réduction
du déficit. Le document ne précise ni hypothèse de croissance pour 2027,
ni multiplicateur.

\textbf{Bénéficiaires et perdants potentiels (premier ordre).}
Bénéficiaires : agents publics, ménages modestes et jeunes (transferts,
SMIC, cantine), retraités (revalorisations, âge), demandeurs d'emploi.
Contributeurs : détenteurs de patrimoines élevés et de revenus du
capital, héritiers, grandes entreprises et multinationales,
bénéficiaires des niches supprimées (composition non publiée). Les
effets indirects (prix, emploi, compétitivité) sont analysés aux
chapitres \ref{sec:transmission} et \ref{sec:resultats}.

\textbf{Traduction dans le modèle.} Le paquet 2027 est entièrement
traduit ; les montants recoupent exactement les totaux publiés. Les
principales hypothèses de traduction sont :

\begin{itemize}
\tightlist
\item
  les recettes sans canal propre (taxe sur les très grands patrimoines,
  multinationales, numérique, superprofits, transactions financières,
  rachats d'actions, crypto-actifs, fraude) sont traitées comme des
  recettes sans effet direct sur la demande intérieure {[}H{]} ;
  l'érosion de leur assiette n'est pas représentée : leurs montants sont
  des bornes hautes ;
\item
  les 22 Md€ de niches, dont la composition n'est pas publiée, sont
  répartis pour moitié sur les ménages et pour moitié sur le coût du
  travail {[}H{]} ;
\item
  la revalorisation de 10 \% (8 Md€) est versée en transferts non ciblés
  : le montant annoncé n'est pas compatible avec une revalorisation
  générale de toutes les pensions (masse de plusieurs centaines de Md€),
  ce qui suggère un périmètre plus restreint, à préciser {[}H{]} ;
\item
  les 100 Md€ du pôle public bancaire sont traités comme un
  investissement hors administrations publiques, appliqué à partir de
  2028, avec une additionnalité totale (borne haute) {[}H{]} ;
\item
  la réforme des retraites est simulée par M9 (retour immédiat à 62 ans,
  calendrier de durée de la loi de 2014) ; le coût annoncé (9 Md€) sert
  de contrôle.
\end{itemize}

\textbf{Tableau : mesures du corpus --- Mélenchon (LFI).} Montants en
Md€ (signe du point de vue du solde public), année de référence entre
parenthèses ; « début » : première année d'application dans le modèle ;
UE : niveau de qualification (chapitre \ref{sec:europe}).

\begin{longtable}[]{@{}
  >{\raggedright\arraybackslash}p{(\columnwidth - 14\tabcolsep) * \real{0.1250}}
  >{\raggedright\arraybackslash}p{(\columnwidth - 14\tabcolsep) * \real{0.1250}}
  >{\raggedright\arraybackslash}p{(\columnwidth - 14\tabcolsep) * \real{0.1250}}
  >{\raggedright\arraybackslash}p{(\columnwidth - 14\tabcolsep) * \real{0.1250}}
  >{\raggedright\arraybackslash}p{(\columnwidth - 14\tabcolsep) * \real{0.1250}}
  >{\raggedright\arraybackslash}p{(\columnwidth - 14\tabcolsep) * \real{0.1250}}
  >{\raggedright\arraybackslash}p{(\columnwidth - 14\tabcolsep) * \real{0.1250}}
  >{\raggedright\arraybackslash}p{(\columnwidth - 14\tabcolsep) * \real{0.1250}}@{}}
\toprule\noalign{}
\begin{minipage}[b]{\linewidth}\raggedright
id
\end{minipage} & \begin{minipage}[b]{\linewidth}\raggedright
mesure
\end{minipage} & \begin{minipage}[b]{\linewidth}\raggedright
montant
\end{minipage} & \begin{minipage}[b]{\linewidth}\raggedright
chiffrage
\end{minipage} & \begin{minipage}[b]{\linewidth}\raggedright
début
\end{minipage} & \begin{minipage}[b]{\linewidth}\raggedright
traduction dans le modèle
\end{minipage} & \begin{minipage}[b]{\linewidth}\raggedright
quant.
\end{minipage} & \begin{minipage}[b]{\linewidth}\raggedright
UE
\end{minipage} \\
\midrule\noalign{}
\endhead
\bottomrule\noalign{}
\endlastfoot
L01 & Suppression du prélèvement forfaitaire unique (flat tax) & +9,0
(2027) & candidat & 2027 & \texttt{pfu} +15,00 & oui & 0 \\
L02 & Rétablissement d'un impôt sur la fortune renforcé & +10,0 (2027) &
candidat & 2027 & \texttt{isf} +1,31 & oui & 0 \\
L03 & Taxe Zucman (patrimoines supérieurs à 100 M€) & +15,0 (2027) &
candidat & 2027 & \texttt{recettes\_neutres} +15,00 & partiel & --- \\
L04 & Réforme des droits de succession & +20,0 (2027) & candidat & 2027
& \texttt{successions} +111,11 & oui & 0 \\
L05 & Impôt sur le revenu à 14 tranches & --- & aucun & 2027 & --- &
partiel & 0 \\
L06 & Rétablissement de la CVAE sur les plus grandes entreprises & +3,0
(2027) & candidat & 2027 & \texttt{impots\_production} +3,00 & oui &
0 \\
L07 & Réforme de l'impôt sur les sociétés & +3,0 (2027) & candidat &
2027 & \texttt{taux\_is} +3,0 Md€ & oui & 0 \\
L08 & Taxe universelle sur les multinationales & +5,0 (2027) & candidat
& 2027 & \texttt{recettes\_neutres} +5,00 & partiel & 1 \\
L09 & Renforcement de la taxe sur les services numériques (GAFAM) & +3,9
(2027) & candidat & 2027 & \texttt{recettes\_neutres} +3,90 & partiel &
0 \\
L10 & Taxe sur les superprofits & +5,0 (2027) & candidat & 2027 &
\texttt{recettes\_neutres} +5,00 & partiel & 0 \\
L11 & Rabot et suppression des niches fiscales et sociales « les plus
injustes » & +22,0 (2027) & candidat & 2027 &
\texttt{transferts\_menages} −11,0 Md€; \texttt{cotisations\_employeurs}
+11,00 & partiel & 0 \\
L12 & Fin des exonérations de cotisations au-dessus de 2 SMIC & --- &
aucun & 2027 & --- & partiel & 0 \\
L13 & Taxes sur les transactions financières & +5,0 (2027) & candidat &
2027 & \texttt{recettes\_neutres} +5,00 & partiel & 0 \\
L14 & Taxe sur les rachats d'actions & +2,5 (2027) & candidat & 2027 &
\texttt{recettes\_neutres} +2,50 & partiel & 0 \\
L15 & Taxe sur les crypto-actifs & +1,0 (2027) & candidat & 2027 &
\texttt{recettes\_neutres} +1,00 & partiel & 0 \\
L16 & Lutte contre la fraude fiscale & +2,0 (2027) & candidat & 2027 &
\texttt{recettes\_neutres} +2,00 & partiel & 0 \\
L20 & Dégel et hausse de 10 \% du point d'indice & −17,0 (2027) &
candidat & 2027 & \texttt{point\_indice} +10,00 & oui & 0 \\
L21 & Revalorisation de 10 \% des pensions, allocations familiales et
APL & −8,0 (2027) & candidat & 2027 & \texttt{transferts\_menages} +8,0
Md€ & partiel & --- \\
L22 & Logements d'urgence et réquisition de logements vacants & −3,0
(2027) & candidat & 2027 & \texttt{autres\_depenses} +3,00 & oui &
--- \\
L23 & Garantie d'autonomie pour les jeunes & −9,5 (2027) & candidat &
2027 & \texttt{transferts\_menages} +9,5 Md€ & oui & --- \\
L24 & SMIC à 1 700 € net & --- & aucun & 2027 & \texttt{smic} +15,03 &
oui & 0 \\
L25 & Blocage du prix des carburants à 1,70 €/l & --- & aucun & 2027 &
--- & partiel & 1 \\
L26 & 10 000 postes dans l'Éducation nationale & −0,5 (2027) & candidat
& 2027 & \texttt{education} +0,50 & oui & --- \\
L27 & Gratuité de la cantine scolaire & −3,6 (2027) & candidat & 2027 &
\texttt{education} +3,60 & oui & --- \\
L28 & 250 000 places à l'université & −3,5 (2027) & candidat & 2027 &
\texttt{education} +3,50 & oui & --- \\
L29 & Hôpital public (+10 000 postes soignants) & −3,0 (2027) & candidat
& 2027 & \texttt{sante} +3,00 & oui & --- \\
L30 & Violences sexistes et sexuelles ; violences faites aux enfants &
−4,6 (2027) & candidat & 2027 & \texttt{autres\_depenses} +4,60 & oui &
--- \\
L31 & Indexation des budgets 2027 des services supprimés en 2026 & −7,0
(2027) & candidat & 2027 & \texttt{autres\_depenses} +7,00 & oui &
--- \\
L32 & Adaptation au dérèglement climatique & −10,0 (2027) & candidat &
2027 & \texttt{investissement\_public} +10,00 & oui & --- \\
L33 & Investissements de planification écologique portés par le pôle
public bancaire & --- & candidat & 2028 &
\texttt{investissement\_hors\_apu} +100,00 & partiel & 1 \\
L34 & Abrogation de la réforme de 2023 (retour à 62 ans) & −9,0 (2027) &
candidat & 2027 & \texttt{retraites} : age=62, mode=immediat, duree=2014
& oui & --- \\
L35 & Abrogation des réformes de l'assurance chômage & −3,0 (2027) &
candidat & 2027 & \texttt{assurance\_chomage} +7,50 & oui & --- \\
L36 & Retour à 60 ans avec 40 annuités & --- & aucun & 2028 &
\texttt{retraites} : age=60, mode=progressif, duree=160,
duree\_mode=progressif & partiel & --- \\
L37 & Plafond du livret A relevé à 30 000 € ; centralisation portée à 80
\% ; livret B & --- & aucun & 2027 & --- & non & 1 \\
L38 & Contrôle des capitaux ; fin de la cotation continue & --- & aucun
& 2027 & --- & non & 4 \\
L39 & Réduction du déficit de 25 Md€ par an ; stabilisation de la dette
en un quinquennat & --- & candidat & 2027 & régime (objectif) & non &
--- \\
\end{longtable}

\hypertarget{sec:prog-rn}{%
\section{Marine Le Pen (Rassemblement national)}\label{sec:prog-rn}}

\textbf{Vision économique et sociale.} Le programme combine une ligne de
défense du pouvoir d'achat (baisse de la TVA sur l'énergie, mesures pour
les jeunes et les familles, retour partiel sur la réforme des
retraites), une attention aux entreprises (suppression de la CFE et des
impôts de production « qui nuisent à la relocalisation »), une priorité
nationale (logement, emploi, prestations) et un objectif de redressement
des comptes fondé sur des économies ciblées : immigration, contribution
au budget européen, organisation de l'État et des collectivités. La
presse a souligné l'évolution du discours vers une ligne plus libérale
qu'en 2022 tout en restant attachée à la redistribution (franceinfo).

\textbf{Principaux objectifs chiffrés.} 140 Md€ d'économies nettes d'ici
2032 (présentées le 6 octobre 2026) ; dette ramenée de 121 \% du PIB en
2027 à 112 \% en 2032 ; déficit à 3 \% du PIB en 2030 ; 30 Md€ de
baisses d'impôts ; règle d'or constitutionnelle par référendum dans les
douze mois, plafonnant le déficit au niveau qui stabilise la dette.

\textbf{Mesures fiscales.} TVA ramenée de 20 \% à 5,5 \% sur
l'électricité, le gaz, le fioul et les carburants (coût non publié par
le candidat ; ordre de grandeur de l'Institut Montaigne en 2022 : une
dizaine de Md€ ; d'autres estimations de presse vont jusqu'à 15-18 Md€)
; exonération d'impôt sur le revenu des moins de 30 ans jusqu'à quatre
fois le SMIC (2 Md€) ; part fiscale pleine dès le deuxième enfant (3,5
Md€ en 2027) ; défiscalisation des dons familiaux ; suppression de la
CFE et des impôts de production (non chiffré).

\textbf{Dépenses publiques.} Économies sur l'immigration (15 Md€ en
2027, 29 Md€ en année pleine), réorganisation de l'État (20,7 Md€ dès
2027, avec centralisation des recettes de TVA et transferts votés par le
Parlement aux collectivités et à la Sécurité sociale), contribution des
collectivités (8,2 Md€, dont suppression des EPCI et des métropoles),
agences de l'État (3,4 Md€ ; une commission sénatoriale évaluait ces
économies à environ 0,5 Md€), baisse de la contribution au budget de
l'Union (9 Md€).

\textbf{Protection sociale et retraites.} Âge légal ramené à 62 ans, 60
ans pour les carrières commencées avant 20 ans, 42 annuités, application
au 1er janvier 2028, coût brut de 18,6 Md€ en 2032 selon le RN ; «
pilier de capitalisation » obligatoire, dont le Fonds de réserve pour
les retraites serait le socle ; 15 à 20 Md€ d'économies « à terme » par
la « correction de mesures injustes » (non détaillées).

\textbf{Emploi, salaires, jeunesse, familles.} Permis de conduire
gratuit (700 000 jeunes par an, 0,7 Md€), prime d'activité étudiante
(0,6 Md€), suppression de la CVEC, prêt à taux zéro au troisième enfant.

\textbf{Entreprises, industrie, énergie, commerce, souveraineté.}
Suppression de la CFE et allègement des impôts de production ; prix
national régulé de l'électricité calculé sur les coûts de production
français ; neutralité carbone « avant 2050 » ; priorité nationale dans
l'accès au logement, à l'emploi et aux prestations.

\textbf{Politique budgétaire et financement.} Le budget est présenté
comme « réversible » ; la trajectoire repose sur des économies dont une
partie seulement est ventilée.

\textbf{Bénéficiaires et perdants potentiels (premier ordre).}
Bénéficiaires : consommateurs d'énergie (davantage, en montant, les
ménages aisés qui consomment plus d'énergie, selon plusieurs analyses de
presse), jeunes actifs, familles, retraités ayant commencé à travailler
tôt, entreprises redevables des impôts de production. Contributeurs :
étrangers (prestations et services publics, selon des modalités non
publiées), collectivités territoriales et leurs agents, opérateurs de
l'État, budget européen.

\textbf{Traduction dans le modèle.} Hypothèses principales :

\begin{itemize}
\tightlist
\item
  économies « immigration » (15 Md€ en 2027) traduites en baisse de
  prestations {[}H{]} ; la part fondée sur la nationalité est
  incompatible avec le droit de l'Union pour les ressortissants
  européens, mais cette part n'est pas connue : la mesure est conservée
  et le risque signalé ;
\item
  « réorganisation de l'État » (20,7 Md€) traduite en baisse de
  consommation publique ; les contributions des collectivités (8,2 Md€)
  et des agences (3,4 Md€) sont traitées comme des composantes de cette
  enveloppe et ne sont pas additionnées {[}H, faute de précision dans
  les sources{]} ;
\item
  TVA sur l'énergie traduite en baisse équivalente du taux normal (le
  modèle ne distingue pas l'énergie) pour un coût de 10 Md€ {[}tiers,
  H{]} ;
\item
  réforme des retraites simulée par M9 : 62 ans, 168 trimestres, départ
  à 60 ans pour 17 \% de chaque génération (part des départs pour
  carrière longue au régime général) {[}H{]} ;
\item
  économies sur les retraites (15-20 Md€) et pilier de capitalisation :
  non simulés faute de mesures identifiées.
\end{itemize}

\textbf{Tableau : mesures du corpus --- Le Pen (RN).} Montants en Md€
(signe du point de vue du solde public), année de référence entre
parenthèses ; « début » : première année d'application dans le modèle ;
UE : niveau de qualification (chapitre \ref{sec:europe}).

\begin{longtable}[]{@{}
  >{\raggedright\arraybackslash}p{(\columnwidth - 14\tabcolsep) * \real{0.1250}}
  >{\raggedright\arraybackslash}p{(\columnwidth - 14\tabcolsep) * \real{0.1250}}
  >{\raggedright\arraybackslash}p{(\columnwidth - 14\tabcolsep) * \real{0.1250}}
  >{\raggedright\arraybackslash}p{(\columnwidth - 14\tabcolsep) * \real{0.1250}}
  >{\raggedright\arraybackslash}p{(\columnwidth - 14\tabcolsep) * \real{0.1250}}
  >{\raggedright\arraybackslash}p{(\columnwidth - 14\tabcolsep) * \real{0.1250}}
  >{\raggedright\arraybackslash}p{(\columnwidth - 14\tabcolsep) * \real{0.1250}}
  >{\raggedright\arraybackslash}p{(\columnwidth - 14\tabcolsep) * \real{0.1250}}@{}}
\toprule\noalign{}
\begin{minipage}[b]{\linewidth}\raggedright
id
\end{minipage} & \begin{minipage}[b]{\linewidth}\raggedright
mesure
\end{minipage} & \begin{minipage}[b]{\linewidth}\raggedright
montant
\end{minipage} & \begin{minipage}[b]{\linewidth}\raggedright
chiffrage
\end{minipage} & \begin{minipage}[b]{\linewidth}\raggedright
début
\end{minipage} & \begin{minipage}[b]{\linewidth}\raggedright
traduction dans le modèle
\end{minipage} & \begin{minipage}[b]{\linewidth}\raggedright
quant.
\end{minipage} & \begin{minipage}[b]{\linewidth}\raggedright
UE
\end{minipage} \\
\midrule\noalign{}
\endhead
\bottomrule\noalign{}
\endlastfoot
P01 & 140 Md€ d'économies nettes d'ici 2032 ; dette de 121 \% (2027) à
112 \% (2032) ; déficit à 3 \% en 2030 ; 30 Md€ de baisses d'impôts &
+140,0 (2032) & candidat & 2027 & régime (objectif) & non & --- \\
P02 & Règle d'or constitutionnelle par référendum (déficit plafonné au
niveau stabilisant la dette) & --- & candidat & 2027 & régime
(R1\_actif) & non & 0 \\
P03 & Économies sur le « coût de l'immigration » & +15,0 (2027) &
candidat & 2027 & \texttt{transferts\_menages} −15,0 Md€ & partiel &
4 \\
P04 & Baisse de la contribution au budget de l'Union & +9,0 (2027) &
candidat & 2027 & \texttt{recettes\_neutres} +9,00 & oui & 2 \\
P05 & Réorganisation de l'État (centralisation de la TVA, transferts aux
collectivités et à la Sécurité sociale votés par le Parlement) & +20,7
(2027) & candidat & 2027 & \texttt{autres\_depenses} −20,70 & partiel &
--- \\
P06 & Contribution des collectivités (dont suppression des EPCI et
métropoles 0,25 Md€, effectifs 0,8 Md€) & +8,2 (2027) & candidat & 2027
& --- & partiel & --- \\
P07 & Agences de l'État & +3,4 (2027) & candidat & 2027 & --- & partiel
& --- \\
P08 & Âge légal ramené à 62 ans (60 ans pour les carrières commencées
avant 20 ans), 42 annuités ; pilier de capitalisation obligatoire (socle
: Fonds de réserve pour les retraites) & −18,6 (2032) & candidat & 2028
& \texttt{retraites} : age=62, mode=immediat, duree=168, groupe=60,
groupe\_part=0,17 & partiel & --- \\
P09 & Économies sur les retraites (« corrections de mesures injustes »)
& +17,5 (2032) & candidat & 2027 & --- & partiel & --- \\
P10 & TVA ramenée de 20 \% à 5,5 \% sur l'électricité, le gaz, le fioul
et les carburants & −10,0 (2022) & tiers & 2027 & \texttt{tva} −1,22 &
partiel & 2 \\
P11 & Exonération d'impôt sur le revenu des moins de 30 ans (jusqu'à 4
SMIC) & −2,0 (2027) & candidat & 2027 & \texttt{transferts\_menages}
+2,0 Md€ & oui & --- \\
P12 & Permis de conduire gratuit (18 ans, casier vierge) ; prime
d'activité étudiante & −1,3 (2027) & candidat & 2027 &
\texttt{transferts\_menages} +1,3 Md€ & oui & --- \\
P13 & Part fiscale pleine dès le deuxième enfant & −3,5 (2027) &
candidat & 2027 & \texttt{transferts\_menages} +3,5 Md€ & oui & --- \\
P14 & Dons familiaux défiscalisés ; prêt à taux zéro au troisième enfant
; suppression de la CVEC & --- & aucun & 2027 & --- & partiel & --- \\
P15 & Suppression de la CFE et des impôts de production « qui nuisent à
la relocalisation » & --- & aucun & 2027 & \texttt{impots\_production}
(montant à préciser) & partiel & 0 \\
P16 & Priorité nationale (logement, emploi, prestations) & --- & aucun &
2027 & --- & non & 4 \\
\end{longtable}

\hypertarget{sec:prog-horizons}{%
\section{Édouard Philippe (Horizons)}\label{sec:prog-horizons}}

\textbf{Vision économique et sociale.} Le programme se présente comme «
massif », « sérieux » et « en rupture » avec la décennie écoulée :
réduction de la dépense publique (« l'argent magique n'existe pas »),
allongement de la durée du travail (retraite à 65 ans, fin des 35
heures), réforme de l'État (250 000 postes en moins, statut limité au
régalien et à l'enseignement), soutien à la compétitivité (baisse de 50
Md€ des impôts de production) et à la natalité, moratoire fiscal et
normatif. Un « volume 2 » est annoncé pour janvier 2027.

\textbf{Principaux objectifs chiffrés.} Près de 150 Md€ d'économies pour
ramener le déficit sous 2,5 \% du PIB à la fin du quinquennat (le
chiffre de 2 \% circule également dans la presse) ; la moitié de
l'effort sur les dépenses sociales, un quart sur l'État et les
collectivités, un quart « sur la croissance » ; règle d'or adoptée par
référendum dès le début du quinquennat ; aucun nouvel impôt pendant cinq
ans.

\textbf{Mesures fiscales.} Baisse de 50 Md€ des impôts de production
financée par 50 Md€ de baisse des aides aux entreprises (« pacte »
pluriannuel, dès la première loi de finances) ; heures supplémentaires
défiscalisées et deux heures hebdomadaires exonérées de cotisations ;
exonération d'impôt sur le revenu pendant les trois premières années de
l'enfant pour les classes moyennes ; suppression des frais de notaire
pour le premier achat avant 35 ans.

\textbf{Dépenses publiques.} 250 000 postes de fonctionnaires en moins
sur le quinquennat (un départ sur trois non remplacé dans les fonctions
de bureau) ; hausse de 20 \% du salaire des enseignants ; fusion,
suppression ou privatisation d'une trentaine d'agences et d'opérateurs
(dont France Travail) ; fonds d'adaptation climatique de 2 Md€ par an.

\textbf{Protection sociale et retraites.} Retraite à 65 ans et 45
annuités, départ précoce pour un tiers des assurés, capitalisation
collective ; indemnisation chômage limitée à 12 mois pour les moins de
50 ans ; baisse de 30 \% du nombre d'allocataires du RSA ; trois jours
de carence pour tous ; allocation familiale unique de 350 € par mois dès
le premier enfant, dégressive.

\textbf{Emploi et salaires ; énergie ; commerce.} Fin des 35 heures ;
relance du nucléaire (14 réacteurs) ; suppression du « zéro
artificialisation nette ».

\textbf{Bénéficiaires et perdants potentiels (premier ordre).}
Bénéficiaires : entreprises redevables des impôts de production (au
détriment de celles qui bénéficient des aides supprimées), salariés
effectuant des heures supplémentaires, familles avec enfants,
enseignants, primo-accédants jeunes. Contributeurs : générations nées
après le milieu des années 1960 (âge et durée), demandeurs d'emploi de
moins de 50 ans, allocataires du RSA, agents publics (effectifs),
bénéficiaires des aides aux entreprises supprimées.

\textbf{Traduction dans le modèle.} Peu de mesures sont chiffrées une à
une ; la lecture par objectifs budgétaires (lecture B) est donc
essentielle. Hypothèses principales :

\begin{itemize}
\tightlist
\item
  250 000 postes en moins atteints en cinq ans (coût moyen de 56 000 €
  par agent du simulateur) ;
\item
  retraites : 65 ans et 180 trimestres, un trimestre par génération, un
  tiers des assurés aux conditions actuelles {[}H{]} ;
\item
  fin des 35 heures : durée de référence portée à 37 heures (deux heures
  supplémentaires annoncées) {[}H{]} ;
\item
  RSA : baisse de 30 \% du montant forfaitaire comme approximation
  budgétaire d'une baisse des effectifs {[}H{]} ;
\item
  pacte impôts de production / aides : moitié des aides assimilée à des
  subventions (effet net nul dans le modèle), moitié à des allègements
  de cotisations (hausse du coût du travail) {[}H{]}, avec deux
  variantes ; le modèle ne distingue pas les distorsions propres aux
  impôts de production de celles des subventions, argument central du
  candidat ;
\item
  non simulés faute de barème : allocations familiales, assurance
  chômage (montant), jours de carence, mesures familiales.
\end{itemize}

\textbf{Tableau : mesures du corpus --- Philippe (Horizons).} Montants
en Md€ (signe du point de vue du solde public), année de référence entre
parenthèses ; « début » : première année d'application dans le modèle ;
UE : niveau de qualification (chapitre \ref{sec:europe}).

\begin{longtable}[]{@{}
  >{\raggedright\arraybackslash}p{(\columnwidth - 14\tabcolsep) * \real{0.1250}}
  >{\raggedright\arraybackslash}p{(\columnwidth - 14\tabcolsep) * \real{0.1250}}
  >{\raggedright\arraybackslash}p{(\columnwidth - 14\tabcolsep) * \real{0.1250}}
  >{\raggedright\arraybackslash}p{(\columnwidth - 14\tabcolsep) * \real{0.1250}}
  >{\raggedright\arraybackslash}p{(\columnwidth - 14\tabcolsep) * \real{0.1250}}
  >{\raggedright\arraybackslash}p{(\columnwidth - 14\tabcolsep) * \real{0.1250}}
  >{\raggedright\arraybackslash}p{(\columnwidth - 14\tabcolsep) * \real{0.1250}}
  >{\raggedright\arraybackslash}p{(\columnwidth - 14\tabcolsep) * \real{0.1250}}@{}}
\toprule\noalign{}
\begin{minipage}[b]{\linewidth}\raggedright
id
\end{minipage} & \begin{minipage}[b]{\linewidth}\raggedright
mesure
\end{minipage} & \begin{minipage}[b]{\linewidth}\raggedright
montant
\end{minipage} & \begin{minipage}[b]{\linewidth}\raggedright
chiffrage
\end{minipage} & \begin{minipage}[b]{\linewidth}\raggedright
début
\end{minipage} & \begin{minipage}[b]{\linewidth}\raggedright
traduction dans le modèle
\end{minipage} & \begin{minipage}[b]{\linewidth}\raggedright
quant.
\end{minipage} & \begin{minipage}[b]{\linewidth}\raggedright
UE
\end{minipage} \\
\midrule\noalign{}
\endhead
\bottomrule\noalign{}
\endlastfoot
H01 & Près de 150 Md€ d'économies ; déficit sous 2,5 \% du PIB en fin de
quinquennat ; moitié de l'effort sur les dépenses sociales, un quart sur
l'État et les collectivités, un quart sur la croissance & +150,0 (2032)
& candidat & 2027 & régime (objectif) & non & --- \\
H02 & Règle d'or constitutionnelle par référendum ; aucun nouvel impôt
pendant cinq ans & --- & candidat & 2027 & régime (R2\_actif\_depenses)
& non & 0 \\
H03 & 250 000 postes de fonctionnaires en moins (un départ sur trois non
remplacé dans les fonctions de bureau) & --- & aucun & 2027 &
\texttt{emploi\_public} −250,00 & oui & --- \\
H04 & Hausse de 20 \% du salaire des enseignants & --- & aucun & 2027 &
--- & partiel & --- \\
H05 & Âge de départ à 65 ans et 45 annuités ; départ précoce pour un
tiers des assurés ; capitalisation collective & --- & aucun & 2028 &
\texttt{retraites} : age=65, mode=progressif, duree=180,
duree\_mode=progressif, groupe=ref, groupe\_part=0,33 & partiel & --- \\
H06 & Fin des 35 heures ; deux heures supplémentaires hebdomadaires
exonérées de cotisations & --- & aucun & 2027 & \texttt{temps\_travail}
+37,00 & partiel & --- \\
H07 & Indemnisation limitée à 12 mois pour les moins de 50 ans & --- &
aucun & 2027 & \texttt{assurance\_chomage} (montant à préciser) &
partiel & --- \\
H08 & Baisse de 30 \% du nombre d'allocataires du RSA (accompagnement,
métiers en tension) & --- & aucun & 2027 & \texttt{rsa} −30,00 & partiel
& --- \\
H09 & Trois jours de carence pour tous (public et privé) & --- & aucun &
2027 & --- & non & --- \\
H10 & Allocation de 350 € par mois dès le premier enfant, dégressive ;
exonération d'impôt les trois premières années de l'enfant ; suppression
des frais de notaire au premier achat (\textless{} 35 ans) ; prêt à taux
zéro famille & --- & aucun & 2027 & --- & partiel & --- \\
H11 & Baisse de 50 Md€ des impôts de production financée par 50 Md€ de
baisse des aides aux entreprises & +0,0 (2028) & candidat & 2028 &
\texttt{impots\_production} −25,00; \texttt{cotisations\_employeurs}
+25,00 & partiel & 0 \\
H12 & Fonds unique d'adaptation au changement climatique & −2,0 (2027) &
candidat & 2027 & \texttt{investissement\_public} +2,00 & oui & --- \\
H13 & Fusion, suppression ou privatisation d'une trentaine d'agences
(dont France Travail) & --- & aucun & 2027 & --- & partiel & --- \\
H14 & Huit réacteurs nucléaires supplémentaires ; suppression du ZAN &
--- & aucun & 2027 & --- & non & --- \\
\end{longtable}

\hypertarget{sec:prog-regimes}{%
\section{Éléments de régime : règles d'or et objectifs de
solde}\label{sec:prog-regimes}}

Les trois candidats annoncent des objectifs de finances publiques qui ne
sont pas des mesures mais des \textbf{règles} : stabilisation de la
dette en un quinquennat (LFI), déficit à 3 \% en 2030 et règle d'or par
référendum (RN), déficit sous 2,5 \% en fin de quinquennat et règle d'or
(Horizons). Le modèle les représente dans la lecture B par une
trajectoire cible du solde public, atteinte par un effort réparti selon
les priorités déclarées :

\begin{longtable}[]{@{}
  >{\raggedright\arraybackslash}p{(\columnwidth - 6\tabcolsep) * \real{0.2500}}
  >{\raggedright\arraybackslash}p{(\columnwidth - 6\tabcolsep) * \real{0.2500}}
  >{\raggedright\arraybackslash}p{(\columnwidth - 6\tabcolsep) * \real{0.2500}}
  >{\raggedright\arraybackslash}p{(\columnwidth - 6\tabcolsep) * \real{0.2500}}@{}}
\toprule\noalign{}
\begin{minipage}[b]{\linewidth}\raggedright
Candidat
\end{minipage} & \begin{minipage}[b]{\linewidth}\raggedright
Objectif annoncé
\end{minipage} & \begin{minipage}[b]{\linewidth}\raggedright
Cible du modèle
\end{minipage} & \begin{minipage}[b]{\linewidth}\raggedright
Instruments de l'effort non ventilé {[}H{]}
\end{minipage} \\
\midrule\noalign{}
\endhead
\bottomrule\noalign{}
\endlastfoot
Mélenchon (LFI) & réduction du déficit de 24,7 Md€ en 2027 (mesures)
puis stabilisation du ratio dette/PIB en un quinquennat & solde
stabilisant la dette en 2032 (calculé dans le modèle) & impots menages
100 \% \\
Le Pen (RN) & déficit à 3 \% du PIB en 2030 ; dette de 121 \% (2027) à
112 \% (2032) ; 140 Md€ d'économies & −3,0 \% du PIB en 2030 & conso
publique 60 \%, prestations 40 \% \\
Philippe (Horizons) & déficit sous 2,5 \% du PIB en fin de quinquennat ;
moitié de l'effort sur les dépenses sociales, un quart sur l'État et les
collectivités, un quart « sur la croissance » ; aucun nouvel impôt &
−2,5 \% du PIB en 2032 & prestations 67 \%, conso publique 33 \% \\
\end{longtable}

Ces compositions sont des lectures prudentes des déclarations, non des
engagements : LFI annonce un financement par des recettes nouvelles
(d'où les impôts des ménages, sans pouvoir cibler les hauts revenus dans
ce canal agrégé) ; le RN des économies sur l'État, les collectivités,
l'immigration et les retraites ; Horizons la moitié de l'effort sur les
dépenses sociales et un quart sur l'État et les collectivités, sans
hausse d'impôt (le quart « croissance » n'est pas un instrument : il est
laissé au modèle).

\hypertarget{sec:incertitude-programmes}{%
\section{Degré d'incertitude des
programmes}\label{sec:incertitude-programmes}}

L'incertitude sur les programmes eux-mêmes est de trois ordres :

\begin{enumerate}
\def\labelenumi{\arabic{enumi}.}
\tightlist
\item
  \textbf{Incertitude de contenu} : mesures non chiffrées, enveloppes
  non ventilées, barèmes non publiés (impôt sur le revenu à 14 tranches,
  allocation familiale de 350 €, durée d'indemnisation du chômage) ;
\item
  \textbf{Incertitude de chiffrage} : rendements annoncés supérieurs aux
  rendements simulés (impôts sur le capital et les patrimoines), coûts
  sans chiffrage du candidat (TVA sur l'énergie), estimations contestées
  (économies sur les agences de l'État) ;
\item
  \textbf{Incertitude d'exécution} : calendrier parlementaire,
  contentieux constitutionnels et européens, réactions des acteurs
  (chapitre \ref{sec:europe}).
\end{enumerate}

Le programme de LFI est le mieux documenté pour 2027 mais ne l'est pas
au-delà ; ceux du RN et d'Horizons fixent des trajectoires de moyen
terme mais ventilent peu leurs économies. Aucun des trois n'est donc «
complètement » simulable : la comparaison doit respecter ces différences
de précision (chapitre \ref{sec:resultats}).

\hypertarget{sec:transmission}{%
\chapter{Transmission macroéconomique des
mesures}\label{sec:transmission}}

Ce chapitre explique comment chaque grande famille de mesures agit sur
l'économie dans le modèle, en distinguant les effets de demande
(rapides), les effets d'offre (lents), les effets de financement et de
dette, et l'incidence des prélèvements. Il s'appuie sur des chocs
standardisés simulés pour l'occasion (section
\ref{sec:multiplicateurs}), qui servent de clé de lecture aux résultats
des programmes (chapitre \ref{sec:resultats}).

\hypertarget{sec:demande-ct}{%
\section{Effets de demande à court terme}\label{sec:demande-ct}}

Une mesure qui augmente le revenu des ménages ou la dépense publique
augmente la demande adressée aux entreprises. Le premier tour dépend de
la nature de la mesure :

\begin{itemize}
\tightlist
\item
  \textbf{dépense publique} (consommation, investissement) : elle entre
  directement dans la demande finale ; une partie est importée (selon
  les contenus en importations du TES) ;
\item
  \textbf{transferts et baisses d'impôts des ménages} : seule la part
  consommée entre dans la demande (85 \% la première année en moyenne,
  plus pour les ménages modestes, moins pour les ménages aisés) ;
\item
  \textbf{baisses d'impôts sur les entreprises} : elles passent par les
  prix (répercussion), les salaires (incidence) et les marges ; l'effet
  de demande est indirect et plus lent.
\end{itemize}

Les tours suivants passent par la production, l'emploi, les revenus et
la consommation : c'est le multiplicateur keynésien. Les fuites sont les
importations, l'épargne et les prélèvements obligatoires.

\hypertarget{sec:offre-lt}{%
\section{Effets d'offre à moyen et long terme}\label{sec:offre-lt}}

L'offre (le PIB potentiel) dépend de la population active, du chômage
structurel, du capital et de la productivité. Dans le modèle, les
mesures agissent sur l'offre par cinq canaux :

\begin{enumerate}
\def\labelenumi{\arabic{enumi}.}
\tightlist
\item
  \textbf{Population active} : âge de la retraite et durée d'assurance
  (M9), incitations au travail des prestations (M15) ;
\item
  \textbf{Chômage structurel} : générosité de l'assurance chômage,
  formation ;
\item
  \textbf{Capital} : investissement des entreprises (accélérateur, coût
  d'usage), capital public ;
\item
  \textbf{Productivité} : capital public, R\&D (exploratoire), durée du
  travail (productivité horaire) ;
\item
  \textbf{Coût relatif des facteurs} : cotisations, impôts de
  production, salaires (substitution capital-travail, \(\sigma\)).
\end{enumerate}

La fermeture d'offre ramène progressivement la demande vers ce potentiel
: un programme qui accroît le potentiel (davantage d'actifs, de capital)
a des effets positifs croissants dans le temps, et inversement.

\hypertarget{sec:multiplicateurs}{%
\section{Multiplicateurs budgétaires
simulés}\label{sec:multiplicateurs}}

Le tableau suivant donne les effets de chocs permanents d'environ 0,5
point de PIB, appliqués à partir de 2028, sous le régime principal (taux
exogènes, règle figée) et sous le régime avec réaction de la BCE et
prime souveraine. Le multiplicateur est le rapport entre la variation du
PIB et le coût ex ante du choc (en part du PIB).

\textbf{Tableau : multiplicateurs simulés (variation du PIB / coût ex
ante, en part du PIB) et effet sur le ratio de dette.} Chocs permanents
à partir de 2028 ; environnement P4 ; source :
\texttt{phase7/multiplicateurs.py}.

\begin{longtable}[]{@{}
  >{\raggedright\arraybackslash}p{(\columnwidth - 14\tabcolsep) * \real{0.1250}}
  >{\raggedright\arraybackslash}p{(\columnwidth - 14\tabcolsep) * \real{0.1250}}
  >{\raggedright\arraybackslash}p{(\columnwidth - 14\tabcolsep) * \real{0.1250}}
  >{\raggedright\arraybackslash}p{(\columnwidth - 14\tabcolsep) * \real{0.1250}}
  >{\raggedright\arraybackslash}p{(\columnwidth - 14\tabcolsep) * \real{0.1250}}
  >{\raggedright\arraybackslash}p{(\columnwidth - 14\tabcolsep) * \real{0.1250}}
  >{\raggedright\arraybackslash}p{(\columnwidth - 14\tabcolsep) * \real{0.1250}}
  >{\raggedright\arraybackslash}p{(\columnwidth - 14\tabcolsep) * \real{0.1250}}@{}}
\toprule\noalign{}
\begin{minipage}[b]{\linewidth}\raggedright
choc --- régime
\end{minipage} & \begin{minipage}[b]{\linewidth}\raggedright
multiplicateur à 1 an
\end{minipage} & \begin{minipage}[b]{\linewidth}\raggedright
multiplicateur à 3 ans
\end{minipage} & \begin{minipage}[b]{\linewidth}\raggedright
multiplicateur à 5 ans
\end{minipage} & \begin{minipage}[b]{\linewidth}\raggedright
multiplicateur à 10 ans
\end{minipage} & \begin{minipage}[b]{\linewidth}\raggedright
dette à 1 an (pt)
\end{minipage} & \begin{minipage}[b]{\linewidth}\raggedright
dette à 5 ans (pt)
\end{minipage} & \begin{minipage}[b]{\linewidth}\raggedright
dette à 10 ans (pt)
\end{minipage} \\
\midrule\noalign{}
\endhead
\bottomrule\noalign{}
\endlastfoot
TVA taux normal (−2 pt) --- BCE + prime & +0,20 & +0,29 & +0,26 & +0,10
& +0,51 & +2,54 & +4,96 \\
TVA taux normal (−2 pt) --- principal (taux exogènes) & +0,20 & +0,29 &
+0,26 & +0,11 & +0,51 & +2,54 & +4,80 \\
consommation publique (+16 Md€) --- BCE + prime & +1,24 & +1,04 & +0,73
& +0,27 & −0,51 & −0,46 & +0,24 \\
consommation publique (+16 Md€) --- principal (taux exogènes) & +1,24 &
+1,04 & +0,75 & +0,30 & −0,51 & −0,56 & −0,21 \\
cotisations employeurs (−16 Md€) --- BCE + prime & +0,69 & +0,65 & +0,55
& +0,28 & +0,90 & +2,44 & +4,47 \\
cotisations employeurs (−16 Md€) --- principal (taux exogènes) & +0,69 &
+0,65 & +0,56 & +0,30 & +0,90 & +2,40 & +4,11 \\
impôts de production (−16 Md€) --- BCE + prime & +0,09 & +0,07 & +0,03 &
−0,05 & +0,80 & +3,25 & +6,88 \\
impôts de production (−16 Md€) --- principal (taux exogènes) & +0,09 &
+0,07 & +0,03 & −0,04 & +0,80 & +3,23 & +6,66 \\
investissement public (+16 Md€) --- BCE + prime & +0,66 & +0,68 & +0,76
& +1,11 & +0,20 & +2,15 & +5,55 \\
investissement public (+16 Md€) --- principal (taux exogènes) & +0,66 &
+0,69 & +0,77 & +1,13 & +0,20 & +2,08 & +5,10 \\
transferts aux ménages (+0,5 pt) --- BCE + prime & +0,54 & +0,58 & +0,45
& +0,22 & +0,11 & +1,17 & +3,02 \\
transferts aux ménages (+0,5 pt) --- principal (taux exogènes) & +0,54 &
+0,58 & +0,46 & +0,24 & +0,11 & +1,11 & +2,70 \\
\end{longtable}

Lecture : une hausse permanente de la consommation publique a un
multiplicateur de 1,24 la première année, 0,75 à cinq ans et 0,30 à dix
ans (régime principal) ; sous le régime BCE + prime, 1,24, 0,73 et 0,27.
Les transferts aux ménages ont un multiplicateur plus faible (0,54 la
première année), une partie étant épargnée. Les multiplicateurs
diminuent avec l'horizon sous l'effet de la fermeture d'offre et de la
hausse des prix. Les baisses de prélèvements sur les entreprises ont des
effets d'activité plus faibles et plus lents : 0,69 la première année
pour les cotisations employeurs (effet sur le coût du travail et
l'emploi), 0,09 pour les impôts de production (dont la moitié est
répercutée sur les prix et une partie conservée en marges) ; la baisse
du taux normal de TVA a un multiplicateur de 0,20, parce que les baisses
ne sont répercutées qu'à moitié sur les prix. L'investissement public a
un effet de première année plus faible que la consommation publique
(0,66, contenu en importations plus élevé) mais croissant (1,13 à dix
ans) grâce à l'accumulation de capital public.

\begin{figure}[htbp]\centering\includegraphics[width=0.9\textwidth]{figures/multiplicateurs.pdf}\caption{Multiplicateurs simulés selon l'horizon (régime principal).}\end{figure}

Ces valeurs sont à comparer aux estimations empiriques (chapitre
\ref{sec:theorie}) : les multiplicateurs de dépense de première année du
modèle se situent dans le haut de la fourchette (0,6-1,2 pour la plupart
des estimations, davantage en récession), parce que les taux sont
exogènes dans le régime principal et que la fermeture est lente. Sous le
régime BCE + prime, la hausse des taux réduit légèrement l'effet ; avec
une fermeture rapide (\(\rho = 0{,}3\)), il s'estompe plus vite
(chapitre \ref{sec:robustesse}).

\hypertarget{sec:financement}{%
\section{Effets de financement et de dette}\label{sec:financement}}

Une mesure non financée augmente la dette ; la charge d'intérêts
correspondante dépend du taux de marché auquel la dette nouvelle est
émise et de la vitesse de renouvellement (cohortes, M17). Avec un taux
de 4,89 \% et une croissance nominale de 2 à 3 \%, chaque point de PIB
de déficit supplémentaire et permanent alourdit la dette de façon
cumulative : l'effet sur le ratio de dette à 20 ans est de l'ordre de 20
points pour un déficit permanent de 1 point, avant effets d'activité.
Inversement, une mesure qui améliore le solde ne réduit le ratio de
dette qu'après avoir compensé l'effet de la baisse d'activité sur le
dénominateur (encadré 1.1).

Dans le régime BCE + prime, une dégradation du solde et de la dette
augmente l'écart OAT-Bund (1,3 point de base par point de dette hors
crise, jusqu'à 11,8 points de base par point de déficit avec la crise),
ce qui renchérit le crédit aux ménages et aux entreprises et amplifie la
dynamique de la dette.

\hypertarget{sec:incidence}{%
\section{Incidence fiscale sur ménages, salariés et
entreprises}\label{sec:incidence}}

L'incidence d'un prélèvement --- qui le supporte réellement --- diffère
de son redevable légal.

\begin{longtable}[]{@{}
  >{\raggedright\arraybackslash}p{(\columnwidth - 4\tabcolsep) * \real{0.3333}}
  >{\raggedright\arraybackslash}p{(\columnwidth - 4\tabcolsep) * \real{0.3333}}
  >{\raggedright\arraybackslash}p{(\columnwidth - 4\tabcolsep) * \real{0.3333}}@{}}
\toprule\noalign{}
\begin{minipage}[b]{\linewidth}\raggedright
Prélèvement
\end{minipage} & \begin{minipage}[b]{\linewidth}\raggedright
Hypothèse du modèle
\end{minipage} & \begin{minipage}[b]{\linewidth}\raggedright
Source, statut
\end{minipage} \\
\midrule\noalign{}
\endhead
\bottomrule\noalign{}
\endlastfoot
TVA, taux normal & répercussion de 95 \% des hausses, 50 \% des baisses
sur les prix à la consommation & Benzarti et al.~(2020) {[}B{]} \\
Taux réduits de TVA & répercussion 50 \% (hausses), 25 \% (baisses) &
idem, extrapolé {[}C{]} \\
Cotisations employeurs & transmission de 50 \% aux salaires à long terme
; le reste sur le coût du travail & Bozio, Breda et Grenet (2019) ;
Carbonnier et al.~(2019) {[}C{]} \\
Impôts de production & 50 \% répercutés sur les prix de la valeur
ajoutée, 20 \% sur les salaires, 30 \% sur les marges & CAE (2019),
hypothèse {[}D{]} \\
Impôt sur les sociétés & 30 \% supportés par les salariés à long terme ;
semi-élasticité des bénéfices déclarés 0,8 (transfert de bénéfices des
groupes, la moitié de l'IS) & Fuest, Peichl et Siegloch (2018) ;
méta-analyses {[}C/B{]} \\
Impôts sur le capital des ménages (PFU, ISF, successions) & érosion
d'assiette selon M13 (fourchette basse, centrale, haute) ; pas d'effet
sur l'investissement des entreprises & littérature {[}C{]} \\
\end{longtable}

\textbf{Exemple chiffré : rendement annoncé et rendement simulé.} Pour
la suppression du prélèvement forfaitaire unique (9 Md€ annoncés), le
module M13 donne un rendement net de 5,3 Md€ en valeur centrale (7,5 Md€
et 1,6 Md€ dans les hypothèses basse et haute de réaction), du fait de
la baisse des dividendes distribués ; pour l'impôt sur la fortune (10
Md€ annoncés), 7,1 Md€ (8,8 et 5,0) ; pour les successions (20 Md€), 18
Md€ (20 et 15). Ces écarts ne mesurent pas la « sincérité » d'un
chiffrage : ils traduisent l'incertitude sur les comportements des
contribuables les plus mobiles.

\hypertarget{sec:cout-capital}{%
\section{Effets sur le coût du capital et
l'investissement}\label{sec:cout-capital}}

L'investissement des entreprises réagit d'abord à la demande
(accélérateur), puis au coût d'usage du capital (taux du crédit, impôt
sur les sociétés, prix de l'investissement) et à la trésorerie. Une
hausse de l'impôt sur les sociétés réduit l'investissement par le coût
d'usage (élasticité −0,18, fragile) et par la trésorerie (0,10) ; une
baisse des impôts de production l'augmente surtout par la demande
induite et la marge. Le modèle n'a pas d'effet de l'incertitude sur
l'investissement, ni d'effet spécifique des impôts de production sur la
localisation des activités : ces deux canaux, invoqués respectivement
contre les programmes jugés risqués et en faveur de la baisse des impôts
de production, ne sont donc pas représentés (chapitre
\ref{sec:comportements}).

\hypertarget{sec:invest-pub-priv}{%
\section{Investissement public et privé}\label{sec:invest-pub-priv}}

L'investissement public a un effet de demande (comme toute dépense) et
un effet d'offre par le capital public (élasticité de 0,08) ;
l'investissement hors administrations (pôle public bancaire) n'a, dans
le modèle, que l'effet de demande, sans effet d'offre ni effet
d'éviction par le crédit : c'est une borne haute de son effet, qui
suppose que les projets financés sont additionnels (ils n'auraient pas
été réalisés sans le pôle public) et que leur financement n'évince pas
d'autres crédits.

\hypertarget{sec:productivite}{%
\section{Productivité et innovation}\label{sec:productivite}}

La productivité tendancielle étant exogène, le modèle ne représente pas
les effets de long terme de l'éducation, de la qualité des services
publics, de l'organisation de l'État, de l'intelligence artificielle
(invoquée par Horizons pour le non-remplacement des agents) ou de la
concurrence. Deux mesures modifient la productivité dans le modèle : la
durée du travail (productivité horaire) et l'investissement public
(capital public). C'est une limite importante pour les programmes qui
revendiquent des gains de productivité de long terme.

\hypertarget{sec:travail-transmission}{%
\section{Marché du travail}\label{sec:travail-transmission}}

L'emploi réagit à la production (à court terme, élasticité de 0,26 puis
convergence) et au coût relatif du travail (\(\sigma\)). Le chômage
dépend en plus de la population active : une réforme qui augmente le
nombre d'actifs (report de l'âge de la retraite) augmente d'abord le
chômage, tant que la demande de travail ne s'est pas ajustée, puis la
fermeture d'offre absorbe progressivement ces actifs supplémentaires ;
une réforme qui réduit le nombre d'actifs (abaissement de l'âge) réduit
d'abord le chômage puis l'emploi. Les hausses du SMIC et du point
d'indice augmentent le coût du travail et les revenus : l'effet net sur
l'emploi dépend de l'élasticité de substitution et de l'effet de
demande.

\hypertarget{sec:inflation-pa}{%
\section{Inflation et pouvoir d'achat}\label{sec:inflation-pa}}

Le pouvoir d'achat (revenu disponible réel) dépend des revenus nominaux
(salaires, emploi, prestations, impôts) et des prix. Une hausse de
salaires augmente le pouvoir d'achat à court terme, mais se diffuse aux
prix (transmission du coût salarial unitaire aux prix de la valeur
ajoutée, puis prix de Leontief) ; une baisse de TVA le relève par les
prix, partiellement (répercussion de 50 \% pour une baisse). Les mesures
d'économie sur les prestations le réduisent directement.

\hypertarget{sec:competitivite}{%
\section{Compétitivité et commerce extérieur}\label{sec:competitivite}}

La compétitivité-prix dépend des prix à l'exportation, eux-mêmes
déterminés par les coûts (salaires, impôts de production, consommations
intermédiaires). Les mesures qui augmentent les coûts (SMIC, point
d'indice indirectement, cotisations) détériorent les exportations
(élasticité-prix −0,60) ; celles qui les réduisent (impôts de
production, cotisations) les améliorent. Les relances de la demande
augmentent les importations (élasticité de 1,47 à court terme), ce qui
dégrade le solde extérieur.

\hypertarget{sec:anticipations-transmission}{%
\section{Anticipations et
comportements}\label{sec:anticipations-transmission}}

Les anticipations du modèle sont adaptatives ; les annonces n'ont pas
d'effet avant leur mise en œuvre. Le modèle ne représente pas : les
effets de confiance (positifs ou négatifs) d'une alternance politique ;
les réactions des marchés financiers au-delà de la prime estimée et des
chocs de scénario ; l'attentisme des entreprises face à l'incertitude
réglementaire ; les comportements d'anticipation des ménages (achats
avant une hausse de TVA, départs à la retraite avant une réforme). Le
chapitre \ref{sec:comportements} discute ces limites.

\hypertarget{sec:contradictoires}{%
\section{Mécanismes contradictoires}\label{sec:contradictoires}}

Une même mesure peut produire des effets de sens opposés selon l'horizon
et le canal. Exemples tirés des programmes :

\begin{longtable}[]{@{}
  >{\raggedright\arraybackslash}p{(\columnwidth - 4\tabcolsep) * \real{0.3333}}
  >{\raggedright\arraybackslash}p{(\columnwidth - 4\tabcolsep) * \real{0.3333}}
  >{\raggedright\arraybackslash}p{(\columnwidth - 4\tabcolsep) * \real{0.3333}}@{}}
\toprule\noalign{}
\begin{minipage}[b]{\linewidth}\raggedright
Mesure
\end{minipage} & \begin{minipage}[b]{\linewidth}\raggedright
Effet favorable
\end{minipage} & \begin{minipage}[b]{\linewidth}\raggedright
Effet défavorable
\end{minipage} \\
\midrule\noalign{}
\endhead
\bottomrule\noalign{}
\endlastfoot
Hausse du point d'indice (+10 \%) & demande (revenus des agents),
pouvoir d'achat & coût budgétaire, diffusion aux salaires privés et aux
prix, compétitivité \\
Abaissement de l'âge de la retraite & revenus des seniors, chômage plus
faible à court terme & population active et potentiel plus faibles,
pensions plus élevées, dette \\
Report de l'âge et allongement de la durée & potentiel, finances
publiques à long terme & chômage des seniors à court terme, pensions
plus faibles pour qui ne reporte pas \\
Baisse de TVA sur l'énergie & prix et pouvoir d'achat & coût budgétaire,
avantage proportionnellement plus élevé en euros pour les gros
consommateurs, signal-prix carbone \\
Baisse de la dépense publique et des prestations & solde public, dette à
moyen terme & demande, emploi, revenus des ménages modestes ; ratio de
dette à court terme \\
Hausse de l'imposition du capital & recettes, redistribution & érosion
d'assiette, éventuels départs, épargne et dividendes \\
Baisse des impôts de production contre baisse des aides & baisse des
coûts de production, prix, marges & perte des bénéficiaires des aides ;
effet net nul si les aides étaient des subventions équivalentes \\
\end{longtable}

\hypertarget{sec:europe}{%
\chapter{Contraintes européennes et risques
institutionnels}\label{sec:europe}}

La France exerce sa politique économique dans un cadre partagé : la
politique monétaire est unique, les règles budgétaires sont communes,
plusieurs domaines (politique commerciale, TVA, aides d'État, libre
circulation) relèvent du droit de l'Union. Ce chapitre distingue
soigneusement trois registres : le \textbf{droit positif} (ce que disent
les textes, statut {[}D{]}), les \textbf{interprétations juridiques}
(qualification d'une mesure au regard de ces textes, qui relève de
l'analyse et peut être contestée), et les \textbf{hypothèses de
scénario} (durées, issues et conséquences financières des conflits
éventuels, statut {[}H{]}). Un conflit institutionnel hypothétique n'est
jamais présenté comme une sanction certaine.

\hypertarget{sec:competences}{%
\section{Compétences nationales et européennes}\label{sec:competences}}

Le traité sur le fonctionnement de l'Union européenne (TFUE) distingue
les compétences exclusives de l'Union (union douanière, politique
commerciale commune, politique monétaire de la zone euro : art. 3), les
compétences partagées (marché intérieur, cohésion, environnement,
énergie\ldots) et les compétences des États membres (fiscalité directe,
protection sociale, organisation de l'État, retraites, droit du travail,
sous réserve du respect des libertés et des règles de concurrence). En
pratique :

\begin{itemize}
\tightlist
\item
  \textbf{Compétence nationale} : impôt sur le revenu, impôt sur la
  fortune, droits de succession, cotisations, prestations, retraites,
  salaire minimum, point d'indice, emploi public, organisation
  territoriale. Le seul canal européen de ces mesures est budgétaire
  (trajectoire de dépenses nettes, procédure de déficit excessif).
\item
  \textbf{Compétence encadrée} : TVA (directive harmonisée), accises sur
  l'énergie (minima de la directive sur la taxation de l'énergie), impôt
  sur les sociétés (taux effectif minimal de 15 \% pour les grands
  groupes), aides aux entreprises (régime des aides d'État), marchés
  publics (non-discrimination).
\item
  \textbf{Compétence exclusive de l'Union} : droits de douane et
  politique commerciale ; un État membre ne peut instaurer seul de
  droits de douane, ni sur les importations extra-européennes (art. 207
  TFUE), ni dans les échanges intra-européens (art. 30 TFUE).
\end{itemize}

\hypertarget{sec:regles-budgetaires-ue}{%
\section{Règles budgétaires
européennes}\label{sec:regles-budgetaires-ue}}

Le cadre de gouvernance économique a été réformé en 2024 par trois
textes entrés en vigueur le 30 avril 2024 : le règlement (UE) 2024/1263
(volet préventif, qui remplace le règlement 1466/97), le règlement (UE)
2024/1264 (volet correctif, qui modifie le règlement 1467/97) et la
directive (UE) 2024/1265 (cadres budgétaires nationaux).

\textbf{Volet préventif.} Chaque État présente un plan budgétaire et
structurel à moyen terme fixant une trajectoire de \textbf{dépenses
nettes} (dépenses hors intérêts, hors mesures nouvelles de recettes et
hors composante conjoncturelle du chômage). Les écarts sont cumulés dans
un compte de contrôle ; une procédure pour déficit excessif fondée sur
la dette peut être ouverte lorsque les écarts dépassent 0,3 point de PIB
par an ou 0,6 point cumulé, la dette étant supérieure à 60 \% du PIB.

\textbf{Volet correctif : la procédure de déficit excessif (PDE).} Le
déficit public ne doit pas dépasser 3 \% du PIB (sauf dépassement
exceptionnel et temporaire) et la dette 60 \% (ou s'en rapprocher à un
rythme satisfaisant). La France est sous procédure de déficit excessif
depuis 2024. Le règlement modifié prévoit :

\begin{itemize}
\tightlist
\item
  que le Conseil fixe un délai de \textbf{six mois au plus} pour engager
  une action suivie d'effets ;
\item
  qu'en cas d'absence d'action suivie d'effets, une \textbf{amende
  pouvant atteindre 0,05 \% du PIB} de l'année précédente \textbf{par
  période de six mois} soit due, payée tous les six mois jusqu'à ce que
  le Conseil constate une action suivie d'effets.
\end{itemize}

Aucune amende n'a jamais été effectivement infligée à un État membre au
titre de la PDE ; l'application de sanctions relève d'une décision
politique du Conseil. C'est pourquoi les amendes ne figurent que dans la
variante « tension », et pour une durée hypothétique.

\hypertarget{sec:tva}{%
\section{Règles fiscales et TVA}\label{sec:tva}}

La directive 2006/112/CE modifiée par la directive (UE) 2022/542 encadre
les taux réduits de TVA : deux taux réduits d'au moins 5 \% au plus,
applicables à 24 catégories au plus de biens et services de l'annexe III
; des taux inférieurs à 5 \% ou nuls pour sept catégories au plus ; et
la \textbf{fin programmée des taux réduits sur les combustibles fossiles
au 1er janvier 2030}. L'électricité et le gaz distribués par réseau
figurent à l'annexe III (taux réduit possible, avec la restriction sur
les combustibles fossiles pour le gaz à partir de 2030) ; les carburants
routiers n'y figurent pas. Le RN a d'ailleurs évoqué une solution de
repli combinant un taux intermédiaire et une baisse d'accise pour les
carburants (selon la presse spécialisée).

\emph{Interprétation retenue.} La baisse de la TVA à 5,5 \% est
réalisable pour l'électricité ; pour les carburants, elle devrait
prendre la forme d'une baisse d'accise (respectant les minima de la
directive sur la taxation de l'énergie) ; pour le gaz, elle serait
remise en cause en 2030. La variante « conforme » suppose que le même
coût budgétaire est atteint par ces voies ; la variante « contentieux »
suppose qu'une moitié de la mesure (la part des carburants, hypothèse
par défaut) n'est pas appliquée et que le reste est reporté de deux ans.

\hypertarget{sec:autres-ue}{%
\section{Autres compétences encadrées}\label{sec:autres-ue}}

\begin{longtable}[]{@{}
  >{\raggedright\arraybackslash}p{(\columnwidth - 6\tabcolsep) * \real{0.2500}}
  >{\raggedright\arraybackslash}p{(\columnwidth - 6\tabcolsep) * \real{0.2500}}
  >{\raggedright\arraybackslash}p{(\columnwidth - 6\tabcolsep) * \real{0.2500}}
  >{\raggedright\arraybackslash}p{(\columnwidth - 6\tabcolsep) * \real{0.2500}}@{}}
\toprule\noalign{}
\begin{minipage}[b]{\linewidth}\raggedright
Domaine
\end{minipage} & \begin{minipage}[b]{\linewidth}\raggedright
Droit positif {[}D{]}
\end{minipage} & \begin{minipage}[b]{\linewidth}\raggedright
Mesures concernées
\end{minipage} & \begin{minipage}[b]{\linewidth}\raggedright
Qualification {[}interprétation{]}
\end{minipage} \\
\midrule\noalign{}
\endhead
\bottomrule\noalign{}
\endlastfoot
Ressources propres de l'Union & décision ressources propres adoptée à
l'unanimité (art. 311 TFUE) & baisse de la contribution (RN, 9 Md€) &
une réduction unilatérale est un manquement ; la voie conforme est la
renégociation du prochain cadre financier pluriannuel (unanimité) \\
Non-discrimination & art. 18 et 45 TFUE ; règlement (CE) 883/2004 ;
directive 2014/24/UE & priorité nationale (RN), économies fondées sur la
nationalité & incompatible pour les ressortissants de l'Union ; à
l'égard des ressortissants de pays tiers, la compatibilité dépend des
directives sectorielles et de la Constitution (non instruit ici) \\
Libre circulation des capitaux & art. 63 TFUE ; sauvegarde vis-à-vis des
pays tiers par le Conseil pour six mois au plus (art. 66) & contrôle des
capitaux (LFI) & incompatible entre États membres hors mesures
prudentielles ou fiscales non discriminatoires \\
Aides d'État & art. 107-108 TFUE & baisse générale des impôts de
production (RN, Horizons), pôle public bancaire et livret A (LFI),
blocage du prix des carburants (LFI) & les mesures générales ne sont pas
des aides ; les mesures sélectives doivent être notifiées ; un
établissement public de crédit doit respecter le principe de
l'investisseur avisé \\
Fiscalité des sociétés & directive (UE) 2022/2523 (taux effectif minimal
de 15 \%) ; conventions fiscales & taxe universelle sur les
multinationales (LFI) & compatible dans son principe ; conformité aux
conventions à vérifier \\
\end{longtable}

\hypertarget{sec:contentieux}{%
\section{Risques de contentieux et délais de mise en
œuvre}\label{sec:contentieux}}

En cas de mesure contraire au droit de l'Union, la Commission peut
engager une \textbf{procédure en manquement} (art. 258 TFUE : mise en
demeure, avis motivé, saisine de la Cour de justice), puis, si l'arrêt
n'est pas exécuté, demander des sanctions financières (art. 260 TFUE) :
une somme forfaitaire et une astreinte. Pour la France, selon la
communication de la Commission de 2025 telle que présentée par la
direction des affaires juridiques du ministère de l'économie, le facteur
de capacité de paiement est de 4,45, la somme forfaitaire minimale de
12,46 M€ et l'astreinte forfaitaire de base de 3 000 € par jour, avant
application des coefficients de gravité et de durée. Ces montants sont
négligeables à l'échelle macroéconomique ; la conséquence principale
d'un contentieux est \textbf{le délai} (une procédure dure souvent
plusieurs années) et \textbf{l'incertitude} (risque de devoir retirer ou
modifier la mesure), non l'amende.

\hypertarget{sec:variantes-ue}{%
\section{Les trois variantes institutionnelles}\label{sec:variantes-ue}}

\begin{longtable}[]{@{}
  >{\raggedright\arraybackslash}p{(\columnwidth - 4\tabcolsep) * \real{0.3333}}
  >{\raggedright\arraybackslash}p{(\columnwidth - 4\tabcolsep) * \real{0.3333}}
  >{\raggedright\arraybackslash}p{(\columnwidth - 4\tabcolsep) * \real{0.3333}}@{}}
\toprule\noalign{}
\begin{minipage}[b]{\linewidth}\raggedright
Variante
\end{minipage} & \begin{minipage}[b]{\linewidth}\raggedright
Contenu
\end{minipage} & \begin{minipage}[b]{\linewidth}\raggedright
Statut
\end{minipage} \\
\midrule\noalign{}
\endhead
\bottomrule\noalign{}
\endlastfoot
\textbf{Conforme} (référence de comparaison) & mesures de niveau 4
retirées pour leur part non conforme (si cette part est connue) ;
mesures de niveaux 1-2 appliquées sous leur forme conforme, au même coût
& {[}interprétation{]} \\
\textbf{Contentieux} & en plus : part non conforme des mesures de
niveaux 1-2 retirée (50 \% pour la TVA sur les carburants, 100 \% pour
la baisse de la contribution européenne) et mesure suspendue pendant le
délai de contentieux (1 à 2 ans) & {[}H{]} \\
\textbf{Tension} & en plus : amendes de la PDE au plafond pendant quatre
semestres (0,1 point de PIB en 2028 et 2029), écart OAT-Bund accru de 50
points de base de 2027 à 2029, coefficients de prime souveraine estimés
en incluant la crise de 2010-2013 & {[}H{]} pour l'amende et le choc ;
{[}E{]} pour les coefficients \\
\end{longtable}

Pour que la comparaison entre variantes isole l'effet des réactions
européennes, les trois variantes sont simulées sous le \textbf{même
régime monétaire} (BCE et prime souveraine endogènes). Le choc de 50
points de base est un ordre de grandeur des variations de l'écart
OAT-Bund lors des épisodes politiques récents ; il n'est pas estimé et
peut être remplacé.

\hypertarget{sec:taux-souverains}{%
\section{Effets possibles sur les taux souverains, le financement et
l'incertitude}\label{sec:taux-souverains}}

La prime de risque est le canal principal par lequel un conflit
institutionnel ou une perte de crédibilité budgétaire affecterait
l'économie. Trois remarques encadrent son interprétation :

\begin{enumerate}
\def\labelenumi{\arabic{enumi}.}
\tightlist
\item
  \textbf{Les coefficients estimés sont des moyennes historiques.} Hors
  crise, la prime réagit peu aux fondamentaux (1,3 point de base par
  point de dette) ; pendant la crise de la zone euro, la réaction au
  déficit a été près de huit fois plus forte. La variante « tension »
  utilise ces coefficients de crise, sans prétendre prévoir une crise.
\item
  \textbf{La BCE dispose d'instruments} (instrument de protection de la
  transmission, programmes d'achats) susceptibles de limiter les écarts
  injustifiés par les fondamentaux, sous conditions (respect du cadre
  budgétaire notamment) ; leur activation n'est pas modélisée. Un pays
  en conflit ouvert avec ce cadre pourrait ne pas en bénéficier.
\item
  \textbf{L'incertitude elle-même a un coût} (attentisme des
  investisseurs, épargne de précaution), non représenté dans le modèle
  (chapitre \ref{sec:comportements}).
\end{enumerate}

\hypertarget{sec:cartographie-ue}{%
\section{Cartographie des incompatibilités du
corpus}\label{sec:cartographie-ue}}

\begin{longtable}[]{@{}
  >{\raggedright\arraybackslash}p{(\columnwidth - 8\tabcolsep) * \real{0.2000}}
  >{\raggedright\arraybackslash}p{(\columnwidth - 8\tabcolsep) * \real{0.2000}}
  >{\raggedright\arraybackslash}p{(\columnwidth - 8\tabcolsep) * \real{0.2000}}
  >{\raggedright\arraybackslash}p{(\columnwidth - 8\tabcolsep) * \real{0.2000}}
  >{\raggedright\arraybackslash}p{(\columnwidth - 8\tabcolsep) * \real{0.2000}}@{}}
\toprule\noalign{}
\begin{minipage}[b]{\linewidth}\raggedright
Mesure
\end{minipage} & \begin{minipage}[b]{\linewidth}\raggedright
Candidat
\end{minipage} & \begin{minipage}[b]{\linewidth}\raggedright
Règle
\end{minipage} & \begin{minipage}[b]{\linewidth}\raggedright
Niveau
\end{minipage} & \begin{minipage}[b]{\linewidth}\raggedright
Conséquence modélisée
\end{minipage} \\
\midrule\noalign{}
\endhead
\bottomrule\noalign{}
\endlastfoot
TVA à 5,5 \% sur les carburants (P10) & Le Pen & directive TVA & 2 &
conforme : même coût par une baisse d'accise ; contentieux : moitié
retirée, report de deux ans \\
Baisse de la contribution à l'UE (P04) & Le Pen & décision ressources
propres & 2 & contentieux : économie non réalisée \\
Économies « immigration » fondées sur la nationalité (P03) & Le Pen &
non-discrimination & 4 & part non conforme inconnue : risque signalé,
montant conservé \\
Priorité nationale (P16) & Le Pen & non-discrimination & 4 & non
quantifiée \\
Contrôle des capitaux (L38) & Mélenchon & art. 63 TFUE & 4 & non
quantifié \\
Taxe sur les multinationales (L08) & Mélenchon & conventions fiscales &
1 & contentieux : report d'un an \\
Pôle public bancaire, livret A (L33, L37) & Mélenchon & aides d'État,
règles prudentielles & 1 & contentieux : investissement hors budget
reporté d'un an \\
Blocage du prix des carburants (L25) & Mélenchon & aides d'État & 1 &
non quantifié \\
Impôts de production (H11, P15) & Philippe, Le Pen & mesure générale & 0
& aucune \\
Règles d'or (P02, H02) & Le Pen, Philippe & cohérentes avec le volet
préventif & 0 & lecture B \\
Trajectoire de déficit & tous & PDE, volet préventif & 3 & variante «
tension » \\
\end{longtable}

\textbf{Limites de l'analyse.} La qualification est sommaire et ne
remplace pas une expertise juridique ; les mesures dont la conformité
dépend de modalités non publiées (priorité nationale à l'égard des
ressortissants de pays tiers, blocage des prix) restent « à instruire ».
Les conditionnalités des fonds européens (cohésion, facilité pour la
reprise) ne sont pas représentées faute de montants vérifiés. Les
renégociations (cadre financier pluriannuel, dérogations de TVA) sont
traitées comme des issues binaires.

\hypertarget{sec:resultats}{%
\chapter{Résultats des simulations comparatives}\label{sec:resultats}}

Ce chapitre présente les résultats disponibles pour chaque programme et
chaque scénario. Il respecte quatre règles : (i) tous les chiffres
proviennent des simulations archivées (\texttt{phase7/resultats}),
générés automatiquement dans les tableaux ; (ii) les résultats de la
\textbf{lecture A} (mesures documentées) et de la \textbf{lecture B}
(objectifs budgétaires) sont présentés séparément et ne doivent pas être
confondus ; (iii) les horizons de 5 et 10 ans sont les horizons
principaux, celui de 20 ans est conditionnel à des hypothèses de long
terme fortes ; (iv) les simulations non exécutées sont signalées comme
telles.

\hypertarget{sec:lire}{%
\section{Lire les résultats}\label{sec:lire}}

Tous les effets sont des \textbf{écarts à une trajectoire de référence}
sans le programme, simulée dans le même environnement et sous le même
régime : un PIB à « +1 \% » à cinq ans signifie que le PIB de 2031 est
supérieur de 1 \% à ce qu'il serait sans le programme, pas que la
croissance de 2031 est de 1 \%. Les niveaux de solde et de dette sont
donnés séparément, car la référence elle-même se dégrade fortement
(chapitre \ref{sec:donnees}) : avec l'OAT à 4,89 \%, la dette de
référence atteint environ 133 \% du PIB en 2031, 156 \% en 2036 et 225
\% en 2046 sous la règle R2 figée.

Trois avertissements spécifiques s'imposent avant toute lecture.

\begin{enumerate}
\def\labelenumi{\arabic{enumi}.}
\tightlist
\item
  \textbf{Les ordres de grandeur du chômage ne sont pas fiables pour les
  grands programmes.} Le chômage du modèle est la différence entre une
  population active en partie exogène et un emploi déterminé par la
  production ; les réformes qui modifient fortement la population active
  (âge de la retraite) ou la demande produisent des écarts de chômage
  beaucoup plus amples que ce que suggérerait la loi d'Okun. Les signes
  sont informatifs, les niveaux beaucoup moins (chapitre
  \ref{sec:robustesse}).
\item
  \textbf{La lecture B mesure l'ampleur d'un effort, pas une trajectoire
  probable.} Comme la référence dérive (effet boule de neige), atteindre
  un déficit de 2,5 ou 3 \% du PIB exige un effort croissant, que le
  modèle traduit en consolidation budgétaire à multiplicateurs élevés.
  Les effets récessifs qui en résultent sont des ordres de grandeur
  conditionnels à la référence et à l'hypothèse de taux.
\item
  \textbf{Les recettes sans canal propre ont été traitées sans effet sur
  la demande} (taxes sur les très grands patrimoines, les
  multinationales, les transactions financières, la fraude ; économies
  sur la contribution européenne) : leurs montants sont des bornes
  hautes de rendement et leurs effets macroéconomiques indirects ne sont
  pas représentés.
\end{enumerate}

\hypertarget{sec:lecture-a}{%
\section{Lecture A : mesures documentées}\label{sec:lecture-a}}

\textbf{Tableau : lecture A (mesures documentées), principaux
indicateurs.} Écarts à la référence du même environnement et du même
régime (R2 figé, taux exogènes) ; horizons 2031, 2036, 2046. Source :
modèle du projet (configuration P4, convention « part constante »),
campagne phase7/campagne.py.

\begin{longtable}[]{@{}llll@{}}
\toprule\noalign{}
candidat --- indicateur & 5 ans & 10 ans & 20 ans \\
\midrule\noalign{}
\endhead
\bottomrule\noalign{}
\endlastfoot
Mélenchon (LFI) --- PIB (\%) & +1,8 & −0,5 & −1,1 \\
Mélenchon (LFI) --- emploi (milliers) & +606,6 & −48,8 & −539,3 \\
Mélenchon (LFI) --- chômage (pt) & −4,5 & −2,8 & −1,4 \\
Mélenchon (LFI) --- pouvoir d'achat (\%) & +3,4 & +4,3 & +5,3 \\
Mélenchon (LFI) --- solde (pt PIB) & +3,3 & +3,7 & +5,6 \\
Mélenchon (LFI) --- dette (pt PIB) & −20,1 & −38,1 & −73,6 \\
Le Pen (RN) --- PIB (\%) & −2,0 & −2,9 & −3,6 \\
Le Pen (RN) --- emploi (milliers) & −440,3 & −754,3 & −1049,7 \\
Le Pen (RN) --- chômage (pt) & −1,8 & −1,6 & −1,1 \\
Le Pen (RN) --- pouvoir d'achat (\%) & −0,9 & −0,6 & −0,0 \\
Le Pen (RN) --- solde (pt PIB) & +0,3 & +0,0 & +0,2 \\
Le Pen (RN) --- dette (pt PIB) & −0,3 & −1,6 & −3,3 \\
Philippe (Horizons) --- PIB (\%) & +1,0 & +2,4 & +4,1 \\
Philippe (Horizons) --- emploi (milliers) & −622,4 & −267,5 & +255,9 \\
Philippe (Horizons) --- chômage (pt) & +2,6 & +2,6 & +1,9 \\
Philippe (Horizons) --- pouvoir d'achat (\%) & −1,1 & −1,0 & −1,9 \\
Philippe (Horizons) --- solde (pt PIB) & −0,6 & −0,5 & −0,7 \\
Philippe (Horizons) --- dette (pt PIB) & +5,3 & +8,4 & +13,2 \\
\end{longtable}

\textbf{Mélenchon (LFI).} Le paquet 2027 combine une relance (point
d'indice, transferts, services publics, SMIC, abaissement de l'âge de la
retraite, investissement hors budget de 100 Md€ par an) et des hausses
de prélèvements concentrées sur le capital, les patrimoines et les
grandes entreprises. À cinq ans, l'activité est supérieure à la
référence, le chômage nettement inférieur, les prix plus élevés
(diffusion du SMIC, du point d'indice et de la moindre offre de
travail), et le solde public amélioré de plus de 3 points de PIB, ce qui
réduit la dette d'une vingtaine de points. Trois éléments expliquent
l'essentiel de ces résultats : (a) les hausses de prélèvements, pour
environ 3 points de PIB, dont une large part est traitée sans effet sur
la demande (hypothèse {[}H{]} la plus favorable au programme) ; (b)
l'investissement hors budget de 100 Md€ par an, qui ajoute de la demande
sans coût budgétaire (borne haute) ; (c) le retour à 62 ans, qui réduit
la population active et le potentiel. Au-delà de cinq ans, la baisse de
l'offre de travail l'emporte : le PIB passe sous la référence à dix et
vingt ans, l'emploi aussi, tandis que les prix restent supérieurs.

\textbf{Le Pen (RN).} Les mesures documentées sont dominées par la
réforme des retraites (62 ans, 42 annuités, carrières longues à 60 ans),
qui réduit la population active et le PIB potentiel de manière
croissante : le PIB est inférieur à la référence d'environ 2 \% à cinq
ans et 3,6 \% à vingt ans, l'emploi d'environ 440 000 à 1 million. Les
économies identifiées (immigration, État, contribution européenne) et
les baisses d'impôts (TVA sur l'énergie, jeunes, familles) se compensent
à peu près sur le solde ; la dette varie peu par rapport à la référence.
Le pouvoir d'achat moyen baisse légèrement à court terme (baisse des
prestations ciblées par les économies « immigration », moindre emploi),
malgré la baisse de la TVA sur l'énergie.

\textbf{Philippe (Horizons).} Les mesures documentées sont celles qui
augmentent l'offre (65 ans et 45 annuités, fin des 35 heures) et
réduisent l'emploi public (250 000 postes), plus le pacte impôts de
production / aides. Le PIB est supérieur à la référence et croissant (de
l'ordre de 1 \% à cinq ans et 4 \% à vingt ans) ; mais l'augmentation de
la population active et de la durée du travail précède celle de la
demande de travail, d'où un chômage plus élevé à cinq et dix ans, qui ne
se résorbe que lentement. Le solde public se dégrade légèrement : les
mesures documentées ne contiennent pas l'essentiel des économies
annoncées (près de 150 Md€), dont la composition n'est pas publiée ;
c'est précisément l'objet de la lecture B.

\hypertarget{sec:decomposition}{%
\section{Décomposition par mesure}\label{sec:decomposition}}

Chaque mesure a été simulée isolément. Le tableau suivant donne les
effets isolés, leur somme et l'écart entre la somme et le paquet complet
(« interactions »), qui provient des non-linéarités (prix, fermeture,
plafonds).

\textbf{Tableau : effets des mesures isolées --- Mélenchon (LFI)}
(régime principal ; chaque mesure simulée seule). Source : modèle du
projet (configuration P4, convention « part constante »), campagne
phase7/campagne.py.

\begin{longtable}[]{@{}
  >{\raggedright\arraybackslash}p{(\columnwidth - 10\tabcolsep) * \real{0.1667}}
  >{\raggedright\arraybackslash}p{(\columnwidth - 10\tabcolsep) * \real{0.1667}}
  >{\raggedright\arraybackslash}p{(\columnwidth - 10\tabcolsep) * \real{0.1667}}
  >{\raggedright\arraybackslash}p{(\columnwidth - 10\tabcolsep) * \real{0.1667}}
  >{\raggedright\arraybackslash}p{(\columnwidth - 10\tabcolsep) * \real{0.1667}}
  >{\raggedright\arraybackslash}p{(\columnwidth - 10\tabcolsep) * \real{0.1667}}@{}}
\toprule\noalign{}
\begin{minipage}[b]{\linewidth}\raggedright
mesure
\end{minipage} & \begin{minipage}[b]{\linewidth}\raggedright
PIB 5 ans (\%)
\end{minipage} & \begin{minipage}[b]{\linewidth}\raggedright
PIB 10 ans (\%)
\end{minipage} & \begin{minipage}[b]{\linewidth}\raggedright
emploi 5 ans (k)
\end{minipage} & \begin{minipage}[b]{\linewidth}\raggedright
solde 5 ans (pt)
\end{minipage} & \begin{minipage}[b]{\linewidth}\raggedright
dette 10 ans (pt)
\end{minipage} \\
\midrule\noalign{}
\endhead
\bottomrule\noalign{}
\endlastfoot
L01 Suppression du prélèvement forfaitaire unique (flat tax & −0,09 &
−0,08 & −16,79 & +0,12 & −0,98 \\
L02 Rétablissement d'un impôt sur la fortune renforcé & −0,11 & −0,09 &
−22,46 & +0,17 & −1,37 \\
L03 Taxe Zucman (patrimoines supérieurs à 100 M€) & +0,00 & +0,00 &
+0,00 & +0,61 & −5,68 \\
L04 Réforme des droits de succession & +0,00 & +0,00 & +0,00 & +0,73 &
−6,82 \\
L06 Rétablissement de la CVAE sur les plus grandes entrepri & −0,02 &
−0,04 & +2,11 & +0,11 & −1,01 \\
L07 Réforme de l'impôt sur les sociétés & −0,04 & −0,07 & −3,40 & +0,05
& −0,29 \\
L08 Taxe universelle sur les multinationales & +0,00 & +0,00 & +0,00 &
+0,20 & −1,89 \\
L09 Renforcement de la taxe sur les services numériques (GA & +0,00 &
+0,00 & +0,00 & +0,16 & −1,48 \\
L10 Taxe sur les superprofits & +0,00 & +0,00 & +0,00 & +0,20 & −1,89 \\
L11 Rabot et suppression des niches fiscales et sociales « & −0,39 &
−0,25 & −109,99 & +0,46 & −4,37 \\
L13 Taxes sur les transactions financières & +0,00 & +0,00 & +0,00 &
+0,20 & −1,89 \\
L14 Taxe sur les rachats d'actions & +0,00 & +0,00 & +0,00 & +0,10 &
−0,95 \\
L15 Taxe sur les crypto-actifs & +0,00 & +0,00 & +0,00 & +0,04 &
−0,38 \\
L16 Lutte contre la fraude fiscale & +0,00 & +0,00 & +0,00 & +0,08 &
−0,76 \\
L20 Dégel et hausse de 10 \% du point d'indice & +0,73 & +0,42 & +175,96
& −0,66 & +5,22 \\
L21 Revalorisation de 10 \% des pensions, allocations famili & +0,11 &
+0,04 & +30,94 & −0,16 & +1,44 \\
L22 Logements d'urgence et réquisition de logements vacants & +0,08 &
+0,03 & +35,70 & −0,02 & −0,03 \\
L23 Garantie d'autonomie pour les jeunes & +0,14 & +0,05 & +36,71 &
−0,19 & +1,68 \\
L24 SMIC à 1 700 € net & −0,00 & +0,02 & −39,31 & −0,16 & +0,70 \\
L26 10 000 postes dans l'Éducation nationale & +0,01 & +0,00 & +5,95 &
−0,00 & −0,00 \\
L27 Gratuité de la cantine scolaire & +0,09 & +0,04 & +42,85 & −0,02 &
−0,03 \\
L28 250 000 places à l'université & +0,09 & +0,03 & +41,66 & −0,02 &
−0,03 \\
L29 Hôpital public (+10 000 postes soignants) & +0,08 & +0,03 & +35,70 &
−0,02 & −0,03 \\
L30 Violences sexistes et sexuelles ; violences faites aux & +0,12 &
+0,05 & +54,75 & −0,02 & −0,04 \\
L31 Indexation des budgets 2027 des services supprimés en 2 & +0,18 &
+0,07 & +83,30 & −0,04 & −0,06 \\
L32 Adaptation au dérèglement climatique & +0,26 & +0,39 & +6,55 & −0,32
& +3,14 \\
L33 Investissements de planification écologique portés par & +1,70 &
+0,65 & +427,99 & +1,21 & −12,75 \\
L34 Abrogation de la réforme de 2023 (retour à 62 ans) & −1,35 & −2,08 &
−211,02 & +0,04 & −0,47 \\
L35 Abrogation des réformes de l'assurance chômage & +0,02 & −0,00 &
+6,21 & −0,07 & +0,53 \\
Σ mesures isolées & +1,59 & −0,79 & +583,41 & +2,79 & −30,49 \\
paquet complet (A) & +1,84 & −0,54 & +606,60 & +3,32 & −38,14 \\
interactions (A − Σ) & +0,25 & +0,25 & +23,19 & +0,53 & −7,65 \\
\end{longtable}

\textbf{Tableau : effets des mesures isolées --- Le Pen (RN)} (régime
principal ; chaque mesure simulée seule). Source : modèle du projet
(configuration P4, convention « part constante »), campagne
phase7/campagne.py.

\begin{longtable}[]{@{}
  >{\raggedright\arraybackslash}p{(\columnwidth - 10\tabcolsep) * \real{0.1667}}
  >{\raggedright\arraybackslash}p{(\columnwidth - 10\tabcolsep) * \real{0.1667}}
  >{\raggedright\arraybackslash}p{(\columnwidth - 10\tabcolsep) * \real{0.1667}}
  >{\raggedright\arraybackslash}p{(\columnwidth - 10\tabcolsep) * \real{0.1667}}
  >{\raggedright\arraybackslash}p{(\columnwidth - 10\tabcolsep) * \real{0.1667}}
  >{\raggedright\arraybackslash}p{(\columnwidth - 10\tabcolsep) * \real{0.1667}}@{}}
\toprule\noalign{}
\begin{minipage}[b]{\linewidth}\raggedright
mesure
\end{minipage} & \begin{minipage}[b]{\linewidth}\raggedright
PIB 5 ans (\%)
\end{minipage} & \begin{minipage}[b]{\linewidth}\raggedright
PIB 10 ans (\%)
\end{minipage} & \begin{minipage}[b]{\linewidth}\raggedright
emploi 5 ans (k)
\end{minipage} & \begin{minipage}[b]{\linewidth}\raggedright
solde 5 ans (pt)
\end{minipage} & \begin{minipage}[b]{\linewidth}\raggedright
dette 10 ans (pt)
\end{minipage} \\
\midrule\noalign{}
\endhead
\bottomrule\noalign{}
\endlastfoot
P03 Économies sur le « coût de l'immigration » & −0,26 & −0,18 & −57,24
& +0,27 & −2,16 \\
P04 Baisse de la contribution au budget de l'Union & +0,00 & +0,00 &
+0,00 & +0,36 & −3,41 \\
P05 Réorganisation de l'État (centralisation de la TVA, tra & −0,52 &
−0,20 & −246,51 & +0,11 & +0,18 \\
P08 Âge légal ramené à 62 ans (60 ans pour les carrières co & −1,39 &
−2,67 & −193,95 & +0,10 & −1,69 \\
P10 TVA ramenée de 20 \% à 5,5 \% sur l'électricité, le gaz, & +0,09 &
+0,02 & +26,49 & −0,29 & +2,98 \\
P11 Exonération d'impôt sur le revenu des moins de 30 ans ( & +0,02 &
−0,02 & +7,86 & −0,05 & +0,50 \\
P12 Permis de conduire gratuit (18 ans, casier vierge) ; pr & +0,00 &
−0,02 & +5,17 & −0,03 & +0,40 \\
P13 Part fiscale pleine dès le deuxième enfant & +0,04 & −0,00 & +13,62
& −0,08 & +0,74 \\
Σ mesures isolées & −2,02 & −3,07 & −444,57 & +0,39 & −2,46 \\
paquet complet (A) & −1,95 & −2,89 & −440,35 & +0,33 & −1,60 \\
interactions (A − Σ) & +0,07 & +0,18 & +4,22 & −0,06 & +0,85 \\
\end{longtable}

\textbf{Tableau : effets des mesures isolées --- Philippe (Horizons)}
(régime principal ; chaque mesure simulée seule). Source : modèle du
projet (configuration P4, convention « part constante »), campagne
phase7/campagne.py.

\begin{longtable}[]{@{}
  >{\raggedright\arraybackslash}p{(\columnwidth - 10\tabcolsep) * \real{0.1667}}
  >{\raggedright\arraybackslash}p{(\columnwidth - 10\tabcolsep) * \real{0.1667}}
  >{\raggedright\arraybackslash}p{(\columnwidth - 10\tabcolsep) * \real{0.1667}}
  >{\raggedright\arraybackslash}p{(\columnwidth - 10\tabcolsep) * \real{0.1667}}
  >{\raggedright\arraybackslash}p{(\columnwidth - 10\tabcolsep) * \real{0.1667}}
  >{\raggedright\arraybackslash}p{(\columnwidth - 10\tabcolsep) * \real{0.1667}}@{}}
\toprule\noalign{}
\begin{minipage}[b]{\linewidth}\raggedright
mesure
\end{minipage} & \begin{minipage}[b]{\linewidth}\raggedright
PIB 5 ans (\%)
\end{minipage} & \begin{minipage}[b]{\linewidth}\raggedright
PIB 10 ans (\%)
\end{minipage} & \begin{minipage}[b]{\linewidth}\raggedright
emploi 5 ans (k)
\end{minipage} & \begin{minipage}[b]{\linewidth}\raggedright
solde 5 ans (pt)
\end{minipage} & \begin{minipage}[b]{\linewidth}\raggedright
dette 10 ans (pt)
\end{minipage} \\
\midrule\noalign{}
\endhead
\bottomrule\noalign{}
\endlastfoot
H03 250 000 postes de fonctionnaires en moins (un départ su & −0,48 &
−0,21 & −169,99 & +0,08 & +0,36 \\
H05 Âge de départ à 65 ans et 45 annuités ; départ précoce & +0,21 &
+0,83 & +32,18 & +0,06 & +0,37 \\
H06 Fin des 35 heures ; deux heures supplémentaires hebdoma & +1,59 &
+1,96 & −336,35 & −0,28 & +3,70 \\
H08 Baisse de 30 \% du nombre d'allocataires du RSA (accompa & +0,03 &
+0,08 & −0,19 & +0,03 & +0,00 \\
H11 Baisse de 50 Md€ des impôts de production financée par & −0,45 &
−0,29 & −151,66 & −0,39 & +3,30 \\
H12 Fonds unique d'adaptation au changement climatique & +0,05 & +0,08 &
+1,34 & −0,06 & +0,63 \\
Σ mesures isolées & +0,96 & +2,46 & −624,68 & −0,56 & +8,36 \\
paquet complet (A) & +0,96 & +2,45 & −622,43 & −0,57 & +8,42 \\
interactions (A − Σ) & +0,01 & −0,01 & +2,25 & −0,01 & +0,05 \\
\end{longtable}

Les décompositions montrent que \textbf{les mesures qui modifient
l'offre de travail sont les plus structurantes pour l'activité à moyen
terme}. Pour LFI et le RN, c'est la réforme des retraites (abaissement
de l'âge et de la durée), qui explique l'essentiel de l'effet négatif à
dix ans sur le PIB. Pour Horizons, c'est d'abord la fin des 35 heures
(effet sur le PIB le plus important à cinq et dix ans, mais emploi en
baisse, la production supplémentaire étant obtenue par des heures plus
longues), puis le relèvement de l'âge de départ, dont l'effet croît avec
le temps. Pour LFI, l'investissement hors budget est la mesure la plus
favorable à l'activité à cinq ans, et les taxes sur le patrimoine les
premières par l'effet sur le solde ; pour le RN, la baisse de la
contribution européenne et les économies sur l'immigration portent
l'essentiel de l'amélioration du solde. Pour Horizons, les suppressions
de postes publics et le pacte impôts de production / aides réduisent
l'activité à court terme dans sa traduction centrale (la moitié des
aides supprimées étant des allègements de cotisations, la hausse du coût
du travail l'emporte sur la baisse des impôts de production). Les
interactions entre mesures (paquet complet moins somme des effets
isolés) sont faibles pour le RN et Horizons ; elles sont plus
importantes pour LFI, dont les mesures se renforcent mutuellement par
les prix et la fermeture.

\hypertarget{sec:lecture-b}{%
\section{Lecture B : objectifs budgétaires}\label{sec:lecture-b}}

La lecture B ajoute aux mesures documentées l'effort nécessaire pour
suivre la trajectoire de solde annoncée. Elle répond à la question :
\emph{si le programme atteignait l'objectif de déficit qu'il affiche,
avec la composition de l'effort qu'il revendique, quelles en seraient
les conséquences macroéconomiques selon ce modèle ?}

\textbf{Tableau : lecture B (objectifs budgétaires).} Mesures
documentées + effort supplémentaire pour suivre la trajectoire de solde
annoncée, réparti selon les compositions du tableau du chapitre
\ref{sec:prog-regimes}. Source : modèle du projet (configuration P4,
convention « part constante »), campagne phase7/campagne.py.

\begin{longtable}[]{@{}
  >{\raggedright\arraybackslash}p{(\columnwidth - 6\tabcolsep) * \real{0.2500}}
  >{\raggedright\arraybackslash}p{(\columnwidth - 6\tabcolsep) * \real{0.2500}}
  >{\raggedright\arraybackslash}p{(\columnwidth - 6\tabcolsep) * \real{0.2500}}
  >{\raggedright\arraybackslash}p{(\columnwidth - 6\tabcolsep) * \real{0.2500}}@{}}
\toprule\noalign{}
\begin{minipage}[b]{\linewidth}\raggedright
candidat --- indicateur
\end{minipage} & \begin{minipage}[b]{\linewidth}\raggedright
5 ans
\end{minipage} & \begin{minipage}[b]{\linewidth}\raggedright
10 ans
\end{minipage} & \begin{minipage}[b]{\linewidth}\raggedright
20 ans
\end{minipage} \\
\midrule\noalign{}
\endhead
\bottomrule\noalign{}
\endlastfoot
Mélenchon (LFI) --- PIB (\%) & +2,7 & −1,2 & −2,7 \\
Mélenchon (LFI) --- emploi (milliers) & +775,9 & −77,4 & −962,3 \\
Mélenchon (LFI) --- chômage (pt) & −5,0 & −2,7 & −0,1 \\
Mélenchon (LFI) --- pouvoir d'achat (\%) & +5,6 & +3,5 & −2,0 \\
Mélenchon (LFI) --- solde (pt PIB) & +2,8 & +4,1 & +7,9 \\
Mélenchon (LFI) --- dette (pt PIB) & −18,7 & −36,1 & −83,3 \\
Mélenchon (LFI) --- effort objectif (pt PIB) & −1,2 & +0,5 & +3,9 \\
Mélenchon (LFI) --- niveau solde (\% PIB) & −4,2 & −4,0 & −4,1 \\
Mélenchon (LFI) --- niveau dette (\% PIB) & +114,1 & +119,7 & +141,4 \\
Le Pen (RN) --- PIB (\%) & −9,6 & −11,1 & −6,6 \\
Le Pen (RN) --- emploi (milliers) & −2560,6 & −4066,8 & −3091,8 \\
Le Pen (RN) --- chômage (pt) & +4,7 & +8,5 & +5,3 \\
Le Pen (RN) --- pouvoir d'achat (\%) & −8,3 & −13,4 & −17,2 \\
Le Pen (RN) --- solde (pt PIB) & +3,2 & +5,3 & +8,9 \\
Le Pen (RN) --- dette (pt PIB) & +2,8 & −5,6 & −65,7 \\
Le Pen (RN) --- effort objectif (pt PIB) & +7,8 & +13,3 & +15,1 \\
Le Pen (RN) --- niveau solde (\% PIB) & −3,8 & −2,9 & −3,1 \\
Le Pen (RN) --- niveau dette (\% PIB) & +135,6 & +150,1 & +159,0 \\
Philippe (Horizons) --- PIB (\%) & −5,2 & −3,5 & +1,1 \\
Philippe (Horizons) --- emploi (milliers) & −2114,9 & −2450,6 &
−1155,8 \\
Philippe (Horizons) --- chômage (pt) & +6,9 & +8,8 & +6,0 \\
Philippe (Horizons) --- pouvoir d'achat (\%) & −11,1 & −16,7 & −21,6 \\
Philippe (Horizons) --- solde (pt PIB) & +3,3 & +5,7 & +9,4 \\
Philippe (Horizons) --- dette (pt PIB) & +3,4 & −11,6 & −77,1 \\
Philippe (Horizons) --- effort objectif (pt PIB) & +8,0 & +12,1 &
+13,7 \\
Philippe (Horizons) --- niveau solde (\% PIB) & −3,6 & −2,4 & −2,6 \\
Philippe (Horizons) --- niveau dette (\% PIB) & +136,2 & +144,2 &
+147,5 \\
\end{longtable}

\textbf{Tableau : atteinte des cibles de la lecture B} (solde simulé −
cible annoncée).

\begin{longtable}[]{@{}
  >{\raggedright\arraybackslash}p{(\columnwidth - 6\tabcolsep) * \real{0.2500}}
  >{\raggedright\arraybackslash}p{(\columnwidth - 6\tabcolsep) * \real{0.2500}}
  >{\raggedright\arraybackslash}p{(\columnwidth - 6\tabcolsep) * \real{0.2500}}
  >{\raggedright\arraybackslash}p{(\columnwidth - 6\tabcolsep) * \real{0.2500}}@{}}
\toprule\noalign{}
\begin{minipage}[b]{\linewidth}\raggedright
candidat
\end{minipage} & \begin{minipage}[b]{\linewidth}\raggedright
cible annoncée
\end{minipage} & \begin{minipage}[b]{\linewidth}\raggedright
écart moyen à la cible 2028-2032 (pt)
\end{minipage} & \begin{minipage}[b]{\linewidth}\raggedright
écart maximal (pt)
\end{minipage} \\
\midrule\noalign{}
\endhead
\bottomrule\noalign{}
\endlastfoot
Mélenchon (LFI) & réduction du déficit de 24,7 Md€ en 2027 (mesures)
puis stabilisation du ratio dette/PIB en un quinquennat & −0,00 &
+0,06 \\
Le Pen (RN) & déficit à 3 \% du PIB en 2030 ; dette de 121 \% (2027) à
112 \% (2032) ; 140 Md€ d'économies & −0,52 & −0,91 \\
Philippe (Horizons) & déficit sous 2,5 \% du PIB en fin de quinquennat ;
moitié de l'effort sur les dépenses sociales, un quart sur l'État et les
collectivités, un quart « sur la croissance » ; aucun nouvel impôt &
−0,34 & −0,45 \\
\end{longtable}

\textbf{Tableau : niveaux de solde et de dette publics (\% du PIB),
lectures A et B, sous deux hypothèses de taux longs} (OAT à 4,89 \%,
dernier point observé, ou à 3,82 \%, moyenne des douze derniers mois).
Source : modèle du projet (configuration P4, convention « part constante
»), campagne phase7/campagne.py.

\begin{longtable}[]{@{}
  >{\raggedright\arraybackslash}p{(\columnwidth - 12\tabcolsep) * \real{0.1429}}
  >{\raggedright\arraybackslash}p{(\columnwidth - 12\tabcolsep) * \real{0.1429}}
  >{\raggedright\arraybackslash}p{(\columnwidth - 12\tabcolsep) * \real{0.1429}}
  >{\raggedright\arraybackslash}p{(\columnwidth - 12\tabcolsep) * \real{0.1429}}
  >{\raggedright\arraybackslash}p{(\columnwidth - 12\tabcolsep) * \real{0.1429}}
  >{\raggedright\arraybackslash}p{(\columnwidth - 12\tabcolsep) * \real{0.1429}}
  >{\raggedright\arraybackslash}p{(\columnwidth - 12\tabcolsep) * \real{0.1429}}@{}}
\toprule\noalign{}
\begin{minipage}[b]{\linewidth}\raggedright
candidat --- lecture
\end{minipage} & \begin{minipage}[b]{\linewidth}\raggedright
solde 5 ans (\% PIB)
\end{minipage} & \begin{minipage}[b]{\linewidth}\raggedright
solde 10 ans (\% PIB)
\end{minipage} & \begin{minipage}[b]{\linewidth}\raggedright
dette 5 ans (\% PIB)
\end{minipage} & \begin{minipage}[b]{\linewidth}\raggedright
dette 10 ans (\% PIB)
\end{minipage} & \begin{minipage}[b]{\linewidth}\raggedright
dette 20 ans (\% PIB)
\end{minipage} & \begin{minipage}[b]{\linewidth}\raggedright
effort B 10 ans (pt)
\end{minipage} \\
\midrule\noalign{}
\endhead
\bottomrule\noalign{}
\endlastfoot
Mélenchon (LFI) --- A & −3,7 & −4,4 & 112,7 & 117,7 & 151,0 & 0,0 \\
Mélenchon (LFI) --- B & −4,2 & −4,0 & 114,1 & 119,7 & 141,4 & 0,5 \\
Mélenchon (LFI) --- A, OAT 3,82 \% & −2,8 & −3,3 & 111,1 & 111,3 & 127,6
& 0,0 \\
Mélenchon (LFI) --- B, OAT 3,82 \% & −4,1 & −4,0 & 113,2 & 116,7 & 135,9
& −2,0 \\
Le Pen (RN) --- A & −6,7 & −8,1 & 132,5 & 154,2 & 221,3 & 0,0 \\
Le Pen (RN) --- B & −3,8 & −2,9 & 135,6 & 150,1 & 159,0 & 13,3 \\
Le Pen (RN) --- A, OAT 3,82 \% & −5,7 & −6,6 & 130,7 & 146,3 & 191,2 &
0,0 \\
Le Pen (RN) --- B, OAT 3,82 \% & −3,1 & −2,9 & 133,5 & 141,8 & 146,1 &
9,8 \\
Philippe (Horizons) --- A & −7,5 & −8,7 & 138,1 & 164,2 & 237,8 & 0,0 \\
Philippe (Horizons) --- B & −3,6 & −2,4 & 136,2 & 144,2 & 147,5 &
12,1 \\
Philippe (Horizons) --- A, OAT 3,82 \% & −6,6 & −7,1 & 136,2 & 156,0 &
206,2 & 0,0 \\
Philippe (Horizons) --- B, OAT 3,82 \% & −3,1 & −2,4 & 134,3 & 138,8 &
139,5 & 9,4 \\
\end{longtable}

\textbf{Lecture.} L'effort supplémentaire nécessaire est très différent
selon les candidats, parce que leurs mesures documentées ne partent pas
du même point et que leurs cibles diffèrent :

\begin{itemize}
\tightlist
\item
  pour \textbf{LFI}, dont les mesures documentées améliorent déjà
  fortement le solde, la cible de stabilisation de la dette en 2032 est
  atteinte sans effort supplémentaire à court terme (l'effort B est même
  négatif jusqu'à la fin du quinquennat, ce qui revient à autoriser des
  dépenses ou des baisses d'impôts supplémentaires) ; il devient positif
  ensuite, pour compenser la dérive de la référence ;
\item
  pour le \textbf{RN} et \textbf{Horizons}, l'effort non ventilé est
  considérable : plusieurs points de PIB dès cinq ans, plus de dix
  points à dix ans dans l'hypothèse de taux de 4,89 \%. Il s'agit d'une
  consolidation budgétaire d'ampleur historique, dont le modèle prévoit
  des effets récessifs importants (PIB inférieur de plusieurs points à
  la référence, chômage nettement plus élevé, pouvoir d'achat en baisse)
  ;
\item
  l'\textbf{hypothèse de taux} modifie fortement l'effort : avec une OAT
  à 3,82 \%, la dérive de la référence est moindre et l'effort
  nécessaire plus faible.
\end{itemize}

Ces résultats ne sont pas des prévisions : ils mesurent l'écart entre
les objectifs annoncés et la dynamique spontanée de la dette dans ce
modèle. Ils dépendent de la référence (convention de dépense, taux), des
multiplicateurs (élevés dans le régime principal) et de la composition
supposée de l'effort. Un gouvernement qui annoncerait un tel effort
bénéficierait peut-être d'une baisse de la prime de risque (non
représentée dans le régime principal) ; il pourrait aussi ne pas le
réaliser. La comparaison A/B rend visible le fait que \textbf{les
mesures documentées du RN et d'Horizons ne représentent qu'une faible
part de l'effort nécessaire à leurs objectifs}, alors que les mesures de
LFI dépassent, à court terme, l'effort requis par son objectif --- sous
les hypothèses favorables de traduction de ses recettes.

\hypertarget{sec:effets-regimes}{%
\section{Effets des régimes budgétaires et
monétaires}\label{sec:effets-regimes}}

\textbf{Tableau : lecture A sous les dix régimes de robustesse}
(fourchette des effets, régimes extrêmes, stabilité du signe). Source :
modèle du projet (configuration P4, convention « part constante »),
campagne phase7/campagne.py.

\begin{longtable}[]{@{}
  >{\raggedright\arraybackslash}p{(\columnwidth - 12\tabcolsep) * \real{0.1429}}
  >{\raggedright\arraybackslash}p{(\columnwidth - 12\tabcolsep) * \real{0.1429}}
  >{\raggedright\arraybackslash}p{(\columnwidth - 12\tabcolsep) * \real{0.1429}}
  >{\raggedright\arraybackslash}p{(\columnwidth - 12\tabcolsep) * \real{0.1429}}
  >{\raggedright\arraybackslash}p{(\columnwidth - 12\tabcolsep) * \real{0.1429}}
  >{\raggedright\arraybackslash}p{(\columnwidth - 12\tabcolsep) * \real{0.1429}}
  >{\raggedright\arraybackslash}p{(\columnwidth - 12\tabcolsep) * \real{0.1429}}@{}}
\toprule\noalign{}
\begin{minipage}[b]{\linewidth}\raggedright
candidat --- indicateur --- horizon
\end{minipage} & \begin{minipage}[b]{\linewidth}\raggedright
principal
\end{minipage} & \begin{minipage}[b]{\linewidth}\raggedright
minimum
\end{minipage} & \begin{minipage}[b]{\linewidth}\raggedright
maximum
\end{minipage} & \begin{minipage}[b]{\linewidth}\raggedright
régime du minimum
\end{minipage} & \begin{minipage}[b]{\linewidth}\raggedright
régime du maximum
\end{minipage} & \begin{minipage}[b]{\linewidth}\raggedright
signe robuste
\end{minipage} \\
\midrule\noalign{}
\endhead
\bottomrule\noalign{}
\endlastfoot
Mélenchon (LFI) --- PIB (\%) --- 5 ans & +1,84 & +0,96 & +2,84 & R2,
fermeture rapide & R2, fermeture lente & oui \\
Mélenchon (LFI) --- PIB (\%) --- 10 ans & −0,54 & −1,47 & +0,91 & R2,
fermeture rapide & R2, fermeture lente & non \\
Mélenchon (LFI) --- PIB (\%) --- 20 ans & −1,07 & −2,13 & +0,40 & R2
active symétrique & R2, fermeture lente & non \\
Mélenchon (LFI) --- chômage (pt) --- 5 ans & −4,51 & −4,96 & −4,08 & R2,
fermeture lente & R2, fermeture rapide & oui \\
Mélenchon (LFI) --- chômage (pt) --- 10 ans & −2,78 & −3,79 & −2,09 &
R2, fermeture lente & R2, fermeture rapide & oui \\
Mélenchon (LFI) --- chômage (pt) --- 20 ans & −1,39 & −2,37 & −0,28 &
R2, fermeture lente & R2 active symétrique & oui \\
Mélenchon (LFI) --- dette (pt PIB) --- 5 ans & −20,11 & −26,85 & −16,63
& R2, ancrage faible & R2, ancrage fort & oui \\
Mélenchon (LFI) --- dette (pt PIB) --- 10 ans & −38,14 & −56,96 & −27,86
& R2, ancrage faible & R2, ancrage fort & oui \\
Mélenchon (LFI) --- dette (pt PIB) --- 20 ans & −73,61 & −118,84 &
−48,09 & R2, ancrage faible & R2, ancrage fort & oui \\
Le Pen (RN) --- PIB (\%) --- 5 ans & −1,95 & −2,03 & −1,27 & R2 active
symétrique & R2, fermeture sur le chômage & oui \\
Le Pen (RN) --- PIB (\%) --- 10 ans & −2,89 & −3,38 & −2,66 & R2 active
symétrique & R2, fermeture lente & oui \\
Le Pen (RN) --- PIB (\%) --- 20 ans & −3,64 & −4,50 & −3,21 & R2 active
symétrique & R2, fermeture lente & oui \\
Le Pen (RN) --- chômage (pt) --- 5 ans & −1,84 & −2,20 & −1,82 & R2,
fermeture sur le chômage & R2, fermeture rapide & oui \\
Le Pen (RN) --- chômage (pt) --- 10 ans & −1,65 & −1,78 & −1,30 & R2,
fermeture lente & R2 active symétrique & oui \\
Le Pen (RN) --- chômage (pt) --- 20 ans & −1,08 & −1,39 & −0,15 & R2,
fermeture lente & R2 active symétrique & oui \\
Le Pen (RN) --- dette (pt PIB) --- 5 ans & −0,34 & −1,91 & +0,10 & R2,
fermeture sur le chômage & R2, ancrage fort & non \\
Le Pen (RN) --- dette (pt PIB) --- 10 ans & −1,60 & −6,21 & +0,74 & R2,
ancrage faible & R2, ancrage fort & non \\
Le Pen (RN) --- dette (pt PIB) --- 20 ans & −3,34 & −18,03 & +4,56 & R2,
ancrage faible & R2, ancrage fort & non \\
Philippe (Horizons) --- PIB (\%) --- 5 ans & +0,96 & +0,19 & +2,09 & R2,
fermeture sur le chômage & R2 active symétrique & oui \\
Philippe (Horizons) --- PIB (\%) --- 10 ans & +2,45 & +1,67 & +4,47 &
R2, fermeture lente & R2 active symétrique & oui \\
Philippe (Horizons) --- PIB (\%) --- 20 ans & +4,08 & +2,97 & +5,55 &
R2, fermeture lente & R2 active symétrique & oui \\
Philippe (Horizons) --- chômage (pt) --- 5 ans & +2,56 & +1,70 & +3,04 &
R2 active symétrique & R2, fermeture sur le chômage & oui \\
Philippe (Horizons) --- chômage (pt) --- 10 ans & +2,56 & +0,69 & +3,05
& R2 active symétrique & R2, fermeture lente & oui \\
Philippe (Horizons) --- chômage (pt) --- 20 ans & +1,93 & +0,07 & +2,70
& R2 active symétrique & R2, fermeture lente & oui \\
Philippe (Horizons) --- dette (pt PIB) --- 5 ans & +5,28 & +3,91 & +8,01
& R2, ancrage fort & R2, ancrage faible & oui \\
Philippe (Horizons) --- dette (pt PIB) --- 10 ans & +8,42 & +3,93 &
+17,50 & R2, ancrage fort & R2, ancrage faible & oui \\
Philippe (Horizons) --- dette (pt PIB) --- 20 ans & +13,21 & −1,48 &
+42,78 & R2, ancrage fort & R2, ancrage faible & non \\
\end{longtable}

Les régimes modifient surtout l'ampleur des effets sur la dette (règle
active, réaction de la BCE et prime souveraine) et sur le chômage
(vitesse de fermeture), beaucoup moins le signe des effets sur le PIB.
Le régime « sans règle budgétaire » (R0) et les régimes « ancrage faible
» ou « BCE + prime » amplifient les effets sur la dette ; « fermeture
sur le chômage » et « R2 active symétrique » modifient la dynamique du
chômage.

\hypertarget{sec:effets-ue-moe}{%
\section{Effets des réactions européennes et de la mise en
œuvre}\label{sec:effets-ue-moe}}

\textbf{Tableau : variantes européennes (régime BCE + prime pour les
trois) et mise en œuvre partielle (50 \%) ou retardée (2 ans) ; pour
LFI, additionnalité de l'investissement hors budget (100 \% dans la
lecture A ; régime principal)} (écarts à la référence du régime
correspondant). Source : modèle du projet (configuration P4, convention
« part constante »), campagne phase7/campagne.py.

\begin{longtable}[]{@{}
  >{\raggedright\arraybackslash}p{(\columnwidth - 12\tabcolsep) * \real{0.1429}}
  >{\raggedright\arraybackslash}p{(\columnwidth - 12\tabcolsep) * \real{0.1429}}
  >{\raggedright\arraybackslash}p{(\columnwidth - 12\tabcolsep) * \real{0.1429}}
  >{\raggedright\arraybackslash}p{(\columnwidth - 12\tabcolsep) * \real{0.1429}}
  >{\raggedright\arraybackslash}p{(\columnwidth - 12\tabcolsep) * \real{0.1429}}
  >{\raggedright\arraybackslash}p{(\columnwidth - 12\tabcolsep) * \real{0.1429}}
  >{\raggedright\arraybackslash}p{(\columnwidth - 12\tabcolsep) * \real{0.1429}}@{}}
\toprule\noalign{}
\begin{minipage}[b]{\linewidth}\raggedright
candidat --- variante
\end{minipage} & \begin{minipage}[b]{\linewidth}\raggedright
PIB 5 ans (\%)
\end{minipage} & \begin{minipage}[b]{\linewidth}\raggedright
PIB 10 ans (\%)
\end{minipage} & \begin{minipage}[b]{\linewidth}\raggedright
chômage 5 ans (pt)
\end{minipage} & \begin{minipage}[b]{\linewidth}\raggedright
solde 5 ans (pt)
\end{minipage} & \begin{minipage}[b]{\linewidth}\raggedright
dette 10 ans (pt)
\end{minipage} & \begin{minipage}[b]{\linewidth}\raggedright
dette 20 ans (pt)
\end{minipage} \\
\midrule\noalign{}
\endhead
\bottomrule\noalign{}
\endlastfoot
Mélenchon (LFI) --- conforme & +1,79 & −0,61 & −4,49 & +3,18 & −36,90 &
−74,54 \\
Mélenchon (LFI) --- contentieux & +2,03 & −0,45 & −4,36 & +3,00 & −35,07
& −72,28 \\
Mélenchon (LFI) --- tension & +2,04 & −0,45 & −4,36 & +2,88 & −34,45 &
−73,52 \\
Mélenchon (LFI) --- partielle 50 & +0,96 & −0,29 & −2,27 & +1,52 &
−16,86 & −31,91 \\
Mélenchon (LFI) --- retard 2ans & +3,00 & +0,10 & −4,33 & +2,72 & −31,77
& −67,22 \\
Mélenchon (LFI) --- investissement hors APU additionnel à 50 \% & +0,93
& −0,89 & −3,82 & +2,61 & −30,13 & −55,98 \\
Mélenchon (LFI) --- investissement hors APU additionnel à 25 \% & +0,47
& −1,06 & −3,47 & +2,24 & −26,02 & −47,25 \\
Mélenchon (LFI) --- investissement hors APU additionnel à 0 \% & +0,02 &
−1,24 & −3,12 & +1,87 & −21,85 & −38,50 \\
Le Pen (RN) --- conforme & −1,93 & −2,84 & −1,84 & +0,40 & −2,79 &
−11,57 \\
Le Pen (RN) --- contentieux & −1,99 & −2,86 & −1,79 & +0,17 & −0,94 &
−6,83 \\
Le Pen (RN) --- tension & −1,99 & −2,86 & −1,79 & −0,05 & +0,84 &
−3,76 \\
Le Pen (RN) --- partielle 50 & −0,89 & −1,32 & −0,82 & +0,17 & −0,58 &
−1,10 \\
Le Pen (RN) --- retard 2ans & −1,41 & −2,66 & −1,73 & +0,39 & −1,10 &
−3,00 \\
Philippe (Horizons) --- conforme & +0,97 & +2,43 & +2,56 & −0,56 & +8,75
& +20,56 \\
Philippe (Horizons) --- contentieux & +0,97 & +2,43 & +2,56 & −0,56 &
+8,75 & +20,56 \\
Philippe (Horizons) --- tension & +0,97 & +2,43 & +2,56 & −0,84 & +11,17
& +26,38 \\
Philippe (Horizons) --- partielle 50 & +0,63 & +1,58 & +1,56 & −0,24 &
+4,36 & +6,67 \\
Philippe (Horizons) --- retard 2ans & +0,38 & +2,01 & +2,53 & −0,73 &
+8,05 & +13,96 \\
\end{longtable}

\textbf{Réactions européennes.} Pour LFI, les variantes « contentieux »
et « tension » retardent d'un an l'investissement hors budget et la taxe
sur les multinationales (mesures de niveau 1), ce qui déplace l'effet
d'activité sans en changer l'ordre de grandeur. Pour le RN, la variante
« contentieux » supprime l'économie sur la contribution européenne et
réduit de moitié la baisse de TVA sur l'énergie : le solde est un peu
moins favorable et la dette plus élevée. Les amendes de la procédure de
déficit excessif (0,1 point de PIB pendant deux ans) et le choc de prime
de 50 points de base de la variante « tension » ont des effets modestes
sur la dette à cinq ans. Pour Horizons, dont les mesures documentées ne
posent pas de problème de compatibilité, « conforme » et « contentieux »
sont identiques ; seules les amendes et le choc de prime distinguent la
variante « tension ».

\textbf{Additionnalité de l'investissement hors budget (LFI).} La
lecture A suppose que les 100 Md€ annuels d'investissement portés par le
pôle public bancaire s'ajoutent entièrement à l'investissement qui
aurait eu lieu sans eux (borne haute). Les variantes à 50 \%, 25 \% et 0
\% montrent que \textbf{cette hypothèse détermine à elle seule le signe
de l'effet d'activité de court terme du paquet LFI} : l'écart de PIB à
cinq ans, positif dans la lecture A, devient proche de zéro sans
additionnalité, tandis que l'amélioration du solde et la baisse du ratio
de dette sont réduites (le PIB nominal plus élevé et les recettes
induites expliquaient une partie de la baisse du ratio). Le tableau
précédent donne les valeurs ; ni le financement de ces investissements
(épargne réglementée, crédit) ni leur éventuelle requalification en
dette publique ne sont représentés.

\textbf{Mise en œuvre partielle ou retardée.} Une mise en œuvre à 50 \%
divise approximativement les effets par deux (les effets sont presque
linéaires à cette échelle) ; un retard de deux ans décale les effets de
court terme sans changer les effets à vingt ans de manière importante.

\hypertarget{sec:effets-environnements}{%
\section{Environnements alternatifs}\label{sec:effets-environnements}}

\textbf{Tableau : lecture A dans les environnements alternatifs} (σ =
0,3 et 0,6 ; équation de prix v0.9). Source : modèle du projet
(configuration P4, convention « part constante »), campagne
phase7/campagne.py.

\begin{longtable}[]{@{}
  >{\raggedright\arraybackslash}p{(\columnwidth - 10\tabcolsep) * \real{0.1667}}
  >{\raggedright\arraybackslash}p{(\columnwidth - 10\tabcolsep) * \real{0.1667}}
  >{\raggedright\arraybackslash}p{(\columnwidth - 10\tabcolsep) * \real{0.1667}}
  >{\raggedright\arraybackslash}p{(\columnwidth - 10\tabcolsep) * \real{0.1667}}
  >{\raggedright\arraybackslash}p{(\columnwidth - 10\tabcolsep) * \real{0.1667}}
  >{\raggedright\arraybackslash}p{(\columnwidth - 10\tabcolsep) * \real{0.1667}}@{}}
\toprule\noalign{}
\begin{minipage}[b]{\linewidth}\raggedright
candidat --- environnement
\end{minipage} & \begin{minipage}[b]{\linewidth}\raggedright
PIB 10 ans (\%)
\end{minipage} & \begin{minipage}[b]{\linewidth}\raggedright
emploi 10 ans (k)
\end{minipage} & \begin{minipage}[b]{\linewidth}\raggedright
chômage 10 ans (pt)
\end{minipage} & \begin{minipage}[b]{\linewidth}\raggedright
dette 10 ans (pt)
\end{minipage} & \begin{minipage}[b]{\linewidth}\raggedright
dette 20 ans (pt)
\end{minipage} \\
\midrule\noalign{}
\endhead
\bottomrule\noalign{}
\endlastfoot
Mélenchon (LFI) --- P4 (référence) & −0,54 & −48,77 & −2,78 & −38,14 &
−73,61 \\
Mélenchon (LFI) --- p4\_sigma03 & −0,61 & −7,09 & −2,91 & −39,63 &
−77,72 \\
Mélenchon (LFI) --- p4\_sigma06 & −0,50 & −72,54 & −2,71 & −36,59 &
−69,83 \\
Mélenchon (LFI) --- p4\_prix\_v09 & −0,59 & −53,59 & −2,77 & −38,32 &
−75,94 \\
Le Pen (RN) --- P4 (référence) & −2,89 & −754,30 & −1,65 & −1,60 &
−3,34 \\
Le Pen (RN) --- p4\_sigma03 & −2,90 & −774,49 & −1,58 & −1,19 & −2,28 \\
Le Pen (RN) --- p4\_sigma06 & −2,90 & −708,66 & −1,79 & −2,51 & −5,75 \\
Le Pen (RN) --- p4\_prix\_v09 & −2,88 & −760,36 & −1,63 & −1,15 &
−2,15 \\
Philippe (Horizons) --- P4 (référence) & +2,45 & −267,47 & +2,56 & +8,42
& +13,21 \\
Philippe (Horizons) --- p4\_sigma03 & +2,43 & −132,25 & +2,17 & +6,38 &
+6,57 \\
Philippe (Horizons) --- p4\_sigma06 & +2,47 & −414,68 & +2,98 & +10,30 &
+21,21 \\
Philippe (Horizons) --- p4\_prix\_v09 & +2,43 & −256,45 & +2,52 & +7,82
& +11,97 \\
\end{longtable}

L'élasticité de substitution capital-travail (\(\sigma\)) et l'équation
des prix de la valeur ajoutée modifient l'ampleur des effets sur
l'emploi et la dette, sans changer leur sens.

\hypertarget{sec:trajectoires}{%
\section{Trajectoires}\label{sec:trajectoires}}

\begin{figure}[htbp]\centering\includegraphics[width=0.95\textwidth]{figures/f9_pib.pdf}\caption{PIB : écart à la référence (lectures A et B).}\end{figure}

\begin{figure}[htbp]\centering\includegraphics[width=0.95\textwidth]{figures/f9_chomage.pdf}\caption{Chômage : écart à la référence (points).}\end{figure}

\begin{figure}[htbp]\centering\includegraphics[width=0.95\textwidth]{figures/f9_solde.pdf}\caption{Solde public en niveau (lectures A et B ; référence en gris).}\end{figure}

\begin{figure}[htbp]\centering\includegraphics[width=0.95\textwidth]{figures/f9_dette.pdf}\caption{Dette publique en niveau, OAT à 4,89 \% (référence en gris).}\end{figure}

\begin{figure}[htbp]\centering\includegraphics[width=0.95\textwidth]{figures/f9_dette_oat.pdf}\caption{Dette publique en niveau, OAT à 3,82 \%.}\end{figure}

\hypertarget{sec:indicateurs-detail}{%
\section{Indicateurs détaillés}\label{sec:indicateurs-detail}}

Les tableaux complets (dix-sept indicateurs, trois horizons, deux
lectures) figurent en annexe \ref{sec:annexe-e}. Les indicateurs
couverts sont : PIB, emploi, chômage, salaire réel, prix à la
consommation, pouvoir d'achat (revenu disponible réel), consommation,
taux d'épargne, investissement, taux de marge, exportations,
importations, solde public, dette publique, charge d'intérêts, recettes
et dépenses publiques (en points de PIB).

\hypertarget{sec:non-simule}{%
\section{Ce qui n'a pas été simulé}\label{sec:non-simule}}

\begin{itemize}
\tightlist
\item
  Les mesures non chiffrées ou sans canal (liste au chapitre
  \ref{sec:programmes}) : impôt sur le revenu à 14 tranches, allocation
  familiale de 350 €, durée d'indemnisation du chômage, jours de
  carence, pilier de capitalisation, contrôle des capitaux, priorité
  nationale, planification industrielle, nucléaire.
\item
  La variante de mandat de LFI (retraite à 60 ans et 40 annuités) n'est
  pas incluse dans le paquet de base.
\item
  Les effets distributifs dynamiques (M16) et les effets territoriaux.
\item
  Le financement de l'investissement hors budget (centralisation de
  l'épargne réglementée, crédit du pôle public) et le risque de
  requalification de sa dette en dette publique.
\end{itemize}

\hypertarget{sec:distribution}{%
\chapter{Distribution des effets économiques et
sociaux}\label{sec:distribution}}

Ce chapitre examine qui gagne et qui perd, au moins au premier ordre. Il
distingue trois niveaux, de précision décroissante : les \textbf{effets
statiques} des mesures socio-fiscales, calculés ménage par ménage par la
microsimulation ; les \textbf{effets générationnels moyens} des réformes
des retraites, calculés par génération et par sexe par le module M9 ; et
les \textbf{effets macroéconomiques moyens} (emploi, salaires, prix),
qui ne sont pas répartis entre ménages dans cette étude et ne peuvent
être commentés que qualitativement.

\hypertarget{sec:moyens-individuels}{%
\section{Effets moyens et effets
individuels}\label{sec:moyens-individuels}}

Un effet moyen par décile peut masquer des situations individuelles très
diverses : au sein d'un même décile, certains ménages gagnent et
d'autres perdent selon leur composition familiale, leur statut
d'activité, leur logement ou leurs revenus du capital. Les tableaux
donnent donc, pour chaque décile, la part de gagnants et de perdants
(ménages dont le revenu disponible varie de plus d'un euro par an). Par
ailleurs, l'effet statique ne tient compte ni des changements de
comportement (offre de travail, optimisation fiscale), ni des effets
macroéconomiques induits (emploi, salaires, prix), qui peuvent être du
même ordre de grandeur que l'effet direct pour certaines catégories.

\hypertarget{sec:deciles}{%
\section{Effets statiques par décile de niveau de
vie}\label{sec:deciles}}

\textbf{Couverture.} Seules les mesures dont l'incidence ménage par
ménage est représentée entrent dans ces tableaux : barèmes et
prestations calculés par OpenFisca (RSA, prime d'activité, allocations
familiales, aides au logement, impôt sur le revenu, CSG) et mesures
réparties selon des profils (TVA selon la structure de consommation par
quintile, impôts sur le capital sur le haut de la distribution). Les
mesures traitées en transferts non ciblés, en recettes sans canal propre
ou en dépenses publiques n'y figurent pas ; les réformes des retraites
sont traitées à la section suivante.

\textbf{Tableau : effets statiques par décile de niveau de vie ---
Mélenchon (LFI)} (mesures dont l'incidence ménage par ménage est
représentée ; revenus de 2024, comportements inchangés). Montants
agrégés pris en compte par la microsimulation : impôts +12,4 Md€,
prestations +6,0 Md€ (euros de 2024).

\begin{longtable}[]{@{}
  >{\raggedright\arraybackslash}p{(\columnwidth - 10\tabcolsep) * \real{0.1667}}
  >{\raggedright\arraybackslash}p{(\columnwidth - 10\tabcolsep) * \real{0.1667}}
  >{\raggedright\arraybackslash}p{(\columnwidth - 10\tabcolsep) * \real{0.1667}}
  >{\raggedright\arraybackslash}p{(\columnwidth - 10\tabcolsep) * \real{0.1667}}
  >{\raggedright\arraybackslash}p{(\columnwidth - 10\tabcolsep) * \real{0.1667}}
  >{\raggedright\arraybackslash}p{(\columnwidth - 10\tabcolsep) * \real{0.1667}}@{}}
\toprule\noalign{}
\begin{minipage}[b]{\linewidth}\raggedright
décile
\end{minipage} & \begin{minipage}[b]{\linewidth}\raggedright
niveau de vie moyen (€/UC)
\end{minipage} & \begin{minipage}[b]{\linewidth}\raggedright
variation (€/UC/an)
\end{minipage} & \begin{minipage}[b]{\linewidth}\raggedright
variation (\%)
\end{minipage} & \begin{minipage}[b]{\linewidth}\raggedright
gagnants (\%)
\end{minipage} & \begin{minipage}[b]{\linewidth}\raggedright
perdants (\%)
\end{minipage} \\
\midrule\noalign{}
\endhead
\bottomrule\noalign{}
\endlastfoot
1 & 5447,7 & 0,0 & 0,0 & 0,0 & 0,0 \\
2 & 10873,5 & 0,0 & 0,0 & 0,0 & 0,0 \\
3 & 14578,5 & 0,0 & 0,0 & 0,0 & 0,0 \\
4 & 17455,8 & 0,0 & 0,0 & 0,0 & 0,0 \\
5 & 20070,7 & 0,0 & 0,0 & 0,0 & 0,0 \\
6 & 22761,5 & 0,0 & 0,0 & 0,0 & 0,0 \\
7 & 25956,9 & 0,0 & 0,0 & 0,0 & 0,0 \\
8 & 30075,9 & 0,0 & 0,0 & 0,0 & 0,0 \\
9 & 36383,4 & −0,8 & −0,0 & 0,0 & 0,4 \\
10 & 71573,1 & −2292,7 & −3,2 & 0,0 & 99,6 \\
\end{longtable}

Indicateurs (avant → après, effet statique) : indice de Gini 0,3493 →
0,3425 ; rapport interdécile D9/D1 4,104 → 4,081 ; taux de pauvreté
(seuil à 60 \% de la médiane) 17,30 \% → 17,30 \%.

\begin{longtable}[]{@{}
  >{\raggedright\arraybackslash}p{(\columnwidth - 8\tabcolsep) * \real{0.2000}}
  >{\raggedright\arraybackslash}p{(\columnwidth - 8\tabcolsep) * \real{0.2000}}
  >{\raggedright\arraybackslash}p{(\columnwidth - 8\tabcolsep) * \real{0.2000}}
  >{\raggedright\arraybackslash}p{(\columnwidth - 8\tabcolsep) * \real{0.2000}}
  >{\raggedright\arraybackslash}p{(\columnwidth - 8\tabcolsep) * \real{0.2000}}@{}}
\toprule\noalign{}
\begin{minipage}[b]{\linewidth}\raggedright
âge de la personne de référence
\end{minipage} & \begin{minipage}[b]{\linewidth}\raggedright
variation (€/UC/an)
\end{minipage} & \begin{minipage}[b]{\linewidth}\raggedright
variation (\%)
\end{minipage} & \begin{minipage}[b]{\linewidth}\raggedright
gagnants (\%)
\end{minipage} & \begin{minipage}[b]{\linewidth}\raggedright
poids (\%)
\end{minipage} \\
\midrule\noalign{}
\endhead
\bottomrule\noalign{}
\endlastfoot
30-44 ans & −358,8 & −1,3 & 0,0 & 23,3 \\
45-59 ans & −290,9 & −1,0 & 0,0 & 28,7 \\
60-74 ans & −205,2 & −0,9 & 0,0 & 22,6 \\
75 ans et plus & −16,1 & −0,1 & 0,0 & 14,5 \\
moins de 30 ans & −125,4 & −0,6 & 0,0 & 10,9 \\
\end{longtable}

\textbf{Tableau : effets statiques par décile de niveau de vie --- Le
Pen (RN)} (mesures dont l'incidence ménage par ménage est représentée ;
revenus de 2024, comportements inchangés). Montants agrégés pris en
compte par la microsimulation : impôts +0,0 Md€, prestations −7,6 Md€
(euros de 2024).

\begin{longtable}[]{@{}
  >{\raggedright\arraybackslash}p{(\columnwidth - 10\tabcolsep) * \real{0.1667}}
  >{\raggedright\arraybackslash}p{(\columnwidth - 10\tabcolsep) * \real{0.1667}}
  >{\raggedright\arraybackslash}p{(\columnwidth - 10\tabcolsep) * \real{0.1667}}
  >{\raggedright\arraybackslash}p{(\columnwidth - 10\tabcolsep) * \real{0.1667}}
  >{\raggedright\arraybackslash}p{(\columnwidth - 10\tabcolsep) * \real{0.1667}}
  >{\raggedright\arraybackslash}p{(\columnwidth - 10\tabcolsep) * \real{0.1667}}@{}}
\toprule\noalign{}
\begin{minipage}[b]{\linewidth}\raggedright
décile
\end{minipage} & \begin{minipage}[b]{\linewidth}\raggedright
niveau de vie moyen (€/UC)
\end{minipage} & \begin{minipage}[b]{\linewidth}\raggedright
variation (€/UC/an)
\end{minipage} & \begin{minipage}[b]{\linewidth}\raggedright
variation (\%)
\end{minipage} & \begin{minipage}[b]{\linewidth}\raggedright
gagnants (\%)
\end{minipage} & \begin{minipage}[b]{\linewidth}\raggedright
perdants (\%)
\end{minipage} \\
\midrule\noalign{}
\endhead
\bottomrule\noalign{}
\endlastfoot
1 & 5447,7 & 27,6 & 0,5 & 99,3 & 0,0 \\
2 & 10873,5 & 53,2 & 0,5 & 100,0 & 0,0 \\
3 & 14578,5 & 64,6 & 0,4 & 100,0 & 0,0 \\
4 & 17455,8 & 77,3 & 0,4 & 100,0 & 0,0 \\
5 & 20070,7 & 85,8 & 0,4 & 100,0 & 0,0 \\
6 & 22761,5 & 97,2 & 0,4 & 100,0 & 0,0 \\
7 & 25956,9 & 106,6 & 0,4 & 100,0 & 0,0 \\
8 & 30075,9 & 123,2 & 0,4 & 100,0 & 0,0 \\
9 & 36383,4 & 134,9 & 0,4 & 100,0 & 0,0 \\
10 & 71573,1 & 265,0 & 0,4 & 100,0 & 0,0 \\
\end{longtable}

Indicateurs (avant → après, effet statique) : indice de Gini 0,3493 →
0,3491 ; rapport interdécile D9/D1 4,104 → 4,099 ; taux de pauvreté
(seuil à 60 \% de la médiane) 17,30 \% → 17,29 \%.

\begin{longtable}[]{@{}
  >{\raggedright\arraybackslash}p{(\columnwidth - 8\tabcolsep) * \real{0.2000}}
  >{\raggedright\arraybackslash}p{(\columnwidth - 8\tabcolsep) * \real{0.2000}}
  >{\raggedright\arraybackslash}p{(\columnwidth - 8\tabcolsep) * \real{0.2000}}
  >{\raggedright\arraybackslash}p{(\columnwidth - 8\tabcolsep) * \real{0.2000}}
  >{\raggedright\arraybackslash}p{(\columnwidth - 8\tabcolsep) * \real{0.2000}}@{}}
\toprule\noalign{}
\begin{minipage}[b]{\linewidth}\raggedright
âge de la personne de référence
\end{minipage} & \begin{minipage}[b]{\linewidth}\raggedright
variation (€/UC/an)
\end{minipage} & \begin{minipage}[b]{\linewidth}\raggedright
variation (\%)
\end{minipage} & \begin{minipage}[b]{\linewidth}\raggedright
gagnants (\%)
\end{minipage} & \begin{minipage}[b]{\linewidth}\raggedright
poids (\%)
\end{minipage} \\
\midrule\noalign{}
\endhead
\bottomrule\noalign{}
\endlastfoot
30-44 ans & 113,8 & 0,4 & 100,0 & 23,3 \\
45-59 ans & 118,8 & 0,4 & 100,0 & 28,7 \\
60-74 ans & 96,5 & 0,4 & 100,0 & 22,6 \\
75 ans et plus & 83,9 & 0,4 & 99,9 & 14,5 \\
moins de 30 ans & 82,4 & 0,4 & 99,6 & 10,9 \\
\end{longtable}

\textbf{Tableau : effets statiques par décile de niveau de vie ---
Philippe (Horizons)} (mesures dont l'incidence ménage par ménage est
représentée ; revenus de 2024, comportements inchangés). Montants
agrégés pris en compte par la microsimulation : impôts +0,0 Md€,
prestations −5,1 Md€ (euros de 2024).

\begin{longtable}[]{@{}
  >{\raggedright\arraybackslash}p{(\columnwidth - 10\tabcolsep) * \real{0.1667}}
  >{\raggedright\arraybackslash}p{(\columnwidth - 10\tabcolsep) * \real{0.1667}}
  >{\raggedright\arraybackslash}p{(\columnwidth - 10\tabcolsep) * \real{0.1667}}
  >{\raggedright\arraybackslash}p{(\columnwidth - 10\tabcolsep) * \real{0.1667}}
  >{\raggedright\arraybackslash}p{(\columnwidth - 10\tabcolsep) * \real{0.1667}}
  >{\raggedright\arraybackslash}p{(\columnwidth - 10\tabcolsep) * \real{0.1667}}@{}}
\toprule\noalign{}
\begin{minipage}[b]{\linewidth}\raggedright
décile
\end{minipage} & \begin{minipage}[b]{\linewidth}\raggedright
niveau de vie moyen (€/UC)
\end{minipage} & \begin{minipage}[b]{\linewidth}\raggedright
variation (€/UC/an)
\end{minipage} & \begin{minipage}[b]{\linewidth}\raggedright
variation (\%)
\end{minipage} & \begin{minipage}[b]{\linewidth}\raggedright
gagnants (\%)
\end{minipage} & \begin{minipage}[b]{\linewidth}\raggedright
perdants (\%)
\end{minipage} \\
\midrule\noalign{}
\endhead
\bottomrule\noalign{}
\endlastfoot
1 & 5447,7 & −451,4 & −8,3 & 0,0 & 34,7 \\
2 & 10873,5 & −383,0 & −3,5 & 0,0 & 31,7 \\
3 & 14578,5 & −66,2 & −0,5 & 0,0 & 7,0 \\
4 & 17455,8 & −54,8 & −0,3 & 0,0 & 5,6 \\
5 & 20070,7 & −43,1 & −0,2 & 0,0 & 4,5 \\
6 & 22761,5 & −34,2 & −0,2 & 0,0 & 3,8 \\
7 & 25956,9 & −31,5 & −0,1 & 0,0 & 3,5 \\
8 & 30075,9 & −28,9 & −0,1 & 0,0 & 3,3 \\
9 & 36383,4 & −22,8 & −0,1 & 0,0 & 2,6 \\
10 & 71573,1 & −17,8 & −0,0 & 0,0 & 1,8 \\
\end{longtable}

Indicateurs (avant → après, effet statique) : indice de Gini 0,3493 →
0,3532 ; rapport interdécile D9/D1 4,104 → 4,222 ; taux de pauvreté
(seuil à 60 \% de la médiane) 17,30 \% → 17,52 \%.

\begin{longtable}[]{@{}
  >{\raggedright\arraybackslash}p{(\columnwidth - 8\tabcolsep) * \real{0.2000}}
  >{\raggedright\arraybackslash}p{(\columnwidth - 8\tabcolsep) * \real{0.2000}}
  >{\raggedright\arraybackslash}p{(\columnwidth - 8\tabcolsep) * \real{0.2000}}
  >{\raggedright\arraybackslash}p{(\columnwidth - 8\tabcolsep) * \real{0.2000}}
  >{\raggedright\arraybackslash}p{(\columnwidth - 8\tabcolsep) * \real{0.2000}}@{}}
\toprule\noalign{}
\begin{minipage}[b]{\linewidth}\raggedright
âge de la personne de référence
\end{minipage} & \begin{minipage}[b]{\linewidth}\raggedright
variation (€/UC/an)
\end{minipage} & \begin{minipage}[b]{\linewidth}\raggedright
variation (\%)
\end{minipage} & \begin{minipage}[b]{\linewidth}\raggedright
gagnants (\%)
\end{minipage} & \begin{minipage}[b]{\linewidth}\raggedright
poids (\%)
\end{minipage} \\
\midrule\noalign{}
\endhead
\bottomrule\noalign{}
\endlastfoot
30-44 ans & −144,2 & −0,5 & 0,0 & 23,3 \\
45-59 ans & −153,5 & −0,5 & 0,0 & 28,7 \\
60-74 ans & −72,9 & −0,3 & 0,0 & 22,6 \\
75 ans et plus & −15,4 & −0,1 & 0,0 & 14,5 \\
moins de 30 ans & −156,4 & −0,8 & 0,0 & 10,9 \\
\end{longtable}

\textbf{Lecture.} Pour LFI, les mesures représentées (imposition du
capital et des patrimoines) pèsent sur le haut de la distribution ; les
transferts (garantie d'autonomie, revalorisations) et la hausse du SMIC,
qui bénéficieraient au bas et au milieu de la distribution, ne sont pas
répartis par la microsimulation dans cette version (transferts non
ciblés, SMIC traité par le canal salarial). Pour le RN, la baisse de la
TVA sur l'énergie, répartie selon la structure de consommation,
bénéficie à tous les déciles, un peu plus en proportion du revenu aux
ménages modestes et davantage en euros aux ménages aisés ; les économies
sur l'immigration, traitées en transferts non ciblés, n'y sont pas
réparties. Pour Horizons, la baisse du RSA concentre les pertes sur les
premiers déciles.

\hypertarget{sec:generations}{%
\section{Actifs, retraités, jeunes et
générations}\label{sec:generations}}

Les réformes des retraites ont des effets très différenciés selon les
générations : elles ne concernent pas les retraités actuels, ni les
générations qui atteignent l'âge légal avant leur entrée en vigueur ;
elles modifient l'âge de départ, la durée de retraite et le niveau de
pension des générations suivantes.

\textbf{Tableau : effets générationnels moyens des réformes des
retraites (module M9).} Âge d'ouverture des droits (AOD), durée requise,
écart d'âge pivot du profil d'activité (≈ report moyen de la cessation
d'activité ; l'écart de durée de retraite attendue en est l'opposé, à
espérance de vie inchangée) et facteur de niveau de pension dû à la
durée requise. Effets moyens, non individuels.

\begin{longtable}[]{@{}
  >{\raggedright\arraybackslash}p{(\columnwidth - 12\tabcolsep) * \real{0.1429}}
  >{\raggedright\arraybackslash}p{(\columnwidth - 12\tabcolsep) * \real{0.1429}}
  >{\raggedright\arraybackslash}p{(\columnwidth - 12\tabcolsep) * \real{0.1429}}
  >{\raggedright\arraybackslash}p{(\columnwidth - 12\tabcolsep) * \real{0.1429}}
  >{\raggedright\arraybackslash}p{(\columnwidth - 12\tabcolsep) * \real{0.1429}}
  >{\raggedright\arraybackslash}p{(\columnwidth - 12\tabcolsep) * \real{0.1429}}
  >{\raggedright\arraybackslash}p{(\columnwidth - 12\tabcolsep) * \real{0.1429}}@{}}
\toprule\noalign{}
\begin{minipage}[b]{\linewidth}\raggedright
candidat --- génération --- sexe
\end{minipage} & \begin{minipage}[b]{\linewidth}\raggedright
AOD actuelle
\end{minipage} & \begin{minipage}[b]{\linewidth}\raggedright
AOD réforme
\end{minipage} & \begin{minipage}[b]{\linewidth}\raggedright
durée actuelle (trim.)
\end{minipage} & \begin{minipage}[b]{\linewidth}\raggedright
durée réforme (trim.)
\end{minipage} & \begin{minipage}[b]{\linewidth}\raggedright
écart d'âge pivot (ans)
\end{minipage} & \begin{minipage}[b]{\linewidth}\raggedright
facteur de pension
\end{minipage} \\
\midrule\noalign{}
\endhead
\bottomrule\noalign{}
\endlastfoot
Mélenchon (LFI) --- 1960 --- F & 62,00 & 62,00 & 167,00 & 167,00 & 0,00
& 1,00 \\
Mélenchon (LFI) --- 1960 --- M & 62,00 & 62,00 & 167,00 & 167,00 & 0,00
& 1,00 \\
Mélenchon (LFI) --- 1965 --- F & 63,00 & 62,00 & 171,00 & 169,00 & −0,99
& 1,01 \\
Mélenchon (LFI) --- 1965 --- M & 63,00 & 62,00 & 171,00 & 169,00 & −1,13
& 1,00 \\
Mélenchon (LFI) --- 1970 --- F & 64,00 & 62,00 & 172,00 & 171,00 & −1,66
& 1,01 \\
Mélenchon (LFI) --- 1970 --- M & 64,00 & 62,00 & 172,00 & 171,00 & −1,69
& 1,00 \\
Mélenchon (LFI) --- 1975 --- F & 64,00 & 62,00 & 172,00 & 172,00 & −1,56
& 1,00 \\
Mélenchon (LFI) --- 1975 --- M & 64,00 & 62,00 & 172,00 & 172,00 & −1,50
& 1,00 \\
Mélenchon (LFI) --- 1980 --- F & 64,00 & 62,00 & 172,00 & 172,00 & −1,56
& 1,00 \\
Mélenchon (LFI) --- 1980 --- M & 64,00 & 62,00 & 172,00 & 172,00 & −1,50
& 1,00 \\
Mélenchon (LFI) --- 1990 --- F & 64,00 & 62,00 & 172,00 & 172,00 & −1,56
& 1,00 \\
Mélenchon (LFI) --- 1990 --- M & 64,00 & 62,00 & 172,00 & 172,00 & −1,50
& 1,00 \\
Le Pen (RN) --- 1960 --- F & 62,00 & 62,00 & 167,00 & 167,00 & 0,00 &
1,00 \\
Le Pen (RN) --- 1960 --- M & 62,00 & 62,00 & 167,00 & 167,00 & 0,00 &
1,00 \\
Le Pen (RN) --- 1965 --- F & 63,00 & 63,00 & 171,00 & 168,00 & −0,45 &
1,02 \\
Le Pen (RN) --- 1965 --- M & 63,00 & 63,00 & 171,00 & 168,00 & −0,69 &
1,01 \\
Le Pen (RN) --- 1970 --- F & 64,00 & 62,00 & 172,00 & 168,00 & −1,98 &
1,02 \\
Le Pen (RN) --- 1970 --- M & 64,00 & 62,00 & 172,00 & 168,00 & −2,25 &
1,01 \\
Le Pen (RN) --- 1975 --- F & 64,00 & 62,00 & 172,00 & 168,00 & −1,98 &
1,02 \\
Le Pen (RN) --- 1975 --- M & 64,00 & 62,00 & 172,00 & 168,00 & −2,25 &
1,01 \\
Le Pen (RN) --- 1980 --- F & 64,00 & 62,00 & 172,00 & 168,00 & −1,98 &
1,02 \\
Le Pen (RN) --- 1980 --- M & 64,00 & 62,00 & 172,00 & 168,00 & −2,25 &
1,01 \\
Le Pen (RN) --- 1990 --- F & 64,00 & 62,00 & 172,00 & 168,00 & −1,98 &
1,02 \\
Le Pen (RN) --- 1990 --- M & 64,00 & 62,00 & 172,00 & 168,00 & −2,25 &
1,01 \\
Philippe (Horizons) --- 1960 --- F & 62,00 & 62,00 & 167,00 & 167,00 &
0,00 & 1,00 \\
Philippe (Horizons) --- 1960 --- M & 62,00 & 62,00 & 167,00 & 167,00 &
0,00 & 1,00 \\
Philippe (Horizons) --- 1965 --- F & 63,00 & 63,25 & 171,00 & 172,00 &
0,20 & 0,99 \\
Philippe (Horizons) --- 1965 --- M & 63,00 & 63,25 & 171,00 & 172,00 &
0,25 & 1,00 \\
Philippe (Horizons) --- 1970 --- F & 64,00 & 64,50 & 172,00 & 177,00 &
0,61 & 0,97 \\
Philippe (Horizons) --- 1970 --- M & 64,00 & 64,50 & 172,00 & 177,00 &
0,88 & 0,99 \\
Philippe (Horizons) --- 1975 --- F & 64,00 & 65,00 & 172,00 & 180,00 &
1,08 & 0,96 \\
Philippe (Horizons) --- 1975 --- M & 64,00 & 65,00 & 172,00 & 180,00 &
1,50 & 0,98 \\
Philippe (Horizons) --- 1980 --- F & 64,00 & 65,00 & 172,00 & 180,00 &
1,08 & 0,96 \\
Philippe (Horizons) --- 1980 --- M & 64,00 & 65,00 & 172,00 & 180,00 &
1,50 & 0,98 \\
Philippe (Horizons) --- 1990 --- F & 64,00 & 65,00 & 172,00 & 180,00 &
1,08 & 0,96 \\
Philippe (Horizons) --- 1990 --- M & 64,00 & 65,00 & 172,00 & 180,00 &
1,50 & 0,98 \\
\end{longtable}

\textbf{Lecture.} Les écarts d'âge pivot mesurent le décalage moyen de
la cessation d'activité d'une génération ; l'écart de durée de retraite
attendue en est approximativement l'opposé (à espérance de vie
inchangée). Pour LFI et le RN, les générations nées à partir du milieu
des années 1960 partiraient en moyenne plus tôt (jusqu'à deux ans
environ, davantage pour les carrières longues dans la réforme du RN) ;
pour Horizons, les générations nées après 1970 partiraient en moyenne
plus d'un an plus tard (davantage pour les hommes, dont le report en
réponse à la durée requise est plus élevé dans la littérature), avec un
niveau de pension légèrement réduit pour la fraction qui ne reporte pas
son départ. Le financement de ces réformes concerne, lui, toutes les
générations (cotisations, impôts, dette) ; sa répartition dépend des
choix budgétaires (lecture B) et de la dynamique de la dette, qui
reporte une partie de la charge sur les générations futures.

\hypertarget{sec:categories}{%
\section{Catégories socioprofessionnelles, entreprises, secteurs et
territoires}\label{sec:categories}}

La microsimulation ne distingue pas les catégories socioprofessionnelles
de manière fiable (score latent construit à partir des grilles de
salaires) ; les effets par catégorie ne sont donc pas publiés. Le modèle
macroéconomique distingue 25 branches marchandes : les effets sectoriels
(industrie, services, construction) sont disponibles dans les sorties
détaillées mais ne sont pas commentés ici, faute de validation
suffisante des équations de branche à cet horizon. Les entreprises ne
sont pas distinguées par taille (la baisse des impôts de production et
la suppression des aides concernent pourtant des entreprises différentes
: la CVAE et la C3S pèsent surtout sur les grandes entreprises, la CFE
sur toutes). Les effets territoriaux ne sont pas représentés.

Qualitativement, les mesures étudiées ont des incidences sectorielles et
territoriales connues : les hausses du SMIC et du point d'indice
concernent davantage les territoires où la part de l'emploi public et
des bas salaires est élevée ; la suppression d'emplois publics et la
réorganisation des collectivités (RN, Horizons) concernent davantage les
territoires peu denses ; la baisse des impôts de production bénéficie
d'abord à l'industrie.

\hypertarget{sec:limites-micro}{%
\section{Limites de la microsimulation}\label{sec:limites-micro}}

\begin{enumerate}
\def\labelenumi{\arabic{enumi}.}
\tightlist
\item
  \textbf{Échantillon} : 2 \% des ménages pour l'impôt, deux
  sous-échantillons de 0,5 \% pour les prestations ; erreur
  d'échantillonnage de l'ordre de 2 \% sur les montants.
\item
  \textbf{Revenus de 2024}, non revalorisés avec les salaires, l'emploi
  et les prix simulés (pas de retour macro → micro).
\item
  \textbf{Patrimoine} : approché par les revenus du patrimoine ; les
  ménages à patrimoine peu productif de revenu imposable sont
  sous-représentés, d'où une répartition approximative des impôts sur le
  patrimoine.
\item
  \textbf{Non-recours} : supposé constant (taux de recours apparent de
  2024) ; une hausse des prestations peut modifier le recours.
\item
  \textbf{Catégories mal représentées} : très hauts revenus et
  patrimoines (agrégats), indépendants, ménages complexes, personnes
  hors ménages ordinaires (institutions).
\item
  \textbf{Comportements} : l'offre de travail réagit en moyenne (M15) ;
  les réactions individuelles (seuils, effets de palier) ne sont pas
  microsimulées.
\end{enumerate}

\hypertarget{sec:comportements}{%
\chapter{Comportements individuels et collectifs : limites des théories
économiques}\label{sec:comportements}}

Un modèle macroéconométrique résume des millions de décisions par
quelques équations moyennes, estimées sur une période donnée. Il suppose
implicitement que les comportements observés de 1997 à 2019 resteront
valables face à des politiques qui, pour certaines, s'écartent fortement
de l'expérience historique. Ce chapitre examine, acteur par acteur, ce
que le modèle suppose, ce que disent les données et la littérature, et
ce qui pourrait se passer différemment. Lorsque c'est possible, il
renvoie à des variantes quantifiées ; lorsque ce ne l'est pas, il
propose une analyse qualitative.

La grille de lecture est la suivante : pour chaque comportement, on
distingue (i) \textbf{l'hypothèse du modèle}, (ii) \textbf{ce que dit la
littérature}, (iii) \textbf{les situations où l'hypothèse peut être mise
en défaut} par les programmes étudiés, et (iv) \textbf{le sens probable
du biais} sur les résultats.

\hypertarget{muxe9nages}{%
\section{Ménages}\label{muxe9nages}}

\hypertarget{consommation-uxe9pargne-et-duxe9sendettement}{%
\subsection{Consommation, épargne et
désendettement}\label{consommation-uxe9pargne-et-duxe9sendettement}}

\emph{Hypothèse du modèle.} La consommation suit le revenu disponible
réel avec une élasticité de court terme de 0,85 et converge vers une
propension moyenne stable ; l'épargne de précaution augmente avec le
chômage (0,3 \% de consommation par point) ; le taux réel du crédit à
l'habitat réduit la consommation (0,2).

\emph{Littérature.} Les propensions marginales à consommer varient
fortement selon la liquidité des ménages : élevées pour les ménages sans
épargne disponible, faibles pour les ménages aisés (Jappelli et
Pistaferri, 2010 ; Kaplan, Moll et Violante, 2018). Les ménages
réagissent davantage aux chocs de revenu jugés permanents (Blundell,
Pistaferri et Preston, 2008). L'épargne a fortement augmenté pendant la
crise sanitaire et n'est redescendue que lentement en France, ce qui
suggère une épargne de précaution plus forte que celle estimée sur la
période antérieure.

\emph{Situations critiques.} (a) Une hausse forte et durable de la dette
publique, ou un débat public sur sa soutenabilité, peut accroître
l'épargne de précaution (anticipation d'impôts futurs ou de baisses de
prestations : comportement « ricardien » partiel). (b) Des baisses de
prestations ciblées sur les ménages modestes ont un effet dépressif plus
fort que la moyenne. (c) Des transferts présentés comme temporaires sont
davantage épargnés.

\emph{Sens du biais.} Le modèle surestime probablement l'effet de
demande des relances financées par la dette dans un contexte de dette
élevée, et sous-estime l'effet dépressif des économies portant sur les
ménages modestes s'il ne passe pas par la microsimulation (le canal
agrégé des prestations n'est pas ciblé).

\emph{Variante disponible.} Épargne de précaution basse/haute (0,1 /
0,6) : effet à 20 ans de moins de 0,03 point de PIB et d'au plus 0,55
point de dette sur les stratégies de test (chapitre
\ref{sec:robustesse}).

\hypertarget{anticipations-de-hausse-ou-de-baisse-des-revenus}{%
\subsection{Anticipations de hausse ou de baisse des
revenus}\label{anticipations-de-hausse-ou-de-baisse-des-revenus}}

\emph{Hypothèse.} Anticipations adaptatives : une mesure annoncée n'a
d'effet qu'à son entrée en vigueur.

\emph{Littérature.} Les ménages anticipent partiellement les réformes
annoncées : achats avant une hausse de TVA (effet observé en Allemagne
en 2007, au Japon en 2014), départs à la retraite avant un relèvement de
l'âge, consommation anticipée d'une hausse de revenu certaine.

\emph{Situations critiques.} Les réformes des retraites annoncées
longtemps à l'avance (Horizons : 65 ans progressivement) peuvent
modifier les comportements d'épargne et de fin de carrière dès l'annonce
; une baisse de TVA annoncée peut retarder certains achats d'énergie.

\emph{Sens du biais.} Effets de calendrier sur la première année, sans
changement probable des effets à 5 ans.

\hypertarget{ruxe9actions-aux-impuxf4ts-et-aux-transferts}{%
\subsection{Réactions aux impôts et aux
transferts}\label{ruxe9actions-aux-impuxf4ts-et-aux-transferts}}

\emph{Hypothèse.} Élasticités du revenu imposable des hauts revenus (0,2
central, 0,1-0,4), du capital (0,5 ; 0,2-1,0), pertes d'assiette de
l'impôt sur la fortune (25 \% puis 35 \% à cinq ans) ; participation au
marché du travail (élasticité 0,2) et heures (0,1).

\emph{Littérature.} Les réactions des très hauts revenus et patrimoines
passent surtout par l'optimisation (choix de la forme de rémunération,
distribution ou non des dividendes, localisation des actifs) plutôt que
par l'effort ; la réforme de 2018 (prélèvement forfaitaire unique) a été
suivie d'une forte hausse des dividendes distribués aux ménages (constat
des rapports d'évaluation du comité d'évaluation des réformes de la
fiscalité du capital, France Stratégie), ce qui suggère une forte
élasticité des distributions au taux. Les départs liés à l'ancien impôt
de solidarité sur la fortune étaient peu nombreux mais concernaient des
patrimoines élevés.

\emph{Situations critiques.} Le programme de LFI cumule plusieurs
hausses d'impôt sur les mêmes contribuables (PFU, impôt sur la fortune,
taxe sur les très grands patrimoines, successions) : les élasticités ne
sont pas additives, et l'effet cumulé peut dépasser la somme des effets
isolés (taux marginaux cumulés très élevés, incitations accrues à
l'expatriation et à l'optimisation). À l'inverse, la mobilité
internationale des patrimoines est contrainte par les dispositifs d'«
exit tax » et les échanges d'information fiscale, renforcés depuis 2014.

\emph{Sens du biais.} Incertain ; l'étude donne les rendements nets dans
les trois hypothèses (basse, centrale, haute) et traite les taxes sans
canal propre (très grands patrimoines, multinationales) comme des bornes
hautes.

\hypertarget{arbitrages-travail-loisir}{%
\subsection{Arbitrages travail-loisir}\label{arbitrages-travail-loisir}}

\emph{Hypothèse.} Les réformes des prestations modifient la
participation selon des élasticités moyennes (M15) ; la durée du travail
est un paramètre légal avec une absorption de 70 \% en productivité
horaire.

\emph{Littérature.} La marge extensive (travailler ou non) réagit
davantage que la marge intensive, en particulier pour les seconds
apporteurs de revenu et les seniors (Chetty et al., 2011). Les effets
des 35 heures sur l'emploi restent discutés (fourchette de quelques
dizaines de milliers à 350 000 emplois selon les évaluations). Les
heures supplémentaires défiscalisées (2007-2012, puis 2019) ont surtout
conduit à déclarer des heures auparavant non déclarées ou à substituer
des heures à des embauches, selon plusieurs évaluations.

\emph{Situations critiques.} La fin des 35 heures (Horizons) est
modélisée de manière symétrique à leur introduction, ce qui est une
hypothèse forte : l'allongement de la durée du travail dépend de la
négociation de branche et d'entreprise et de la demande.

\emph{Sens du biais.} L'effet sur la productivité horaire et sur
l'emploi de la mesure H06 est très incertain ; il est signalé comme tel
dans les résultats.

\hypertarget{pruxe9caution-face-uxe0-lincertitude-et-ruxe9actions-diffuxe9renciuxe9es-selon-le-patrimoine}{%
\subsection{Précaution face à l'incertitude et réactions différenciées
selon le
patrimoine}\label{pruxe9caution-face-uxe0-lincertitude-et-ruxe9actions-diffuxe9renciuxe9es-selon-le-patrimoine}}

Les ménages sans patrimoine liquide n'ont pas de marge pour lisser leurs
dépenses ; les ménages propriétaires endettés à taux fixe sont peu
sensibles aux taux ; les ménages patrimoniaux réagissent à la
valorisation de leurs actifs (effets de richesse, absents du modèle).
Une remontée des taux longs (variante « tension ») réduirait la valeur
des actifs immobiliers et financiers et pèserait sur la consommation des
ménages aisés, un canal non représenté.

\hypertarget{entreprises}{%
\section{Entreprises}\label{entreprises}}

\hypertarget{investissement-et-attentisme}{%
\subsection{Investissement et
attentisme}\label{investissement-et-attentisme}}

\emph{Hypothèse.} Accélérateur, coût d'usage (élasticité −0,18),
trésorerie (0,10).

\emph{Littérature.} L'incertitude économique et politique réduit
l'investissement et l'embauche par un effet d'« option d'attendre »
(Bloom, 2009) ; les indices d'incertitude de politique économique
(Baker, Bloom et Davis, 2016) sont associés à des baisses
d'investissement. Les périodes électorales et les changements de
majorité s'accompagnent souvent d'un attentisme temporaire.

\emph{Situations critiques.} Les programmes de rupture (fiscale,
institutionnelle ou européenne) augmentent l'incertitude ; les
programmes de consolidation réduisent la demande attendue. Aucun des
deux effets n'est représenté au-delà de la prime souveraine.

\emph{Sens du biais.} Le modèle sous-estime probablement l'effet négatif
de court terme des programmes qui accroissent l'incertitude, quelle que
soit leur orientation.

\hypertarget{embauches-licenciements-marges-et-prix}{%
\subsection{Embauches, licenciements, marges et
prix}\label{embauches-licenciements-marges-et-prix}}

\emph{Hypothèse.} Ajustement partiel de l'emploi à la production
(élasticité de 0,26 à court terme) ; prix de la valeur ajoutée vers une
frontière de coûts ; marges en résultante.

\emph{Littérature.} Les entreprises françaises ajustent l'emploi plus
lentement que la production (rétention de main-d'œuvre), et ont absorbé
une partie des chocs de coûts par leurs marges (2022-2023) ; la
transmission des baisses de coûts aux prix est souvent incomplète et
asymétrique.

\emph{Situations critiques.} Une hausse forte du SMIC (LFI : +15 \%) se
diffuse dans la grille salariale (effet « tassement ») plus que ne le
suppose le modèle (11 \% de la masse salariale) ; elle pourrait aussi
accroître la substitution vers le capital ou l'externalisation. La
baisse de 50 Md€ des impôts de production (Horizons) pourrait être
davantage conservée en marges que répercutée en prix.

\hypertarget{duxe9localisations-et-relocalisations}{%
\subsection{Délocalisations et
relocalisations}\label{duxe9localisations-et-relocalisations}}

Le modèle représente la compétitivité-prix et, pour l'énergie, une
délocalisation des industries énergo-intensives (élasticité 0,15). Il ne
représente ni les décisions de localisation des sièges et des
investissements directs étrangers en fonction de la fiscalité et de la
stabilité réglementaire, ni les relocalisations induites par des
politiques industrielles (sauf comme hypothèse exogène). Les arguments
des programmes sur la relocalisation (RN, Horizons) ou sur la
planification industrielle (LFI) ne peuvent donc être évalués que
qualitativement.

\hypertarget{innovation-et-productivituxe9}{%
\subsection{Innovation et
productivité}\label{innovation-et-productivituxe9}}

La productivité est exogène (chapitre \ref{sec:transmission}). Les
réformes qui revendiquent des gains de productivité (IA dans
l'administration, réforme de l'école, planification écologique,
concurrence) ne sont pas récompensées par le modèle, et celles qui
pourraient en coûter (baisse de l'investissement dans l'éducation ou la
recherche, hausse de la fiscalité du capital) ne sont pas pénalisées par
ce canal.

\hypertarget{ruxe9actions-uxe0-lincertitude-politique-et-ruxe9glementaire}{%
\subsection{Réactions à l'incertitude politique et
réglementaire}\label{ruxe9actions-uxe0-lincertitude-politique-et-ruxe9glementaire}}

Les entreprises internationales arbitrent entre pays selon la fiscalité,
le coût du travail, la stabilité des règles et l'accès au marché
européen. Un moratoire fiscal et normatif (Horizons) vise explicitement
à réduire cette incertitude ; des mesures en conflit avec le droit
européen (RN, LFI pour certaines) l'augmentent. Ces effets ne sont pas
quantifiés ; leur ordre de grandeur est inconnu, mais leur sens est
opposé selon les programmes.

\hypertarget{institutions-et-marchuxe9s}{%
\section{Institutions et marchés}\label{institutions-et-marchuxe9s}}

\hypertarget{ruxe9actions-des-marchuxe9s-financiers-et-anticipations-des-investisseurs}{%
\subsection{Réactions des marchés financiers et anticipations des
investisseurs}\label{ruxe9actions-des-marchuxe9s-financiers-et-anticipations-des-investisseurs}}

\emph{Hypothèse.} Prime souveraine linéaire en dette et déficit
(estimée), plus chocs de scénario.

\emph{Littérature.} Les marchés réagissent de manière non linéaire :
longtemps peu sensibles aux fondamentaux, puis brutalement (crise de la
zone euro). Les épisodes de 2022 au Royaume-Uni (budget non financé
suivi d'une forte hausse des taux longs) et les tensions sur l'OAT lors
des crises politiques françaises récentes montrent que la crédibilité
budgétaire perçue peut déplacer les taux de plusieurs dizaines de points
de base en quelques jours.

\emph{Sens du biais.} Le modèle, dans son régime principal, ne
représente aucune réaction de marché ; dans le régime BCE + prime, une
réaction moyenne ; dans la variante « tension », un choc de 50 points de
base. Une crise de confiance plus grave est possible mais non
quantifiable de manière crédible ; elle concernerait en priorité les
programmes qui dégradent le solde.

\hypertarget{ruxe9actions-de-la-bce}{%
\subsection{Réactions de la BCE}\label{ruxe9actions-de-la-bce}}

La règle estimée décrit une BCE qui réagit à l'inflation et à l'activité
de l'ensemble de la zone ; elle ne réagit pas spécifiquement à la
France. Dans une situation de tension sur la dette française, la BCE
pourrait intervenir par ses instruments de protection de la
transmission, sous conditions de respect du cadre budgétaire : la
réaction serait alors asymétrique selon que le programme respecte ou non
ce cadre.

\hypertarget{ruxe9actions-des-partenaires-commerciaux}{%
\subsection{Réactions des partenaires
commerciaux}\label{ruxe9actions-des-partenaires-commerciaux}}

Le modèle ne représente pas de représailles commerciales (sauf pour le
curseur des droits de douane, hors programme) ni de réactions des
partenaires européens à des mesures unilatérales (non-paiement de la
contribution, priorité nationale). Ces réactions peuvent prendre des
formes non économiques (contentieux, blocages politiques) dont les
conséquences économiques sont indirectes.

\hypertarget{confiance-incertitude-uxe9conomie-informelle-et-comportements-duxe9vitement}{%
\subsection{Confiance, incertitude, économie informelle et comportements
d'évitement}\label{confiance-incertitude-uxe9conomie-informelle-et-comportements-duxe9vitement}}

La confiance des ménages et des entreprises n'est pas une variable du
modèle. L'économie informelle et l'évitement fiscal réagissent aux taux
d'imposition et aux contrôles : la lutte contre la fraude (LFI : 2 Md€ ;
RN : sans montant dans les sources) a des rendements difficiles à
mesurer et à anticiper, comme le relèvent régulièrement les travaux de
la Cour des comptes. Ces recettes sont traitées comme des bornes hautes.

\hypertarget{sec:synthese-biais}{%
\section{Synthèse : sens probable des biais par
programme}\label{sec:synthese-biais}}

\begin{longtable}[]{@{}
  >{\raggedright\arraybackslash}p{(\columnwidth - 2\tabcolsep) * \real{0.5000}}
  >{\raggedright\arraybackslash}p{(\columnwidth - 2\tabcolsep) * \real{0.5000}}@{}}
\toprule\noalign{}
\begin{minipage}[b]{\linewidth}\raggedright
Comportement non représenté ou simplifié
\end{minipage} & \begin{minipage}[b]{\linewidth}\raggedright
Effet sur la lecture des résultats
\end{minipage} \\
\midrule\noalign{}
\endhead
\bottomrule\noalign{}
\endlastfoot
Épargne de précaution liée à la dette publique & surestimation de
l'effet des relances financées par la dette (LFI dans la lecture A) \\
Effets de richesse et taux longs & sous-estimation de l'effet dépressif
d'une hausse des taux (tous, variante « tension ») \\
Incertitude politique et attentisme & sous-estimation des coûts de court
terme des programmes de rupture (LFI, RN) \\
Gains de productivité de long terme (éducation, État, IA, planification)
& absence de récompense des réformes structurelles revendiquées
(Horizons, LFI) \\
Diffusion du SMIC dans la grille salariale & sous-estimation des effets
sur les coûts et les prix d'une forte hausse (LFI) \\
Réactions cumulées des contribuables aisés & surestimation probable des
rendements fiscaux cumulés (LFI) \\
Rendement de la lutte contre la fraude et des économies sur les agences
& surestimation probable des économies annoncées (RN, LFI) \\
Négociation collective sur la durée du travail & incertitude forte sur
l'effet de la fin des 35 heures (Horizons) \\
\end{longtable}

Aucun de ces biais ne peut être corrigé par un paramètre unique. Ils
justifient deux attitudes : présenter les résultats sous plusieurs
régimes et hypothèses (chapitre \ref{sec:robustesse}), et ne pas
interpréter les écarts faibles entre programmes comme significatifs.

\hypertarget{sec:robustesse}{%
\chapter{Robustesse, incertitudes et scénarios
alternatifs}\label{sec:robustesse}}

Ce chapitre rassemble les analyses de sensibilité : sur des stratégies
de test standardisées (S1 à S5) et sur les programmes eux-mêmes, sous
dix régimes de politique économique, trois environnements alternatifs et
des variantes de paramètres structurels. Il distingue les types
d'incertitude et indique ce qui est robuste et ce qui ne l'est pas.

\hypertarget{sec:typologie-incertitudes}{%
\section{Typologie des incertitudes}\label{sec:typologie-incertitudes}}

\begin{longtable}[]{@{}
  >{\raggedright\arraybackslash}p{(\columnwidth - 4\tabcolsep) * \real{0.3333}}
  >{\raggedright\arraybackslash}p{(\columnwidth - 4\tabcolsep) * \real{0.3333}}
  >{\raggedright\arraybackslash}p{(\columnwidth - 4\tabcolsep) * \real{0.3333}}@{}}
\toprule\noalign{}
\begin{minipage}[b]{\linewidth}\raggedright
Type
\end{minipage} & \begin{minipage}[b]{\linewidth}\raggedright
Définition
\end{minipage} & \begin{minipage}[b]{\linewidth}\raggedright
Traitement dans l'étude
\end{minipage} \\
\midrule\noalign{}
\endhead
\bottomrule\noalign{}
\endlastfoot
\textbf{Statistique} & précision des paramètres estimés (écarts types) &
intervalles disponibles pour les paramètres {[}E{]} ; non propagés
systématiquement (le module M12 permet des tirages) \\
\textbf{Paramétrique} & valeur des paramètres calibrés ou non identifiés
& variantes basse / haute de 17 paramètres structurels ; jeux « ancrage
», « fermeture » ; environnements \(\sigma\) et prix \\
\textbf{Structurelle} & forme des équations et des mécanismes
(fermeture, anticipations, politique monétaire) & dix régimes ;
comparaison avec les familles de modèles (chapitre \ref{sec:theorie}) ;
mécanismes manquants (annexe \ref{sec:annexe-h}) \\
\textbf{Chocs futurs} & environnement international, prix de l'énergie,
taux & variante de taux longs (OAT 3,82 \% / 4,89 \%) ; chocs du régime
monétaire (non utilisés dans les résultats de base) \\
\textbf{Mise en œuvre} & ampleur et calendrier effectifs des mesures &
variantes 50 \% et retard de deux ans \\
\textbf{Politique et institutionnelle} & réactions européennes,
contentieux, crédibilité & variantes « conforme », « contentieux », «
tension » \\
\end{longtable}

\hypertarget{sec:s1s5}{%
\section{Stratégies de test S1 à S5 sous dix régimes}\label{sec:s1s5}}

Les cinq stratégies de test (S1 « protection », S2 « emploi », S3 «
productivité », S4 « redressement », S5 « global ») sont des
combinaisons de curseurs conçues lors des phases antérieures pour
couvrir des orientations contrastées ; elles servent de banc d'essai de
la robustesse du modèle, indépendamment des programmes.

\textbf{Tableau : stratégies de test S1-S5 sous les dix régimes}
(principal, minimum, maximum, stabilité du signe). Source :
phase5/resultats/lot53\_synthese\_S1\_S5.csv.

\begin{longtable}[]{@{}
  >{\raggedright\arraybackslash}p{(\columnwidth - 8\tabcolsep) * \real{0.2000}}
  >{\raggedright\arraybackslash}p{(\columnwidth - 8\tabcolsep) * \real{0.2000}}
  >{\raggedright\arraybackslash}p{(\columnwidth - 8\tabcolsep) * \real{0.2000}}
  >{\raggedright\arraybackslash}p{(\columnwidth - 8\tabcolsep) * \real{0.2000}}
  >{\raggedright\arraybackslash}p{(\columnwidth - 8\tabcolsep) * \real{0.2000}}@{}}
\toprule\noalign{}
\begin{minipage}[b]{\linewidth}\raggedright
stratégie --- horizon --- indicateur
\end{minipage} & \begin{minipage}[b]{\linewidth}\raggedright
principal
\end{minipage} & \begin{minipage}[b]{\linewidth}\raggedright
min
\end{minipage} & \begin{minipage}[b]{\linewidth}\raggedright
max
\end{minipage} & \begin{minipage}[b]{\linewidth}\raggedright
signe robuste
\end{minipage} \\
\midrule\noalign{}
\endhead
\bottomrule\noalign{}
\endlastfoot
S1\_protection --- 10 ans --- PIB (\%) & +1,00 & +0,32 & +1,49 & oui \\
S1\_protection --- 10 ans --- emploi (milliers) & +420,65 & +222,34 &
+531,10 & oui \\
S1\_protection --- 10 ans --- chômage (pt) & −0,70 & −1,02 & −0,12 &
oui \\
S1\_protection --- 10 ans --- dette (pt PIB) & +3,64 & +0,58 & +5,32 &
oui \\
S1\_protection --- 20 ans --- PIB (\%) & +0,32 & +0,10 & +0,78 & oui \\
S1\_protection --- 20 ans --- emploi (milliers) & +178,68 & +65,56 &
+295,64 & oui \\
S1\_protection --- 20 ans --- chômage (pt) & +0,01 & −0,34 & +0,34 &
non \\
S1\_protection --- 20 ans --- dette (pt PIB) & +16,92 & +9,31 & +23,43 &
oui \\
S2\_emploi --- 10 ans --- PIB (\%) & +1,48 & +1,14 & +1,66 & oui \\
S2\_emploi --- 10 ans --- emploi (milliers) & +203,63 & +136,59 &
+251,75 & oui \\
S2\_emploi --- 10 ans --- chômage (pt) & −0,15 & −0,29 & +0,05 & non \\
S2\_emploi --- 10 ans --- dette (pt PIB) & +8,43 & +7,34 & +9,65 &
oui \\
S2\_emploi --- 20 ans --- PIB (\%) & +2,39 & +1,89 & +2,93 & oui \\
S2\_emploi --- 20 ans --- emploi (milliers) & +102,73 & −28,24 & +300,54
& non \\
S2\_emploi --- 20 ans --- chômage (pt) & +0,15 & −0,44 & +0,53 & non \\
S2\_emploi --- 20 ans --- dette (pt PIB) & +24,22 & +22,28 & +32,06 &
oui \\
S3\_productivite --- 10 ans --- PIB (\%) & +2,76 & +2,22 & +3,25 &
oui \\
S3\_productivite --- 10 ans --- emploi (milliers) & +458,22 & +346,55 &
+531,84 & oui \\
S3\_productivite --- 10 ans --- chômage (pt) & +0,74 & +0,53 & +1,06 &
oui \\
S3\_productivite --- 10 ans --- dette (pt PIB) & +14,14 & +12,93 &
+16,34 & oui \\
S3\_productivite --- 20 ans --- PIB (\%) & +5,19 & +4,53 & +6,66 &
oui \\
S3\_productivite --- 20 ans --- emploi (milliers) & +758,84 & +581,10 &
+1293,94 & oui \\
S3\_productivite --- 20 ans --- chômage (pt) & +1,38 & −0,14 & +1,90 &
non \\
S3\_productivite --- 20 ans --- dette (pt PIB) & +37,75 & +33,20 &
+52,36 & oui \\
S4\_redressement --- 10 ans --- PIB (\%) & −0,79 & −0,93 & −0,45 &
oui \\
S4\_redressement --- 10 ans --- emploi (milliers) & −69,03 & −100,12 &
+44,16 & non \\
S4\_redressement --- 10 ans --- chômage (pt) & +0,20 & −0,13 & +0,29 &
non \\
S4\_redressement --- 10 ans --- dette (pt PIB) & −16,58 & −17,55 &
−15,30 & oui \\
S4\_redressement --- 20 ans --- PIB (\%) & −1,00 & −1,43 & −0,38 &
oui \\
S4\_redressement --- 20 ans --- emploi (milliers) & −14,00 & −120,72 &
+133,77 & non \\
S4\_redressement --- 20 ans --- chômage (pt) & +0,04 & −0,40 & +0,36 &
non \\
S4\_redressement --- 20 ans --- dette (pt PIB) & −38,54 & −44,07 &
−35,49 & oui \\
S5\_global --- 10 ans --- PIB (\%) & +0,64 & +0,57 & +0,94 & oui \\
S5\_global --- 10 ans --- emploi (milliers) & −47,89 & −61,37 & +56,40 &
non \\
S5\_global --- 10 ans --- chômage (pt) & +0,47 & +0,16 & +0,51 & oui \\
S5\_global --- 10 ans --- dette (pt PIB) & +3,22 & +2,35 & +4,92 &
oui \\
S5\_global --- 20 ans --- PIB (\%) & +1,33 & +1,12 & +1,56 & oui \\
S5\_global --- 20 ans --- emploi (milliers) & −92,05 & −141,64 & +57,93
& non \\
S5\_global --- 20 ans --- chômage (pt) & +0,60 & +0,16 & +0,75 & oui \\
S5\_global --- 20 ans --- dette (pt PIB) & +11,17 & +7,65 & +18,07 &
oui \\
\end{longtable}

\textbf{Tableau : classements des stratégies S1-S5 par régime} (du plus
favorable au moins favorable). Source :
phase5/resultats/lot53\_classements\_S1\_S5.csv.

\begin{longtable}[]{@{}
  >{\raggedright\arraybackslash}p{(\columnwidth - 16\tabcolsep) * \real{0.1807}}
  >{\raggedright\arraybackslash}p{(\columnwidth - 16\tabcolsep) * \real{0.0964}}
  >{\raggedright\arraybackslash}p{(\columnwidth - 16\tabcolsep) * \real{0.0964}}
  >{\raggedright\arraybackslash}p{(\columnwidth - 16\tabcolsep) * \real{0.1044}}
  >{\raggedright\arraybackslash}p{(\columnwidth - 16\tabcolsep) * \real{0.1044}}
  >{\raggedright\arraybackslash}p{(\columnwidth - 16\tabcolsep) * \real{0.1124}}
  >{\raggedright\arraybackslash}p{(\columnwidth - 16\tabcolsep) * \real{0.1124}}
  >{\raggedright\arraybackslash}p{(\columnwidth - 16\tabcolsep) * \real{0.0964}}
  >{\raggedright\arraybackslash}p{(\columnwidth - 16\tabcolsep) * \real{0.0964}}@{}}
\toprule\noalign{}
\begin{minipage}[b]{\linewidth}\raggedright
régime
\end{minipage} & \begin{minipage}[b]{\linewidth}\raggedright
(`PIB (\%)', `10')
\end{minipage} & \begin{minipage}[b]{\linewidth}\raggedright
(`PIB (\%)', `20')
\end{minipage} & \begin{minipage}[b]{\linewidth}\raggedright
(`chômage (pt)', `10')
\end{minipage} & \begin{minipage}[b]{\linewidth}\raggedright
(`chômage (pt)', `20')
\end{minipage} & \begin{minipage}[b]{\linewidth}\raggedright
(`dette (pt PIB)', `10')
\end{minipage} & \begin{minipage}[b]{\linewidth}\raggedright
(`dette (pt PIB)', `20')
\end{minipage} & \begin{minipage}[b]{\linewidth}\raggedright
(`prix (\%)', `10')
\end{minipage} & \begin{minipage}[b]{\linewidth}\raggedright
(`prix (\%)', `20')
\end{minipage} \\
\midrule\noalign{}
\endhead
\bottomrule\noalign{}
\endlastfoot
R2 + BCE et prime & S3 \textgreater{} S2 \textgreater{} S1
\textgreater{} S5 \textgreater{} S4 & S3 \textgreater{} S2
\textgreater{} S5 \textgreater{} S1 \textgreater{} S4 & S1
\textgreater{} S2 \textgreater{} S4 \textgreater{} S5 \textgreater{} S3
& S4 \textgreater{} S1 \textgreater{} S2 \textgreater{} S5
\textgreater{} S3 & S4 \textgreater{} S5 \textgreater{} S1
\textgreater{} S2 \textgreater{} S3 & S4 \textgreater{} S5
\textgreater{} S1 \textgreater{} S2 \textgreater{} S3 & S5
\textgreater{} S2 \textgreater{} S3 \textgreater{} S4 \textgreater{} S1
& S3 \textgreater{} S5 \textgreater{} S2 \textgreater{} S4
\textgreater{} S1 \\
R2 active symétrique & S3 \textgreater{} S2 \textgreater{} S5
\textgreater{} S1 \textgreater{} S4 & S3 \textgreater{} S2
\textgreater{} S5 \textgreater{} S1 \textgreater{} S4 & S4
\textgreater{} S1 \textgreater{} S2 \textgreater{} S5 \textgreater{} S3
& S2 \textgreater{} S3 \textgreater{} S5 \textgreater{} S1
\textgreater{} S4 & S4 \textgreater{} S1 \textgreater{} S5
\textgreater{} S2 \textgreater{} S3 & S4 \textgreater{} S1
\textgreater{} S5 \textgreater{} S2 \textgreater{} S3 & S5
\textgreater{} S3 \textgreater{} S2 \textgreater{} S4 \textgreater{} S1
& S3 \textgreater{} S5 \textgreater{} S2 \textgreater{} S4
\textgreater{} S1 \\
R2, ancrage faible & S3 \textgreater{} S2 \textgreater{} S1
\textgreater{} S5 \textgreater{} S4 & S3 \textgreater{} S2
\textgreater{} S5 \textgreater{} S1 \textgreater{} S4 & S1
\textgreater{} S2 \textgreater{} S4 \textgreater{} S5 \textgreater{} S3
& S4 \textgreater{} S1 \textgreater{} S2 \textgreater{} S5
\textgreater{} S3 & S4 \textgreater{} S1 \textgreater{} S5
\textgreater{} S2 \textgreater{} S3 & S4 \textgreater{} S1
\textgreater{} S5 \textgreater{} S2 \textgreater{} S3 & S5
\textgreater{} S3 \textgreater{} S2 \textgreater{} S4 \textgreater{} S1
& S3 \textgreater{} S5 \textgreater{} S2 \textgreater{} S4
\textgreater{} S1 \\
R2, ancrage fort & S3 \textgreater{} S2 \textgreater{} S1 \textgreater{}
S5 \textgreater{} S4 & S3 \textgreater{} S2 \textgreater{} S5
\textgreater{} S1 \textgreater{} S4 & S1 \textgreater{} S2
\textgreater{} S4 \textgreater{} S5 \textgreater{} S3 & S1
\textgreater{} S4 \textgreater{} S2 \textgreater{} S5 \textgreater{} S3
& S4 \textgreater{} S5 \textgreater{} S1 \textgreater{} S2
\textgreater{} S3 & S4 \textgreater{} S5 \textgreater{} S1
\textgreater{} S2 \textgreater{} S3 & S5 \textgreater{} S2
\textgreater{} S3 \textgreater{} S4 \textgreater{} S1 & S5
\textgreater{} S3 \textgreater{} S2 \textgreater{} S1 \textgreater{}
S4 \\
R2, fermeture lente & S3 \textgreater{} S2 \textgreater{} S1
\textgreater{} S5 \textgreater{} S4 & S3 \textgreater{} S2
\textgreater{} S5 \textgreater{} S1 \textgreater{} S4 & S1
\textgreater{} S2 \textgreater{} S4 \textgreater{} S5 \textgreater{} S3
& S1 \textgreater{} S4 \textgreater{} S2 \textgreater{} S5
\textgreater{} S3 & S4 \textgreater{} S1 \textgreater{} S5
\textgreater{} S2 \textgreater{} S3 & S4 \textgreater{} S1
\textgreater{} S5 \textgreater{} S2 \textgreater{} S3 & S5
\textgreater{} S2 \textgreater{} S3 \textgreater{} S4 \textgreater{} S1
& S3 \textgreater{} S5 \textgreater{} S2 \textgreater{} S4
\textgreater{} S1 \\
R2, fermeture rapide & S3 \textgreater{} S2 \textgreater{} S5
\textgreater{} S1 \textgreater{} S4 & S3 \textgreater{} S2
\textgreater{} S5 \textgreater{} S1 \textgreater{} S4 & S1
\textgreater{} S2 \textgreater{} S4 \textgreater{} S5 \textgreater{} S3
& S4 \textgreater{} S2 \textgreater{} S1 \textgreater{} S5
\textgreater{} S3 & S4 \textgreater{} S5 \textgreater{} S1
\textgreater{} S2 \textgreater{} S3 & S4 \textgreater{} S5
\textgreater{} S1 \textgreater{} S2 \textgreater{} S3 & S5
\textgreater{} S2 \textgreater{} S3 \textgreater{} S4 \textgreater{} S1
& S3 \textgreater{} S5 \textgreater{} S2 \textgreater{} S4
\textgreater{} S1 \\
R2, fermeture sur le chômage & S3 \textgreater{} S2 \textgreater{} S1
\textgreater{} S5 \textgreater{} S4 & S3 \textgreater{} S2
\textgreater{} S5 \textgreater{} S1 \textgreater{} S4 & S1
\textgreater{} S4 \textgreater{} S2 \textgreater{} S5 \textgreater{} S3
& S4 \textgreater{} S1 \textgreater{} S5 \textgreater{} S2
\textgreater{} S3 & S4 \textgreater{} S5 \textgreater{} S1
\textgreater{} S2 \textgreater{} S3 & S4 \textgreater{} S5
\textgreater{} S1 \textgreater{} S2 \textgreater{} S3 & S5
\textgreater{} S3 \textgreater{} S2 \textgreater{} S4 \textgreater{} S1
& S3 \textgreater{} S5 \textgreater{} S2 \textgreater{} S4
\textgreater{} S1 \\
principal (R2 figé, taux exogènes, central) & S3 \textgreater{} S2
\textgreater{} S1 \textgreater{} S5 \textgreater{} S4 & S3
\textgreater{} S2 \textgreater{} S5 \textgreater{} S1 \textgreater{} S4
& S1 \textgreater{} S2 \textgreater{} S4 \textgreater{} S5
\textgreater{} S3 & S1 \textgreater{} S4 \textgreater{} S2
\textgreater{} S5 \textgreater{} S3 & S4 \textgreater{} S5
\textgreater{} S1 \textgreater{} S2 \textgreater{} S3 & S4
\textgreater{} S5 \textgreater{} S1 \textgreater{} S2 \textgreater{} S3
& S5 \textgreater{} S2 \textgreater{} S3 \textgreater{} S4
\textgreater{} S1 & S3 \textgreater{} S5 \textgreater{} S2
\textgreater{} S4 \textgreater{} S1 \\
sans règle budgétaire (R0) & S3 \textgreater{} S2 \textgreater{} S1
\textgreater{} S5 \textgreater{} S4 & S3 \textgreater{} S2
\textgreater{} S5 \textgreater{} S1 \textgreater{} S4 & S1
\textgreater{} S2 \textgreater{} S4 \textgreater{} S5 \textgreater{} S3
& S1 \textgreater{} S4 \textgreater{} S2 \textgreater{} S5
\textgreater{} S3 & S4 \textgreater{} S1 \textgreater{} S5
\textgreater{} S2 \textgreater{} S3 & S4 \textgreater{} S5
\textgreater{} S1 \textgreater{} S2 \textgreater{} S3 & S5
\textgreater{} S3 \textgreater{} S2 \textgreater{} S4 \textgreater{} S1
& S3 \textgreater{} S5 \textgreater{} S2 \textgreater{} S4
\textgreater{} S1 \\
stabilisation (R1) & S3 \textgreater{} S2 \textgreater{} S1
\textgreater{} S5 \textgreater{} S4 & S3 \textgreater{} S2
\textgreater{} S5 \textgreater{} S1 \textgreater{} S4 & S1
\textgreater{} S2 \textgreater{} S4 \textgreater{} S5 \textgreater{} S3
& S1 \textgreater{} S4 \textgreater{} S2 \textgreater{} S5
\textgreater{} S3 & S4 \textgreater{} S5 \textgreater{} S1
\textgreater{} S2 \textgreater{} S3 & S4 \textgreater{} S5
\textgreater{} S1 \textgreater{} S2 \textgreater{} S3 & S5
\textgreater{} S2 \textgreater{} S3 \textgreater{} S4 \textgreater{} S1
& S3 \textgreater{} S5 \textgreater{} S2 \textgreater{} S4
\textgreater{} S1 \\
\end{longtable}

\textbf{Constats.} (1) Les signes des effets sur le PIB et la dette sont
identiques dans les dix régimes pour les cinq stratégies, aux horizons
de 10 et 20 ans. (2) Les classements extrêmes sont stables : S3 a
toujours l'effet le plus favorable sur le PIB et le plus défavorable sur
la dette, S4 l'inverse. (3) L'ordre de S1 et S5 en PIB change selon le
régime. (4) Les effets sur l'emploi et le chômage changent de signe dans
11 cas sur 40 : régime actif symétrique, variantes de fermeture et
environnements \(\sigma\) en sont les principales causes.

\hypertarget{sec:pourquoi}{%
\section{Pourquoi le PIB et la dette sont plus robustes que l'emploi et
le chômage --- et quand ce n'est pas le cas}\label{sec:pourquoi}}

Pour les \textbf{stratégies de test S1-S5}, les simulations confirment
que les résultats sur le PIB et la dette sont plus robustes que ceux sur
l'emploi et le chômage. Trois raisons l'expliquent :

\begin{enumerate}
\def\labelenumi{\arabic{enumi}.}
\tightlist
\item
  \textbf{Le PIB est contraint par la demande à court terme et par le
  potentiel à long terme}, deux mécanismes communs à tous les régimes ;
  l'emploi dépend en plus du partage de l'ajustement entre emploi et
  productivité (élasticité de court terme, \(\sigma\)), qui varie selon
  l'environnement.
\item
  \textbf{Le chômage est un solde} entre deux grandeurs (population
  active et emploi) qui réagissent toutes deux aux mesures ; une petite
  différence relative sur chacune produit une grande différence sur le
  solde.
\item
  \textbf{La dette cumule les soldes} : ses écarts croissent avec le
  temps et dominent les variations des régimes, sauf quand le régime
  lui-même agit sur la dette (règle active, prime souveraine), auquel
  cas l'ampleur change mais pas le signe.
\end{enumerate}

Pour les \textbf{programmes des candidats}, le constat est différent et
doit être énoncé tel quel (tableau suivant). Les signes des effets sur
le \textbf{chômage} sont stables dans toutes les variantes, parce qu'ils
sont dominés par les réformes des retraites, qui modifient la population
active dans le même sens quel que soit le régime ; leur \emph{ampleur}
reste peu fiable. À l'inverse, certains effets sur le \textbf{PIB} et la
\textbf{dette} changent de signe : l'effet de LFI sur le PIB à dix et
vingt ans (proche de zéro et encadré par des valeurs des deux signes),
l'effet du RN sur la dette (petit à toutes les échéances), l'effet
d'Horizons sur la dette à vingt ans. La règle générale qui ressort des
deux exercices est que \textbf{la robustesse du signe dépend du rapport
entre l'ampleur de l'effet et l'étendue des variantes}, plus que de la
nature de l'indicateur : un effet important est robuste, un effet proche
de zéro ne l'est pas, quelle que soit la variable.

\hypertarget{sec:robustesse-programmes}{%
\section{Robustesse des résultats des
programmes}\label{sec:robustesse-programmes}}

\textbf{Tableau : robustesse de la lecture A des programmes aux dix
régimes et aux trois environnements alternatifs} (signe identique dans
les 13 variantes ; étendue max − min). Source : modèle du projet
(configuration P4, convention « part constante »), campagne
phase7/campagne.py.

\begin{longtable}[]{@{}
  >{\raggedright\arraybackslash}p{(\columnwidth - 12\tabcolsep) * \real{0.1429}}
  >{\raggedright\arraybackslash}p{(\columnwidth - 12\tabcolsep) * \real{0.1429}}
  >{\raggedright\arraybackslash}p{(\columnwidth - 12\tabcolsep) * \real{0.1429}}
  >{\raggedright\arraybackslash}p{(\columnwidth - 12\tabcolsep) * \real{0.1429}}
  >{\raggedright\arraybackslash}p{(\columnwidth - 12\tabcolsep) * \real{0.1429}}
  >{\raggedright\arraybackslash}p{(\columnwidth - 12\tabcolsep) * \real{0.1429}}
  >{\raggedright\arraybackslash}p{(\columnwidth - 12\tabcolsep) * \real{0.1429}}@{}}
\toprule\noalign{}
\begin{minipage}[b]{\linewidth}\raggedright
candidat --- indicateur
\end{minipage} & \begin{minipage}[b]{\linewidth}\raggedright
signe robuste à 5 ans
\end{minipage} & \begin{minipage}[b]{\linewidth}\raggedright
signe robuste à 10 ans
\end{minipage} & \begin{minipage}[b]{\linewidth}\raggedright
signe robuste à 20 ans
\end{minipage} & \begin{minipage}[b]{\linewidth}\raggedright
étendue à 5 ans
\end{minipage} & \begin{minipage}[b]{\linewidth}\raggedright
étendue à 10 ans
\end{minipage} & \begin{minipage}[b]{\linewidth}\raggedright
étendue à 20 ans
\end{minipage} \\
\midrule\noalign{}
\endhead
\bottomrule\noalign{}
\endlastfoot
Le Pen (RN) --- PIB (\%) & oui & oui & oui & +0,76 & +0,72 & +1,30 \\
Le Pen (RN) --- chômage (pt) & oui & oui & oui & +0,39 & +0,49 &
+1,24 \\
Le Pen (RN) --- dette (pt PIB) & non & non & non & +2,01 & +6,95 &
+22,59 \\
Le Pen (RN) --- emploi (milliers) & oui & oui & oui & +127,92 & +157,84
& +400,68 \\
Le Pen (RN) --- pouvoir d'achat (\%) & oui & non & non & +0,37 & +1,22 &
+3,89 \\
Le Pen (RN) --- solde (pt PIB) & oui & non & non & +0,42 & +0,80 &
+2,08 \\
Mélenchon (LFI) --- PIB (\%) & oui & non & non & +1,88 & +2,38 &
+2,53 \\
Mélenchon (LFI) --- chômage (pt) & oui & oui & oui & +0,88 & +1,70 &
+2,09 \\
Mélenchon (LFI) --- dette (pt PIB) & oui & oui & oui & +10,22 & +29,10 &
+70,76 \\
Mélenchon (LFI) --- emploi (milliers) & oui & non & oui & +290,01 &
+560,03 & +681,77 \\
Mélenchon (LFI) --- pouvoir d'achat (\%) & oui & oui & oui & +1,38 &
+3,70 & +6,26 \\
Mélenchon (LFI) --- solde (pt PIB) & oui & oui & oui & +1,56 & +3,19 &
+6,03 \\
Philippe (Horizons) --- PIB (\%) & oui & oui & oui & +1,90 & +2,80 &
+2,58 \\
Philippe (Horizons) --- chômage (pt) & oui & oui & oui & +1,34 & +2,37 &
+2,63 \\
Philippe (Horizons) --- dette (pt PIB) & oui & oui & non & +4,10 &
+13,57 & +44,26 \\
Philippe (Horizons) --- emploi (milliers) & oui & non & non & +458,15 &
+820,85 & +913,36 \\
Philippe (Horizons) --- pouvoir d'achat (\%) & non & non & non & +2,26 &
+5,76 & +9,44 \\
Philippe (Horizons) --- solde (pt PIB) & oui & oui & non & +0,82 & +2,04
& +4,17 \\
\end{longtable}

Part des couples (candidat, horizon) dont le signe est robuste aux
treize variantes, par indicateur : PIB (\%) : 78 \% ; chômage (pt) : 100
\% ; dette (pt PIB) : 56 \% ; emploi (milliers) : 67 \% ; pouvoir
d'achat (\%) : 44 \% ; solde (pt PIB) : 67 \%.

\hypertarget{sec:robustesse-parametres}{%
\section{Paramètres structurels}\label{sec:robustesse-parametres}}

Amplitude maximale de l'effet à 20 ans sur les stratégies S1-S5 quand un
paramètre passe de sa valeur basse à sa valeur haute (régime principal ;
paramètres monétaires sous le régime BCE + prime) :

\begin{longtable}[]{@{}
  >{\raggedright\arraybackslash}p{(\columnwidth - 4\tabcolsep) * \real{0.3333}}
  >{\raggedright\arraybackslash}p{(\columnwidth - 4\tabcolsep) * \real{0.3333}}
  >{\raggedright\arraybackslash}p{(\columnwidth - 4\tabcolsep) * \real{0.3333}}@{}}
\toprule\noalign{}
\begin{minipage}[b]{\linewidth}\raggedright
paramètre
\end{minipage} & \begin{minipage}[b]{\linewidth}\raggedright
amplitude PIB 20 ans (pt de \%)
\end{minipage} & \begin{minipage}[b]{\linewidth}\raggedright
amplitude dette 20 ans (pt)
\end{minipage} \\
\midrule\noalign{}
\endhead
\bottomrule\noalign{}
\endlastfoot
PSI\_CHOMAGE & 0,02 & 0,55 \\
PSI\_TAUX & 0,02 & 0,14 \\
SEMI\_ELAST\_IS & 0,00 & 1,29 \\
THETA\_TRESORERIE & 0,01 & 0,73 \\
bund\_transmission & 0,00 & 0,00 \\
div\_lissage & 0,01 & 1,41 \\
mu\_taux & 0,00 & 0,02 \\
prime\_dette & 0,00 & 0,00 \\
\end{longtable}

\begin{longtable}[]{@{}
  >{\raggedright\arraybackslash}p{(\columnwidth - 4\tabcolsep) * \real{0.3333}}
  >{\raggedright\arraybackslash}p{(\columnwidth - 4\tabcolsep) * \real{0.3333}}
  >{\raggedright\arraybackslash}p{(\columnwidth - 4\tabcolsep) * \real{0.3333}}@{}}
\toprule\noalign{}
\begin{minipage}[b]{\linewidth}\raggedright
paramètre
\end{minipage} & \begin{minipage}[b]{\linewidth}\raggedright
amplitude PIB 20 ans (pt de \%)
\end{minipage} & \begin{minipage}[b]{\linewidth}\raggedright
amplitude dette 20 ans (pt)
\end{minipage} \\
\midrule\noalign{}
\endhead
\bottomrule\noalign{}
\endlastfoot
bce\_activite & 0,13 & 6,47 \\
bce\_inflation & 0,00 & 0,27 \\
bund\_transmission & 0,00 & 2,32 \\
entreprises\_vitesse & 0,08 & 0,69 \\
menages\_vitesse & 0,00 & 0,13 \\
prime\_dette & 0,02 & 7,09 \\
prime\_solde & 0,01 & 7,63 \\
\end{longtable}

Aucun paramètre structurel testé ne modifie le classement des stratégies
en PIB ou en dette ; les paramètres de prime souveraine sont les plus
influents sur la dette des stratégies déficitaires.

\hypertarget{sec:robustesse-env}{%
\section{Environnements alternatifs}\label{sec:robustesse-env}}

Les environnements \(\sigma = 0{,}3\) et \(\sigma = 0{,}6\) et
l'équation de prix v0.9 ont chacun été calibrés sur les mêmes cibles.
Pour les stratégies de test, \(\sigma\) modifie l'effet sur le PIB de
moins de 0,05 point à 20 ans mais l'emploi et la dette davantage
(jusqu'à 7 points de dette) ; l'équation de prix modifie la dette de 2,6
points au plus. Les résultats des programmes dans ces environnements
figurent au chapitre \ref{sec:resultats}.

\hypertarget{sec:robuste-fragile}{%
\section{Ce qui est robuste, ce qui est
fragile}\label{sec:robuste-fragile}}

\textbf{Résultats robustes (signe stable dans toutes les variantes,
ampleur comparable).} - Le sens des effets des réformes des retraites
sur la population active, le PIB potentiel et le PIB à 10 et 20 ans. -
Le sens des effets des programmes sur la dette relativement à la
référence, dans la lecture A, pour LFI (baisse) et pour Horizons jusqu'à
dix ans (hausse) ; pas pour le RN, dont l'effet sur la dette est proche
de zéro. - Le sens des effets des programmes sur le chômage (dominés par
la population active), mais pas leur ampleur. - L'ordre de grandeur de
l'effort budgétaire implicite des objectifs de déficit (lecture B),
relativement aux mesures documentées, sous une hypothèse de taux donnée.

\textbf{Résultats fragiles (signe ou ampleur dépendant des hypothèses).}
- L'ampleur des effets sur le chômage à toutes les échéances ; leur
signe pour les stratégies de test S1-S5. - Les effets proches de zéro :
PIB de LFI à dix et vingt ans, dette du RN, dette d'Horizons à vingt
ans, pouvoir d'achat du RN et d'Horizons. - Les effets de court terme
des mesures de demande, qui dépendent des multiplicateurs (régime
monétaire, fermeture). - Les rendements nets des impôts sur le capital
et les patrimoines (fourchette des réactions des contribuables). - Tous
les résultats à 20 ans qui dépendent de la dynamique de la dette (taux
longs, prime souveraine). - L'effet de l'investissement hors budget,
dont l'additionnalité et le financement ne sont pas modélisés.

\hypertarget{sec:discussion}{%
\chapter{Discussion générale et comparaison des philosophies
économiques}\label{sec:discussion}}

Ce chapitre compare les logiques économiques des trois programmes. Il ne
propose pas de classement unique : un programme peut être préférable
selon un critère et défavorable selon un autre, et le poids accordé à
chaque critère relève du choix politique, non de l'analyse économique.
L'objectif est d'expliciter les arbitrages, de distinguer les critères
et de séparer les conclusions qui résistent aux variantes de celles qui
dépendent fortement des hypothèses.

\hypertarget{sec:philosophies}{%
\section{Trois philosophies économiques}\label{sec:philosophies}}

\textbf{Une logique de demande et de redistribution (LFI).} Le programme
de La France insoumise s'inscrit dans la tradition keynésienne et
post-keynésienne : la faiblesse de l'activité tient à l'insuffisance de
la demande et à la répartition des revenus en faveur du capital ; la
relance passe par les revenus des ménages modestes (salaires,
prestations, emploi public) et par l'investissement public et parapublic
; son financement par l'imposition des hauts patrimoines et des grandes
entreprises est jugé peu dommageable pour l'activité, ces contribuables
ayant une faible propension à consommer. Le programme assume une
intervention publique forte dans le financement de l'économie (pôle
public bancaire, épargne réglementée) et une rupture avec certaines
règles européennes.

\textbf{Une logique de pouvoir d'achat ciblé, de souveraineté et de
redressement par la dépense (RN).} Le programme du Rassemblement
national combine des transferts ciblés (TVA sur l'énergie, jeunes,
familles), le maintien d'un âge de départ plus précoce, et un
redressement des comptes fondé sur des économies sur des postes désignés
(immigration, contribution européenne, organisation de l'État et des
collectivités). Sa philosophie est moins doctrinale que celle des deux
autres : elle emprunte à l'économie de l'offre (impôts de production,
relocalisation) et à la redistribution (retraites, pouvoir d'achat),
avec une forte dimension de priorité nationale et de souveraineté.

\textbf{Une logique d'offre et de consolidation (Horizons).} Le
programme d'Édouard Philippe s'inscrit dans la tradition de l'économie
de l'offre et de la consolidation budgétaire par la dépense : augmenter
la quantité de travail (âge, durée, temps de travail), réduire la
dépense publique et la taille de l'administration, alléger les
prélèvements pesant sur la production, stabiliser l'environnement fiscal
et réglementaire, et faire porter la moitié de l'effort sur les dépenses
sociales. Il revendique un effet de croissance (« un quart de l'effort
sur la croissance »).

\hypertarget{sec:redistribution-incitations}{%
\section{Redistribution et
incitations}\label{sec:redistribution-incitations}}

Le programme de LFI redistribue des hauts patrimoines et des grandes
entreprises vers les ménages modestes, les agents publics et les
retraités ; son risque principal, du point de vue des incitations, est
l'érosion des assiettes (contribuables mobiles, optimisation) et l'effet
sur l'épargne et l'investissement des entreprises. Le modèle donne des
rendements nets inférieurs aux rendements annoncés pour les impôts sur
le capital (chapitre \ref{sec:transmission}) mais ne représente pas
d'effet négatif de ces impôts sur l'investissement productif ; la
littérature est partagée sur ce point. Le programme d'Horizons réduit
les prestations (RSA, assurance chômage, dépenses sociales) pour
renforcer les incitations au travail ; son risque principal est l'effet
sur les ménages modestes et sur la demande. Le RN combine des transferts
universels (TVA sur l'énergie, dont le bénéfice en euros croît avec la
consommation) et des économies ciblées sur une catégorie de population
définie par la nationalité, dont la compatibilité juridique est en
partie douteuse.

\hypertarget{sec:demande-offre}{%
\section{Demande et offre}\label{sec:demande-offre}}

Les résultats illustrent l'arbitrage temporel entre demande et offre
(chapitre \ref{sec:resultats}). Les mesures de demande (LFI) produisent
des effets favorables rapides sur l'activité et l'emploi, qui
s'estompent avec la fermeture d'offre et la hausse des prix ; les
mesures d'offre (Horizons : âge, durée, temps de travail) produisent des
effets favorables lents sur le PIB potentiel, précédés d'une hausse du
chômage tant que la demande de travail ne suit pas l'augmentation de la
population active. Les mesures d'offre de travail dominent cet arbitrage
dans les trois programmes, dans des sens opposés : l'abaissement de
l'âge de départ (LFI, RN) réduit le potentiel de manière croissante ; la
fin des 35 heures et le relèvement de l'âge (Horizons) l'augmentent.
C'est le résultat le plus robuste de l'étude : il ne dépend ni du
régime, ni de l'environnement, seulement de l'hypothèse --- conforme à
l'expérience des réformes de 2010 et 2023 --- selon laquelle un report
de l'âge légal maintient une partie des seniors en activité.

\hypertarget{sec:soutenabilite}{%
\section{Dépenses publiques et soutenabilité
budgétaire}\label{sec:soutenabilite}}

Dans un contexte où le taux d'intérêt de marché dépasse la croissance
nominale, toute dégradation durable du solde primaire se cumule
rapidement dans la dette, et toute amélioration aussi. La comparaison
des deux lectures montre que les \textbf{objectifs budgétaires affichés
par le RN et Horizons impliquent un effort non ventilé beaucoup plus
important que leurs mesures documentées}, et que cet effort, s'il était
réalisé par la dépense, aurait selon le modèle des effets récessifs
importants. Pour LFI, l'amélioration du solde dans la lecture A repose
sur des recettes nouvelles dont les montants sont, pour plusieurs
d'entre elles, des bornes hautes (taxes sans canal propre, érosion non
représentée) ; l'investissement hors budget de 100 Md€ par an n'affecte
pas le solde public dans le modèle, mais il crée une dette du pôle
public dont la requalification éventuelle en dette publique n'est pas
examinée.

\hypertarget{sec:protection-emploi}{%
\section{Protection sociale et emploi}\label{sec:protection-emploi}}

Les programmes opposent deux conceptions de la protection sociale :
extension (LFI : revalorisations, garantie d'autonomie, retraite plus
précoce) ou recentrage (Horizons : durée d'indemnisation, RSA, dépenses
sociales). Dans le modèle, les effets sur l'emploi passent surtout par
la demande (extension) et par l'offre de travail (recentrage) ; les
effets d'incitation au retour à l'emploi des réformes de l'assurance
chômage et du RSA sont représentés de manière partielle (participation,
chômage structurel), et leur ampleur est incertaine.

\hypertarget{sec:souverainete}{%
\section{Souveraineté économique et intégration
commerciale}\label{sec:souverainete}}

Les programmes de LFI et du RN comportent des mesures qui supposent de
s'écarter du droit européen (contrôle des capitaux, priorité nationale,
contribution européenne, TVA sur les carburants) ; celui d'Horizons
s'inscrit dans le cadre existant. Le modèle ne peut évaluer que les
conséquences directes et chiffrables de ces écarts (retrait, report,
amendes, prime de risque) ; les conséquences indirectes (perte de
crédibilité, réactions des partenaires, accès aux instruments de la BCE)
relèvent de l'analyse qualitative (chapitres \ref{sec:europe} et
\ref{sec:comportements}).

\hypertarget{sec:invest}{%
\section{Investissement public et privé}\label{sec:invest}}

LFI mise sur l'investissement public et parapublic (adaptation
climatique, pôle public bancaire) ; Horizons et le RN sur
l'investissement privé par la baisse des prélèvements sur la production.
Le modèle attribue à l'investissement public un multiplicateur croissant
dans le temps (capital public) et à la baisse des impôts de production
un effet d'activité faible (une grande partie est répercutée sur les
prix ou conservée en marges) ; il ne représente pas les effets de
localisation des activités que l'on attend de la baisse des impôts de
production, ni les effets d'éviction possibles d'un pôle public
d'investissement de grande taille.

\hypertarget{sec:ct-lt}{%
\section{Court terme et long terme}\label{sec:ct-lt}}

Les horizons modifient la lecture des programmes : à cinq ans, les
mesures de demande et la baisse de l'âge de la retraite (qui réduit le
chômage en réduisant la population active) paraissent favorables ; à dix
et vingt ans, le potentiel et la dette dominent. Le choix d'un horizon
de référence n'est pas neutre ; l'étude présente les trois.

\hypertarget{sec:synthese-criteres}{%
\section{Synthèse par critère}\label{sec:synthese-criteres}}

\begin{longtable}[]{@{}
  >{\raggedright\arraybackslash}p{(\columnwidth - 4\tabcolsep) * \real{0.3333}}
  >{\raggedright\arraybackslash}p{(\columnwidth - 4\tabcolsep) * \real{0.3333}}
  >{\raggedright\arraybackslash}p{(\columnwidth - 4\tabcolsep) * \real{0.3333}}@{}}
\toprule\noalign{}
\begin{minipage}[b]{\linewidth}\raggedright
Critère
\end{minipage} & \begin{minipage}[b]{\linewidth}\raggedright
Ce que montrent les simulations
\end{minipage} & \begin{minipage}[b]{\linewidth}\raggedright
Robustesse
\end{minipage} \\
\midrule\noalign{}
\endhead
\bottomrule\noalign{}
\endlastfoot
Activité à 5 ans (lecture A) & LFI au-dessus de la référence (demande,
investissement hors budget) ; RN au-dessous (retraites) ; Horizons
au-dessus (offre de travail), avec un chômage plus élevé & signes
robustes ; ampleurs dépendantes des multiplicateurs et de la traduction
des recettes \\
Activité à 10-20 ans (lecture A) & Horizons au-dessus et croissant ; RN
au-dessous ; LFI légèrement au-dessous, du fait de la baisse de l'offre
de travail & robuste pour Horizons et le RN ; non robuste pour LFI
(effet proche de zéro, signe variable selon le régime) ; dépend de
l'hypothèse d'exogénéité de la productivité \\
Chômage & effets amples et de signe parfois instable ; ne pas
interpréter les niveaux & fragile \\
Prix et pouvoir d'achat & LFI : prix et pouvoir d'achat plus élevés ; RN
: pouvoir d'achat légèrement inférieur à cinq ans, proche de la
référence à vingt ans ; Horizons : pouvoir d'achat moyen légèrement plus
faible & robuste en signe à 5 ans \\
Finances publiques (lecture A) & LFI : solde amélioré (recettes en
bornes hautes) ; RN : quasi neutre ; Horizons : légère dégradation
(mesures documentées) & dépend des rendements fiscaux et des taux \\
Objectifs affichés (lecture B) & effort implicite important et récessif
pour le RN et Horizons ; faible à court terme pour LFI & ordre de
grandeur robuste à hypothèse de taux donnée \\
Distribution (statique) & LFI : prélèvements sur le haut de la
distribution ; RN : gains diffus (TVA énergie) ; Horizons : pertes
concentrées sur les bénéficiaires du RSA & couverture partielle \\
Générations & baisse de l'âge (LFI, RN) favorable aux générations
proches de la retraite, coûteuse pour les finances publiques ; hausse
(Horizons) défavorable aux générations nées après 1970 & robuste en
sens \\
Contraintes européennes & mesures en tension pour LFI et le RN, aucune
pour les mesures documentées d'Horizons & qualification juridique
sommaire \\
\end{longtable}

\hypertarget{sec:conclusions-robustes}{%
\section{Conclusions robustes et conclusions dépendantes des
hypothèses}\label{sec:conclusions-robustes}}

\textbf{Résistent aux variantes :} (1) le sens des effets des réformes
des retraites sur l'activité à moyen terme ; (2) le fait que les mesures
documentées du RN et d'Horizons ne représentent qu'une faible part de
l'effort nécessaire à leurs objectifs de déficit ; (3) le sens des
effets sur la dette relativement à la référence lorsqu'ils sont
importants (lecture A : baisse avec LFI, hausse avec Horizons jusqu'à
dix ans ; l'effet du RN, proche de zéro, n'a pas de signe robuste) ; (4)
la qualification européenne des mesures les plus problématiques.

\textbf{Dépendent fortement des hypothèses :} (1) l'ampleur des effets
de court terme des mesures de demande (multiplicateurs, fermeture) ; (2)
les rendements des impôts sur le capital et les patrimoines ; (3)
l'effet de l'investissement hors budget ; (4) l'ampleur des effets sur
le chômage ; (5) les effets à vingt ans sur la dette (taux longs, prime
souveraine) ; (6) les effets récessifs de la lecture B (référence,
composition de l'effort, crédibilité).

\hypertarget{sec:conclusion}{%
\chapter{Conclusion générale}\label{sec:conclusion}}

Cette étude posait une question de méthode avant d'être une question
politique : \emph{comment comparer rigoureusement les conséquences
économiques et sociales de programmes présidentiels qui ne parlent pas
le même langage, ne sont pas chiffrés avec la même précision et ne
visent pas les mêmes objectifs ?} La réponse proposée tient en trois
principes, qu'il faut rappeler avant les résultats, parce qu'ils
conditionnent ce que ces résultats veulent dire.

\begin{enumerate}
\def\labelenumi{\arabic{enumi}.}
\tightlist
\item
  \textbf{Comparer des écarts, pas des niveaux.} Chaque programme est
  évalué par l'écart entre deux simulations du même modèle, dans le même
  environnement économique et sous le même régime de politique
  économique, avec et sans le programme. Les niveaux absolus de dette ou
  de chômage dépendent d'une trajectoire de référence elle-même très
  incertaine ; les écarts en dépendent beaucoup moins.
\item
  \textbf{Séparer ce que les candidats ont documenté de ce qu'ils ont
  seulement annoncé.} La lecture A simule les mesures identifiées et
  chiffrées ; la lecture B simule l'effort nécessaire pour atteindre les
  objectifs de déficit affichés, avec une composition qui est une
  hypothèse de l'étude. Les deux lectures ne sont pas de même nature et
  ne sont jamais additionnées ni confondues.
\item
  \textbf{Présenter des fourchettes et des conditions plutôt que des
  points.} Chaque résultat est accompagné de sa sensibilité au régime
  budgétaire et monétaire, à l'environnement (élasticité de
  substitution, équation des prix), à l'hypothèse de taux longs, aux
  réactions européennes et à une mise en œuvre partielle ou retardée.
\end{enumerate}

Les simulations sont des \textbf{scénarios conditionnels}. Elles
répondent à la question « que se passerait-il, selon ce modèle et sous
ces hypothèses, si ces mesures étaient appliquées ? », non à la question
« que se passera-t-il ? ».

\hypertarget{sec:reponse}{%
\section{Réponse à la problématique}\label{sec:reponse}}

\textbf{Tableau : synthèse des résultats, horizons de 5 et 10 ans.}
Lecture A (mesures documentées) avec la fourchette de ses treize
variantes (dix régimes, trois environnements) ; lecture B (objectifs
budgétaires). Écarts à la référence. Source : modèle du projet
(configuration P4, convention « part constante »), campagne
phase7/campagne.py.

\begin{longtable}[]{@{}
  >{\raggedright\arraybackslash}p{(\columnwidth - 8\tabcolsep) * \real{0.2000}}
  >{\raggedright\arraybackslash}p{(\columnwidth - 8\tabcolsep) * \real{0.2000}}
  >{\raggedright\arraybackslash}p{(\columnwidth - 8\tabcolsep) * \real{0.2000}}
  >{\raggedright\arraybackslash}p{(\columnwidth - 8\tabcolsep) * \real{0.2000}}
  >{\raggedright\arraybackslash}p{(\columnwidth - 8\tabcolsep) * \real{0.2000}}@{}}
\toprule\noalign{}
\begin{minipage}[b]{\linewidth}\raggedright
candidat --- indicateur --- horizon
\end{minipage} & \begin{minipage}[b]{\linewidth}\raggedright
lecture A
\end{minipage} & \begin{minipage}[b]{\linewidth}\raggedright
min. variantes A
\end{minipage} & \begin{minipage}[b]{\linewidth}\raggedright
max. variantes A
\end{minipage} & \begin{minipage}[b]{\linewidth}\raggedright
lecture B
\end{minipage} \\
\midrule\noalign{}
\endhead
\bottomrule\noalign{}
\endlastfoot
Mélenchon (LFI) --- PIB (\%) --- 5 ans & +1,84 & +0,96 & +2,84 &
+2,66 \\
Mélenchon (LFI) --- PIB (\%) --- 10 ans & −0,54 & −1,47 & +0,91 &
−1,19 \\
Mélenchon (LFI) --- chômage (pt) --- 5 ans & −4,51 & −4,96 & −4,08 &
−5,02 \\
Mélenchon (LFI) --- chômage (pt) --- 10 ans & −2,78 & −3,79 & −2,09 &
−2,70 \\
Mélenchon (LFI) --- pouvoir d'achat (\%) --- 5 ans & +3,35 & +2,84 &
+4,22 & +5,58 \\
Mélenchon (LFI) --- pouvoir d'achat (\%) --- 10 ans & +4,33 & +3,05 &
+6,75 & +3,51 \\
Mélenchon (LFI) --- dette (pt PIB) --- 5 ans & −20,11 & −26,85 & −16,63
& −18,74 \\
Mélenchon (LFI) --- dette (pt PIB) --- 10 ans & −38,14 & −56,96 & −27,86
& −36,08 \\
Le Pen (RN) --- PIB (\%) --- 5 ans & −1,95 & −2,03 & −1,27 & −9,59 \\
Le Pen (RN) --- PIB (\%) --- 10 ans & −2,89 & −3,38 & −2,66 & −11,12 \\
Le Pen (RN) --- chômage (pt) --- 5 ans & −1,84 & −2,20 & −1,80 &
+4,65 \\
Le Pen (RN) --- chômage (pt) --- 10 ans & −1,65 & −1,79 & −1,30 &
+8,52 \\
Le Pen (RN) --- pouvoir d'achat (\%) --- 5 ans & −0,90 & −0,96 & −0,59 &
−8,30 \\
Le Pen (RN) --- pouvoir d'achat (\%) --- 10 ans & −0,62 & −1,19 & +0,03
& −13,38 \\
Le Pen (RN) --- dette (pt PIB) --- 5 ans & −0,34 & −1,91 & +0,10 &
+2,79 \\
Le Pen (RN) --- dette (pt PIB) --- 10 ans & −1,60 & −6,21 & +0,74 &
−5,64 \\
Philippe (Horizons) --- PIB (\%) --- 5 ans & +0,96 & +0,19 & +2,09 &
−5,18 \\
Philippe (Horizons) --- PIB (\%) --- 10 ans & +2,45 & +1,67 & +4,47 &
−3,45 \\
Philippe (Horizons) --- chômage (pt) --- 5 ans & +2,56 & +1,70 & +3,04 &
+6,94 \\
Philippe (Horizons) --- chômage (pt) --- 10 ans & +2,56 & +0,69 & +3,05
& +8,85 \\
Philippe (Horizons) --- pouvoir d'achat (\%) --- 5 ans & −1,08 & −1,42 &
+0,84 & −11,12 \\
Philippe (Horizons) --- pouvoir d'achat (\%) --- 10 ans & −1,05 & −2,20
& +3,56 & −16,67 \\
Philippe (Horizons) --- dette (pt PIB) --- 5 ans & +5,28 & +3,91 & +8,01
& +3,40 \\
Philippe (Horizons) --- dette (pt PIB) --- 10 ans & +8,42 & +3,93 &
+17,50 & −11,59 \\
\end{longtable}

Le tableau précédent résume l'essentiel. Il appelle quatre constats
généraux.

\textbf{Premier constat : les trois programmes ne se distinguent pas
d'abord par leurs multiplicateurs, mais par ce qu'ils font de l'offre de
travail.} La décomposition mesure par mesure (chapitre
\ref{sec:resultats}) montre que les mesures d'offre de travail dominent
les effets à moyen terme sur l'activité. L'abaissement de l'âge
d'ouverture des droits et de la durée d'assurance (LFI, RN) réduit la
population active et le PIB potentiel ; la fin des 35 heures et le
relèvement de l'âge de départ (Horizons) les augmentent, la première par
les heures travaillées (avec moins d'emplois), le second par la
population active. Les effets de demande des autres mesures --- relance
par les revenus et l'investissement pour LFI, baisses ciblées d'impôts
pour le RN, pacte sur les impôts de production pour Horizons ---
dominent à cinq ans mais s'estompent ensuite sous l'effet de la
fermeture d'offre et de la hausse des prix. C'est pourquoi les
classements par le PIB à cinq ans et à dix ans ne coïncident pas, et
pourquoi aucun horizon unique ne peut résumer un programme.

\textbf{Deuxième constat : la précision documentaire inégale des
programmes se traduit par des incertitudes de nature différente.} Le
paquet 2027 de LFI est chiffré mesure par mesure par le candidat et
recoupe exactement ses totaux ; son incertitude principale tient à la
\emph{traduction} de recettes sans précédent (taxe sur les très grands
patrimoines, fiscalité des multinationales, superprofits) et à
l'additionnalité des 100 Md€ annuels d'investissement hors budget, dont
la variante sans additionnalité ramène l'effet sur le PIB à cinq ans à
zéro. Les programmes du RN et d'Horizons sont d'abord des
\emph{trajectoires} (140 Md€ d'économies et un déficit de 3 \% en 2030 ;
près de 150 Md€ et un déficit sous 2,5 \% en fin de quinquennat) dont la
composition n'est que partiellement publiée ; leur incertitude
principale tient au contenu des économies, que la lecture B ne peut que
supposer.

\textbf{Troisième constat : dans ce modèle, l'écart entre les objectifs
de déficit annoncés et la dynamique spontanée des finances publiques est
considérable.} La trajectoire de référence à politique inchangée dérive
(taux apparent de la dette supérieur à la croissance, déficit primaire),
si bien qu'atteindre un déficit de 3 \% ou de 2,5 \% du PIB exige un
effort qui croît avec le temps : plusieurs points de PIB dès cinq ans,
davantage à dix ans dans l'hypothèse d'une OAT à 4,89 \%, moins avec une
OAT à 3,82 \%. Avec les multiplicateurs du régime principal, cet effort
réduit l'activité et l'emploi, et la baisse du ratio de dette est
beaucoup plus lente que celle du déficit, le dénominateur (le PIB
nominal) étant lui-même réduit. Ce résultat porte sur l'\emph{ampleur}
de l'ajustement implicite dans les objectifs : il ne dit pas que ces
objectifs sont inatteignables (une consolidation crédible peut réduire
la prime de risque, ce que le régime principal ne représente pas, et les
multiplicateurs de consolidation sont débattus, chapitre
\ref{sec:theorie}) ; il dit qu'ils supposent un effort d'une ampleur
sans précédent récent en France, que les mesures documentées ne
contiennent pas encore.

\textbf{Quatrième constat : les contraintes européennes pèsent plus sur
la faisabilité de certaines mesures que sur les agrégats.} Les mesures
les plus exposées (TVA à taux réduit sur les carburants, réduction
unilatérale de la contribution au budget de l'Union, priorité nationale,
contrôle des capitaux) sont soit modestes à l'échelle macroéconomique,
soit non quantifiables. Les variantes « contentieux » et « tension »
modifient peu les effets à cinq ans ; la prime souveraine pèse surtout
sur la dette à long terme. Ces variantes ne sont pas des prévisions de
sanctions : elles bornent un risque.

\hypertarget{sec:robustes}{%
\section{Résultats robustes}\label{sec:robustes}}

Sont qualifiés de robustes les résultats dont le signe ne change dans
aucune des treize variantes de la lecture A (dix régimes, trois
environnements) et qui ne reposent pas sur une hypothèse de traduction
isolée :

\begin{itemize}
\tightlist
\item
  \textbf{le sens des effets des mesures d'offre de travail} (retraites,
  durée du travail) sur la population active, les heures travaillées et
  le PIB à moyen terme ;
\item
  \textbf{le sens des effets importants} : PIB de LFI à cinq ans, PIB du
  RN et d'Horizons à cinq et dix ans, baisse de la dette avec LFI,
  hausse avec Horizons jusqu'à dix ans, et sens des effets sur le
  chômage, dominés par la population active (chapitre
  \ref{sec:robustesse}) ; à l'inverse, les effets proches de zéro (PIB
  de LFI à dix ans, dette du RN) n'ont pas de signe robuste ;
\item
  \textbf{l'écart entre mesures documentées et objectifs budgétaires}
  pour le RN et Horizons : quelle que soit l'hypothèse de taux, les
  mesures identifiées ne représentent qu'une faible part de l'effort
  nécessaire ;
\item
  \textbf{le profil temporel} des effets des mesures de demande (maximum
  à court terme, atténuation ensuite) ;
\item
  \textbf{la qualification européenne} des mesures les plus
  problématiques, qui relève du droit positif.
\end{itemize}

\hypertarget{sec:incertains}{%
\section{Résultats incertains}\label{sec:incertains}}

\begin{itemize}
\tightlist
\item
  \textbf{L'ampleur des effets sur le chômage.} Le chômage du modèle est
  la différence entre une population active en partie exogène et un
  emploi déterminé par la production ; les réformes qui modifient
  fortement la population active produisent des écarts de chômage plus
  amples que ce que suggérerait la loi d'Okun. Pour les programmes, le
  signe est stable (il suit la population active) ; pour les stratégies
  de test, il change dans plusieurs variantes. Les signes sont
  informatifs, les ordres de grandeur beaucoup moins.
\item
  \textbf{Les rendements des impôts sur le capital et le patrimoine},
  qui dépendent de comportements d'évitement et de mobilité mal mesurés
  (chapitre \ref{sec:comportements}).
\item
  \textbf{L'effet de l'investissement hors budget}, entièrement
  conditionné par son additionnalité et son financement.
\item
  \textbf{Les effets à vingt ans}, qui dépendent de la trajectoire de
  référence, des taux longs et de l'hypothèse de productivité exogène.
\item
  \textbf{Les effets récessifs de la lecture B}, qui dépendent de la
  référence, de la composition supposée de l'effort et de l'absence
  d'effets de crédibilité.
\end{itemize}

\hypertarget{sec:arbitrages}{%
\section{Arbitrages économiques}\label{sec:arbitrages}}

Les trois programmes incarnent des arbitrages différents, que le modèle
permet de rendre explicites sans les trancher :

\begin{itemize}
\tightlist
\item
  \textbf{demande ou offre} : LFI mise sur la demande (revenus, services
  publics, investissement) et accepte une baisse de l'offre de travail
  (retraites) ; Horizons mise sur l'offre (âge de départ, durée du
  travail, impôts de production) et accepte une phase de chômage plus
  élevé, le temps que la demande de travail rejoigne l'offre ; le RN
  combine une baisse de l'offre de travail (retraites) et des baisses
  ciblées de prélèvements ;
\item
  \textbf{redistribution ou incitations} : les prélèvements de LFI
  portent sur le haut de la distribution et le capital, avec des risques
  d'érosion d'assiette ; les mesures d'Horizons sur le RSA et
  l'assurance chômage visent les incitations au travail, avec des pertes
  concentrées sur les ménages modestes ;
\item
  \textbf{court terme ou long terme} : les mesures favorables à
  l'activité à cinq ans ne sont pas toujours celles qui améliorent le
  potentiel à dix ou vingt ans, et réciproquement ;
\item
  \textbf{dépense publique ou soutenabilité} : la lecture B montre le
  prix macroéconomique, dans ce modèle, d'un retour rapide vers 3 \% de
  déficit à partir d'une référence dégradée.
\end{itemize}

\hypertarget{sec:effets-sociaux-conclusion}{%
\section{Effets sociaux}\label{sec:effets-sociaux-conclusion}}

Les effets distributifs statiques (chapitre \ref{sec:distribution})
diffèrent surtout par leur concentration : prélèvements concentrés sur
le dernier décile pour LFI (les transferts et la hausse du SMIC, qui
bénéficieraient au bas et au milieu de la distribution, ne sont pas
répartis ménage par ménage dans cette version), gains diffus et
proportionnels pour le RN (TVA sur l'énergie), pertes concentrées sur
les deux premiers déciles pour Horizons (RSA). Les effets générationnels
des réformes des retraites sont nets : l'abaissement de l'âge profite
aux générations proches de la retraite et pèse sur les finances
publiques, donc sur les générations suivantes ; le relèvement à 65 ans
et 45 annuités réduit la durée de retraite des générations nées après
1970. Ces résultats sont des moyennes : la microsimulation statique ne
représente ni les trajectoires individuelles, ni les territoires, ni les
effets de second tour sur les revenus.

\hypertarget{sec:limites-conclusion}{%
\section{Limites de l'exercice}\label{sec:limites-conclusion}}

Les limites principales, détaillées en annexe \ref{sec:annexe-h}, sont :
(i) un modèle annuel à anticipations adaptatives, sans effets de
crédibilité ni de confiance ; (ii) une trajectoire de référence
explosive, qui impose de lire les niveaux à vingt ans avec prudence ;
(iii) des comportements calibrés hors de leur domaine d'estimation pour
les chocs les plus amples (lecture B) ; (iv) une microsimulation
statique ; (v) des mesures non représentées (capitalisation, durée
d'indemnisation du chômage, priorité nationale, contrôle des capitaux,
planification industrielle) ; (vi) un corpus provisoire, arrêté au 24
novembre 2026, avant les chiffrages complets annoncés par deux des trois
candidats. Les tests automatisés vérifient la cohérence comptable, la
reproductibilité et la non-régression ; ils ne constituent pas une
validation scientifique des comportements, qui repose sur des
estimations économétriques imparfaites et sur des rétrosimulations
(chapitre \ref{sec:donnees}).

\hypertarget{sec:questions-ouvertes}{%
\section{Questions ouvertes et éléments à
actualiser}\label{sec:questions-ouvertes}}

\begin{longtable}[]{@{}
  >{\raggedright\arraybackslash}p{(\columnwidth - 4\tabcolsep) * \real{0.3333}}
  >{\raggedright\arraybackslash}p{(\columnwidth - 4\tabcolsep) * \real{0.3333}}
  >{\raggedright\arraybackslash}p{(\columnwidth - 4\tabcolsep) * \real{0.3333}}@{}}
\toprule\noalign{}
\begin{minipage}[b]{\linewidth}\raggedright
Élément attendu
\end{minipage} & \begin{minipage}[b]{\linewidth}\raggedright
Candidat
\end{minipage} & \begin{minipage}[b]{\linewidth}\raggedright
Effet sur l'étude
\end{minipage} \\
\midrule\noalign{}
\endhead
\bottomrule\noalign{}
\endlastfoot
Chiffrage complet de \emph{L'Avenir en commun} (annoncé pour début
2027), barème de l'impôt sur le revenu à 14 tranches, périmètre de la
revalorisation de 10 \% & Mélenchon & nouvelles mesures et révision de
L05 et L21 ; lecture B \\
« Volume 2 » du programme (janvier 2027) : ventilation des économies,
allocation de 350 €, durée d'indemnisation du chômage, rémunération des
enseignants & Philippe & passage d'une partie de la lecture B à la
lecture A \\
Composition des 140 Md€ d'économies, économies sur les retraites, coût
de la baisse de TVA sur l'énergie, impôts de production & Le Pen &
idem \\
Programme détaillé de Gabriel Attal ou d'un autre candidat du centre &
--- & réévaluation du choix du candidat du centre (ajout au corpus sans
changement de code) \\
Comptes nationaux 2025, prévisions du FMI d'avril 2027, taux de marché &
--- & recalibration des environnements \\
\end{longtable}

Les questions scientifiques ouvertes les plus importantes pour ce type
d'exercice sont : la taille des multiplicateurs de consolidation dans
une économie très endettée de la zone euro ; la réaction de la prime
souveraine aux annonces budgétaires ; l'élasticité de l'activité des
seniors à la durée d'assurance pour les générations récentes ;
l'élasticité des bases taxables des très hauts patrimoines ;
l'additionnalité de l'investissement des banques publiques.

\hypertarget{sec:enseignements}{%
\section{Enseignements méthodologiques}\label{sec:enseignements}}

\begin{enumerate}
\def\labelenumi{\arabic{enumi}.}
\tightlist
\item
  \textbf{Un chiffrage statique ne suffit pas, mais un modèle dynamique
  non plus} : la valeur ajoutée de l'exercice tient moins à un chiffre
  central qu'à la décomposition des effets (mesures, régimes,
  environnements, institutions) et à la séparation explicite entre
  faits, hypothèses de traduction et résultats.
\item
  \textbf{La symétrie de traitement exige deux lectures} : simuler
  seulement les mesures documentées favoriserait les programmes les
  moins détaillés ; simuler seulement les objectifs reviendrait à
  remplacer les programmes par les hypothèses de l'évaluateur.
\item
  \textbf{La robustesse se démontre par variantes}, pas par une
  validation unique : les conclusions retenues sont celles qui survivent
  aux régimes, aux environnements et aux hypothèses de taux.
\item
  \textbf{La traçabilité est une condition de la discussion publique} :
  chaque nombre de cette étude est produit par un script archivé, à
  partir de sources datées, avec des hypothèses nommées (annexe
  \ref{sec:annexe-k}).
\end{enumerate}

L'étude ne désigne pas de programme « meilleur » : les critères
économiques (activité, emploi, prix), sociaux (distribution,
générations), budgétaires (solde, dette) et institutionnels
(compatibilité européenne) ne se hiérarchisent pas sans jugement de
valeur, et leurs résultats dépendent d'hypothèses que le lecteur peut
faire varier.

\appendix

\hypertarget{sec:annexe-a}{%
\chapter{Registre des paramètres, valeurs, unités et
sources}\label{sec:annexe-a}}

Statuts : {[}D{]} donnée ou règle publiée ; {[}E{]} estimé par le projet
; {[}C{]} calibré (littérature ou cible) ; {[}H{]} hypothèse ; {[}X{]}
exploratoire ; confiance A à D (registre de la version 0.8). Les
coefficients des équations à correction d'erreur (consommation,
investissement, commerce, emploi, salaires, prix) sont donnés au
chapitre \ref{sec:equations} ; les effets fixes de branche et les
constantes de calage sont dans le code (\texttt{src/m7\_equations.py},
cache des environnements).

\hypertarget{paramuxe8tres-structurels-modifiables-sans-recalibration-m23}{%
\subsection{Paramètres structurels modifiables sans recalibration
(M23)}\label{paramuxe8tres-structurels-modifiables-sans-recalibration-m23}}

\begin{longtable}[]{@{}
  >{\raggedright\arraybackslash}p{(\columnwidth - 12\tabcolsep) * \real{0.1429}}
  >{\raggedright\arraybackslash}p{(\columnwidth - 12\tabcolsep) * \real{0.1429}}
  >{\raggedright\arraybackslash}p{(\columnwidth - 12\tabcolsep) * \real{0.1429}}
  >{\raggedright\arraybackslash}p{(\columnwidth - 12\tabcolsep) * \real{0.1429}}
  >{\raggedright\arraybackslash}p{(\columnwidth - 12\tabcolsep) * \real{0.1429}}
  >{\raggedright\arraybackslash}p{(\columnwidth - 12\tabcolsep) * \real{0.1429}}
  >{\raggedright\arraybackslash}p{(\columnwidth - 12\tabcolsep) * \real{0.1429}}@{}}
\toprule\noalign{}
\begin{minipage}[b]{\linewidth}\raggedright
paramètre
\end{minipage} & \begin{minipage}[b]{\linewidth}\raggedright
catégorie
\end{minipage} & \begin{minipage}[b]{\linewidth}\raggedright
valeur centrale
\end{minipage} & \begin{minipage}[b]{\linewidth}\raggedright
basse
\end{minipage} & \begin{minipage}[b]{\linewidth}\raggedright
haute
\end{minipage} & \begin{minipage}[b]{\linewidth}\raggedright
statut
\end{minipage} & \begin{minipage}[b]{\linewidth}\raggedright
source, justification
\end{minipage} \\
\midrule\noalign{}
\endhead
\bottomrule\noalign{}
\endlastfoot
\texttt{PSI\_TAUX} & comportements & 0.2 & 0.1 & 0.4 & {[}L{]} &
consommation : −\% par point de taux réel (FR-BDF, Mésange : 0,1-0,4) \\
\texttt{PSI\_CHOMAGE} & comportements & 0.3 & 0.1 & 0.5 & {[}H/L{]} &
épargne de précaution : −\% de consommation par point de chômage \\
\texttt{SEMI\_ELAST\_IS} & comportements & 0.8 & 0.2 & 1.5 & {[}L{]} &
transfert de bénéfices : semi-élasticité de la base de l'IS (Heckemeyer
et Overesch 2017) \\
\texttt{PART\_IS\_MULTINAT} & comportements & 0.5 & 0.3 & 0.7 & {[}C{]}
& part de l'IS acquittée par les groupes multinationaux \\
\texttt{THETA\_TRESORERIE} & comportements & 0.1 & 0.0 & 0.3 & {[}H{]} &
part des profits conservés réinvestie (contrainte de trésorerie) \\
\texttt{THETA\_ANTICIPATION} & comportements & 0.75 & 0.5 & 1.0 &
{[}E{]} & anticipation du potentiel (phase 3.1, inférence indirecte) \\
\texttt{mu\_taux} & societes & estimé (M24) & 0.35 & 0.83 & {[}E{]} &
vitesse de renouvellement du taux apparent des SNF (0,60 estimé
1999-2024) \\
\texttt{phi\_dette} & societes & estimé (M24) & 0.5 & 1.5 & {[}E{]} &
dette des SNF / besoin de financement (1,06 ± 0,31) \\
\texttt{div\_lissage} & societes & 1.0 & 0.3 & None & {[}H{]} & vitesse
d'ajustement des dividendes au taux de distribution (Lintner non
identifié) \\
\texttt{fiscalite\_dividendes} & societes & False & None & True & {[}E
fragile{]} & effets de la contribution de 3 \% et du PFU sur la
distribution \\
\texttt{bund\_transmission} & monetaire & estimé (M24) & 0.2 & 0.6 &
{[}E{]} & Δ Bund / Δ Euribor (0,39 ± 0,09) \\
\texttt{menages\_vitesse} & monetaire & estimé (M24) & 0.25 & 0.6 &
{[}E{]} & vitesse d'ajustement des taux des crédits à l'habitat (0,43 ±
0,07) \\
\texttt{entreprises\_vitesse} & monetaire & estimé (M24) & 0.6 & 1.0 &
{[}E{]} & vitesse d'ajustement du coût du crédit aux SNF (0,93 ±
0,02) \\
\texttt{bce\_inflation} & monetaire & estimé (M24) & 0.2 & 0.6 & {[}E{]}
& réponse de court terme de la BCE à l'inflation (0,39 ± 0,18) \\
\texttt{bce\_activite} & monetaire & estimé (M24) & 0.1 & 0.5 & {[}E{]}
& réponse de court terme de la BCE à l'activité (0,30 ± 0,09) \\
\texttt{prime\_dette} & monetaire & estimé (M24) & 0.007 & 0.03 &
{[}E{]}/{[}L{]} & prime souveraine par point de dette (1,3 pb ; France
seule 0,7 pb ; littérature 3 pb) \\
\texttt{prime\_solde} & monetaire & estimé (M24) & 0.0 & -0.118 & {[}E
fragile{]} & prime par point de solde (−1,5 pb hors crise ; −11,8 pb
avec la crise) \\
\end{longtable}

\hypertarget{politique-monuxe9taire-et-conditions-financiuxe8res-m24-estimuxe9s}{%
\subsection{Politique monétaire et conditions financières (M24,
estimés)}\label{politique-monuxe9taire-et-conditions-financiuxe8res-m24-estimuxe9s}}

\begin{longtable}[]{@{}lll@{}}
\toprule\noalign{}
paramètre (M24) & valeur & statut \\
\midrule\noalign{}
\endhead
\bottomrule\noalign{}
\endlastfoot
\texttt{bce\_lissage} & 0,7278 & {[}E{]} \\
\texttt{bce\_inflation} & 0,3928 & {[}E{]} \\
\texttt{bce\_activite} & 0,3022 & {[}E{]} \\
\texttt{bce\_forme} & contemporaine appliquée à t − 1 {[}E*{]} &
{[}E{]} \\
\texttt{bund\_transmission} & 0,3858 & {[}E{]} \\
\texttt{menages\_vitesse} & 0,4294 & {[}E{]} \\
\texttt{menages\_court} & 0,2167 & {[}E{]} \\
\texttt{menages\_long} & 0,1679 & {[}E{]} \\
\texttt{menages\_transmission\_lt} & 0,8957 & {[}E{]} \\
\texttt{entreprises\_vitesse} & 0,9263 & {[}E{]} \\
\texttt{entreprises\_court} & 0,6781 & {[}E{]} \\
\texttt{entreprises\_long} & 0,0440 & {[}E{]} \\
\texttt{entreprises\_transmission\_lt} & 0,7796 & {[}E{]} \\
\texttt{prime\_dette} & 0,0130 & {[}E{]} \\
\texttt{prime\_solde} & −0,0148 & {[}E{]} \\
\texttt{prime\_dette\_avec\_crise} & 0,0122 & {[}E{]} \\
\texttt{prime\_solde\_avec\_crise} & −0,1184 & {[}E{]} \\
\end{longtable}

\hypertarget{retraites-m9-finalisation}{%
\subsection{Retraites (M9,
finalisation)}\label{retraites-m9-finalisation}}

\begin{longtable}[]{@{}
  >{\raggedright\arraybackslash}p{(\columnwidth - 6\tabcolsep) * \real{0.2500}}
  >{\raggedright\arraybackslash}p{(\columnwidth - 6\tabcolsep) * \real{0.2500}}
  >{\raggedright\arraybackslash}p{(\columnwidth - 6\tabcolsep) * \real{0.2500}}
  >{\raggedright\arraybackslash}p{(\columnwidth - 6\tabcolsep) * \real{0.2500}}@{}}
\toprule\noalign{}
\begin{minipage}[b]{\linewidth}\raggedright
paramètre (M9)
\end{minipage} & \begin{minipage}[b]{\linewidth}\raggedright
valeur
\end{minipage} & \begin{minipage}[b]{\linewidth}\raggedright
statut
\end{minipage} & \begin{minipage}[b]{\linewidth}\raggedright
source
\end{minipage} \\
\midrule\noalign{}
\endhead
\bottomrule\noalign{}
\endlastfoot
report par année de durée requise, hommes & 0.75 an & {[}E,
littérature{]} & Bozio (2011) \\
report par année de durée requise, femmes & 0,417 an & {[}E,
littérature{]} & Bozio (2011) \\
part de la pension soumise à proratisation et décote & 0.5 & {[}H{]} &
ordre de grandeur \\
décote par année manquante & 0.05 & {[}D{]} & code de la sécurité
sociale \\
part des carrières longues & 0.17 & {[}D{]} & CNAV, données 2024 \\
coefficient sur l'écart de masse des pensions & 0.55 & {[}C{]} & étude
d'impact 2023, Institut Montaigne \\
\end{longtable}

\hypertarget{ruxe8gle-budguxe9taire-r2-m23}{%
\subsection{Règle budgétaire R2
(M23)}\label{ruxe8gle-budguxe9taire-r2-m23}}

\begin{longtable}[]{@{}
  >{\raggedright\arraybackslash}p{(\columnwidth - 4\tabcolsep) * \real{0.3333}}
  >{\raggedright\arraybackslash}p{(\columnwidth - 4\tabcolsep) * \real{0.3333}}
  >{\raggedright\arraybackslash}p{(\columnwidth - 4\tabcolsep) * \real{0.3333}}@{}}
\toprule\noalign{}
\begin{minipage}[b]{\linewidth}\raggedright
paramètre de régime
\end{minipage} & \begin{minipage}[b]{\linewidth}\raggedright
valeur
\end{minipage} & \begin{minipage}[b]{\linewidth}\raggedright
statut
\end{minipage} \\
\midrule\noalign{}
\endhead
\bottomrule\noalign{}
\endlastfoot
\texttt{debut} & 2027 & {[}H{]} \\
\texttt{alpha} & 0.5 & {[}H{]} \\
\texttt{beta} & 0.02 & {[}H{]} \\
\texttt{dette\_cible} & 90.0 & {[}H{]} \\
\texttt{plafond} & 0.5 & {[}H{]} \\
\texttt{ecart\_suspension} & 1.5 & {[}H{]} \\
\texttt{epsilon} & 0.5 & {[}H{]} \\
répartition de l'effort & \{`conso\_publique': 0.3, `prestations': 0.2,
`impots\_menages': 0.3, `impots\_produits': 0.2\} & {[}H{]} \\
\end{longtable}

\hypertarget{registre-duxe9tailluxe9-de-la-version-0.8-toujours-en-vigueur-pour-ces-paramuxe8tres}{%
\subsection{Registre détaillé de la version 0.8 (toujours en vigueur
pour ces
paramètres)}\label{registre-duxe9tailluxe9-de-la-version-0.8-toujours-en-vigueur-pour-ces-paramuxe8tres}}

\hypertarget{bouclage-macrouxe9conomique}{%
\subsection{Bouclage
macroéconomique}\label{bouclage-macrouxe9conomique}}

\begin{longtable}[]{@{}
  >{\raggedright\arraybackslash}p{(\columnwidth - 12\tabcolsep) * \real{0.1429}}
  >{\raggedright\arraybackslash}p{(\columnwidth - 12\tabcolsep) * \real{0.1429}}
  >{\raggedright\arraybackslash}p{(\columnwidth - 12\tabcolsep) * \real{0.1429}}
  >{\raggedright\arraybackslash}p{(\columnwidth - 12\tabcolsep) * \real{0.1429}}
  >{\raggedright\arraybackslash}p{(\columnwidth - 12\tabcolsep) * \real{0.1429}}
  >{\raggedright\arraybackslash}p{(\columnwidth - 12\tabcolsep) * \real{0.1429}}
  >{\raggedright\arraybackslash}p{(\columnwidth - 12\tabcolsep) * \real{0.1429}}@{}}
\toprule\noalign{}
\begin{minipage}[b]{\linewidth}\raggedright
Paramètre
\end{minipage} & \begin{minipage}[b]{\linewidth}\raggedright
Définition
\end{minipage} & \begin{minipage}[b]{\linewidth}\raggedright
Valeur
\end{minipage} & \begin{minipage}[b]{\linewidth}\raggedright
Intervalle
\end{minipage} & \begin{minipage}[b]{\linewidth}\raggedright
Source et méthode
\end{minipage} & \begin{minipage}[b]{\linewidth}\raggedright
Confiance
\end{minipage} & \begin{minipage}[b]{\linewidth}\raggedright
Date
\end{minipage} \\
\midrule\noalign{}
\endhead
\bottomrule\noalign{}
\endlastfoot
α & Part du capital (fonction de production) & 0,35 & 0,30-0,40 &
Partage de la valeur ajoutée, comptes nationaux & A & 2024 \\
ρ & Rappel annuel de la demande vers le potentiel & 0,2 & 0,1-0,3 &
Inférence indirecte : durée des écarts de production (AMECO, FMI
1982-2019) & B & 2026 \\
θ & Anticipation des révisions du potentiel par la demande & 0,75 &
0,5-1,0 & Même méthode + ancrage sur la réponse de la consommation à un
choc permanent (Blundell, Pistaferri, Preston 2008) & C & 2026 \\
a & Terme proportionnel de la fermeture & 0 & 0-0,15 & Retiré :
oscillations non observées & C & 2026 \\
μ & Effet des tensions sur les importations & 0 & 0-0,3 & Retiré :
oscillations non observées ; non estimé & D & 2026 \\
Part du capital dans l'anticipation & Part de l'écart de potentiel dû au
capital prise en compte & 1 & 0-1 & Non identifiable ; incertitude de
structure & D & 2026 \\
κ & Pente de Phillips sur les prix hors industrie & 0,10 & 0,05-0,18 &
Banque de France (États-Unis 0,18) ; ancrage des anticipations & C &
2024 \\
u* & Chômage structurel & 7,5 → 7,0 \% & --- & Hypothèses de la
référence (FMI, COR) & C & 2026 \\
\end{longtable}

\hypertarget{muxe9nages-1}{%
\subsection{Ménages}\label{muxe9nages-1}}

\begin{longtable}[]{@{}
  >{\raggedright\arraybackslash}p{(\columnwidth - 12\tabcolsep) * \real{0.1429}}
  >{\raggedright\arraybackslash}p{(\columnwidth - 12\tabcolsep) * \real{0.1429}}
  >{\raggedright\arraybackslash}p{(\columnwidth - 12\tabcolsep) * \real{0.1429}}
  >{\raggedright\arraybackslash}p{(\columnwidth - 12\tabcolsep) * \real{0.1429}}
  >{\raggedright\arraybackslash}p{(\columnwidth - 12\tabcolsep) * \real{0.1429}}
  >{\raggedright\arraybackslash}p{(\columnwidth - 12\tabcolsep) * \real{0.1429}}
  >{\raggedright\arraybackslash}p{(\columnwidth - 12\tabcolsep) * \real{0.1429}}@{}}
\toprule\noalign{}
\begin{minipage}[b]{\linewidth}\raggedright
Paramètre
\end{minipage} & \begin{minipage}[b]{\linewidth}\raggedright
Définition
\end{minipage} & \begin{minipage}[b]{\linewidth}\raggedright
Valeur
\end{minipage} & \begin{minipage}[b]{\linewidth}\raggedright
Intervalle
\end{minipage} & \begin{minipage}[b]{\linewidth}\raggedright
Source et méthode
\end{minipage} & \begin{minipage}[b]{\linewidth}\raggedright
Confiance
\end{minipage} & \begin{minipage}[b]{\linewidth}\raggedright
Date
\end{minipage} \\
\midrule\noalign{}
\endhead
\bottomrule\noalign{}
\endlastfoot
η\_ext & Élasticité de participation au gain à l'emploi & 0,2 & 0,1-0,4
& Revues OFCE, DARES ; 27 études sur 35 \textgreater{} 0,1 & B & 2023 \\
η\_int & Élasticité des heures au taux marginal net & 0,1 & 0,05-0,2 &
Cabannes, Houdré, Landais (2014) ; revues & B & 2014 \\
ψ\_u & Épargne de précaution (consommation, \% par point de chômage) &
0,3 & 0,1-0,6 & Littérature ; épisode 2020-2021 & C & --- \\
ψ\_r & Effet du taux réel sur la consommation & 0,2 & 0-0,5 & Havranek
(2015) & C & 2015 \\
Rigidité salariale à la baisse & Multiplicateur de la courbe de Phillips
quand u \textgreater{} u* & 0,5 & 0,25-1 & Littérature sur la rigidité
nominale & C & --- \\
Effet de l'assurance chômage sur u* & Points par 10 \% d'allocation &
0,063 & 0-0,13 & Le Barbanchon, Rathelot, Roulet (2017) & B & 2017 \\
Salaire d'embauche (M16) & Relatif au salaire moyen & 0,75 & 0,6-0,9 &
Insee, salaires d'embauche & C & --- \\
Allocation perdue au retour à l'emploi (M16) & € par an, en espérance &
5 200 & 3 000-8 000 & Indemnisation moyenne × part indemnisée (Unédic,
ordres de grandeur) & C & 2024 \\
Masse indiciaire nette (M16, point d'indice) & M€ pour +100 \% & 187 000
& 170 000-200 000 & Masse indiciaire brute ≈ 220 Md€ (DGAFP, 1 \% ≈ 2,2
Md€) × 0,85 & C & 2025 \\
Âges des ménages concernés par l'âge légal (M16) & Âge de la personne de
référence & 58-67 ans & --- & Hypothèse de pseudo-cohorte & D & --- \\
Prime de risque souveraine (M11) & Point de taux long par point de dette
& 0,03 & 0-0,05 & Littérature (2 à 5 points de base par point de dette)
; non estimée par le projet & C & --- \\
Taux net d'un supplément de salaire (M16) & Part conservée après impôt &
0,88 & 0,8-0,95 & Taux moyen d'IR & D & 2024 \\
\end{longtable}

\hypertarget{entreprises-et-incidence}{%
\subsection{Entreprises et incidence}\label{entreprises-et-incidence}}

\begin{longtable}[]{@{}
  >{\raggedright\arraybackslash}p{(\columnwidth - 12\tabcolsep) * \real{0.1429}}
  >{\raggedright\arraybackslash}p{(\columnwidth - 12\tabcolsep) * \real{0.1429}}
  >{\raggedright\arraybackslash}p{(\columnwidth - 12\tabcolsep) * \real{0.1429}}
  >{\raggedright\arraybackslash}p{(\columnwidth - 12\tabcolsep) * \real{0.1429}}
  >{\raggedright\arraybackslash}p{(\columnwidth - 12\tabcolsep) * \real{0.1429}}
  >{\raggedright\arraybackslash}p{(\columnwidth - 12\tabcolsep) * \real{0.1429}}
  >{\raggedright\arraybackslash}p{(\columnwidth - 12\tabcolsep) * \real{0.1429}}@{}}
\toprule\noalign{}
\begin{minipage}[b]{\linewidth}\raggedright
Paramètre
\end{minipage} & \begin{minipage}[b]{\linewidth}\raggedright
Définition
\end{minipage} & \begin{minipage}[b]{\linewidth}\raggedright
Valeur
\end{minipage} & \begin{minipage}[b]{\linewidth}\raggedright
Intervalle
\end{minipage} & \begin{minipage}[b]{\linewidth}\raggedright
Source et méthode
\end{minipage} & \begin{minipage}[b]{\linewidth}\raggedright
Confiance
\end{minipage} & \begin{minipage}[b]{\linewidth}\raggedright
Date
\end{minipage} \\
\midrule\noalign{}
\endhead
\bottomrule\noalign{}
\endlastfoot
Semi-élasticité des bénéfices à l'IS & Transfert de bénéfices des
groupes & 0,8 & 0,2-1,5 & Heckemeyer-Overesch ; Beer, de Mooij, Liu ;
Merlo-Wamser & B & 2017-2020 \\
Part des groupes dans l'IS & --- & 0,5 & --- & Concentration de l'IS & D
& --- \\
Trésorerie → investissement & Part des bénéfices conservés investie &
0,10 & 0-0,3 & CICE (France Stratégie 2020) ; dividendes 2013 (Boissel,
Matray) & C & 2020-2022 \\
Incidence salariale de l'IS & Part supportée par les salariés à long
terme & 0,3 & 0-0,5 & Fuest, Peichl, Siegloch (AER 2018) : ≈ 0,5 (impôt
local allemand) & C & 2018 \\
Incidence salariale des impôts de production & idem & 0,2 & 0-0,4 & Pas
d'estimation causale française & D & --- \\
Répercussion des impôts de production sur les prix & --- & 0,5 & 0,3-0,7
& CAE (2019), hypothèse & D & 2019 \\
Transmission des cotisations employeurs aux salaires & Part transmise à
long terme & 0,5 & 0,2-1,0 & Bozio, Breda, Grenet (2019) ; Carbonnier et
al.~(2019) & C & 2019 \\
Répercussion de la TVA (hausse / baisse, taux normal) & --- & 0,95 /
0,50 & --- & Benzarti et al.~(JPE 2020) ; restauration en France & B &
2020 \\
Délocalisation énergo-intensive & Élasticité des importations au prix
relatif de l'énergie & 0,15 & 0-0,4 & Dussaux (2020) ; Marin et Vona
(2021) & C & 2021 \\
Additionnalité de l'aide à la R\&D & € de R\&D privée par € d'aide & 0,7
& 0,3-1,2 & CNEPI (2021) : 0,9-1,5 PME, ≈ 0 grandes & C & 2021 \\
Élasticité au stock de R\&D & --- & 0,08 & 0,03-0,15 & Hall, Mairesse,
Mohnen (2010) & C/X & 2010 \\
Formation → u* & Points par Md€ et par an & 0,02 & 0-0,05 & Évaluations
des formations de demandeurs d'emploi & X & --- \\
Durée du travail & Part des heures supplémentaires convertie en
production & 0,3 & --- & Hypothèse D de la v0.6 (35 heures) & D & --- \\
\end{longtable}

\hypertarget{finances-publiques}{%
\subsection{Finances publiques}\label{finances-publiques}}

\begin{longtable}[]{@{}
  >{\raggedright\arraybackslash}p{(\columnwidth - 12\tabcolsep) * \real{0.1429}}
  >{\raggedright\arraybackslash}p{(\columnwidth - 12\tabcolsep) * \real{0.1429}}
  >{\raggedright\arraybackslash}p{(\columnwidth - 12\tabcolsep) * \real{0.1429}}
  >{\raggedright\arraybackslash}p{(\columnwidth - 12\tabcolsep) * \real{0.1429}}
  >{\raggedright\arraybackslash}p{(\columnwidth - 12\tabcolsep) * \real{0.1429}}
  >{\raggedright\arraybackslash}p{(\columnwidth - 12\tabcolsep) * \real{0.1429}}
  >{\raggedright\arraybackslash}p{(\columnwidth - 12\tabcolsep) * \real{0.1429}}@{}}
\toprule\noalign{}
\begin{minipage}[b]{\linewidth}\raggedright
Paramètre
\end{minipage} & \begin{minipage}[b]{\linewidth}\raggedright
Définition
\end{minipage} & \begin{minipage}[b]{\linewidth}\raggedright
Valeur
\end{minipage} & \begin{minipage}[b]{\linewidth}\raggedright
Intervalle
\end{minipage} & \begin{minipage}[b]{\linewidth}\raggedright
Source
\end{minipage} & \begin{minipage}[b]{\linewidth}\raggedright
Confiance
\end{minipage} & \begin{minipage}[b]{\linewidth}\raggedright
Date
\end{minipage} \\
\midrule\noalign{}
\endhead
\bottomrule\noalign{}
\endlastfoot
Renouvellement de la dette & Part de la dette renouvelée chaque année au
taux de marché & 1/8 & 1/10-1/6 & Durée moyenne de la dette (AFT) & C ;
le backtest montre une baisse plus rapide du taux apparent en 2012-2019
que celle simulée & 2024 \\
Dépense d'indemnisation chômage & \% du PIB à chômage de base & 1,5 &
1,3-1,7 & Unédic & A & 2024 \\
\end{longtable}

\hypertarget{donnuxe9es-indisponibles-et-substitutions}{%
\subsection{Données indisponibles et
substitutions}\label{donnuxe9es-indisponibles-et-substitutions}}

\begin{longtable}[]{@{}
  >{\raggedright\arraybackslash}p{(\columnwidth - 4\tabcolsep) * \real{0.3333}}
  >{\raggedright\arraybackslash}p{(\columnwidth - 4\tabcolsep) * \real{0.3333}}
  >{\raggedright\arraybackslash}p{(\columnwidth - 4\tabcolsep) * \real{0.3333}}@{}}
\toprule\noalign{}
\begin{minipage}[b]{\linewidth}\raggedright
Donnée manquante
\end{minipage} & \begin{minipage}[b]{\linewidth}\raggedright
Substitution
\end{minipage} & \begin{minipage}[b]{\linewidth}\raggedright
Perte de précision
\end{minipage} \\
\midrule\noalign{}
\endhead
\bottomrule\noalign{}
\endlastfoot
Potentiel en temps réel (millésimes successifs) & Potentiel reconstruit
en temps réel par filtrage jusqu'à l'origine & Le potentiel « vintage »
d'époque aurait pu être différent \\
Patrimoine des ménages par ménage dans l'échantillon & Revenus non
salariaux et non pensions du revenu fiscal & Les ménages à patrimoine
non productif de revenu imposable sont sous-représentés \\
Héritages reçus par ménage & Assiette de patrimoine pondérée par l'âge &
Distribution des droits de succession approximative \\
Transitions individuelles emploi-chômage (panel) & Probabilités en
espérance par risque relatif âge × catégorie & Pas de distribution des
pertes individuelles \\
Effets de la crise sanitaire & Période exclue du backtest & Pas de
validation sur 2020-2021 \\
\end{longtable}

\hypertarget{sec:annexe-b}{%
\chapter{Équations complètes du modèle (forme
compacte)}\label{sec:annexe-b}}

Cette annexe récapitule le système dans l'ordre de résolution d'une
année (chapitre \ref{sec:annee}). Les numéros renvoient aux sections du
chapitre \ref{sec:equations}. Variables en logarithme en minuscules ;
\(\Delta\) : variation annuelle ; \(b\) : branche (\(25\) branches
marchandes) ; \(j\) : produit (\(64\)).

\textbf{1. Crochets de régime (\(t\), information \(\le t-1\)).}
\(\Delta i_t = 0{,}728\,\Delta i_{t-1} + 0{,}393\,\Delta\pi^{ZE}_{t-1} + 0{,}302\,\Delta y^{ZE}_{t-1}\)
;
\(\Delta B_t = \Delta B_{t-1} + 0{,}386\,(\Delta i_t - \Delta i_{t-1})\)
;
\(\Delta \text{écart}_t = \gamma_d\,\Delta d_{t-1} + \gamma_s\,\Delta s_{t-1} + \text{chocs}\)
;
\(\Delta i^m_t = (1 - 0{,}43)\,\Delta i^m_{t-1} + b^c_m \Delta i_t + b^l_m \Delta i^{OAT}_t\)
;
\(\Delta i^e_t = (1 - 0{,}93)\,\Delta i^e_{t-1} + b^c_e \Delta i_t + b^l_e \Delta i^{OAT}_t\)
; règle budgétaire
\(\Delta e_t = \text{clip}(\alpha(s^* - s^{ca}), 0, \bar e)\) ; objectif
(lecture B) \(\Delta e^B_t = \text{clip}(c_t - s_{t-1}, -2, 2)\).

\textbf{2. Prix d'importation.}
\(\Delta p^m_{j,t} = \theta\,\Delta p^{*}_{j,t} + (1-\theta)\,\Delta p^{*}_{j,t-1}\),
\(\theta \approx 0{,}5\).

\textbf{3. Salaires et prix.}
\(\Delta w_t = 0{,}39\,\Delta p^c_t - 0{,}09\,(w - p_{va} - \text{prod} - \mu_w)_{t-1} - \phi(u_t - u^*_t) + \dots\)
;
\(p^*_{va,b} = \theta_b \ln \text{CSU}_b + (1-\theta_b) \ln(p_I K_{b,t-1}/Q_b)\)
; correction d'erreur des prix de VA ;
\(p_j = \sum_i a^d_{ij} p_i + \sum_i a^m_{ij} p^m_i + v_j p_{va,j} + \tau_j\).

\textbf{4. Emploi et investissement.}
\(n_b = q_b - \sigma \ln(c_h/p_{va}) - (1-\sigma)(\ln E + \ln d)\)
(dynamique partielle) ;
\(\Delta i_b = 0{,}67\,\Delta q_b + 0{,}26\,\Delta q_{b,t-1} - \lambda_I(i_b - k^*_b)_{t-1} + \theta_T\,\Delta\text{Ép.}/I\)
; \(k^* = q - \sigma_K \ln c_K\).

\textbf{5. Revenus et demande.} \$\text{RDB} = \$ salaires \(+\) revenus
mixtes \(+\) revenus de la propriété (dividendes M21) \(+\) prestations
(dont pensions M9) \(-\) impôts et cotisations ;
\(\Delta c = 0{,}85\,\Delta r - 0{,}22\,(c - r - \mu_c)_{t-1} - \psi_u \Delta u - \psi_r \Delta i^{m,r}\)
; logement ; exportations
\(\Delta x = 0{,}71\,\Delta d^* - 0{,}60\,\Delta\ln(p_X/p^*)\) ;
fermeture \(S_t = S_{t-1} - \rho\,\tilde y_{t-1} + \theta\,\Delta y^*\).

\textbf{6. Importations et production.}
\(\Delta m = 1{,}47\,\Delta \bar d - \lambda_M(m - 1{,}37\,\bar d - 0{,}006\,t + 0{,}16 \ln p_M/p)_{t-1}\)
; \(q = (I - A_d)^{-1} d_d\).

\textbf{7. Après convergence.} Potentiel CES ;
\(\tilde y = 100\ln(Y/Y^*)\) ; finances publiques (recettes par taux
apparents, dépenses, intérêts par cohortes M17,
\(D_t = D_{t-1} - \text{Solde}_t + \text{AF}_t\)) ; compte des sociétés
(intérêts, IS, dividendes \(0{,}364\,\Pi_{t-1}\), dette
\(+1{,}06 \sum \Delta \text{BF}\)) ; comptes d'agents (RNB, capacités de
financement).

\textbf{8. Démographie et retraites (M1, M9).} Population Insee 2026 ;
activité des seniors
\(a(x,t) = A_s(t)\,\tilde h(x + 0{,}5 - \pi_{g,s})\),
\(\pi_{g,s} = \beta_s \text{AOD}_g + (1-\beta_s)\,62 + \varepsilon_s \Delta D_g / 4\)
; pensions par génération, facteur de durée \(f_{g,s}\) ; liens :
population active, emploi des seniors, écart de masse des pensions
\(\times 0{,}55\).

\hypertarget{sec:annexe-c}{%
\chapter{Registre des mesures politiques et de leur
traduction}\label{sec:annexe-c}}

Les tableaux du chapitre \ref{sec:programmes} donnent, pour chaque
mesure, le montant, le calendrier, le canal et la qualification
européenne. Cette annexe en donne les règles de conversion complètes et
les sources ; le registre est généré à partir du module
\texttt{src/m27\_corpus.py} (corpus daté du 24 novembre 2026).

\begin{itemize}
\tightlist
\item
  \textbf{L01} (Mélenchon (LFI)) --- Suppression du prélèvement
  forfaitaire unique (flat tax). \emph{Sources} : LFI\_BUDGET,
  LFI\_LCP\_BUDGET. \emph{Règle} : Points de PFU = montant / 0,6 Md€ par
  point (assiette 60 Md€).
\item
  \textbf{L02} (Mélenchon (LFI)) --- Rétablissement d'un impôt sur la
  fortune renforcé. \emph{Sources} : LFI\_BUDGET, LFI\_LCP\_BUDGET.
  \emph{Règle} : Taux moyen = (montant + rendement de l'IFI remplacé) /
  assiette ISF 930 Md€.
\item
  \textbf{L03} (Mélenchon (LFI)) --- Taxe Zucman (patrimoines supérieurs
  à 100 M€). \emph{Sources} : LFI\_BUDGET, LFI\_LCP\_BUDGET.
  \emph{Règle} : {[}H{]} Assiette (\textgreater{} 100 M€) distincte de
  celle de l'ISF du modèle : recette sans effet de demande (propension à
  consommer quasi nulle) ; l'érosion d'assiette n'est pas représentée
  (le montant annoncé est une borne haute). Variante : canal isf
  (érosion M13).
\item
  \textbf{L04} (Mélenchon (LFI)) --- Réforme des droits de succession.
  \emph{Sources} : LFI\_BUDGET, LFI\_LCP\_BUDGET. \emph{Règle} :
  Variation du rendement = montant / 18 Md€ (au-delà de la borne de
  l'interface).
\item
  \textbf{L05} (Mélenchon (LFI)) --- Impôt sur le revenu à 14 tranches.
  \emph{Sources} : LFI\_BUDGET, LFI\_LCP\_BUDGET. \emph{Règle} : Barème
  non publié : canal ir\_t1..ir\_t4 identifié, paramètres manquants.
\item
  \textbf{L06} (Mélenchon (LFI)) --- Rétablissement de la CVAE sur les
  plus grandes entreprises. \emph{Sources} : LFI\_BUDGET,
  LFI\_LCP\_BUDGET. \emph{Règle} : Md€ directs.
\item
  \textbf{L07} (Mélenchon (LFI)) --- Réforme de l'impôt sur les
  sociétés. \emph{Sources} : LFI\_BUDGET, LFI\_LCP\_BUDGET. \emph{Règle}
  : Paramètres non publiés : taux\_is calé sur 3 Md€ par le rendement de
  référence de l'IS du modèle (fonction conversion\_is).
\item
  \textbf{L08} (Mélenchon (LFI)) --- Taxe universelle sur les
  multinationales. \emph{Sources} : LFI\_BUDGET, LFI\_LCP\_BUDGET.
  \emph{Règle} : {[}H{]} Recette sans effet de demande intérieure
  (bénéfices en grande partie distribués aux non-résidents).
\item
  \textbf{L09} (Mélenchon (LFI)) --- Renforcement de la taxe sur les
  services numériques (GAFAM). \emph{Sources} : LFI\_BUDGET,
  LFI\_LCP\_BUDGET. \emph{Règle} : {[}H{]} idem L08.
\item
  \textbf{L10} (Mélenchon (LFI)) --- Taxe sur les superprofits.
  \emph{Sources} : LFI\_BUDGET, LFI\_LCP\_BUDGET. \emph{Règle} : {[}H{]}
  Prélèvement exceptionnel sur profits.
\item
  \textbf{L11} (Mélenchon (LFI)) --- Rabot et suppression des niches
  fiscales et sociales « les plus injustes ». \emph{Sources} :
  LFI\_BUDGET, LFI\_LCP\_BUDGET. \emph{Règle} : {[}H{]} Composition non
  publiée : moitié ménages (baisse de revenu disponible), moitié
  entreprises (hausse du coût du travail). Variante : recettes\_neutres.
\item
  \textbf{L12} (Mélenchon (LFI)) --- Fin des exonérations de cotisations
  au-dessus de 2 SMIC. \emph{Sources} : LFI\_BUDGET, LFI\_LCP\_BUDGET.
  \emph{Règle} : Canal cotisations\_employeurs ; montant non publié.
\item
  \textbf{L13} (Mélenchon (LFI)) --- Taxes sur les transactions
  financières. \emph{Sources} : LFI\_BUDGET, LFI\_LCP\_BUDGET.
  \emph{Règle} : ---
\item
  \textbf{L14} (Mélenchon (LFI)) --- Taxe sur les rachats d'actions.
  \emph{Sources} : LFI\_BUDGET, LFI\_LCP\_BUDGET. \emph{Règle} : ---
\item
  \textbf{L15} (Mélenchon (LFI)) --- Taxe sur les crypto-actifs.
  \emph{Sources} : LFI\_BUDGET, LFI\_LCP\_BUDGET. \emph{Règle} : ---
\item
  \textbf{L16} (Mélenchon (LFI)) --- Lutte contre la fraude fiscale.
  \emph{Sources} : LFI\_BUDGET, LFI\_LCP\_BUDGET. \emph{Règle} : ---
\item
  \textbf{L20} (Mélenchon (LFI)) --- Dégel et hausse de 10 \% du point
  d'indice. \emph{Sources} : LFI\_BUDGET, LFI\_LCP\_BUDGET. \emph{Règle}
  : Point d'indice +10 \%.
\item
  \textbf{L21} (Mélenchon (LFI)) --- Revalorisation de 10 \% des
  pensions, allocations familiales et APL. \emph{Sources} : LFI\_BUDGET,
  LFI\_LCP\_BUDGET. \emph{Règle} : {[}H{]} Le montant annoncé (8 Md€)
  n'est pas compatible avec une revalorisation générale de toutes les
  pensions (masse de plusieurs centaines de Md€) : périmètre à préciser
  ; montant annoncé versé en transferts non ciblés. Variante :
  allocations\_familiales et aides\_logement +10 \% (microsimulation).
\item
  \textbf{L22} (Mélenchon (LFI)) --- Logements d'urgence et réquisition
  de logements vacants. \emph{Sources} : LFI\_BUDGET, LFI\_LCP\_BUDGET.
  \emph{Règle} : ---
\item
  \textbf{L23} (Mélenchon (LFI)) --- Garantie d'autonomie pour les
  jeunes. \emph{Sources} : LFI\_BUDGET, LFI\_LCP\_BUDGET. \emph{Règle} :
  {[}H{]} Transfert non ciblé (pas de mesure jeunes dans M4).
\item
  \textbf{L24} (Mélenchon (LFI)) --- SMIC à 1 700 € net. \emph{Sources}
  : LFI\_BUDGET, LFI\_LCP\_BUDGET, SMIC\_2026. \emph{Règle} : Hausse = 1
  700 / 1 477,93 − 1 ≈ 15,0 \% (SMIC net au 1er juin 2026).
\item
  \textbf{L25} (Mélenchon (LFI)) --- Blocage du prix des carburants à
  1,70 €/l. \emph{Sources} : LFI\_BUDGET, LFI\_LCP\_BUDGET. \emph{Règle}
  : Coût dépendant du prix de marché : canal energie\_import +
  subvention, non chiffrable sans hypothèse de prix.
\item
  \textbf{L26} (Mélenchon (LFI)) --- 10 000 postes dans l'Éducation
  nationale. \emph{Sources} : LFI\_BUDGET, LFI\_LCP\_BUDGET.
  \emph{Règle} : ---
\item
  \textbf{L27} (Mélenchon (LFI)) --- Gratuité de la cantine scolaire.
  \emph{Sources} : LFI\_BUDGET, LFI\_LCP\_BUDGET. \emph{Règle} : {[}H{]}
  Dépense de fonctionnement (transfert en nature).
\item
  \textbf{L28} (Mélenchon (LFI)) --- 250 000 places à l'université.
  \emph{Sources} : LFI\_BUDGET, LFI\_LCP\_BUDGET. \emph{Règle} : ---
\item
  \textbf{L29} (Mélenchon (LFI)) --- Hôpital public (+10 000 postes
  soignants). \emph{Sources} : LFI\_BUDGET, LFI\_LCP\_BUDGET.
  \emph{Règle} : ---
\item
  \textbf{L30} (Mélenchon (LFI)) --- Violences sexistes et sexuelles ;
  violences faites aux enfants. \emph{Sources} : LFI\_BUDGET,
  LFI\_LCP\_BUDGET. \emph{Règle} : 2,6 + 2,0 Md€.
\item
  \textbf{L31} (Mélenchon (LFI)) --- Indexation des budgets 2027 des
  services supprimés en 2026. \emph{Sources} : LFI\_BUDGET,
  LFI\_LCP\_BUDGET. \emph{Règle} : ---
\item
  \textbf{L32} (Mélenchon (LFI)) --- Adaptation au dérèglement
  climatique. \emph{Sources} : LFI\_BUDGET, LFI\_LCP\_BUDGET.
  \emph{Règle} : ---
\item
  \textbf{L33} (Mélenchon (LFI)) --- Investissements de planification
  écologique portés par le pôle public bancaire. \emph{Sources} :
  LFI\_BUDGET, LFI\_LCP\_BUDGET. \emph{Règle} : 100 Md€ par an annoncés
  hors budget : investissement hors APU (effet de demande sans coût pour
  le solde public). {[}H{]} Montée en charge non précisée : appliqué dès
  2028 ; additionnalité totale (borne haute). Financement, dette du pôle
  public et effet d'offre non représentés.
\item
  \textbf{L34} (Mélenchon (LFI)) --- Abrogation de la réforme de 2023
  (retour à 62 ans). \emph{Sources} : LFI\_BUDGET, LFI\_LCP\_BUDGET.
  \emph{Règle} : Retour à 62 ans sans transition et au calendrier de
  durée de la loi de 2014 ; coût calculé par M9 (le montant annoncé sert
  de contrôle).
\item
  \textbf{L35} (Mélenchon (LFI)) --- Abrogation des réformes de
  l'assurance chômage. \emph{Sources} : LFI\_BUDGET, LFI\_LCP\_BUDGET.
  \emph{Règle} : Variation = montant / 40 Md€ de dépense.
\item
  \textbf{L36} (Mélenchon (LFI)) --- Retour à 60 ans avec 40 annuités.
  \emph{Sources} : LFI\_BUDGET, LFI\_LCP\_BUDGET. \emph{Règle} :
  Objectif de mandat : variante (non incluse dans le paquet 2027) ; 40
  annuités = 160 trimestres, un trimestre par génération {[}H{]}.
\item
  \textbf{L37} (Mélenchon (LFI)) --- Plafond du livret A relevé à 30 000
  € ; centralisation portée à 80 \% ; livret B. \emph{Sources} :
  LFI\_BUDGET, LFI\_EPARGNE. \emph{Règle} : Pas de canal de
  l'intermédiation financière dans le modèle.
\item
  \textbf{L38} (Mélenchon (LFI)) --- Contrôle des capitaux ; fin de la
  cotation continue. \emph{Sources} : LFI\_EPARGNE. \emph{Règle} :
  Mesures de régulation sans canal.
\item
  \textbf{L39} (Mélenchon (LFI)) --- Réduction du déficit de 25 Md€ par
  an ; stabilisation de la dette en un quinquennat. \emph{Sources} :
  LFI\_BUDGET, LFI\_LCP\_BUDGET. \emph{Règle} : Objectif de solde :
  élément de régime (pas de règle automatique).
\item
  \textbf{P01} (Le Pen (RN)) --- 140 Md€ d'économies nettes d'ici 2032 ;
  dette de 121 \% (2027) à 112 \% (2032) ; déficit à 3 \% en 2030 ; 30
  Md€ de baisses d'impôts. \emph{Sources} : RN\_PS\_BUDGET,
  RN\_MONEYVOX. \emph{Règle} : Enveloppe agrégée : les mesures
  détaillées ci-dessous en couvrent une partie.
\item
  \textbf{P02} (Le Pen (RN)) --- Règle d'or constitutionnelle par
  référendum (déficit plafonné au niveau stabilisant la dette).
  \emph{Sources} : RN\_PS\_BUDGET. \emph{Règle} : Élément de régime :
  règle de stabilisation (M23 R1, mode actif).
\item
  \textbf{P03} (Le Pen (RN)) --- Économies sur le « coût de
  l'immigration ». \emph{Sources} : RN\_PS\_BUDGET. \emph{Règle} :
  {[}H{]} Composition non publiée : baisse de prestations (transferts).
  La part fondée sur la nationalité est incompatible pour les
  ressortissants de l'Union (M26, niveau 4) ; cette part n'est pas
  connue.
\item
  \textbf{P04} (Le Pen (RN)) --- Baisse de la contribution au budget de
  l'Union. \emph{Sources} : RN\_PS\_BUDGET. \emph{Règle} : Recette nette
  sans effet de demande.
\item
  \textbf{P05} (Le Pen (RN)) --- Réorganisation de l'État
  (centralisation de la TVA, transferts aux collectivités et à la
  Sécurité sociale votés par le Parlement). \emph{Sources} :
  RN\_PS\_BUDGET. \emph{Règle} : {[}H{]} Les sources ne disent pas si
  les économies sur les collectivités (P06) et les agences (P07) sont
  incluses : par prudence, P06 et P07 sont traitées comme des
  composantes de P05 (non additionnées).
\item
  \textbf{P06} (Le Pen (RN)) --- Contribution des collectivités (dont
  suppression des EPCI et métropoles 0,25 Md€, effectifs 0,8 Md€).
  \emph{Sources} : RN\_PS\_BUDGET. \emph{Règle} : Composante de P05
  {[}H{]}.
\item
  \textbf{P07} (Le Pen (RN)) --- Agences de l'État. \emph{Sources} :
  RN\_PS\_BUDGET. \emph{Règle} : Composante de P05 {[}H{]} ; une
  commission sénatoriale évaluait ces économies à environ 0,5 Md€.
\item
  \textbf{P08} (Le Pen (RN)) --- Âge légal ramené à 62 ans (60 ans pour
  les carrières commencées avant 20 ans), 42 annuités ; pilier de
  capitalisation obligatoire (socle : Fonds de réserve pour les
  retraites). \emph{Sources} : RN\_RETRAITES\_PS, RN\_RETRAITES\_SUD,
  RN\_RETRAITES\_BRUT. \emph{Règle} : Coût brut annoncé 18,6 Md€ en
  2032. M9 : âge 62 ans, 42 annuités (168 trimestres), départ à 60 ans
  pour une fraction de chaque génération égale à la part des départs
  pour carrière longue au régime général (17 \%, CNAV 2024) {[}H :
  éligibilité « avant 20 ans » approchée par cette part{]}. Pilier de
  capitalisation non chiffré : non représenté.
\item
  \textbf{P09} (Le Pen (RN)) --- Économies sur les retraites («
  corrections de mesures injustes »). \emph{Sources} : RN\_PS\_BUDGET,
  RN\_MONEYVOX. \emph{Règle} : Fourchette 15-20 Md€, mesures non
  détaillées.
\item
  \textbf{P10} (Le Pen (RN)) --- TVA ramenée de 20 \% à 5,5 \% sur
  l'électricité, le gaz, le fioul et les carburants. \emph{Sources} :
  RN\_ENERGIE, MONTAIGNE\_TVA\_2022. \emph{Règle} : Coût : ordre de
  grandeur de l'Institut Montaigne (2022), aucun montant du candidat
  dans les sources. {[}H{]} Baisse équivalente du taux normal (le modèle
  ne distingue pas l'énergie) : la répercussion sur les prix est
  répartie sur toute la consommation.
\item
  \textbf{P11} (Le Pen (RN)) --- Exonération d'impôt sur le revenu des
  moins de 30 ans (jusqu'à 4 SMIC). \emph{Sources} : RN\_LCP\_JEUNES.
  \emph{Règle} : {[}H{]} Baisse d'impôt traitée en transfert non ciblé.
\item
  \textbf{P12} (Le Pen (RN)) --- Permis de conduire gratuit (18 ans,
  casier vierge) ; prime d'activité étudiante. \emph{Sources} :
  RN\_LCP\_JEUNES. \emph{Règle} : 0,7 + 0,6 Md€.
\item
  \textbf{P13} (Le Pen (RN)) --- Part fiscale pleine dès le deuxième
  enfant. \emph{Sources} : RN\_LCP\_JEUNES. \emph{Règle} : {[}H{]}
  Transfert non ciblé.
\item
  \textbf{P14} (Le Pen (RN)) --- Dons familiaux défiscalisés ; prêt à
  taux zéro au troisième enfant ; suppression de la CVEC. \emph{Sources}
  : RN\_LCP\_JEUNES. \emph{Règle} : ---
\item
  \textbf{P15} (Le Pen (RN)) --- Suppression de la CFE et des impôts de
  production « qui nuisent à la relocalisation ». \emph{Sources} :
  RN\_VOTONS. \emph{Règle} : Montant non publié : canal
  impots\_production.
\item
  \textbf{P16} (Le Pen (RN)) --- Priorité nationale (logement, emploi,
  prestations). \emph{Sources} : RN\_LCP\_JEUNES. \emph{Règle} : ---
\item
  \textbf{H01} (Philippe (Horizons)) --- Près de 150 Md€ d'économies ;
  déficit sous 2,5 \% du PIB en fin de quinquennat ; moitié de l'effort
  sur les dépenses sociales, un quart sur l'État et les collectivités,
  un quart sur la croissance. \emph{Sources} : EP\_COMPTES,
  EP\_PS\_VOL1. \emph{Règle} : Enveloppe agrégée.
\item
  \textbf{H02} (Philippe (Horizons)) --- Règle d'or constitutionnelle
  par référendum ; aucun nouvel impôt pendant cinq ans. \emph{Sources} :
  EP\_COMPTES, EP\_PS\_VOL1. \emph{Règle} : Élément de régime :
  trajectoire de rétablissement (M23 R2, mode actif, instruments de
  dépense).
\item
  \textbf{H03} (Philippe (Horizons)) --- 250 000 postes de
  fonctionnaires en moins (un départ sur trois non remplacé dans les
  fonctions de bureau). \emph{Sources} : EP\_COMPTES, EP\_PS\_VOL1.
  \emph{Règle} : Effectif atteint en 5 ans ; coût moyen du simulateur 56
  k€ par agent.
\item
  \textbf{H04} (Philippe (Horizons)) --- Hausse de 20 \% du salaire des
  enseignants. \emph{Sources} : EP\_PS\_VOL1. \emph{Règle} : Masse
  salariale des enseignants non disponible dans le modèle : canal
  education (Md€) à chiffrer.
\item
  \textbf{H05} (Philippe (Horizons)) --- Âge de départ à 65 ans et 45
  annuités ; départ précoce pour un tiers des assurés ; capitalisation
  collective. \emph{Sources} : EP\_COMPTES, EP\_PS\_VOL1, EP\_EURONEWS.
  \emph{Règle} : M9 : âge 65 ans et 45 annuités (180 trimestres), montée
  d'un trimestre par génération {[}H{]} ; le tiers des assurés annoncé
  comme pouvant partir plus tôt garde les conditions actuelles {[}H{]}.
  Capitalisation collective non chiffrée : non représentée.
\item
  \textbf{H06} (Philippe (Horizons)) --- Fin des 35 heures ; deux heures
  supplémentaires hebdomadaires exonérées de cotisations. \emph{Sources}
  : EP\_PS\_VOL1, EP\_EURONEWS. \emph{Règle} : {[}H{]} Durée de
  référence portée à 37 h (2 h supplémentaires) ; exonération non
  chiffrée.
\item
  \textbf{H07} (Philippe (Horizons)) --- Indemnisation limitée à 12 mois
  pour les moins de 50 ans. \emph{Sources} : EP\_COMPTES, EP\_PS\_VOL1,
  EP\_EURONEWS. \emph{Règle} : Le curseur agit sur les montants, pas sur
  la durée : montant à chiffrer.
\item
  \textbf{H08} (Philippe (Horizons)) --- Baisse de 30 \% du nombre
  d'allocataires du RSA (accompagnement, métiers en tension).
  \emph{Sources} : EP\_COMPTES, EP\_PS\_VOL1. \emph{Règle} : {[}H{]} Le
  curseur réduit le montant forfaitaire (−30 \%) : approximation
  budgétaire d'une baisse des effectifs ; la réponse de participation
  passe par M15 (gain à l'emploi accru), pas par l'accompagnement
  annoncé ; effet distributif différent d'une baisse d'effectifs.
\item
  \textbf{H09} (Philippe (Horizons)) --- Trois jours de carence pour
  tous (public et privé). \emph{Sources} : EP\_COMPTES, EP\_PS\_VOL1.
  \emph{Règle} : Pas de canal des indemnités journalières.
\item
  \textbf{H10} (Philippe (Horizons)) --- Allocation de 350 € par mois
  dès le premier enfant, dégressive ; exonération d'impôt les trois
  premières années de l'enfant ; suppression des frais de notaire au
  premier achat (\textless{} 35 ans) ; prêt à taux zéro famille.
  \emph{Sources} : EP\_COMPTES, EP\_PS\_VOL1, EP\_EURONEWS. \emph{Règle}
  : Barèmes non publiés.
\item
  \textbf{H11} (Philippe (Horizons)) --- Baisse de 50 Md€ des impôts de
  production financée par 50 Md€ de baisse des aides aux entreprises.
  \emph{Sources} : EP\_PS\_VOL1, EP\_DEAL. \emph{Règle} : {[}H{]}
  Composition des aides non publiée. Le modèle représente une aide
  supprimée de deux façons : (a) aide assimilable à une moindre
  imposition de la production (subventions, dépenses fiscales) : même
  canal que les impôts de production, de signe opposé, donc effet net
  nul ; (b) allègement de cotisations supprimé : hausse du coût du
  travail (cotisations\_employeurs). Central : moitié (a), moitié (b),
  soit impôts de production −25 Md€ nets et cotisations +25 Md€.
  Variantes : (a) seul (mesure neutre dans le modèle) ; (b) seul (impôts
  de production −50, cotisations +50). Limite : le modèle ne distingue
  pas les distorsions propres aux impôts de production de celles des
  subventions (argument central du candidat). Début : première loi de
  finances du quinquennat (2028).
\item
  \textbf{H12} (Philippe (Horizons)) --- Fonds unique d'adaptation au
  changement climatique. \emph{Sources} : EP\_EURONEWS. \emph{Règle} :
  ---
\item
  \textbf{H13} (Philippe (Horizons)) --- Fusion, suppression ou
  privatisation d'une trentaine d'agences (dont France Travail).
  \emph{Sources} : EP\_COMPTES, EP\_PS\_VOL1. \emph{Règle} : ---
\item
  \textbf{H14} (Philippe (Horizons)) --- Huit réacteurs nucléaires
  supplémentaires ; suppression du ZAN. \emph{Sources} : EP\_PS\_VOL1.
  \emph{Règle} : ---
\end{itemize}

\textbf{Sources du corpus.}

\begin{itemize}
\tightlist
\item
  \texttt{LFI\_BUDGET} :
  https://melenchon2027.fr/wp-content/uploads/2026/10/Comment-nous-allons-remettre-de-l-ordre-dans-les-comptes-publics.pdf
\item
  \texttt{LFI\_AEC} : https://melenchon2027.fr/?p=24158
\item
  \texttt{LFI\_PRESSE} :
  https://melenchon2027.fr/wp-content/uploads/2026/11/DOSSIER-DE-PRESSE\_0611\_PROGRAMME\_ST.pdf
\item
  \texttt{LFI\_LCP\_BUDGET} :
  https://lcp.fr/actualites/contre-budget-lfi-propose-1064-milliards-d-euros-de-recettes-supplementaires-et-817
\item
  \texttt{LFI\_EPARGNE} :
  https://lessentieldeleco.fr/9362-melenchon-epargne/
\item
  \texttt{RN\_PS\_BUDGET} :
  https://www.publicsenat.fr/actualites/politique/140-milliards-deuros-deconomies-dici-2032-leffort-radical-propose-par-marine-le-pen-pour-redresser-les-f
\item
  \texttt{RN\_LCP\_JEUNES} :
  https://lcp.fr/actualites/permis-de-conduire-gratuit-exoneration-d-impot-jusqu-a-30-ans-comment-le-rn-veut
\item
  \texttt{RN\_MONEYVOX} :
  https://www.moneyvox.fr/impot/actualites/110673/140-milliards-economies-ce-quil-faut-retenir-de-intervention-de-marine-le-pen
\item
  \texttt{RN\_RETRAITES\_PS} :
  https://www.publicsenat.fr/actualites/politique/presidentielle-2027-que-contient-la-reforme-des-retraites-presentee-par-marine-le-pen
\item
  \texttt{RN\_RETRAITES\_SUD} :
  https://www.sudradio.fr/sud-radio/age-de-depart-cotisation-et-capitalisation-les-propositions-du-rn-pour-reformer-les-retraites
\item
  \texttt{RN\_RETRAITES\_BRUT} :
  https://www.brut.media/fr/articles/le-pen-affirme-que-sa-reforme-des-retraites-sera-engagee-des-son-arrivee-au-pouvoir-et-sera-appliquee-au-1er-janvier-2028
\item
  \texttt{RN\_ENERGIE} :
  https://www.les-energies-renouvelables.eu/article/actualites/energies/baisse-de-la-tva-arret-du-photovoltaique-ce-que-le-programme-energie-du-rn-changerait-pour-les-menages-2302/
\item
  \texttt{RN\_VOTONS} :
  https://votons-2027.fr/candidats/le-pen/programme
\item
  \texttt{MONTAIGNE\_TVA\_2022} :
  https://www.institutmontaigne.org/presidentielle-2022/marine-le-pen/baisser-la-tva-de-20-a-55-pour-les-carburants-lelectricite-le-gaz-et-le-fioul-domestique
\item
  \texttt{EP\_COMPTES} :
  https://www.edouardphilippe.fr/programme/comptes-publics
\item
  \texttt{EP\_PS\_VOL1} :
  https://www.publicsenat.fr/actualites/politique/capitalisation-suppression-de-250-000-postes-de-fonctionnaires-autonomie-des-ecoles-fin-des-35-heures-aides-pour-les-familles-edouard-philippe-presente-son-projet-massif-et
\item
  \texttt{EP\_EURONEWS} :
  https://fr.euronews.com/2026/10/09/presidentielle-2027-edouard-philippe-mise-sur-le-travail-et-la-reduction-des-depenses-publ
\item
  \texttt{EP\_DEAL} :
  https://www.publicsenat.fr/actualites/economie/entreprises-un-deal-a-250-milliards-deuros-sur-un-quinquennat-que-represente-la-proposition-dedouard-philippe
\item
  \texttt{SMIC\_2026} :
  https://www.info.gouv.fr/public/index.php/actualite/le-smic-revalorise-le-1er-juin-2026
\end{itemize}

\hypertarget{sec:annexe-d}{%
\chapter{Corpus bibliographique et documentaire}\label{sec:annexe-d}}

Les références sont classées en quatre groupes : littérature académique,
institutions et statistique publique, textes juridiques, documents des
candidats et presse. Les documents en ligne ont été consultés entre le
14 et le 24 novembre 2026 (date du corpus). Les références académiques
sont citées sous la forme (auteur, année) dans le texte. Quand une
référence n'a pas pu être vérifiée dans sa forme bibliographique
complète, seuls l'auteur, l'année et le titre sont donnés.

\hypertarget{d.1-littuxe9rature-acaduxe9mique}{%
\section{D.1 Littérature
académique}\label{d.1-littuxe9rature-acaduxe9mique}}

\begin{itemize}
\tightlist
\item
  Aghion, P. et Howitt, P. (1992), « A Model of Growth through Creative
  Destruction », \emph{Econometrica}, 60(2).
\item
  Alesina, A. et Ardagna, S. (2010), « Large Changes in Fiscal Policy:
  Taxes versus Spending », \emph{Tax Policy and the Economy}, 24.
\item
  Auerbach, A. et Gorodnichenko, Y. (2012), « Measuring the Output
  Responses to Fiscal Policy », \emph{American Economic Journal:
  Economic Policy}, 4(2).
\item
  Autor, D., Dorn, D. et Hanson, G. (2013), « The China Syndrome: Local
  Labor Market Effects of Import Competition in the United States »,
  \emph{American Economic Review}, 103(6).
\item
  Baker, S., Bloom, N. et Davis, S. (2016), « Measuring Economic Policy
  Uncertainty », \emph{Quarterly Journal of Economics}, 131(4).
\item
  Barro, R. (1974), « Are Government Bonds Net Wealth? », \emph{Journal
  of Political Economy}, 82(6).
\item
  Benzarti, Y., Carloni, D., Harju, J. et Kosonen, T. (2020), « What
  Goes Up May Not Come Down: Asymmetric Incidence of Value-Added Taxes
  », \emph{Journal of Political Economy}, 128(12).
\item
  Blanchard, O. (2019), « Public Debt and Low Interest Rates »,
  \emph{American Economic Review}, 109(4).
\item
  Blanchard, O. et Leigh, D. (2013), « Growth Forecast Errors and Fiscal
  Multipliers », \emph{American Economic Review}, 103(3).
\item
  Blanchard, O. et Perotti, R. (2002), « An Empirical Characterization
  of the Dynamic Effects of Changes in Government Spending and Taxes on
  Output », \emph{Quarterly Journal of Economics}, 117(4).
\item
  Blanchard, O. et Summers, L. (1986), « Hysteresis and the European
  Unemployment Problem », \emph{NBER Macroeconomics Annual}, 1.
\item
  Bloom, N. (2009), « The Impact of Uncertainty Shocks »,
  \emph{Econometrica}, 77(3).
\item
  Blundell, R., Pistaferri, L. et Preston, I. (2008), « Consumption
  Inequality and Partial Insurance », \emph{American Economic Review},
  98(5).
\item
  Bohn, H. (1998), « The Behavior of U.S. Public Debt and Deficits »,
  \emph{Quarterly Journal of Economics}, 113(3).
\item
  Bom, P. et Ligthart, J. (2014), « What Have We Learned from Three
  Decades of Research on the Productivity of Public Capital? »,
  \emph{Journal of Economic Surveys}, 28(5).
\item
  Bozio, A. (2011), « La réforme des retraites de 1993 : l'impact de
  l'augmentation de la durée d'assurance », \emph{Économie et
  Statistique}, 441-442, Insee.
\item
  Bozio, A., Breda, T. et Grenet, J. (2019), travaux sur l'incidence des
  cotisations sociales selon le lien entre cotisations et prestations
  (cités dans le registre des paramètres du projet).
\item
  Cabannes, P.-Y., Houdré, C. et Landais, C. (2014), travaux sur
  l'élasticité du revenu imposable en France (cités dans le registre des
  paramètres du projet).
\item
  Calvo, G. (1983), « Staggered Prices in a Utility-Maximizing Framework
  », \emph{Journal of Monetary Economics}, 12(3).
\item
  Carroll, C. (1997), « Buffer-Stock Saving and the Life Cycle/Permanent
  Income Hypothesis », \emph{Quarterly Journal of Economics}, 112(1).
\item
  Chetty, R., Guren, A., Manoli, D. et Weber, A. (2011), « Are Micro and
  Macro Labor Supply Elasticities Consistent? A Review of Evidence on
  the Intensive and Extensive Margins », \emph{American Economic
  Review}, 101(3).
\item
  Clarida, R., Galí, J. et Gertler, M. (1999), « The Science of Monetary
  Policy: A New Keynesian Perspective », \emph{Journal of Economic
  Literature}, 37(4).
\item
  Deaton, A. (1991), « Saving and Liquidity Constraints »,
  \emph{Econometrica}, 59(5).
\item
  Engle, R. et Granger, C. (1987), « Co-Integration and Error
  Correction: Representation, Estimation, and Testing »,
  \emph{Econometrica}, 55(2).
\item
  Fisher, I. (1933), « The Debt-Deflation Theory of Great Depressions »,
  \emph{Econometrica}, 1(4).
\item
  Friedman, M. (1957), \emph{A Theory of the Consumption Function},
  Princeton University Press.
\item
  Friedman, M. (1968), « The Role of Monetary Policy », \emph{American
  Economic Review}, 58(1).
\item
  Fuest, C., Peichl, A. et Siegloch, S. (2018), « Do Higher Corporate
  Taxes Reduce Wages? Micro Evidence from Germany », \emph{American
  Economic Review}, 108(2).
\item
  Galí, J. (2015), \emph{Monetary Policy, Inflation, and the Business
  Cycle}, 2\textsuperscript{e} éd., Princeton University Press.
\item
  Giavazzi, F. et Pagano, M. (1990), « Can Severe Fiscal Contractions Be
  Expansionary? Tales of Two Small European Countries », \emph{NBER
  Macroeconomics Annual}, 5.
\item
  Godley, W. et Lavoie, M. (2007), \emph{Monetary Economics: An
  Integrated Approach to Credit, Money, Income, Production and Wealth},
  Palgrave Macmillan.
\item
  Guajardo, J., Leigh, D. et Pescatori, A. (2014), « Expansionary
  Austerity? International Evidence », \emph{Journal of the European
  Economic Association}, 12(4).
\item
  Hall, B., Mairesse, J. et Mohnen, P. (2010), « Measuring the Returns
  to R\&D », in \emph{Handbook of the Economics of Innovation}, vol.~2,
  Elsevier.
\item
  Hall, R. et Jorgenson, D. (1967), « Tax Policy and Investment Behavior
  », \emph{American Economic Review}, 57(3).
\item
  Havranek, T. (2015), « Measuring Intertemporal Substitution: The
  Importance of Method Choices and Selective Reporting », \emph{Journal
  of the European Economic Association}, 13(6).
\item
  Herndon, T., Ash, M. et Pollin, R. (2014), « Does High Public Debt
  Consistently Stifle Economic Growth? A Critique of Reinhart and Rogoff
  », \emph{Cambridge Journal of Economics}, 38(2).
\item
  Hicks, J. (1937), « Mr.~Keynes and the ``Classics''; A Suggested
  Interpretation », \emph{Econometrica}, 5(2).
\item
  Jappelli, T. et Pistaferri, L. (2010), « The Consumption Response to
  Income Changes », \emph{Annual Review of Economics}, 2.
\item
  Jorgenson, D. (1963), « Capital Theory and Investment Behavior »,
  \emph{American Economic Review}, 53(2).
\item
  Kaplan, G., Moll, B. et Violante, G. (2018), « Monetary Policy
  According to HANK », \emph{American Economic Review}, 108(3).
\item
  Keynes, J. M. (1936), \emph{The General Theory of Employment, Interest
  and Money}, Macmillan.
\item
  Le Barbanchon, T., Rathelot, R. et Roulet, A. (2019), « Unemployment
  Insurance and Reservation Wages: Evidence from Administrative Data »,
  \emph{Journal of Public Economics}, 171.
\item
  Leeper, E. (1991), « Equilibria under ``Active'' and ``Passive''
  Monetary and Fiscal Policies », \emph{Journal of Monetary Economics},
  27(1).
\item
  Lucas, R. (1972), « Expectations and the Neutrality of Money »,
  \emph{Journal of Economic Theory}, 4(2).
\item
  Lucas, R. (1976), « Econometric Policy Evaluation: A Critique »,
  \emph{Carnegie-Rochester Conference Series on Public Policy}, 1.
\item
  Lucas, R. (1988), « On the Mechanics of Economic Development »,
  \emph{Journal of Monetary Economics}, 22(1).
\item
  Mankiw, N. G., Romer, D. et Weil, D. (1992), « A Contribution to the
  Empirics of Economic Growth », \emph{Quarterly Journal of Economics},
  107(2).
\item
  Modigliani, F. et Brumberg, R. (1954), « Utility Analysis and the
  Consumption Function: An Interpretation of Cross-Section Data », in
  Kurihara (dir.), \emph{Post-Keynesian Economics}.
\item
  Mortensen, D. et Pissarides, C. (1994), « Job Creation and Job
  Destruction in the Theory of Unemployment », \emph{Review of Economic
  Studies}, 61(3).
\item
  Muth, J. (1961), « Rational Expectations and the Theory of Price
  Movements », \emph{Econometrica}, 29(3).
\item
  North, D. (1990), \emph{Institutions, Institutional Change and
  Economic Performance}, Cambridge University Press.
\item
  Okun, A. (1962), « Potential GNP: Its Measurement and Significance »,
  \emph{Proceedings of the Business and Economic Statistics Section},
  American Statistical Association.
\item
  Phelps, E. (1967), « Phillips Curves, Expectations of Inflation and
  Optimal Unemployment over Time », \emph{Economica}, 34(135).
\item
  Phillips, A. W. (1958), « The Relation between Unemployment and the
  Rate of Change of Money Wage Rates in the United Kingdom, 1861-1957 »,
  \emph{Economica}, 25(100).
\item
  Ramey, V. (2011), « Identifying Government Spending Shocks: It's All
  in the Timing », \emph{Quarterly Journal of Economics}, 126(1).
\item
  Ramey, V. (2019), « Ten Years after the Financial Crisis: What Have We
  Learned from the Renaissance in Fiscal Research? », \emph{Journal of
  Economic Perspectives}, 33(2).
\item
  Ramey, V. et Zubairy, S. (2018), « Government Spending Multipliers in
  Good Times and in Bad: Evidence from US Historical Data »,
  \emph{Journal of Political Economy}, 126(2).
\item
  Reinhart, C. et Rogoff, K. (2010), « Growth in a Time of Debt »,
  \emph{American Economic Review}, 100(2).
\item
  Romer, C. et Romer, D. (2010), « The Macroeconomic Effects of Tax
  Changes: Estimates Based on a New Measure of Fiscal Shocks »,
  \emph{American Economic Review}, 100(3).
\item
  Romer, P. (1990), « Endogenous Technological Change », \emph{Journal
  of Political Economy}, 98(5).
\item
  Saez, E., Slemrod, J. et Giertz, S. (2012), « The Elasticity of
  Taxable Income with Respect to Marginal Tax Rates: A Critical Review
  », \emph{Journal of Economic Literature}, 50(1).
\item
  Sargent, T. et Wallace, N. (1975), « ``Rational'' Expectations, the
  Optimal Monetary Instrument, and the Optimal Money Supply Rule »,
  \emph{Journal of Political Economy}, 83(2).
\item
  Smets, F. et Wouters, R. (2007), « Shocks and Frictions in US Business
  Cycles: A Bayesian DSGE Approach », \emph{American Economic Review},
  97(3).
\item
  Solow, R. (1956), « A Contribution to the Theory of Economic Growth »,
  \emph{Quarterly Journal of Economics}, 70(1).
\item
  Taylor, J. (1980), « Aggregate Dynamics and Staggered Contracts »,
  \emph{Journal of Political Economy}, 88(1).
\item
  Taylor, J. (1993), « Discretion versus Policy Rules in Practice »,
  \emph{Carnegie-Rochester Conference Series on Public Policy}, 39.
\item
  Tobin, J. (1969), « A General Equilibrium Approach to Monetary Theory
  », \emph{Journal of Money, Credit and Banking}, 1(1).
\item
  Woodford, M. (2003), \emph{Interest and Prices: Foundations of a
  Theory of Monetary Policy}, Princeton University Press.
\end{itemize}

\hypertarget{d.2-institutions-et-statistique-publique}{%
\section{D.2 Institutions et statistique
publique}\label{d.2-institutions-et-statistique-publique}}

\begin{itemize}
\tightlist
\item
  Conseil d'orientation des retraites (2026), \emph{Rapport annuel},
  juin 2026 (partie 1, chapitre 2 : aménagement de la montée en charge
  de l'âge d'ouverture des droits et de la durée d'assurance) ---
  https://www.cor-retraites.fr/sites/default/files/2026-06/RA\_Part1\_Chap2.pdf
\item
  Conseil d'orientation des retraites (2024), séance du 14 mars 2024 «
  Départs anticipés » ---
  https://www.cor-retraites.fr/sites/default/files/2024-03/Doc\_01\_Note\%20de\%20pr\%C3\%A9sentation\%20g\%C3\%A9n\%C3\%A9rale.pdf
\item
  DREES (2025), \emph{Les retraités et les retraites}, édition 2025
  (fiches 11, 15, 16, 17) ---
  https://drees.solidarites-sante.gouv.fr/publications-communique-de-presse-documents-de-reference/250731\_PANORAMAS-retraites
\item
  CNAV, \emph{Recueil statistique du régime général}, données 2024
  (nouveaux retraités de droit direct) ---
  https://www.statistiques-recherche.lassuranceretraite.fr/app/uploads/2026/02/2.1-Les-nouveaux-retraites-de-droit-direct-1.pdf
\item
  Insee, comptes nationaux annuels, tableau des entrées-sorties,
  projections de population 2026, recensement 2022 ; Banque de données
  macroéconomiques.
\item
  Eurostat, comptes nationaux (\texttt{nama\_10\_*},
  \texttt{nasa\_10\_*}), finances publiques (\texttt{gov\_10a\_main},
  \texttt{gov\_10dd\_edpt1}), enquête sur les forces de travail
  (\texttt{lfsa\_argan}), protection sociale (\texttt{spr\_*}), taux
  d'intérêt (\texttt{irt\_*}).
\item
  Banque centrale européenne, statistiques des taux d'intérêt des
  institutions financières monétaires (MIR).
\item
  Fonds monétaire international, \emph{World Economic Outlook}
  (millésimes 2010-2026) ; FMI (2010), \emph{World Economic Outlook},
  octobre, chapitre 3 « Will It Hurt? Macroeconomic Effects of Fiscal
  Consolidation ».
\item
  France Stratégie (2020), comité de suivi et d'évaluation du CICE,
  rapport final (synthèse des évaluations).
\item
  Institut Montaigne (2022), chiffrage des programmes présidentiels
  (fiche « Baisser la TVA de 20 à 5,5 \% pour les carburants,
  l'électricité, le gaz et le fioul domestique ») ---
  https://www.institutmontaigne.org/presidentielle-2022/marine-le-pen/baisser-la-tva-de-20-a-55-pour-les-carburants-lelectricite-le-gaz-et-le-fioul-domestique
\item
  Ministère de l'Économie, direction des affaires juridiques, note sur
  le barème des sanctions financières de l'article 260 TFUE ---
  https://www.economie.gouv.fr/node/3219467
\item
  Info.gouv.fr (2026), « Le SMIC revalorisé le 1er juin 2026 » ---
  https://www.info.gouv.fr/public/index.php/actualite/le-smic-revalorise-le-1er-juin-2026
\item
  Modèles de référence : Insee et DG Trésor, modèle Mésange (documents
  de travail) ; Banque de France, modèle FR-BDF (documents de travail) ;
  OFCE et ADEME, modèle ThreeME.
\end{itemize}

\hypertarget{d.3-textes-juridiques}{%
\section{D.3 Textes juridiques}\label{d.3-textes-juridiques}}

\begin{itemize}
\tightlist
\item
  Traité sur le fonctionnement de l'Union européenne, articles 3, 18,
  30, 45, 63, 65, 66, 107, 108, 126, 207, 258, 260, 311.
\item
  Règlement (UE) 2024/1263 du Parlement européen et du Conseil du 29
  avril 2024 (volet préventif).
\item
  Règlement (UE) 2024/1264 du Conseil du 29 avril 2024 modifiant le
  règlement (CE) n° 1467/97 (procédure concernant les déficits
  excessifs) --- résumé :
  https://oeil.europarl.europa.eu/oeil/en/document-summary?id=1784571
\item
  Directive (UE) 2024/1265 du Conseil (cadres budgétaires nationaux).
\item
  Directive 2006/112/CE modifiée par la directive (UE) 2022/542 du 5
  avril 2022 (taux de TVA) ---
  https://eur-lex.europa.eu/eli/dir/2022/542/oj
\item
  Directive (UE) 2022/2523 (niveau minimum mondial d'imposition des
  groupes).
\item
  Communication de la Commission (2025) relative au calcul des sanctions
  financières (art. 260 TFUE) ---
  https://eur-lex.europa.eu/eli/C/2025/1481/oj
\item
  Loi de financement de la sécurité sociale pour 2026 (aménagement du
  calendrier de la réforme de 2023) ; réponses ministérielles aux
  questions écrites n° 14752 et 15271 (Assemblée nationale,
  17\textsuperscript{e} législature).
\item
  Loi n° 2023-270 du 14 avril 2023 de financement rectificative de la
  sécurité sociale pour 2023 ; loi n° 2014-40 du 20 janvier 2014.
\end{itemize}

\hypertarget{d.4-documents-des-candidats-et-presse-consultuxe9s-en-novembre-2026}{%
\section{D.4 Documents des candidats et presse (consultés en novembre
2026)}\label{d.4-documents-des-candidats-et-presse-consultuxe9s-en-novembre-2026}}

Mélenchon (LFI) : - « Comment nous allons remettre de l'ordre dans les
comptes publics » (octobre 2026) ---
https://melenchon2027.fr/wp-content/uploads/2026/10/Comment-nous-allons-remettre-de-l-ordre-dans-les-comptes-publics.pdf
- \emph{L'Avenir en commun}, 5\textsuperscript{e} édition, et dossier de
presse du 7 novembre 2026 --- https://melenchon2027.fr/?p=24158 ;
https://melenchon2027.fr/wp-content/uploads/2026/11/DOSSIER-DE-PRESSE\_0611\_PROGRAMME\_ST.pdf
- LCP, « Contre-budget : LFI propose 106,4 milliards d'euros de recettes
supplémentaires et 81,7 milliards de dépenses nouvelles » (14 octobre
2026).

Le Pen (RN) : - Public Sénat, « 140 milliards d'euros d'économies d'ici
2032 : l'effort radical proposé par Marine Le Pen » (6 octobre 2026). -
LCP, « Permis de conduire gratuit, exonération d'impôt jusqu'à 30
ans\ldots{} » (6 octobre 2026). - Moneyvox, « 140 milliards d'économies
: ce qu'il faut retenir de l'intervention de Marine Le Pen ». - Public
Sénat, « Que contient la réforme des retraites présentée par Marine Le
Pen ? » ; Sud Radio ; Brut (17 novembre 2026). - Les énergies
renouvelables, « Baisse de la TVA, arrêt du photovoltaïque : ce que le
programme énergie du RN changerait pour les ménages ». - votons-2027.fr,
page programme de Marine Le Pen (agrégateur citant les sources
officielles).

Philippe (Horizons) : - Site de campagne, « Finances publiques » ---
https://www.edouardphilippe.fr/programme/comptes-publics - Public Sénat,
« Capitalisation, suppression de 250 000 postes de
fonctionnaires\ldots{} » (18 novembre 2026) ; Euronews (9 octobre 2026)
; Public Sénat, « Entreprises : un deal à 250 milliards d'euros sur un
quinquennat » (6 novembre 2026).

\hypertarget{sec:annexe-e}{%
\chapter{Résultats détaillés des simulations et
sensibilités}\label{sec:annexe-e}}

Écarts à la référence (même environnement, même régime), horizons 5, 10
et 20 ans (2031, 2036, 2046). Unités : \% pour les volumes et les prix,
points pour les taux et les ratios au PIB, milliers pour l'emploi.
Toutes les simulations sont archivées dans
\texttt{phase7/resultats/\textless{}candidat\textgreater{}.pkl}
(trajectoires annuelles complètes) et
\texttt{\textless{}candidat\textgreater{}\_synthese.csv}.

\hypertarget{e.1-lectures-a-et-b-dix-sept-indicateurs}{%
\section{E.1 Lectures A et B, dix-sept
indicateurs}\label{e.1-lectures-a-et-b-dix-sept-indicateurs}}

\hypertarget{muxe9lenchon-lfi-lecture-a}{%
\subsection{Mélenchon (LFI) --- lecture
A}\label{muxe9lenchon-lfi-lecture-a}}

\begin{longtable}[]{@{}llll@{}}
\toprule\noalign{}
indicateur (écart à la référence) & 5 ans & 10 ans & 20 ans \\
\midrule\noalign{}
\endhead
\bottomrule\noalign{}
\endlastfoot
PIB (\%) & +1,84 & −0,54 & −1,07 \\
emploi (milliers) & +606,60 & −48,77 & −539,33 \\
chômage (pt) & −4,51 & −2,78 & −1,39 \\
salaire réel (\%) & +5,20 & +8,91 & +12,12 \\
prix (\%) & +3,96 & +6,99 & +5,63 \\
pouvoir d'achat (\%) & +3,35 & +4,33 & +5,28 \\
consommation (\%) & +4,18 & +4,17 & +5,34 \\
taux d'épargne (pt) & −0,70 & +0,13 & −0,05 \\
investissement (\%) & +18,90 & +17,69 & +17,79 \\
taux de marge (pt) & −1,04 & −1,71 & −1,01 \\
exportations (\%) & −2,66 & −4,47 & −4,05 \\
importations (\%) & +3,33 & +0,39 & −0,83 \\
solde (pt PIB) & +3,32 & +3,71 & +5,61 \\
dette (pt PIB) & −20,11 & −38,14 & −73,61 \\
intérêts (pt PIB) & −0,61 & −1,55 & −3,17 \\
recettes (pt PIB) & +2,74 & +2,96 & +2,99 \\
dépenses (pt PIB) & −0,58 & −0,74 & −2,61 \\
\end{longtable}

\hypertarget{muxe9lenchon-lfi-lecture-b}{%
\subsection{Mélenchon (LFI) --- lecture
B}\label{muxe9lenchon-lfi-lecture-b}}

\begin{longtable}[]{@{}llll@{}}
\toprule\noalign{}
indicateur (écart à la référence) & 5 ans & 10 ans & 20 ans \\
\midrule\noalign{}
\endhead
\bottomrule\noalign{}
\endlastfoot
PIB (\%) & +2,66 & −1,19 & −2,70 \\
emploi (milliers) & +775,89 & −77,38 & −962,34 \\
chômage (pt) & −5,02 & −2,70 & −0,10 \\
salaire réel (\%) & +5,58 & +9,50 & +11,31 \\
prix (\%) & +4,05 & +7,89 & +5,04 \\
pouvoir d'achat (\%) & +5,58 & +3,51 & −1,98 \\
consommation (\%) & +6,35 & +3,57 & −1,71 \\
taux d'épargne (pt) & −0,66 & −0,05 & −0,23 \\
investissement (\%) & +19,77 & +17,04 & +15,84 \\
taux de marge (pt) & −1,00 & −2,09 & −0,93 \\
exportations (\%) & −2,78 & −4,99 & −3,59 \\
importations (\%) & +4,27 & −0,56 & −2,36 \\
solde (pt PIB) & +2,79 & +4,10 & +7,89 \\
dette (pt PIB) & −18,74 & −36,08 & −83,25 \\
intérêts (pt PIB) & −0,53 & −1,44 & −3,56 \\
recettes (pt PIB) & +1,75 & +3,53 & +6,18 \\
dépenses (pt PIB) & −1,03 & −0,57 & −1,71 \\
\end{longtable}

\hypertarget{le-pen-rn-lecture-a}{%
\subsection{Le Pen (RN) --- lecture A}\label{le-pen-rn-lecture-a}}

\begin{longtable}[]{@{}llll@{}}
\toprule\noalign{}
indicateur (écart à la référence) & 5 ans & 10 ans & 20 ans \\
\midrule\noalign{}
\endhead
\bottomrule\noalign{}
\endlastfoot
PIB (\%) & −1,95 & −2,89 & −3,64 \\
emploi (milliers) & −440,35 & −754,30 & −1049,67 \\
chômage (pt) & −1,84 & −1,65 & −1,08 \\
salaire réel (\%) & +0,59 & +1,86 & +4,01 \\
prix (\%) & +0,16 & +1,09 & +1,54 \\
pouvoir d'achat (\%) & −0,90 & −0,62 & −0,00 \\
consommation (\%) & −0,48 & −0,68 & −0,19 \\
taux d'épargne (pt) & −0,35 & +0,05 & +0,16 \\
investissement (\%) & −0,87 & −1,14 & −1,48 \\
taux de marge (pt) & −0,29 & −0,35 & −0,31 \\
exportations (\%) & −0,33 & −0,92 & −1,40 \\
importations (\%) & −2,25 & −3,36 & −3,94 \\
solde (pt PIB) & +0,33 & +0,03 & +0,15 \\
dette (pt PIB) & −0,34 & −1,60 & −3,34 \\
intérêts (pt PIB) & +0,01 & −0,05 & −0,13 \\
recettes (pt PIB) & +0,24 & +0,33 & +0,43 \\
dépenses (pt PIB) & −0,09 & +0,30 & +0,28 \\
\end{longtable}

\hypertarget{le-pen-rn-lecture-b}{%
\subsection{Le Pen (RN) --- lecture B}\label{le-pen-rn-lecture-b}}

\begin{longtable}[]{@{}llll@{}}
\toprule\noalign{}
indicateur (écart à la référence) & 5 ans & 10 ans & 20 ans \\
\midrule\noalign{}
\endhead
\bottomrule\noalign{}
\endlastfoot
PIB (\%) & −9,59 & −11,12 & −6,59 \\
emploi (milliers) & −2560,63 & −4066,79 & −3091,77 \\
chômage (pt) & +4,65 & +8,52 & +5,27 \\
salaire réel (\%) & −1,30 & −5,45 & −15,90 \\
prix (\%) & −0,26 & −5,40 & −8,14 \\
pouvoir d'achat (\%) & −8,30 & −13,38 & −17,15 \\
consommation (\%) & −8,65 & −14,01 & −16,90 \\
taux d'épargne (pt) & +0,30 & +0,53 & −0,22 \\
investissement (\%) & −5,10 & −5,24 & −3,55 \\
taux de marge (pt) & −0,85 & +0,70 & +1,87 \\
exportations (\%) & +0,16 & +3,38 & +5,70 \\
importations (\%) & −6,75 & −3,88 & +1,88 \\
solde (pt PIB) & +3,19 & +5,28 & +8,89 \\
dette (pt PIB) & +2,79 & −5,64 & −65,65 \\
intérêts (pt PIB) & +0,06 & −0,38 & −2,79 \\
recettes (pt PIB) & +0,54 & +0,28 & −0,28 \\
dépenses (pt PIB) & −2,65 & −5,00 & −9,16 \\
\end{longtable}

\hypertarget{philippe-horizons-lecture-a}{%
\subsection{Philippe (Horizons) --- lecture
A}\label{philippe-horizons-lecture-a}}

\begin{longtable}[]{@{}llll@{}}
\toprule\noalign{}
indicateur (écart à la référence) & 5 ans & 10 ans & 20 ans \\
\midrule\noalign{}
\endhead
\bottomrule\noalign{}
\endlastfoot
PIB (\%) & +0,96 & +2,45 & +4,08 \\
emploi (milliers) & −622,43 & −267,47 & +255,87 \\
chômage (pt) & +2,56 & +2,56 & +1,93 \\
salaire réel (\%) & +1,69 & +0,39 & −2,74 \\
prix (\%) & −2,07 & −3,21 & −3,42 \\
pouvoir d'achat (\%) & −1,08 & −1,05 & −1,87 \\
consommation (\%) & −1,60 & −1,35 & −1,96 \\
taux d'épargne (pt) & +0,44 & +0,26 & +0,08 \\
investissement (\%) & +0,66 & +1,31 & +2,08 \\
taux de marge (pt) & +0,20 & +0,32 & +0,33 \\
exportations (\%) & +1,31 & +2,01 & +2,48 \\
importations (\%) & +2,43 & +4,20 & +5,69 \\
solde (pt PIB) & −0,57 & −0,53 & −0,68 \\
dette (pt PIB) & +5,28 & +8,42 & +13,21 \\
intérêts (pt PIB) & +0,15 & +0,33 & +0,57 \\
recettes (pt PIB) & −0,50 & −0,66 & −0,91 \\
dépenses (pt PIB) & +0,07 & −0,13 & −0,23 \\
\end{longtable}

\hypertarget{philippe-horizons-lecture-b}{%
\subsection{Philippe (Horizons) --- lecture
B}\label{philippe-horizons-lecture-b}}

\begin{longtable}[]{@{}llll@{}}
\toprule\noalign{}
indicateur (écart à la référence) & 5 ans & 10 ans & 20 ans \\
\midrule\noalign{}
\endhead
\bottomrule\noalign{}
\endlastfoot
PIB (\%) & −5,18 & −3,45 & +1,13 \\
emploi (milliers) & −2114,92 & −2450,62 & −1155,81 \\
chômage (pt) & +6,94 & +8,85 & +5,99 \\
salaire réel (\%) & +0,25 & −4,55 & −15,56 \\
prix (\%) & −2,13 & −7,32 & −9,05 \\
pouvoir d'achat (\%) & −11,12 & −16,67 & −21,63 \\
consommation (\%) & −11,63 & −16,94 & −21,37 \\
taux d'épargne (pt) & +0,43 & +0,23 & −0,22 \\
investissement (\%) & −4,32 & −3,60 & −1,80 \\
taux de marge (pt) & −0,41 & +0,93 & +1,67 \\
exportations (\%) & +1,55 & +4,77 & +6,67 \\
importations (\%) & −2,47 & +1,85 & +6,97 \\
solde (pt PIB) & +3,34 & +5,73 & +9,42 \\
dette (pt PIB) & +3,40 & −11,59 & −77,09 \\
intérêts (pt PIB) & +0,04 & −0,58 & −3,29 \\
recettes (pt PIB) & −0,77 & −1,66 & −2,60 \\
dépenses (pt PIB) & −4,10 & −7,39 & −12,02 \\
\end{longtable}

\hypertarget{e.2-toutes-les-simulations-de-la-campagne-pib-chuxf4mage-dette}{%
\section{E.2 Toutes les simulations de la campagne (PIB, chômage,
dette)}\label{e.2-toutes-les-simulations-de-la-campagne-pib-chuxf4mage-dette}}

\hypertarget{muxe9lenchon-lfi-toutes-les-simulations}{%
\subsection{Mélenchon (LFI) --- toutes les
simulations}\label{muxe9lenchon-lfi-toutes-les-simulations}}

\begin{longtable}[]{@{}
  >{\raggedright\arraybackslash}p{(\columnwidth - 18\tabcolsep) * \real{0.1000}}
  >{\raggedright\arraybackslash}p{(\columnwidth - 18\tabcolsep) * \real{0.1000}}
  >{\raggedright\arraybackslash}p{(\columnwidth - 18\tabcolsep) * \real{0.1000}}
  >{\raggedright\arraybackslash}p{(\columnwidth - 18\tabcolsep) * \real{0.1000}}
  >{\raggedright\arraybackslash}p{(\columnwidth - 18\tabcolsep) * \real{0.1000}}
  >{\raggedright\arraybackslash}p{(\columnwidth - 18\tabcolsep) * \real{0.1000}}
  >{\raggedright\arraybackslash}p{(\columnwidth - 18\tabcolsep) * \real{0.1000}}
  >{\raggedright\arraybackslash}p{(\columnwidth - 18\tabcolsep) * \real{0.1000}}
  >{\raggedright\arraybackslash}p{(\columnwidth - 18\tabcolsep) * \real{0.1000}}
  >{\raggedright\arraybackslash}p{(\columnwidth - 18\tabcolsep) * \real{0.1000}}@{}}
\toprule\noalign{}
\begin{minipage}[b]{\linewidth}\raggedright
simulation
\end{minipage} & \begin{minipage}[b]{\linewidth}\raggedright
PIB (\%) 5a
\end{minipage} & \begin{minipage}[b]{\linewidth}\raggedright
PIB (\%) 10a
\end{minipage} & \begin{minipage}[b]{\linewidth}\raggedright
PIB (\%) 20a
\end{minipage} & \begin{minipage}[b]{\linewidth}\raggedright
chômage (pt) 5a
\end{minipage} & \begin{minipage}[b]{\linewidth}\raggedright
chômage (pt) 10a
\end{minipage} & \begin{minipage}[b]{\linewidth}\raggedright
chômage (pt) 20a
\end{minipage} & \begin{minipage}[b]{\linewidth}\raggedright
dette (pt PIB) 5a
\end{minipage} & \begin{minipage}[b]{\linewidth}\raggedright
dette (pt PIB) 10a
\end{minipage} & \begin{minipage}[b]{\linewidth}\raggedright
dette (pt PIB) 20a
\end{minipage} \\
\midrule\noalign{}
\endhead
\bottomrule\noalign{}
\endlastfoot
A & +1,84 & −0,54 & −1,07 & −4,51 & −2,78 & −1,39 & −20,11 & −38,14 &
−73,61 \\
B & +2,66 & −1,19 & −2,70 & −5,02 & −2,70 & −0,10 & −18,74 & −36,08 &
−83,25 \\
A:oat\_382 & +1,82 & −0,55 & −1,11 & −4,49 & −2,78 & −1,37 & −19,92 &
−36,65 & −66,67 \\
B:oat\_382 & +3,66 & +0,16 & −1,35 & −5,43 & −3,55 & −0,89 & −17,85 &
−31,22 & −58,40 \\
mesure:L01 & −0,09 & −0,08 & −0,02 & +0,05 & +0,05 & +0,01 & −0,48 &
−0,98 & −2,35 \\
mesure:L02 & −0,11 & −0,09 & −0,03 & +0,07 & +0,06 & +0,01 & −0,66 &
−1,37 & −3,33 \\
mesure:L03 & +0,00 & +0,00 & +0,00 & +0,00 & +0,00 & +0,00 & −2,66 &
−5,68 & −13,24 \\
mesure:L04 & +0,00 & +0,00 & +0,00 & +0,00 & +0,00 & +0,00 & −3,19 &
−6,82 & −15,89 \\
mesure:L06 & −0,02 & −0,04 & +0,00 & −0,01 & +0,00 & −0,01 & −0,56 &
−1,01 & −2,18 \\
mesure:L07 & −0,04 & −0,07 & −0,06 & +0,01 & +0,03 & +0,04 & −0,23 &
−0,29 & −0,16 \\
mesure:L08 & +0,00 & +0,00 & +0,00 & +0,00 & +0,00 & +0,00 & −0,89 &
−1,89 & −4,41 \\
mesure:L09 & +0,00 & +0,00 & +0,00 & +0,00 & +0,00 & +0,00 & −0,69 &
−1,48 & −3,44 \\
mesure:L10 & +0,00 & +0,00 & +0,00 & +0,00 & +0,00 & +0,00 & −0,89 &
−1,89 & −4,41 \\
mesure:L11 & −0,39 & −0,25 & −0,09 & +0,33 & +0,29 & +0,12 & −2,35 &
−4,37 & −10,39 \\
mesure:L13 & +0,00 & +0,00 & +0,00 & +0,00 & +0,00 & +0,00 & −0,89 &
−1,89 & −4,41 \\
mesure:L14 & +0,00 & +0,00 & +0,00 & +0,00 & +0,00 & +0,00 & −0,44 &
−0,95 & −2,21 \\
mesure:L15 & +0,00 & +0,00 & +0,00 & +0,00 & +0,00 & +0,00 & −0,18 &
−0,38 & −0,88 \\
mesure:L16 & +0,00 & +0,00 & +0,00 & +0,00 & +0,00 & +0,00 & −0,35 &
−0,76 & −1,77 \\
mesure:L20 & +0,73 & +0,42 & +0,23 & −0,52 & −0,38 & −0,14 & +2,18 &
+5,22 & +13,90 \\
mesure:L21 & +0,11 & +0,04 & +0,04 & −0,09 & −0,06 & −0,03 & +0,57 &
+1,44 & +3,77 \\
mesure:L22 & +0,08 & +0,03 & +0,00 & −0,11 & −0,07 & −0,03 & −0,10 &
−0,03 & +0,27 \\
mesure:L23 & +0,14 & +0,05 & +0,05 & −0,11 & −0,07 & −0,04 & +0,67 &
+1,68 & +4,38 \\
mesure:L24 & −0,00 & +0,02 & +0,04 & +0,12 & +0,13 & +0,05 & −0,63 &
+0,70 & +3,70 \\
mesure:L26 & +0,01 & +0,00 & +0,00 & −0,02 & −0,01 & −0,01 & −0,02 &
−0,00 & +0,04 \\
mesure:L27 & +0,09 & +0,04 & +0,00 & −0,13 & −0,08 & −0,04 & −0,12 &
−0,03 & +0,32 \\
mesure:L28 & +0,09 & +0,03 & +0,00 & −0,12 & −0,08 & −0,04 & −0,12 &
−0,03 & +0,31 \\
mesure:L29 & +0,08 & +0,03 & +0,00 & −0,11 & −0,07 & −0,03 & −0,10 &
−0,03 & +0,27 \\
mesure:L30 & +0,12 & +0,05 & +0,00 & −0,16 & −0,11 & −0,05 & −0,15 &
−0,04 & +0,41 \\
mesure:L31 & +0,18 & +0,07 & +0,00 & −0,25 & −0,16 & −0,08 & −0,23 &
−0,06 & +0,62 \\
mesure:L32 & +0,26 & +0,39 & +0,72 & −0,02 & +0,07 & +0,12 & +1,28 &
+3,14 & +8,09 \\
mesure:L33 & +1,70 & +0,65 & +0,54 & −1,27 & −0,97 & −0,59 & −6,00 &
−12,75 & −27,10 \\
mesure:L34 & −1,35 & −2,08 & −2,58 & −2,02 & −1,40 & −0,84 & −0,72 &
−0,47 & +1,74 \\
mesure:L35 & +0,02 & −0,00 & −0,02 & −0,02 & +0,00 & +0,02 & +0,22 &
+0,53 & +1,39 \\
regime:principal (R2 figé, taux exogènes, central) & +1,84 & −0,54 &
−1,07 & −4,51 & −2,78 & −1,39 & −20,11 & −38,14 & −73,61 \\
regime:sans règle budgétaire (R0) & +1,89 & −0,45 & −0,85 & −4,58 &
−2,81 & −1,23 & −20,66 & −41,08 & −84,99 \\
regime:stabilisation (R1) & +1,83 & −0,54 & −1,06 & −4,49 & −2,78 &
−1,39 & −20,23 & −38,23 & −73,84 \\
regime:R2 active symétrique & +1,53 & −0,77 & −2,13 & −4,33 & −2,64 &
−0,28 & −20,08 & −38,24 & −78,04 \\
regime:R2 + BCE et prime & +1,79 & −0,61 & −1,12 & −4,49 & −2,73 & −1,35
& −19,76 & −36,90 & −74,54 \\
regime:R2, ancrage faible & +1,72 & −0,68 & −0,70 & −4,42 & −2,33 &
−1,09 & −26,85 & −56,96 & −118,84 \\
regime:R2, ancrage fort & +1,91 & −0,45 & −1,29 & −4,55 & −3,01 & −1,58
& −16,63 & −27,86 & −48,09 \\
regime:R2, fermeture lente & +2,84 & +0,91 & +0,40 & −4,96 & −3,79 &
−2,37 & −21,90 & −45,79 & −99,83 \\
regime:R2, fermeture rapide & +0,96 & −1,47 & −1,82 & −4,08 & −2,09 &
−0,88 & −18,53 & −32,12 & −57,16 \\
regime:R2, fermeture sur le chômage & +2,04 & −1,35 & −1,79 & −4,63 &
−2,40 & −0,75 & −20,64 & −37,37 & −64,21 \\
environnement:p4\_sigma03 & +1,84 & −0,61 & −1,20 & −4,69 & −2,91 &
−1,63 & −20,70 & −39,63 & −77,72 \\
environnement:p4\_sigma06 & +1,84 & −0,50 & −1,00 & −4,29 & −2,71 &
−1,18 & −19,39 & −36,59 & −69,83 \\
environnement:p4\_prix\_v09 & +1,85 & −0,59 & −1,15 & −4,50 & −2,77 &
−1,52 & −20,02 & −38,32 & −75,94 \\
europe:conforme & +1,79 & −0,61 & −1,12 & −4,49 & −2,73 & −1,35 & −19,76
& −36,90 & −74,54 \\
europe:contentieux & +2,03 & −0,45 & −1,09 & −4,36 & −2,86 & −1,40 &
−17,81 & −35,07 & −72,28 \\
europe:tension & +2,04 & −0,45 & −1,08 & −4,36 & −2,87 & −1,41 & −17,03
& −34,45 & −73,52 \\
additionnalite:50 & +0,93 & −0,89 & −1,36 & −3,82 & −2,31 & −1,13 &
−16,73 & −30,13 & −55,98 \\
additionnalite:25 & +0,47 & −1,06 & −1,52 & −3,47 & −2,07 & −1,00 &
−14,99 & −26,02 & −47,25 \\
additionnalite:0 & +0,02 & −1,24 & −1,67 & −3,12 & −1,83 & −0,87 &
−13,21 & −21,85 & −38,50 \\
mise\_en\_oeuvre:partielle\_50 & +0,96 & −0,29 & −0,55 & −2,27 & −1,52 &
−0,74 & −9,23 & −16,86 & −31,91 \\
mise\_en\_oeuvre:retard\_2ans & +3,00 & +0,10 & −0,97 & −4,33 & −3,44 &
−1,61 & −12,71 & −31,77 & −67,22 \\
\end{longtable}

\hypertarget{le-pen-rn-toutes-les-simulations}{%
\subsection{Le Pen (RN) --- toutes les
simulations}\label{le-pen-rn-toutes-les-simulations}}

\begin{longtable}[]{@{}
  >{\raggedright\arraybackslash}p{(\columnwidth - 18\tabcolsep) * \real{0.1000}}
  >{\raggedright\arraybackslash}p{(\columnwidth - 18\tabcolsep) * \real{0.1000}}
  >{\raggedright\arraybackslash}p{(\columnwidth - 18\tabcolsep) * \real{0.1000}}
  >{\raggedright\arraybackslash}p{(\columnwidth - 18\tabcolsep) * \real{0.1000}}
  >{\raggedright\arraybackslash}p{(\columnwidth - 18\tabcolsep) * \real{0.1000}}
  >{\raggedright\arraybackslash}p{(\columnwidth - 18\tabcolsep) * \real{0.1000}}
  >{\raggedright\arraybackslash}p{(\columnwidth - 18\tabcolsep) * \real{0.1000}}
  >{\raggedright\arraybackslash}p{(\columnwidth - 18\tabcolsep) * \real{0.1000}}
  >{\raggedright\arraybackslash}p{(\columnwidth - 18\tabcolsep) * \real{0.1000}}
  >{\raggedright\arraybackslash}p{(\columnwidth - 18\tabcolsep) * \real{0.1000}}@{}}
\toprule\noalign{}
\begin{minipage}[b]{\linewidth}\raggedright
simulation
\end{minipage} & \begin{minipage}[b]{\linewidth}\raggedright
PIB (\%) 5a
\end{minipage} & \begin{minipage}[b]{\linewidth}\raggedright
PIB (\%) 10a
\end{minipage} & \begin{minipage}[b]{\linewidth}\raggedright
PIB (\%) 20a
\end{minipage} & \begin{minipage}[b]{\linewidth}\raggedright
chômage (pt) 5a
\end{minipage} & \begin{minipage}[b]{\linewidth}\raggedright
chômage (pt) 10a
\end{minipage} & \begin{minipage}[b]{\linewidth}\raggedright
chômage (pt) 20a
\end{minipage} & \begin{minipage}[b]{\linewidth}\raggedright
dette (pt PIB) 5a
\end{minipage} & \begin{minipage}[b]{\linewidth}\raggedright
dette (pt PIB) 10a
\end{minipage} & \begin{minipage}[b]{\linewidth}\raggedright
dette (pt PIB) 20a
\end{minipage} \\
\midrule\noalign{}
\endhead
\bottomrule\noalign{}
\endlastfoot
A & −1,95 & −2,89 & −3,64 & −1,84 & −1,65 & −1,08 & −0,34 & −1,60 &
−3,34 \\
B & −9,59 & −11,12 & −6,59 & +4,65 & +8,52 & +5,27 & +2,79 & −5,64 &
−65,65 \\
A:oat\_382 & −1,96 & −2,89 & −3,67 & −1,83 & −1,65 & −1,08 & −0,36 &
−1,57 & −3,03 \\
B:oat\_382 & −9,01 & −8,31 & −4,87 & +4,17 & +5,71 & +2,59 & +2,50 &
−6,09 & −48,12 \\
mesure:P03 & −0,26 & −0,18 & −0,08 & +0,17 & +0,14 & +0,04 & −0,97 &
−2,16 & −5,52 \\
mesure:P04 & +0,00 & +0,00 & +0,00 & +0,00 & +0,00 & +0,00 & −1,59 &
−3,41 & −7,94 \\
mesure:P05 & −0,52 & −0,20 & −0,02 & +0,73 & +0,49 & +0,23 & +0,69 &
+0,18 & −1,87 \\
mesure:P08 & −1,39 & −2,67 & −3,59 & −2,59 & −2,07 & −1,26 & −1,03 &
−1,69 & +0,32 \\
mesure:P10 & +0,09 & +0,02 & +0,01 & −0,08 & −0,06 & −0,03 & +1,56 &
+2,98 & +6,75 \\
mesure:P11 & +0,02 & −0,02 & +0,01 & −0,02 & −0,01 & −0,01 & +0,17 &
+0,50 & +1,35 \\
mesure:P12 & +0,00 & −0,02 & +0,01 & −0,02 & −0,00 & −0,01 & +0,12 &
+0,40 & +1,07 \\
mesure:P13 & +0,04 & −0,00 & +0,02 & −0,04 & −0,02 & −0,02 & +0,27 &
+0,74 & +1,96 \\
regime:principal (R2 figé, taux exogènes, central) & −1,95 & −2,89 &
−3,64 & −1,84 & −1,65 & −1,08 & −0,34 & −1,60 & −3,34 \\
regime:sans règle budgétaire (R0) & −1,95 & −2,91 & −3,61 & −1,85 &
−1,57 & −0,82 & −0,91 & −4,79 & −14,02 \\
regime:stabilisation (R1) & −1,96 & −2,89 & −3,64 & −1,83 & −1,64 &
−1,08 & −0,39 & −1,65 & −3,42 \\
regime:R2 active symétrique & −2,03 & −3,38 & −4,50 & −1,92 & −1,30 &
−0,15 & +0,04 & −1,56 & −6,59 \\
regime:R2 + BCE et prime & −1,93 & −2,84 & −3,46 & −1,84 & −1,68 & −1,21
& −0,52 & −2,79 & −11,57 \\
regime:R2, ancrage faible & −1,95 & −2,91 & −3,52 & −1,83 & −1,53 &
−0,91 & −1,24 & −6,21 & −18,03 \\
regime:R2, ancrage fort & −1,95 & −2,87 & −3,71 & −1,84 & −1,70 & −1,18
& +0,10 & +0,74 & +4,56 \\
regime:R2, fermeture lente & −1,90 & −2,66 & −3,21 & −1,84 & −1,78 &
−1,39 & −0,32 & −2,15 & −7,39 \\
regime:R2, fermeture rapide & −2,02 & −3,07 & −3,91 & −1,82 & −1,53 &
−0,88 & −0,36 & −1,23 & −0,76 \\
regime:R2, fermeture sur le chômage & −1,27 & −2,92 & −4,03 & −2,20 &
−1,78 & −0,73 & −1,91 & −5,00 & −5,46 \\
environnement:p4\_sigma03 & −1,95 & −2,90 & −3,75 & −1,80 & −1,58 &
−1,05 & −0,26 & −1,19 & −2,28 \\
environnement:p4\_sigma06 & −1,96 & −2,90 & −3,58 & −1,90 & −1,79 &
−1,16 & −0,50 & −2,51 & −5,75 \\
environnement:p4\_prix\_v09 & −1,94 & −2,88 & −3,69 & −1,84 & −1,63 &
−1,04 & −0,20 & −1,15 & −2,15 \\
europe:conforme & −1,93 & −2,84 & −3,46 & −1,84 & −1,68 & −1,21 & −0,52
& −2,79 & −11,57 \\
europe:contentieux & −1,99 & −2,86 & −3,47 & −1,79 & −1,65 & −1,21 &
+0,12 & −0,94 & −6,83 \\
europe:tension & −1,99 & −2,86 & −3,47 & −1,79 & −1,65 & −1,21 & +1,13 &
+0,84 & −3,76 \\
mise\_en\_oeuvre:partielle\_50 & −0,89 & −1,32 & −1,79 & −0,82 & −0,88 &
−0,61 & −0,17 & −0,58 & −1,10 \\
mise\_en\_oeuvre:retard\_2ans & −1,41 & −2,66 & −3,60 & −1,73 & −1,82 &
−1,16 & +0,43 & −1,10 & −3,00 \\
\end{longtable}

\hypertarget{philippe-horizons-toutes-les-simulations}{%
\subsection{Philippe (Horizons) --- toutes les
simulations}\label{philippe-horizons-toutes-les-simulations}}

\begin{longtable}[]{@{}
  >{\raggedright\arraybackslash}p{(\columnwidth - 18\tabcolsep) * \real{0.1000}}
  >{\raggedright\arraybackslash}p{(\columnwidth - 18\tabcolsep) * \real{0.1000}}
  >{\raggedright\arraybackslash}p{(\columnwidth - 18\tabcolsep) * \real{0.1000}}
  >{\raggedright\arraybackslash}p{(\columnwidth - 18\tabcolsep) * \real{0.1000}}
  >{\raggedright\arraybackslash}p{(\columnwidth - 18\tabcolsep) * \real{0.1000}}
  >{\raggedright\arraybackslash}p{(\columnwidth - 18\tabcolsep) * \real{0.1000}}
  >{\raggedright\arraybackslash}p{(\columnwidth - 18\tabcolsep) * \real{0.1000}}
  >{\raggedright\arraybackslash}p{(\columnwidth - 18\tabcolsep) * \real{0.1000}}
  >{\raggedright\arraybackslash}p{(\columnwidth - 18\tabcolsep) * \real{0.1000}}
  >{\raggedright\arraybackslash}p{(\columnwidth - 18\tabcolsep) * \real{0.1000}}@{}}
\toprule\noalign{}
\begin{minipage}[b]{\linewidth}\raggedright
simulation
\end{minipage} & \begin{minipage}[b]{\linewidth}\raggedright
PIB (\%) 5a
\end{minipage} & \begin{minipage}[b]{\linewidth}\raggedright
PIB (\%) 10a
\end{minipage} & \begin{minipage}[b]{\linewidth}\raggedright
PIB (\%) 20a
\end{minipage} & \begin{minipage}[b]{\linewidth}\raggedright
chômage (pt) 5a
\end{minipage} & \begin{minipage}[b]{\linewidth}\raggedright
chômage (pt) 10a
\end{minipage} & \begin{minipage}[b]{\linewidth}\raggedright
chômage (pt) 20a
\end{minipage} & \begin{minipage}[b]{\linewidth}\raggedright
dette (pt PIB) 5a
\end{minipage} & \begin{minipage}[b]{\linewidth}\raggedright
dette (pt PIB) 10a
\end{minipage} & \begin{minipage}[b]{\linewidth}\raggedright
dette (pt PIB) 20a
\end{minipage} \\
\midrule\noalign{}
\endhead
\bottomrule\noalign{}
\endlastfoot
A & +0,96 & +2,45 & +4,08 & +2,56 & +2,56 & +1,93 & +5,28 & +8,42 &
+13,21 \\
B & −5,18 & −3,45 & +1,13 & +6,94 & +8,85 & +5,99 & +3,40 & −11,59 &
−77,09 \\
B\_variante\_2pc & −5,18 & −4,10 & +1,10 & +6,94 & +9,39 & +6,13 & +3,40
& −12,24 & −82,41 \\
A:oat\_382 & +0,97 & +2,44 & +4,09 & +2,56 & +2,57 & +1,94 & +5,22 &
+8,10 & +11,94 \\
B:oat\_382 & −4,49 & −1,81 & +2,38 & +6,51 & +7,44 & +4,58 & +3,29 &
−9,10 & −54,74 \\
mesure:H03 & −0,48 & −0,21 & −0,03 & +0,50 & +0,39 & +0,18 & +0,48 &
+0,36 & −0,77 \\
mesure:H05 & +0,21 & +0,83 & +1,82 & +0,41 & +1,10 & +1,10 & +0,11 &
+0,37 & +0,33 \\
mesure:H06 & +1,59 & +1,96 & +2,12 & +0,99 & +0,38 & +0,22 & +3,63 &
+3,70 & +3,95 \\
mesure:H08 & +0,03 & +0,08 & +0,17 & +0,24 & +0,16 & +0,08 & +0,04 &
+0,00 & −0,28 \\
mesure:H11 & −0,45 & −0,29 & −0,10 & +0,45 & +0,52 & +0,29 & +0,73 &
+3,30 & +7,97 \\
mesure:H12 & +0,05 & +0,08 & +0,14 & −0,00 & +0,01 & +0,02 & +0,26 &
+0,63 & +1,61 \\
regime:principal (R2 figé, taux exogènes, central) & +0,96 & +2,45 &
+4,08 & +2,56 & +2,56 & +1,93 & +5,28 & +8,42 & +13,21 \\
regime:sans règle budgétaire (R0) & +0,91 & +2,38 & +4,08 & +2,62 &
+2,66 & +1,98 & +5,34 & +8,66 & +13,71 \\
regime:stabilisation (R1) & +0,98 & +2,45 & +4,07 & +2,55 & +2,56 &
+1,93 & +5,29 & +8,42 & +13,22 \\
regime:R2 active symétrique & +2,09 & +4,47 & +5,55 & +1,70 & +0,69 &
+0,07 & +6,35 & +12,85 & +33,71 \\
regime:R2 + BCE et prime & +0,97 & +2,43 & +3,95 & +2,56 & +2,56 & +2,01
& +5,22 & +8,75 & +20,56 \\
regime:R2, ancrage faible & +1,02 & +2,52 & +3,84 & +2,54 & +2,45 &
+1,75 & +8,01 & +17,50 & +42,78 \\
regime:R2, ancrage fort & +0,93 & +2,41 & +4,21 & +2,57 & +2,61 & +2,00
& +3,91 & +3,93 & −1,48 \\
regime:R2, fermeture lente & +0,48 & +1,67 & +2,97 & +2,77 & +3,05 &
+2,70 & +6,21 & +12,15 & +27,60 \\
regime:R2, fermeture rapide & +1,39 & +2,97 & +4,72 & +2,37 & +2,21 &
+1,46 & +4,46 & +5,61 & +4,09 \\
regime:R2, fermeture sur le chômage & +0,19 & +2,49 & +4,63 & +3,04 &
+2,73 & +1,54 & +7,52 & +12,00 & +14,98 \\
environnement:p4\_sigma03 & +0,96 & +2,43 & +4,14 & +2,39 & +2,17 &
+1,51 & +4,86 & +6,38 & +6,57 \\
environnement:p4\_sigma06 & +0,96 & +2,47 & +4,04 & +2,67 & +2,98 &
+2,53 & +5,50 & +10,30 & +21,21 \\
environnement:p4\_prix\_v09 & +0,94 & +2,43 & +4,10 & +2,56 & +2,52 &
+1,90 & +5,04 & +7,82 & +11,97 \\
europe:conforme & +0,97 & +2,43 & +3,95 & +2,56 & +2,56 & +2,01 & +5,22
& +8,75 & +20,56 \\
europe:contentieux & +0,97 & +2,43 & +3,95 & +2,56 & +2,56 & +2,01 &
+5,22 & +8,75 & +20,56 \\
europe:tension & +0,97 & +2,43 & +3,94 & +2,56 & +2,56 & +2,02 & +6,39 &
+11,17 & +26,38 \\
mise\_en\_oeuvre:partielle\_50 & +0,63 & +1,58 & +2,45 & +1,56 & +1,61 &
+1,01 & +2,72 & +4,36 & +6,67 \\
mise\_en\_oeuvre:retard\_2ans & +0,38 & +2,01 & +3,92 & +2,53 & +2,80 &
+2,07 & +3,28 & +8,05 & +13,96 \\
\end{longtable}

\hypertarget{sec:annexe-f}{%
\chapter{Tests, contrôles comptables et
validations}\label{sec:annexe-f}}

\hypertarget{f.1-contruxf4les-comptables}{%
\section{F.1 Contrôles comptables}\label{f.1-contruxf4les-comptables}}

\begin{longtable}[]{@{}
  >{\raggedright\arraybackslash}p{(\columnwidth - 4\tabcolsep) * \real{0.3333}}
  >{\raggedright\arraybackslash}p{(\columnwidth - 4\tabcolsep) * \real{0.3333}}
  >{\raggedright\arraybackslash}p{(\columnwidth - 4\tabcolsep) * \real{0.3333}}@{}}
\toprule\noalign{}
\begin{minipage}[b]{\linewidth}\raggedright
Contrôle
\end{minipage} & \begin{minipage}[b]{\linewidth}\raggedright
Résultat
\end{minipage} & \begin{minipage}[b]{\linewidth}\raggedright
Source
\end{minipage} \\
\midrule\noalign{}
\endhead
\bottomrule\noalign{}
\endlastfoot
142 contrôles d'intégration (comptes d'agents, TES, flux-stocks) &
142/142 & phase 4, doc 21 §8 \\
Identité dette-solde dans les simulations des programmes (années ≥ 2027)
& écart relatif ≤ 1,1e-16 &
\texttt{phase6/resultats/simulabilite.json} \\
Transfert des programmes entre environnements (historique → part
constante) & écart nul & \texttt{phase5/resultats/lot53\_bilan.json} \\
Recoupement des totaux publiés par LFI (106,4 et 81,7 Md€) & exact &
\texttt{tests/test\_phase6.py} \\
Reproduction de la v0.8, de la v0.9 et de la référence R0 & exacte &
\texttt{tests/test\_phase5.py} \\
Revalidation du cache (empreinte des paramètres calibrés) & refuse un
coefficient modifié & \texttt{tests/test\_lot53.py} \\
\end{longtable}

\hypertarget{f.2-validations-uxe9conomiques}{%
\section{F.2 Validations
économiques}\label{f.2-validations-uxe9conomiques}}

Voir le chapitre \ref{sec:donnees} : tests équation par équation,
rétrosimulations conditionnelles (2010-2019, 2014-2019, 2019-2024),
prévisions ex ante contre les millésimes du FMI, choc identifié du CICE.

\hypertarget{f.3-liste-des-tests-automatisuxe9s}{%
\section{F.3 Liste des tests
automatisés}\label{f.3-liste-des-tests-automatisuxe9s}}

Suites de la finalisation, de la phase 6, des lots 5.2 et 5.3, de la
phase 5 et de la phase 4.3 (collecte pytest au moment de la génération)
:

\begin{itemize}
\tightlist
\item
  \texttt{tests/test\_finalisation.py::test\_duree\_requise\_legislation}
\item
  \texttt{tests/test\_finalisation.py::test\_regle\_simple\_inchangee\_et\_options\_neutres}
\item
  \texttt{tests/test\_finalisation.py::test\_sens\_des\_effets}
\item
  \texttt{tests/test\_finalisation.py::test\_effets\_generationnels}
\item
  \texttt{tests/test\_finalisation.py::test\_objectif\_budgetaire\_interpolation}
\item
  \texttt{tests/test\_finalisation.py::test\_corpus\_retraites\_et\_objectifs}
\item
  \texttt{tests/test\_finalisation.py::test\_variante\_environnement\_transmise\_au\_regime\_monetaire}
\item
  \texttt{tests/test\_finalisation.py::test\_investissement\_hors\_apu\_sans\_double\_compte}
\item
  \texttt{tests/test\_phase6.py::test\_corpus\_integrite}
\item
  \texttt{tests/test\_phase6.py::test\_paquet\_lfi\_recoupe\_les\_totaux\_publies}
\item
  \texttt{tests/test\_phase6.py::test\_conversions\_inversent\_le\_rendement\_statique}
\item
  \texttt{tests/test\_phase6.py::test\_variantes\_institutionnelles}
\item
  \texttt{tests/test\_phase6.py::test\_amendes\_pde\_et\_regime\_monetaire}
\item
  \texttt{tests/test\_phase6.py::test\_qualification\_politique\_commerciale}
\item
  \texttt{tests/test\_lot53.py::test\_robustesse\_programme\_nul\_et\_coherence}
\item
  \texttt{tests/test\_lot53.py::test\_parametres\_structurels\_appliques\_sans\_recalibration}
\item
  \texttt{tests/test\_lot53.py::test\_transfert\_des\_programmes\_exact}
\item
  \texttt{tests/test\_lot53.py::test\_revalidation\_refuse\_un\_coefficient\_modifie}
\item
  \texttt{tests/test\_lot52.py::test\_estimations\_et\_statuts}
\item
  \texttt{tests/test\_lot52.py::test\_regime\_bce\_sans\_choc\_identique\_au\_regime\_exogene}
\item
  \texttt{tests/test\_lot52.py::test\_choc\_permanent\_transmission}
\item
  \texttt{tests/test\_lot52.py::test\_regle\_bce\_lissage}
\item
  \texttt{tests/test\_lot52.py::test\_prime\_souveraine}
\item
  \texttt{tests/test\_phase5.py::test\_regles\_parametrables}
\item
  \texttt{tests/test\_phase5.py::test\_R0\_reproduit\_l\_environnement}
\item
  \texttt{tests/test\_phase5.py::test\_changer\_de\_regime\_sans\_recalibrer\_et\_rapidement}
\item
  \texttt{tests/test\_phase5.py::test\_identites\_et\_regle}
\item
  \texttt{tests/test\_phase5.py::test\_programme\_nul\_et\_regle\_figee}
\item
  \texttt{tests/test\_phase5.py::test\_programme\_figee\_effort\_identique}
\item
  \texttt{tests/test\_phase5.py::test\_conversion\_effort\_conso\_publique}
\item
  \texttt{tests/test\_phase43.py::test\_identites\_compte\_societes\_donnees}
\item
  \texttt{tests/test\_phase43.py::test\_compte\_societes\_isole\_continuite\_et\_identites}
\item
  \texttt{tests/test\_phase43.py::test\_hausse\_des\_taux\_augmente\_les\_interets\_et\_reduit\_les\_dividendes}
\item
  \texttt{tests/test\_phase43.py::test\_millesime\_fmi\_disponible\_et\_ex\_ante\_sans\_information\_future}
\item
  \texttt{tests/test\_phase43.py::test\_v09\_reproduite\_apres\_phase43}
\item
  \texttt{tests/test\_phase43.py::test\_configurations\_sans\_fuite\_d\_etat}
\item
  \texttt{tests/test\_phase43.py::test\_v08\_reproduite\_apres\_un\_essai\_v09}
\item
  \texttt{tests/test\_phase43.py::test\_p4\_identites\_et\_lien\_menages}
\item
  \texttt{tests/test\_phase43.py::test\_p4\_backtest\_proche\_de\_t1}
\item
  \texttt{tests/test\_phase43.py::test\_p4\_choc\_de\_taux\_stable\_sur\_15\_ans}
\end{itemize}

Échecs antérieurs connus et documentés (suite v0.8, configurations
historiques conservées en l'état) :
\texttt{test\_offre\_travail\_heterogene} (clés obsolètes) et
\texttt{test\_choc\_de\_productivite\_defavorable} (signe de la dette
dans la référence explosive).

\hypertarget{f.4-duxe9fauts-corriguxe9s-pendant-la-finalisation}{%
\section{F.4 Défauts corrigés pendant la
finalisation}\label{f.4-duxe9fauts-corriguxe9s-pendant-la-finalisation}}

\begin{longtable}[]{@{}
  >{\raggedright\arraybackslash}p{(\columnwidth - 4\tabcolsep) * \real{0.3333}}
  >{\raggedright\arraybackslash}p{(\columnwidth - 4\tabcolsep) * \real{0.3333}}
  >{\raggedright\arraybackslash}p{(\columnwidth - 4\tabcolsep) * \real{0.3333}}@{}}
\toprule\noalign{}
\begin{minipage}[b]{\linewidth}\raggedright
Défaut
\end{minipage} & \begin{minipage}[b]{\linewidth}\raggedright
Effet
\end{minipage} & \begin{minipage}[b]{\linewidth}\raggedright
Correction
\end{minipage} \\
\midrule\noalign{}
\endhead
\bottomrule\noalign{}
\endlastfoot
Une variante d'environnement portant sur les taux (OAT 3,82 \%) était
écrasée par l'appliqueur monétaire à partir de 2027 & les premières
simulations « OAT 3,82 \% » étaient identiques aux simulations de base &
base des taux du régime monétaire calculée après la variante ; test de
non-régression ajouté ; variantes recalculées \\
Sensibilité des paramètres monétaires évaluée à taux exogènes (lot 5.3)
& effet nul par construction & évaluation sous le régime BCE + prime \\
Pacte impôts de production / aides (Horizons) d'abord traduit en baisse
de consommation publique & effet de demande erroné & traduction par les
canaux des prélèvements sur les entreprises (deux variantes) \\
Variante européenne « tension » : le choc de prime souveraine figurait
aussi dans la simulation de référence & le choc s'annulait dans l'écart
; seules les amendes étaient mesurées & référence simulée sans le choc
(le choc est une réponse au programme) ; variantes européennes
recalculées pour les trois candidats
(\texttt{phase7/complement\_europe.py}) \\
Variantes d'additionnalité de l'investissement hors budget non simulées
& effet déduit approximativement de la mesure isolée & variantes 50 \%,
25 \%, 0 \% ajoutées à la campagne \\
\end{longtable}

\textbf{Reproductibilité constatée.} La campagne complète a été relancée
le 25 novembre 2026 avec le code gelé : les résultats de la lecture A,
de la lecture B et des variantes déjà calculées la veille ont été
reproduits à l'identique (écart maximal nul) pour les trois candidats.

\hypertarget{sec:annexe-g}{%
\chapter{Hypothèses des scénarios européens}\label{sec:annexe-g}}

\begin{longtable}[]{@{}
  >{\raggedright\arraybackslash}p{(\columnwidth - 6\tabcolsep) * \real{0.2500}}
  >{\raggedright\arraybackslash}p{(\columnwidth - 6\tabcolsep) * \real{0.2500}}
  >{\raggedright\arraybackslash}p{(\columnwidth - 6\tabcolsep) * \real{0.2500}}
  >{\raggedright\arraybackslash}p{(\columnwidth - 6\tabcolsep) * \real{0.2500}}@{}}
\toprule\noalign{}
\begin{minipage}[b]{\linewidth}\raggedright
Paramètre
\end{minipage} & \begin{minipage}[b]{\linewidth}\raggedright
Valeur
\end{minipage} & \begin{minipage}[b]{\linewidth}\raggedright
Statut
\end{minipage} & \begin{minipage}[b]{\linewidth}\raggedright
Justification
\end{minipage} \\
\midrule\noalign{}
\endhead
\bottomrule\noalign{}
\endlastfoot
Amende PDE par semestre & 0,05 \% du PIB de l'année précédente (plafond)
& {[}D{]} & règlement (UE) 2024/1264 \\
Nombre de semestres d'amende (tension) & 4 (2028-2029) & {[}H{]} & durée
illustrative ; aucune amende n'a été infligée à ce jour \\
Délai d'action après recommandation & 6 mois & {[}D{]} & idem \\
Choc d'écart OAT-Bund (tension) & +50 pb, 2027-2029 & {[}H{]} & ordre de
grandeur des épisodes politiques récents \\
Coefficients de prime (tension) & 0,0122 (dette), −0,118 (solde) &
{[}E{]} & panel zone euro incluant 2010-2013 \\
Délai de contentieux, niveau 1 & 1 an & {[}H{]} & --- \\
Délai de contentieux, niveau 2 & 2 ans & {[}H{]} & durée typique d'une
procédure en manquement jusqu'à l'arrêt \\
Part retirée, TVA sur les carburants & 50 \% de la mesure & {[}H{]} &
part des carburants dans l'assiette, à préciser \\
Part retirée, contribution européenne & 100 \% & {[}H{]} & réduction
unilatérale non réalisable \\
Part non conforme, économies « immigration » & inconnue (mesure
conservée) & {[}H{]} & modalités non publiées \\
Sanctions de l'article 260 TFUE & non représentées & {[}D{]} & somme
forfaitaire minimale 12,46 M€, astreinte de base 3 000 €/jour :
négligeables \\
Régime monétaire des trois variantes & BCE + prime souveraine en écart à
la référence & {[}H{]} & isoler l'effet des réactions européennes \\
\end{longtable}

Le code (\texttt{src/m26\_europe\_institutions.py}) permet de définir
d'autres scénarios (\texttt{ScenarioInstitutionnel}), par exemple une
durée d'amende différente, un choc de prime plus fort ou plus long, ou
la neutralisation de la prime par une intervention de la BCE.

\hypertarget{sec:annexe-h}{%
\chapter{Limites et mécanismes non modélisés}\label{sec:annexe-h}}

\hypertarget{h.1-muxe9canismes-absents-ou-simplifiuxe9s-liste-priorisuxe9e}{%
\section{H.1 Mécanismes absents ou simplifiés (liste
priorisée)}\label{h.1-muxe9canismes-absents-ou-simplifiuxe9s-liste-priorisuxe9e}}

\begin{longtable}[]{@{}
  >{\raggedright\arraybackslash}p{(\columnwidth - 8\tabcolsep) * \real{0.2000}}
  >{\raggedright\arraybackslash}p{(\columnwidth - 8\tabcolsep) * \real{0.2000}}
  >{\raggedright\arraybackslash}p{(\columnwidth - 8\tabcolsep) * \real{0.2000}}
  >{\raggedright\arraybackslash}p{(\columnwidth - 8\tabcolsep) * \real{0.2000}}
  >{\raggedright\arraybackslash}p{(\columnwidth - 8\tabcolsep) * \real{0.2000}}@{}}
\toprule\noalign{}
\begin{minipage}[b]{\linewidth}\raggedright
Priorité
\end{minipage} & \begin{minipage}[b]{\linewidth}\raggedright
Mécanisme
\end{minipage} & \begin{minipage}[b]{\linewidth}\raggedright
Mesures concernées
\end{minipage} & \begin{minipage}[b]{\linewidth}\raggedright
État
\end{minipage} & \begin{minipage}[b]{\linewidth}\raggedright
Conséquence
\end{minipage} \\
\midrule\noalign{}
\endhead
\bottomrule\noalign{}
\endlastfoot
1 & Paramètres fins des retraites : régimes, surcote, minima, réversion,
durées individuelles & réformes des trois candidats & âge, durée requise
et groupes de départ anticipé en moyenne par génération (finalisation) &
coûts et effets générationnels approchés \\
2 & Piliers de capitalisation & RN, Horizons & non chiffrés par les
candidats ; canal disponible (recettes en moins sans effet de demande) &
non simulés \\
3 & Financement et effet d'offre de l'investissement hors budget & LFI
(100 Md€/an) & effet de demande seul, additionnalité totale & borne
haute \\
4 & TVA et prix de l'énergie par produit & RN (TVA énergie), LFI
(blocage des carburants) & taux normal équivalent & répartition et
répercussion approchées \\
5 & Durée d'indemnisation du chômage, accompagnement des allocataires du
RSA & Horizons, LFI & montants seulement & effets sur le retour à
l'emploi non représentés \\
6 & Composition des économies non ventilées & RN, Horizons & hypothèses
de composition ; lecture B & incertitude forte \\
7 & Heures supplémentaires et allongement de la durée du travail &
Horizons & curseur symétrique de la réduction des 35 heures & effet
incertain \\
8 & Crédibilité des règles d'or (prime) & RN, Horizons & non représentée
& prime peut-être surestimée dans la lecture B \\
9 & Effets de richesse et marché du logement & impôts sur le patrimoine,
taux & absents & sous-estimation des effets des taux \\
10 & Incertitude et confiance & tous & absentes & sous-estimation des
coûts de court terme des ruptures \\
11 & Productivité endogène (éducation, innovation, organisation de
l'État) & tous & capital public et R\&D seulement & réformes
structurelles non récompensées ni pénalisées \\
12 & Économie européenne et change & tous & chocs exogènes, BCE & pas
d'effets d'entraînement des partenaires \\
13 & Retour macro → micro, distribution dynamique & tous & effets
statiques seulement & distribution des effets induits non mesurée \\
14 & Pas trimestriel & calendrier infra-annuel & non retenu & effets de
la première année approchés \\
\end{longtable}

\hypertarget{h.2-limites-de-validation}{%
\section{H.2 Limites de validation}\label{h.2-limites-de-validation}}

Rétrosimulations correctes sur la croissance, moins bonnes sur les prix
(épisode 2022-2023) et l'investissement des entreprises ; paramètres
clés non identifiés (σ, Phillips, fermeture) ; comportements estimés sur
1997-2019 (critique de Lucas) ; référence à politique inchangée
explosive avec les taux actuels.

\hypertarget{h.3-limites-liuxe9es-aux-programmes}{%
\section{H.3 Limites liées aux
programmes}\label{h.3-limites-liuxe9es-aux-programmes}}

Programmes provisoires (chiffrage complet de LFI attendu début 2027, «
volume 2 » d'Horizons en janvier 2027) ; enveloppes non ventilées (RN,
Horizons) ; hypothèses de traduction {[}H{]} documentées mesure par
mesure (annexe \ref{sec:annexe-c}).

\hypertarget{sec:annexe-i}{%
\chapter{Glossaire pédagogique}\label{sec:annexe-i}}

\textbf{Ajustement stock-flux.} Écart entre la variation de la dette
publique et le déficit de l'année (acquisitions d'actifs financiers,
effets de valorisation, décalages de trésorerie).

\textbf{Anticipations adaptatives.} Les agents forment leurs prévisions
à partir des valeurs passées (par exemple l'inflation de l'an dernier).
\textbf{Anticipations rationnelles} : les agents utilisent toute
l'information disponible et connaissent le fonctionnement de l'économie
; leurs erreurs de prévision sont imprévisibles.

\textbf{Assiette fiscale.} Base sur laquelle s'applique un impôt
(revenus, consommation, patrimoine). Une assiette « s'érode » lorsque
les contribuables réduisent ou déplacent la base imposable en réaction à
l'impôt.

\textbf{Boule de neige (effet).} Hausse automatique du ratio de dette
quand le taux d'intérêt apparent dépasse le taux de croissance nominal
du PIB, même à solde primaire nul.

\textbf{CES (fonction).} Fonction de production à élasticité de
substitution constante (\(\sigma\)) entre capital et travail.
\(\sigma < 1\) : les facteurs sont plutôt complémentaires ;
\(\sigma = 1\) : Cobb-Douglas.

\textbf{Charge excédentaire (perte sèche).} Coût économique d'un impôt
au-delà de la recette qu'il procure, dû aux changements de comportements
qu'il induit.

\textbf{Coût d'usage du capital.} Coût annuel de l'utilisation d'un euro
de capital : intérêt, dépréciation, fiscalité, variation du prix des
équipements.

\textbf{Coût salarial unitaire (CSU).} Coût du travail par unité
produite : rémunérations divisées par la production en volume.

\textbf{Correction d'erreur (modèle à).} Équation où la variable revient
chaque année vers une valeur de long terme, à une vitesse estimée, tout
en réagissant à court terme aux variations de ses déterminants.

\textbf{Décile de niveau de vie.} Dixième de la population classée par
revenu disponible par unité de consommation.

\textbf{Décote.} Minoration de la pension pour les assurés qui partent
avant d'avoir la durée requise ou l'âge du taux plein (1,25 \% par
trimestre manquant au régime général, dans la limite de 20 trimestres).

\textbf{DSGE.} Modèle dynamique, stochastique, d'équilibre général, à
fondements microéconomiques et anticipations rationnelles, utilisé par
les banques centrales.

\textbf{Durée d'assurance requise.} Nombre de trimestres validés
nécessaires pour une pension à taux plein (172 pour les générations nées
à partir de 1966 en législation actuelle).

\textbf{Écart de production (output gap).} Écart, en pourcentage, entre
le PIB effectif et le PIB potentiel.

\textbf{Effet d'éviction.} Réduction de la dépense privée provoquée par
une hausse de la dépense publique, par exemple via une hausse des taux
d'intérêt.

\textbf{Élasticité.} Variation en pourcentage d'une variable lorsqu'une
autre varie de 1 \%.

\textbf{Environnement (dans ce modèle).} Économie à politique inchangée,
avec ses équations, paramètres et trajectoire de référence calibrée.

\textbf{Équivalence ricardienne.} Proposition selon laquelle des ménages
prévoyants épargnent intégralement une baisse d'impôt financée par la
dette, anticipant les impôts futurs ; elle annule l'effet de la
politique budgétaire.

\textbf{Fermeture d'offre.} Mécanisme qui ramène la production vers son
potentiel à moyen terme.

\textbf{Incidence fiscale.} Répartition effective de la charge d'un
impôt entre agents (consommateurs, salariés, actionnaires), qui peut
différer du redevable légal.

\textbf{Loi d'Okun.} Relation empirique entre la croissance du PIB (ou
l'écart de production) et les variations du chômage.

\textbf{Microsimulation.} Calcul des effets d'une mesure sur un
échantillon représentatif de ménages, un par un, à partir de leur
situation et de la législation.

\textbf{Multiplicateur budgétaire.} Variation du PIB provoquée par une
variation d'un euro de dépense publique ou d'impôt.

\textbf{NAIRU / chômage structurel.} Taux de chômage compatible avec une
inflation stable ; il dépend des institutions du marché du travail.

\textbf{OAT.} Obligation assimilable du Trésor, titre de dette à long
terme de l'État français ; l'OAT à dix ans sert de référence de taux.

\textbf{OpenFisca.} Moteur libre de calcul de la législation
socio-fiscale, utilisé ici pour calculer impôts et prestations ménage
par ménage.

\textbf{PIB potentiel.} Production compatible avec une utilisation
normale des capacités de production et un chômage à son niveau
structurel.

\textbf{Pivot (âge).} Dans le module retraites, âge autour duquel se
concentrent les sorties d'activité d'une génération ; il dépend de l'âge
légal et de la durée requise.

\textbf{Prime de risque souveraine.} Écart entre le taux d'emprunt d'un
État et celui d'un emprunteur jugé sans risque (ici l'Allemagne).

\textbf{Procédure de déficit excessif (PDE).} Volet correctif du pacte
de stabilité : recommandations, délais et sanctions éventuelles pour un
État dont le déficit dépasse 3 \% du PIB.

\textbf{Propension marginale à consommer.} Part d'un euro de revenu
supplémentaire consacrée à la consommation.

\textbf{Proratisation.} Réduction proportionnelle de la pension de base
quand la durée validée est inférieure à la durée requise.

\textbf{Référence (trajectoire de).} Évolution de l'économie sans le
programme étudié, dans le même environnement et le même régime.

\textbf{Régime (dans ce modèle).} Ensemble des règles de politique
économique appliquées pendant la projection (règle budgétaire, politique
monétaire, paramètres fragiles).

\textbf{Règle de Taylor.} Règle selon laquelle la banque centrale fixe
son taux en fonction de l'inflation et de l'écart de production.

\textbf{Rétrosimulation.} Simulation du modèle sur une période passée,
avec les données observées pour les variables exogènes, afin de comparer
ses résultats à l'histoire.

\textbf{Solde primaire.} Solde public hors charges d'intérêts.
\textbf{Solde stabilisant} : solde qui maintient constant le ratio de
dette.

\textbf{Solde structurel / corrigé du cycle.} Solde public corrigé de
l'effet de la conjoncture (recettes et dépenses cycliques).

\textbf{Tableau des entrées-sorties (TES).} Tableau des échanges entre
branches de production, des emplois finals et des importations, qui
permet de calculer les effets d'entraînement et la formation des prix.

\textbf{Taux apparent de la dette.} Charge d'intérêts de l'année
rapportée à l'encours de dette.

\textbf{Taux de marge.} Excédent brut d'exploitation rapporté à la
valeur ajoutée.

\textbf{VAR structurel.} Modèle statistique multivarié permettant
d'identifier l'effet de chocs (par exemple budgétaires) sous des
hypothèses d'identification.

\hypertarget{sec:annexe-j}{%
\chapter{Guide technique de reproduction des
simulations}\label{sec:annexe-j}}

\hypertarget{j.1-environnement-logiciel}{%
\section{J.1 Environnement logiciel}\label{j.1-environnement-logiciel}}

\begin{itemize}
\tightlist
\item
  Dépôt : \texttt{vjwkjw26wk-dot/Pr-sidentielle-}, branche
  \texttt{phase43} (la branche \texttt{main} conserve la version 0.8,
  immuable). L'étiquette de gel de cette étude est indiquée dans
  \texttt{GEL.md} à la racine du dépôt.
\item
  Python 3.13 ; bibliothèques : numpy, pandas, scipy, statsmodels,
  matplotlib, OpenFisca-France 177.0.0.
\item
  Compilation de l'étude : pandoc (≥ 3), XeLaTeX (TeX Live), polices TeX
  Gyre.
\item
  Données : \texttt{data/raw/} (téléchargées par
  \texttt{src/telechargement.py}, empreintes SHA-256 dans
  \texttt{data/raw/\_registre.json}).
\end{itemize}

\hypertarget{sec:gel}{%
\section{J.2 Configuration de référence (gel technique)}\label{sec:gel}}

\begin{longtable}[]{@{}
  >{\raggedright\arraybackslash}p{(\columnwidth - 4\tabcolsep) * \real{0.3333}}
  >{\raggedright\arraybackslash}p{(\columnwidth - 4\tabcolsep) * \real{0.3333}}
  >{\raggedright\arraybackslash}p{(\columnwidth - 4\tabcolsep) * \real{0.3333}}@{}}
\toprule\noalign{}
\begin{minipage}[b]{\linewidth}\raggedright
Élément
\end{minipage} & \begin{minipage}[b]{\linewidth}\raggedright
Valeur
\end{minipage} & \begin{minipage}[b]{\linewidth}\raggedright
Code
\end{minipage} \\
\midrule\noalign{}
\endhead
\bottomrule\noalign{}
\endlastfoot
Configuration du modèle & P4 : v0.9 + prix de VA « frontière » + sans
effet de tension + compte des sociétés &
\texttt{M8.configurer\_phase4()} \\
Environnement &
\texttt{Environnement.charger("p4",\ "part\_constante")}, base 2022,
horizon 2047 & \texttt{src/m23\_regimes.py} \\
Environnement de traduction &
\texttt{Environnement.charger("p4",\ "historique")} & idem \\
Régime principal & \texttt{regime("R2")} : règle R2 figée, taux
exogènes, jeu central & idem \\
Variante européenne de référence & « conforme » &
\texttt{src/m26\_europe\_institutions.py} \\
Corpus & \texttt{src/m27\_corpus.py}, daté du 24 novembre 2026 & --- \\
\end{longtable}

Le premier chargement d'un environnement le calibre (environ 50 minutes
de calcul par environnement) et le met en cache dans \texttt{outputs/} ;
les chargements suivants vérifient l'empreinte des paramètres et la
trajectoire de contrôle.

\hypertarget{j.3-commandes}{%
\section{J.3 Commandes}\label{j.3-commandes}}

\begin{verbatim}
# environnements (une fois)
python3 phase5/construire_environnement.py p4 part_constante
python3 phase5/construire_environnement.py p4 historique
python3 phase5/construire_environnement.py p4_sigma03 part_constante
python3 phase5/construire_environnement.py p4_sigma06 part_constante
python3 phase5/construire_environnement.py p4_prix_v09 part_constante

# campagne de simulations des programmes (≈ 15 à 20 minutes par candidat)
python3 phase7/campagne.py melenchon
python3 phase7/campagne.py le_pen
python3 phase7/campagne.py philippe
python3 phase7/multiplicateurs.py

# stratégies de test S1-S5 et robustesse (phase 5)
python3 phase5/lot53_strategies.py
python3 phase5/lot53_structurels.py
python3 phase5/lot53_structurels_monetaires.py
python3 phase5/lot53_environnements.py

# tests
python3 -m pytest -q tests/test_finalisation.py tests/test_phase6.py tests/test_lot53.py \
    tests/test_lot52.py tests/test_phase5.py tests/test_phase43.py

# étude : tableaux, figures, chapitres de résultats, puis PDF
python3 etude/generer.py
bash etude/compiler.sh
\end{verbatim}

\hypertarget{j.4-simuler-un-programme-ou-une-mesure}{%
\section{J.4 Simuler un programme ou une
mesure}\label{j.4-simuler-un-programme-ou-une-mesure}}

\begin{Shaded}
\begin{Highlighting}[]
\ImportTok{import}\NormalTok{ m23\_regimes }\ImportTok{as}\NormalTok{ G, m25\_programmes }\ImportTok{as}\NormalTok{ P, m27\_corpus }\ImportTok{as}\NormalTok{ C}
\NormalTok{env\_h }\OperatorTok{=}\NormalTok{ G.Environnement.charger(}\StringTok{"p4"}\NormalTok{, }\StringTok{"historique"}\NormalTok{)}
\NormalTok{env }\OperatorTok{=}\NormalTok{ G.Environnement.charger(}\StringTok{"p4"}\NormalTok{, }\StringTok{"part\_constante"}\NormalTok{)}
\NormalTok{mesures, hors\_modele }\OperatorTok{=}\NormalTok{ C.mesures(}\StringTok{"philippe"}\NormalTok{, env.reference, }\StringTok{"conforme"}\NormalTok{)   }\CommentTok{\# lecture A}
\NormalTok{prog }\OperatorTok{=}\NormalTok{ P.programme(P.traduire([m[:}\DecValTok{3}\NormalTok{] }\ControlFlowTok{for}\NormalTok{ m }\KeywordTok{in}\NormalTok{ mesures], env\_h, }\StringTok{"philippe"}\NormalTok{, microsimulation}\OperatorTok{=}\StringTok{"calcul"}\NormalTok{))}
\NormalTok{effet }\OperatorTok{=}\NormalTok{ G.evaluer(env, G.regime(}\StringTok{"R2"}\NormalTok{), prog)[}\StringTok{"ecart"}\NormalTok{]                     }\CommentTok{\# écarts à la référence}
\CommentTok{\# lecture B : objectif de solde}
\NormalTok{obj }\OperatorTok{=}\NormalTok{ G.ObjectifBudgetaire(cibles}\OperatorTok{=}\NormalTok{\{}\DecValTok{2027}\NormalTok{: }\OperatorTok{{-}}\FloatTok{6.0}\NormalTok{, }\DecValTok{2032}\NormalTok{: }\OperatorTok{{-}}\FloatTok{2.5}\NormalTok{\}, parts}\OperatorTok{=}\NormalTok{\{}\StringTok{"prestations"}\NormalTok{: }\DecValTok{2}\OperatorTok{/}\DecValTok{3}\NormalTok{, }\StringTok{"conso\_publique"}\NormalTok{: }\DecValTok{1}\OperatorTok{/}\DecValTok{3}\NormalTok{\})}
\NormalTok{effet\_B }\OperatorTok{=}\NormalTok{ G.evaluer(env, G.regime(}\StringTok{"R2"}\NormalTok{), prog, objectif}\OperatorTok{=}\NormalTok{obj)[}\StringTok{"ecart"}\NormalTok{]}
\end{Highlighting}
\end{Shaded}

Une mesure isolée se simule en ne gardant que ses éléments dans
\texttt{mesures} ; une mise en œuvre partielle ou retardée avec la
fonction \texttt{modifier} de \texttt{phase7/campagne.py} ; un régime
alternatif avec
\texttt{G.regime(code,\ monetaire=…,\ structurels=…,\ **options)} (voir
\texttt{G.REGIMES\_ROBUSTESSE}).

\hypertarget{j.5-mettre-uxe0-jour-le-corpus-apruxe8s-les-publications-des-candidats}{%
\section{J.5 Mettre à jour le corpus après les publications des
candidats}\label{j.5-mettre-uxe0-jour-le-corpus-apruxe8s-les-publications-des-candidats}}

\begin{enumerate}
\def\labelenumi{\arabic{enumi}.}
\tightlist
\item
  Ajouter ou modifier les mesures dans \texttt{src/m27\_corpus.py}
  (source dans \texttt{SOURCES}, montant, calendrier, canal, règle,
  qualification européenne) et, le cas échéant, l'objectif dans
  \texttt{OBJECTIFS}.
\item
  Lancer
  \texttt{python3\ -m\ pytest\ -q\ tests/test\_phase6.py\ tests/test\_finalisation.py}
  (intégrité du corpus, recoupements).
\item
  Relancer la campagne du candidat
  (\texttt{phase7/campagne.py\ \textless{}candidat\textgreater{}}), puis
  \texttt{etude/generer.py} et \texttt{etude/compiler.sh}.
\end{enumerate}

La liste des éléments à actualiser figure au chapitre
\ref{sec:conclusion} (section « Questions ouvertes »).

\hypertarget{sec:annexe-k}{%
\chapter{Traçabilité des conclusions}\label{sec:annexe-k}}

Chaque conclusion du rapport (chapitres \ref{sec:discussion} et
\ref{sec:conclusion}) est reliée ci-dessous aux simulations qui la
fondent (clés des fichiers
\texttt{phase7/resultats/\textless{}candidat\textgreater{}.pkl},
régénérables par \texttt{phase7/campagne.py}), aux tableaux et sections
où ses chiffres sont présentés, et aux sources documentaires (codes du
corpus \texttt{src/m27\_corpus.py}, annexe \ref{sec:annexe-d}). Les
tableaux sont produits par \texttt{etude/resultats\_etude.py} sans
saisie manuelle de résultats.

\begin{longtable}[]{@{}
  >{\raggedright\arraybackslash}p{(\columnwidth - 6\tabcolsep) * \real{0.2500}}
  >{\raggedright\arraybackslash}p{(\columnwidth - 6\tabcolsep) * \real{0.2500}}
  >{\raggedright\arraybackslash}p{(\columnwidth - 6\tabcolsep) * \real{0.2500}}
  >{\raggedright\arraybackslash}p{(\columnwidth - 6\tabcolsep) * \real{0.2500}}@{}}
\toprule\noalign{}
\begin{minipage}[b]{\linewidth}\raggedright
conclusion
\end{minipage} & \begin{minipage}[b]{\linewidth}\raggedright
simulations (clés de
\texttt{phase7/resultats/\textless{}candidat\textgreater{}.pkl})
\end{minipage} & \begin{minipage}[b]{\linewidth}\raggedright
tableaux et sections
\end{minipage} & \begin{minipage}[b]{\linewidth}\raggedright
sources
\end{minipage} \\
\midrule\noalign{}
\endhead
\bottomrule\noalign{}
\endlastfoot
R1. Les mesures d'offre de travail sont les plus structurantes pour
l'activité à moyen terme : retraites pour LFI et le RN (effet négatif),
fin des 35 heures puis retraites pour Horizons (effet positif). &
mesure:L34, mesure:P08, mesure:H05, mesure:H06 ; A & t9\_mesures ;
t10\_generations ; ch.~9, 10 & LFI\_BUDGET ; RN\_RETRAITES\_PS ;
EP\_COMPTES ; Bozio (2011) ; CNAV 2024 \\
R2. Les signes des effets importants (PIB à cinq ans, dette de LFI et
d'Horizons à dix ans, chômage) sont stables dans les dix régimes et les
trois environnements ; les effets proches de zéro ne le sont pas. &
regime:* ; environnement:* & t12\_candidats ; t12\_part ; ch.~12 &
phase7/resultats/*\_synthese.csv \\
R3. L'ampleur des effets sur le chômage est peu fiable ; pour les
stratégies S1-S5, leur signe est aussi moins robuste que celui des
effets sur le PIB et la dette. & regime:* ; environnement:* ; S1-S5 &
t12\_S ; t12\_candidats ; ch.~12 &
phase5/resultats/lot53\_synthese\_S1\_S5.csv \\
R4. Les mesures documentées du RN et d'Horizons ne représentent qu'une
faible part de l'effort nécessaire à leurs objectifs de déficit ; celles
de LFI dépassent à court terme l'effort requis par son objectif. & A ; B
; B:oat\_382 & t9\_B ; t9\_cible ; t9\_niveaux ; ch.~9 & RN\_MONEYVOX ;
EP\_COMPTES ; LFI\_BUDGET \\
R5. Atteindre les objectifs de déficit du RN et d'Horizons dans ce
modèle demande un effort de plusieurs points de PIB, aux effets
récessifs importants ; l'ampleur dépend de l'hypothèse de taux longs. &
B ; B:oat\_382 ; B\_variante\_2pc & t9\_B ; t9\_niveaux ; ch.~9, 12 &
doc 22 (référence) ; M24 \\
R6. Les résultats de LFI dépendent fortement de la traduction des
recettes sans canal propre et de l'additionnalité de l'investissement
hors budget. & A ; additionnalite:* ; mesure:L03, L08-L16, L33 &
t9\_mesures ; t9\_variantes ; ch.~9 & LFI\_BUDGET \\
R7. Les variantes européennes modifient peu les effets à cinq ans ; la
variante « tension » pèse surtout sur la dette à long terme. &
europe:conforme, contentieux, tension & t9\_variantes ; ch.~8, 9 &
règlement (UE) 2024/1264 ; directive 2022/542 \\
R8. Les effets distributifs statiques diffèrent par leur concentration :
prélèvements sur le haut de la distribution (LFI), gains diffus (RN),
pertes concentrées sur les allocataires du RSA (Horizons). & traductions
A (microsimulation) & t10\_deciles ; ch.~10 & OpenFisca-France ; M4 \\
R9. Une mise en œuvre à 50 \% divise les effets à peu près par deux ; un
retard de deux ans décale les effets de court terme. &
mise\_en\_oeuvre:* & t9\_variantes ; ch.~9 & --- \\
R10. Les multiplicateurs simulés diminuent avec l'horizon ; ils sont
plus élevés pour la consommation publique que pour les transferts et les
baisses de prélèvements sur les entreprises. & multiplicateurs & ch.~7 &
phase7/resultats/multiplicateurs.csv ; Blanchard et Leigh (2013) ; Ramey
(2019) \\
\end{longtable}

\textbf{Statut des conclusions.} R1, R2, R4 et R10 sont robustes au sens
du chapitre \ref{sec:robustesse} ; R3 est une conclusion sur la
fragilité elle-même ; R5, R6 et R8 dépendent d'hypothèses de traduction
explicitées (compositions {[}H{]}, additionnalité, rendements) ; R7
dépend d'hypothèses de scénario européen {[}H{]} ; R9 est une propriété
quasi linéaire du modèle à cette échelle.

\backmatter
\end{document}
